% !TeX program = xelatex
% Compile twice: xelatex -interaction=nonstopmode main.tex
\documentclass[UTF8,10.5pt,a4paper,oneside,fontset=fandol,scheme=plain]{ctexbook}
% ==================== Art-Book Editorial Preamble for AI Series Volume 4 ====================
\usepackage[top=24mm,bottom=24mm,left=23mm,right=23mm,headheight=22pt,headsep=7mm]{geometry}
\usepackage{amsmath,amssymb,mathtools,amsthm}
\usepackage{booktabs,tabularx,array,multirow,longtable}
\usepackage{graphicx}
\usepackage{xcolor}
\usepackage{tikz}
\usepackage{pgfplots}
\pgfplotsset{compat=1.18}
\usetikzlibrary{arrows.meta,positioning,shapes.geometric,calc,decorations.pathreplacing,backgrounds}
\usepackage[most]{tcolorbox}
\usepackage{fancyhdr}
\usepackage{listings}
\usepackage{enumitem}
\usepackage{hyperref}

% ==================== Curated Gallery Color Palette ====================
\definecolor{ink}{HTML}{1B2430}        % Deep charcoal-slate ink
\definecolor{teal}{HTML}{0B5555}       % Refined verdigris teal
\definecolor{tealDark}{HTML}{083B3B}   % Deep gallery spine teal
\definecolor{tealLight}{HTML}{F2F8F8}  % Delicate ice-teal parchment
\definecolor{tealLine}{HTML}{A3C6C5}   % Soft architectural rule
\definecolor{tealMist}{HTML}{DCEBEB}   % Watermark chapter numeral tint
\definecolor{navy}{HTML}{1D3B5E}       % Muted Prussian indigo
\definecolor{navyLight}{HTML}{F2F5FA}  % Soft indigo tint
\definecolor{ochre}{HTML}{9E4A22}      % Burnished terracotta-ochre
\definecolor{gold}{HTML}{B97A24}       % Fine burnished amber-gold accent
\definecolor{ochreLight}{HTML}{FCF5F0} % Warm parchment tint
\definecolor{plum}{HTML}{5B3A6B}       % Muted editorial plum
\definecolor{plumLight}{HTML}{F8F4FB}  % Soft plum tint
\definecolor{slate}{HTML}{526071}      % Muted architectural gray
\definecolor{codebg}{HTML}{F7F9FB}     % Clean light code canvas
\definecolor{codeborder}{HTML}{D4DCE5}

\hypersetup{
  unicode=true,
  colorlinks=true,
  urlcolor=teal,
  linkcolor=teal,
  citecolor=teal,
  pdftitle={天书中的智能 · 第4卷：凸优化与经典机器学习},
  pdfauthor={天书中的智能：人工智能的七卷手记}
}

% ==================== Editorial Chapter Openers & Typography ====================
\newcommand{\chapWatermarkNum}{\ifnum\value{chapter}<10 0\fi\thechapter}
\ctexset{
  contentsname={目\quad 录},
  chapter={
    name={第,章},
    number=\arabic{chapter},
    format=\LARGE\sffamily\bfseries\color{tealDark},
    nameformat=\small\sffamily\bfseries\color{gold},
    aftername=\hfill{\fontsize{38}{38}\selectfont\sffamily\bfseries\color{tealMist}\chapWatermarkNum}\par\vspace{-4mm},
    aftertitle=\par\vspace{2.8mm}{\noindent\color{gold}\rule{26mm}{1.5pt}}{\color{tealLine}\rule{\dimexpr\linewidth-26mm\relax}{0.55pt}}\par,
    beforeskip=6pt,
    afterskip=10pt
  },
  section={
    format=\Large\sffamily\bfseries\color{teal},
    beforeskip=14pt plus 3pt minus 2pt,
    afterskip=6pt plus 1pt
  },
  subsection={
    format=\large\sffamily\bfseries\color{ink},
    beforeskip=10pt plus 2pt minus 1pt,
    afterskip=4pt plus 1pt
  },
  subsubsection={
    format=\normalsize\sffamily\bfseries\color{navy},
    beforeskip=7pt,
    afterskip=2pt
  }
}
\renewcommand{\lstlistingname}{代码}
\renewcommand{\figurename}{图}
\renewcommand{\tablename}{表}

% Artisanal Chapter Epigraph macro (2 args: quote + right-aligned attribution)
\newcommand{\chapEpigraph}[2]{%
  \vspace{-2mm}%
  \begin{tcolorbox}[
    enhanced,
    colback=ochreLight!45!white,
    colframe=ochreLight!45!white,
    borderline west={1.8pt}{0pt}{gold},
    boxrule=0pt,
    arc=0.8mm,
    left=4mm,right=3.5mm,top=2mm,bottom=2mm,
    before skip=0pt,after skip=8pt
  ]
    {\kaishu\normalsize\color{ink!85}“#1”}\par\vspace{1mm}%
    \hfill{\footnotesize\sffamily\color{slate}——\,#2}%
  \end{tcolorbox}%
}

% Exercise tier header macro
\newcommand{\exTier}[2]{%
  \par\vspace{3.8mm}\noindent
  \colorbox{tealLight}{\strut\;\sffamily\bfseries\footnotesize\color{tealDark}#1\;}\;%
  {\small\sffamily\bfseries\color{slate}#2}\par\nopagebreak\vspace{1.2mm}%
}
\newcommand{\exHint}[1]{\par\smallskip\noindent{\footnotesize\sffamily\color{navy}\textbf{【提示】}\,#1}}

\setlength{\parindent}{0pt}
\setlength{\parskip}{5pt plus 1.5pt minus 0.5pt}
\linespread{1.26}
\setlist[itemize]{leftmargin=1.8em,topsep=3pt,itemsep=2.5pt,parsep=1pt}
\setlist[enumerate]{leftmargin=2.0em,topsep=3pt,itemsep=2.5pt,parsep=1pt}
\setlength{\tabcolsep}{5.5pt}
\renewcommand{\arraystretch}{1.22}

% ==================== Running Headers & Footers ====================
\pagestyle{fancy}
\fancyhf{}
\fancyhead[L]{\small\sffamily\color{slate}天书中的智能 \;\textcolor{gold}{$\cdot$}\; 第 4 卷}
\fancyhead[R]{\small\sffamily\color{slate}\nouppercase{\leftmark}}
\fancyfoot[C]{\small\sffamily\color{slate}\textcolor{gold}{$\diamond$}\quad\textcolor{tealDark}{\thepage}\quad\textcolor{gold}{$\diamond$}}
\renewcommand{\headrulewidth}{0.4pt}
\renewcommand{\headrule}{\hbox to\headwidth{\color{tealLine}\leaders\hrule height \headrulewidth\hfill}}
\renewcommand{\chaptermark}[1]{\markboth{第 \thechapter\ 章\quad #1}{}}

\fancypagestyle{plain}{
  \fancyhf{}
  \fancyfoot[C]{\small\sffamily\color{slate}\textcolor{gold}{$\diamond$}\quad\textcolor{tealDark}{\thepage}\quad\textcolor{gold}{$\diamond$}}
  \renewcommand{\headrulewidth}{0pt}
}

% ==================== Theorem & Pedagogical Box Environments ====================
\theoremstyle{definition}
\newtheorem{theorem}{定理}[chapter]
\newtheorem{lemma}[theorem]{引理}
\newtheorem{proposition}[theorem]{命题}
\newtheorem{corollary}[theorem]{推论}
\newtheorem{definition}[theorem]{定义}
\newtheorem{example}[theorem]{例}
\newtheorem{remark}[theorem]{注记}
\renewcommand{\proofname}{\textbf{\sffamily\color{teal}证明}}

\newtcolorbox{thmBox}[1][]{
  enhanced,breakable,
  colback=tealLight,
  colframe=tealLine,
  borderline west={2.4pt}{0pt}{teal},
  boxrule=0.45pt,
  arc=1.5mm,
  left=4mm,right=3.8mm,top=2.8mm,bottom=2.8mm,
  before skip=6.5pt,after skip=6.5pt,
  #1
}

\newtcolorbox{defBox}[1][]{
  enhanced,breakable,
  colback=navyLight,
  colframe=navy!30,
  borderline west={2.4pt}{0pt}{navy},
  boxrule=0.45pt,
  arc=1.5mm,
  left=4mm,right=3.8mm,top=2.8mm,bottom=2.8mm,
  before skip=6.5pt,after skip=6.5pt,
  #1
}

\newtcolorbox{insightBox}[1]{
  enhanced,breakable,
  colback=white,
  colframe=navy!35,
  borderline west={2.4pt}{0pt}{gold},
  coltitle=navy,
  fonttitle=\sffamily\bfseries\small,
  attach title to upper={\quad},
  title={【直觉与下游坐标】#1},
  boxrule=0.45pt,
  arc=1.2mm,
  left=4mm,right=3.8mm,top=2.6mm,bottom=2.6mm,
  before skip=6.5pt,after skip=6.5pt
}

\newtcolorbox{recallBox}[1]{
  enhanced,breakable,
  colback=codebg,
  colframe=codeborder,
  borderline west={2.2pt}{0pt}{slate},
  coltitle=slate,
  fonttitle=\sffamily\bfseries\small,
  attach title to upper={\quad},
  title={【前置速忆】#1},
  boxrule=0.4pt,
  arc=1.2mm,
  left=3.8mm,right=3.5mm,top=2.4mm,bottom=2.4mm,
  before skip=5.5pt,after skip=5.5pt
}

\newtcolorbox{failBox}[1]{
  enhanced,breakable,
  colback=ochreLight,
  colframe=ochre!35,
  borderline west={2.4pt}{0pt}{ochre},
  coltitle=ochre,
  fonttitle=\sffamily\bfseries\small,
  attach title to upper={\quad},
  title={【反例、数值陷阱与排错】#1},
  boxrule=0.45pt,
  arc=1.2mm,
  left=4mm,right=3.8mm,top=2.6mm,bottom=2.6mm,
  before skip=6.5pt,after skip=6.5pt
}

\newtcolorbox{extBox}[1]{
  enhanced,breakable,
  colback=plumLight,
  colframe=plum!30,
  borderline west={2.4pt}{0pt}{plum},
  coltitle=plum,
  fonttitle=\sffamily\bfseries\small,
  attach title to upper={\quad},
  title={【选读拓展 · 初读可跳过】#1},
  boxrule=0.45pt,
  arc=1.2mm,
  left=4mm,right=3.8mm,top=2.6mm,bottom=2.6mm,
  before skip=6.5pt,after skip=6.5pt
}

\newtcolorbox{chapIntro}[1]{
  enhanced,
  colback=tealLight,
  colframe=tealLine,
  borderline west={2.8pt}{0pt}{tealDark},
  boxrule=0.45pt,
  arc=1.5mm,
  left=4.2mm,right=4mm,top=3mm,bottom=3mm,
  before skip=2pt,after skip=10pt,
  title={\sffamily\bfseries\small\color{tealDark}#1},
  colbacktitle=tealLight,
  attach boxed title to top left={yshift=-1mm,xshift=2mm},
  boxed title style={empty,size=minimal}
}

% ==================== Code Listing Style ====================
\lstdefinestyle{pycode}{
  language=Python,
  basicstyle=\footnotesize\ttfamily\color{ink},
  keywordstyle=\bfseries\color{teal},
  commentstyle=\itshape\color{slate},
  stringstyle=\color{ochre},
  backgroundcolor=\color{codebg},
  frame=single,
  rulecolor=\color{codeborder},
  framesep=5pt,
  xleftmargin=4pt,
  xrightmargin=4pt,
  numbers=left,
  numberstyle=\tiny\ttfamily\color{slate!80},
  numbersep=6pt,
  columns=fullflexible,
  keepspaces=true,
  showstringspaces=false,
  breaklines=true,
  aboveskip=6.5pt,
  belowskip=6.5pt
}
\lstset{style=pycode}

% ==================== Math Shortcuts ====================
\newcommand{\R}{\mathbb{R}}
\newcommand{\C}{\mathbb{C}}
\newcommand{\N}{\mathbb{N}}
\newcommand{\E}{\mathbb{E}}
\newcommand{\Prob}{\mathbb{P}}
\newcommand{\Var}{\operatorname{Var}}
\newcommand{\Cov}{\operatorname{Cov}}
\newcommand{\Bias}{\operatorname{Bias}}
\newcommand{\DKL}{D_{\mathrm{KL}}}
\newcommand{\sign}{\operatorname{sign}}
\newcommand{\prox}{\operatorname{prox}}
\newcommand{\diag}{\operatorname{diag}}
\newcommand{\tr}{\operatorname{tr}}
\newcommand{\rank}{\operatorname{rank}}
\newcommand{\argmin}{\operatorname*{arg\,min}}
\newcommand{\argmax}{\operatorname*{arg\,max}}
\newcommand{\cond}{\kappa}
\newcommand{\norm}[1]{\left\|#1\right\|}
\newcommand{\abs}[1]{\left|#1\right|}
\newcommand{\inner}[2]{\left\langle #1,\, #2 \right\rangle}
\newcommand{\had}{\odot}

\makeatletter
\setlength{\@fpsep}{16pt}
\makeatother


\begin{document}
\color{ink}

% ==================== Cover Page (Gallery Art-Book Edition) ====================
\frontmatter
\thispagestyle{empty}
\begin{tikzpicture}[remember picture, overlay]
  % Architectural deep teal spine + burnished gold hairline + vertical monograph label
  \fill[tealDark] (current page.north west) rectangle ([xshift=11mm]current page.south west);
  \fill[gold] ([xshift=11mm]current page.north west) rectangle ([xshift=12.2mm]current page.south west);
  \node[rotate=90, anchor=center, font=\sffamily\bfseries\scriptsize, color=tealMist, opacity=0.85]
    at ([xshift=5.5mm]current page.west) {INTELLIGENCE FROM THE BOOK \quad $\cdot$ \quad VOLUME IV};

  % Bespoke fine-art geometric SVM max-margin & convex contours motif (top-right)
  \begin{scope}[shift={([xshift=-95mm, yshift=-20mm]current page.north east)}, scale=0.52, transform shape]
    % Faint convex level sets
    \begin{scope}[rotate=-18, opacity=0.11]
      \draw[navy, line width=0.8pt] (0.2, 0.0) ellipse (2.5cm and 1.25cm);
      \draw[navy, line width=0.7pt] (0.2, 0.0) ellipse (1.7cm and 0.85cm);
      \draw[navy, line width=0.6pt] (0.2, 0.0) ellipse (0.95cm and 0.48cm);
    \end{scope}
    % Max-margin hyperplane & margin slab
    \fill[teal, opacity=0.07] (-2.3,-1.2) -- (2.1, 1.0) -- (1.6, 1.7) -- (-2.8,-0.5) -- cycle;
    \draw[gold, opacity=0.65, line width=1.4pt] (-2.55,-0.85) -- (1.85, 1.35);
    \draw[teal, opacity=0.38, line width=0.85pt, dashed] (-2.3,-1.2) -- (2.1, 1.0);
    \draw[teal, opacity=0.38, line width=0.85pt, dashed] (-2.8,-0.5) -- (1.6, 1.7);
    % Support vectors (circled) & data points
    \foreach \x/\y in {-1.2/-0.65, 0.6/0.25} {
      \fill[ochre, opacity=0.55] (\x,\y) circle (2.6pt);
      \draw[ochre, opacity=0.45, thick] (\x,\y) circle (5.2pt);
    }
    \foreach \x/\y in {-0.9/0.45, 0.9/1.35} {
      \fill[teal, opacity=0.55] (\x,\y) circle (2.6pt);
      \draw[teal, opacity=0.45, thick] (\x,\y) circle (5.2pt);
    }
    \foreach \x/\y in {-0.2/-1.05, 1.1/-0.45, 1.7/0.2, -1.7/-1.3} {
      \fill[ochre, opacity=0.25] (\x,\y) circle (2.2pt);
    }
    \foreach \x/\y in {-1.8/0.75, -0.3/1.35, 0.2/1.65, -1.2/1.25} {
      \fill[teal, opacity=0.25] (\x,\y) circle (2.2pt);
    }
  \end{scope}
\end{tikzpicture}
\vspace*{1mm}
\noindent{\small\sffamily\bfseries\color{teal}天书中的智能 \;\textcolor{gold}{$\cdot$}\; 人工智能的七卷手记 \hfill 第 4 卷 \quad 2026 年 11 月版}
\vspace{1.5mm}
\noindent{\color{gold}\rule{30mm}{1.8pt}}{\color{tealLine}\rule{\dimexpr\linewidth-30mm\relax}{0.7pt}}
\vspace{6.5mm}

{\fontsize{25}{33}\selectfont\sffamily\bfseries\color{ink}凸优化与经典机器学习}\par
\vspace{2.5mm}
{\large\sffamily\color{teal}从凸分析、拉格朗日对偶与泛化理论，走向 SVM、集成树与 EM 算法}\par
\vspace{4.5mm}
{\small\color{slate}对标 Stanford CS229 / UC Berkeley CS189 \& EECS127 / CMU 07-380 / 南大《最优化方法》与《机器学习导论》}\par

\vspace{5.5mm}

\begin{tcolorbox}[
  enhanced,
  colback=tealLight,
  colframe=tealLine,
  borderline west={2.6pt}{0pt}{tealDark},
  arc=2mm,
  boxrule=0.5pt,
  left=4.5mm,right=4.5mm,top=3mm,bottom=3mm
]
\textbf{\sffamily\color{tealDark}本卷定位：建立统计学习思维与严格最优化功底}\par\vspace{1.2mm}
\small
经过前三卷的数理与智能体算法训练，本卷将正式攻克经典机器学习与最优化的核心理论——不仅知道怎么训练模型，更要从数学上回答“收敛有多快、为什么能泛化、如何突破线性边界以及怎样在无标签数据中发现隐藏结构”：
\begin{itemize}
  \item \textbf{凸优化与对偶理论}：严格证明 $L$-光滑下降引理、$O(1/T)$ 与强凸线性收敛率、弱对偶与 KKT 互补松弛条件，并用近端梯度（ISTA）亲手跑出 $L_1$ Lasso 稀疏路径。
  \item \textbf{统计学习与泛化诊断}：严格推导偏差-方差-噪声分解、有限假设空间 PAC 泛化界，掌握学习曲线与 $K$ 折交叉验证诊断法。
  \item \textbf{三大监督经典支柱}：最大间隔 SVM 与核技巧（RBF 非线性升维）、CART 决策树、随机森林（降方差证明）与梯度提升树 GBDT（函数空间负梯度降偏差）。
  \item \textbf{无监督与隐变量推断}：从 K-Means 球形假设的失效出发，借第 2 卷 Jensen 与 KL 散度完整推导 GMM 与 EM 算法的 ELBO 单调上升定理，直通第 5 卷变分自编码器（VAE）。
\end{itemize}
\end{tcolorbox}

\vspace{4.5mm}

\noindent\textbf{\sffamily\color{tealDark}第 4 卷章节导览}
\vspace{1.5mm}

{\footnotesize\renewcommand{\arraystretch}{1.15}
\begin{tabularx}{\linewidth}{@{}p{15mm}p{47mm}X@{}}
\toprule
\textbf{章节} & \textbf{标题} & \textbf{核心理论、完整证明与实测实验} \\
\midrule
序\quad 章 & 本卷坐标与八校考点覆盖对照 & 优化与 ML 如何耦合、八校对标表与本卷知识依赖 DAG \\
第 1 章 & 凸分析与一阶优化收敛率理论 & 凸集/凸函数、$L$-光滑下降引理、PL 不等式、$O(1/T)$ 与线性收敛、动量加速实测 \\
第 2 章 & 拉格朗日对偶、KKT 与近端梯度 & 弱/强对偶、Slater 条件、KKT 互补松弛、次梯度、软阈值 ISTA 与 Lasso 路径实测 \\
第 3 章 & 统计学习：偏差-方差与泛化界 & 偏差-方差严格证明与实测拆解、$L_1/L_2$ 几何、$K$ 折交叉验证、PAC 界与 VC 维直觉 \\
第 4 章 & 支持向量机（SVM）与核技巧 & 最大间隔、对偶 QP 与支持向量、软间隔合页损失、RBF 核升维与同心圆实测 \\
第 5 章 & 决策树、随机森林与 GBDT & CART 切分（信息增益/Gini）、Bagging 降方差证明、GBDT 负梯度提升与三模型实测 \\
第 6 章 & K-Means、高斯混合与 EM 算法 & K-Means++、GMM 软聚类、ELBO 分解与单调收敛证明、倾斜椭圆实测与第 5 卷衔接 \\
附\quad 录 & 全卷 48 道分层习题完整精解 & 第 1--6 章 A/B/C/D 四层习题分步推导、反例验算与代码自检基准 \\
\bottomrule
\end{tabularx}}

\vfill
\noindent{\color{tealLine}\rule{\linewidth}{0.5pt}}
{\footnotesize\color{slate}配套实验脚本 \texttt{code/vol4\_experiments.py} 与实测数据随本卷提供。上一卷：第 3 卷；下一卷：第 5 卷《深度学习、自动微分与现代架构》。}

% ==================== Table of Contents ====================
\clearpage
\tableofcontents

% ==================== Main Matter ====================
\mainmatter
\chapter*{序章：最优化与统计学习的交汇与八校考点对照}
\addcontentsline{toc}{chapter}{序章：最优化与统计学习的交汇与八校考点对照}
\markboth{序章\quad 本卷坐标与八校考点对照}{}

\chapEpigraph{你不必相信上帝，但你应当相信“天书”的存在。}{Paul Erdős（保罗·埃尔德什，谈论记录最美证明的 \emph{THE BOOK}）}

\begin{chapIntro}{本章定位与丛书题记}
保罗·埃尔德什曾说，宇宙中藏着一本记录最纯粹、最不可避免证明的“天书”（\emph{THE BOOK}）。本套《天书中的智能》七卷手记，便是我们在那本天书的页边空白处，用严密推导与可验证代码一笔一划还原出的智能基石。

本卷（第 4 卷）正式收编了第 1 卷刻意留白的\textbf{凸优化收敛理论与对偶分析}，并将其与\textbf{统计学习泛化思维及四大经典机器学习支柱（SVM/核、树与集成、聚类与 EM）}融为一体。读完本卷，你将拥有判断任何经典与现代机器学习模型“为何有效、何时过拟合、怎样求解”的完整理论视角。
\end{chapIntro}

\section*{0.1\quad 为什么把凸优化与经典机器学习合为一卷？}
\addcontentsline{toc}{section}{0.1\quad 为什么把凸优化与经典机器学习合为一卷？}

在 UC Berkeley（EECS127 + CS189）、Stanford（CS229）与南京大学（《最优化方法导论》+《机器学习导论》）的培养体系中，最优化与机器学习如车之两轮：
\begin{enumerate}
  \item \textbf{不学凸优化，经典模型只剩黑箱公式}：不懂 $L$-光滑与强凸，就只能死背步长经验；不懂拉格朗日对偶与 KKT 互补松弛，就看不懂 SVM 为什么只有少数支持向量起作用、更看不懂核技巧如何自然浮现；不懂次梯度与近端算子，就不知道不可微的 $L_1$ Lasso 怎么求精确零解。
  \item \textbf{不学统计学习，训练误差会误导全局}：偏差-方差分解、PAC 泛化界与交叉验证教会我们区分“拟合训练集”与“真正学到规律”，这一准则在进入第 5 卷深度学习后依然是最高铁律。
  \item \textbf{经典算法至今仍是不可替代的工业与科研基线}：表格数据上 GBDT（XGBoost/LightGBM）仍常胜过深度神经网络；EM 算法的 ELBO 下界推导将在第 5 卷直接演化为变分自编码器（VAE）的损失函数！
\end{enumerate}

\section*{0.2\quad 第 4 卷八校核心考点覆盖对照表}
\addcontentsline{toc}{section}{0.2\quad 第 4 卷八校核心考点覆盖对照表}

\begin{center}
\footnotesize
\begin{tabularx}{\linewidth}{@{}p{34mm}p{34mm}p{16mm}X@{}}
\toprule
\textbf{八校核心考点} & \textbf{主要对标课程} & \textbf{本卷位置} & \textbf{下游直接用途} \\
\midrule
凸分析、$L$-光滑下降引理、强凸与收敛率 & Berkeley EECS127 / 南大最优化 & 第 1 章 & 理解动量/加速、正则化加速优化与非凸训练诊断 \\
拉格朗日对偶、KKT 互补松弛与近端 ISTA & EECS127 / Stanford CS229 & 第 2 章 & 第 4 章 SVM 对偶、$L_1$ 稀疏求解与带约束优化 \\
偏差-方差分解、正则化、交叉验证与 PAC 界 & CS229 / Berkeley CS189 & 第 3 章 & 全系列过拟合诊断、超参选择与泛化误差评估 \\
最大间隔 SVM、合页损失与核技巧（RBF） & CS229 / CMU 07-380 / 南大 ML & 第 4 章 & 非线性分类、再生核直觉与大间隔分类思维 \\
CART 决策树、Bagging/随机森林与 GBDT & CS189 / CMU 07-380 & 第 5 章 & 表格数据强基线、函数空间梯度提升思想 \\
K-Means++、高斯混合模型与 EM（ELBO） & CS229 / 南大 ML & 第 6 章 & 无监督聚类、隐变量建模、直通第 5 卷 VAE \\
\bottomrule
\end{tabularx}
\end{center}

\section*{0.3\quad 第 4 卷知识阶梯（零先修倒挂）}
\addcontentsline{toc}{section}{0.3\quad 第 4 卷知识阶梯（零先修倒挂）}

\begin{center}
\resizebox{0.98\linewidth}{!}{%
\begin{tikzpicture}[
  node distance=6.5mm and 7.5mm,
  ch/.style={
    rectangle, rounded corners=1.8mm, draw=teal, thick, fill=tealLight,
    text width=34mm, minimum height=14.5mm, align=center, font=\sffamily\footnotesize, text=ink, inner sep=4pt
  },
  goal/.style={
    rectangle, rounded corners=1.8mm, draw=ochre, very thick, fill=ochreLight,
    text width=37mm, minimum height=14.5mm, align=center, font=\sffamily\footnotesize, text=ink, inner sep=4pt
  },
  arr/.style={-{Stealth[length=2.2mm]}, thick, teal, preaction={draw=white, line width=2.5pt, arrows=-}}
]
  \node[ch] (c1) {\textbf{第 1 章}\\凸分析与一阶优化\\收敛率理论};
  \node[ch, right=of c1] (c2) {\textbf{第 2 章}\\拉格朗日对偶、KKT\\与近端梯度 ISTA};
  \node[ch, right=of c2] (c3) {\textbf{第 3 章}\\统计学习：偏差-方差\\与 PAC 泛化界};

  \node[ch, below=9mm of c1] (c4) {\textbf{第 4 章}\\支持向量机（SVM）\\与核技巧};
  \node[ch, below=9mm of c2] (c5) {\textbf{第 5 章}\\决策树、随机森林\\与 GBDT};
  \node[goal, below=9mm of c3] (c6) {\textbf{\color{ochre}第 6 章（直通第 5 卷 VAE）}\\K-Means、GMM 与\\EM（ELBO）算法};

  \draw[arr] (c1) -- (c2);
  \draw[arr] (c2) -- (c3);
  \draw[arr] ([xshift=-9mm]c2.south) -- ([xshift=3mm]c4.north);
  \draw[arr] ([xshift=-12mm]c3.south) -- ([xshift=12mm]c4.north);
  \draw[arr] ([xshift=-4mm]c3.south) -- ([xshift=10mm]c5.north);
  \draw[arr] ([xshift=9mm]c1.south) -- ([xshift=-9mm]c5.north);
  \draw[arr] (c3.south) -- (c6.north);
\end{tikzpicture}}
\end{center}

\chapter[凸分析与一阶优化收敛率]{凸分析与一阶优化收敛率理论}\label{ch:v4-convex}

\chapEpigraph{事实上，优化问题的真正分水岭并不在于线性与非线性，而在于凸与非凸。}{R. Tyrrell Rockafellar（1993）}

\begin{chapIntro}{本章导读：从“沿着负梯度走”到“精确预报收敛步数”}
在第 1 卷与第 2 卷中，我们已经多次使用梯度下降（Gradient Descent, GD）去最小化线性回归的均方误差或逻辑回归的交叉熵损失。然而当时的理解仍停留在局部直觉上：“负梯度是最速下降方向，所以只要步长别太大，慢慢走总能走到谷底。”

本章将把这一经验直觉升级为一套严密的定量理论。我们将依次回答四个核心问题：
\begin{enumerate}
  \item \textbf{何时“停下来”就意味着全局最优？}——通过凸集、凸函数与一阶/二阶等价刻画，严格证明凸优化不存在“陷阱局部极小点”。
  \item \textbf{步长 $\eta$ 究竟可以取多大？每走一步至少能下降多少？}——利用 $L$-光滑条件证明\textbf{下降引理}（Descent Lemma），定量导出步长阈值 $\eta = 1/L$ 与经典 $O(1/T)$ 亚线性收敛率。
  \item \textbf{为什么加上一个看似普通的 $L_2$ 正则项 $\frac{\lambda}{2}\norm{w}_2^2$，收敛会从“蜗牛爬”变成指数飞跃？}——建立 $\mu$-强凸与 Polyak-\L{}ojasiewicz（PL）不等式，证明线性收敛率 $(1-\mu/L)^T$，并揭示条件数 $\kappa = L/\mu$ 的统治地位。
  \item \textbf{当条件数 $\kappa \approx 10^3$ 时如何破局？}——剖析动量法（Heavy-Ball）与 Nesterov 加速梯度（NAG）将依赖从 $\kappa$ 压缩至 $\sqrt{\kappa}$ 的几何机理，并在真实病态逻辑回归数据上用纯 NumPy 代码验证。
\end{enumerate}
\end{chapIntro}

% ============================================================
\section{病态地形的警钟：为什么梯度下降会剧烈“之字形”震荡？}
% ============================================================

在引入任何抽象定义之前，我们先看一个在机器学习中最常见、却最令人困惑的失败现象。

\begin{recallBox}{与第 1 卷和第 2 卷的衔接}
在第 1 卷第 5--6 章我们学过：对称正定矩阵 $A \in \R^{d\times d}$ 的特征值 $\lambda_{\max}(A) \ge \dots \ge \lambda_{\min}(A) > 0$ 决定了二次型 $x^\top A x$ 在各个主轴方向上的弯曲程度，其比值 $\kappa(A) = \lambda_{\max}(A)/\lambda_{\min}(A)$ 称为\textbf{条件数}。在第 2 卷第 5 章中，我们通过极大似然导出了逻辑回归的负对数似然损失。今天我们要看看，一旦输入特征的量纲相差一个数量级（比如身高以米计、年薪以元计），梯度下降会遭遇怎样的灾难。
\end{recallBox}

\begin{example}[二维狭长山谷上的之字形停滞]\label{ex:v4-ch1-zigzag}
考虑极简的二维强凸二次函数
\[
  f(x_1, x_2) = \frac{1}{2}\left(x_1^2 + 100 x_2^2\right),
\]
其唯一全局极小点显然为 $x^* = (0,0)^\top$，$f^* = 0$。梯度为 $\nabla f(x) = (x_1,\, 100x_2)^\top$，海森矩阵为 $\nabla^2 f(x) = \diag(1, 100)$，最小曲率 $\mu = 1$，最大曲率 $L = 100$，条件数 $\kappa = L/\mu = 100$。

从起点 $x_0 = (10, 1)^\top$ 出发，执行固定步长梯度下降 $x_{t+1} = x_t - \eta \nabla f(x_t)$，两个坐标完全解耦：
\[
  x_{t,1} = (1 - \eta)^t \cdot 10, \qquad x_{t,2} = (1 - 100\eta)^t \cdot 1.
\]
为了避免 $x_{t,2}$ 爆炸发散，我们必须要求 $|1 - 100\eta| < 1$，即步长存在刚性上限 $\eta < \frac{2}{100} = 0.02$：
\begin{itemize}
  \item 若贪心地取 $\eta = 0.025$，则 $x_{t,2} = (-1.5)^t$，仅 10 步后 $|x_{10,2}| \approx 57.7$，迅速数值溢出！
  \item 若保守地取理论安全步长 $\eta = 1/L = 0.01$，则 $x_{t,2}$ 在第 1 步就直接归零，但平坦方向 $x_{t,1} = 0.99^t \cdot 10$ 每步只收缩 $1\%$；要让 $x_{t,1}$ 衰减到 $0.01$，需要 $0.99^t \le 10^{-3}$，即 $t \ge \frac{\ln 1000}{0.01} \approx 688$ 步！
  \item 若取接近上限的 $\eta = 0.019$，则 $1 - 100\eta = -0.9$，$x_{t,2}$ 在山谷两侧剧烈正负震荡，呈现典型的“之字形（zig-zag）”轨迹，而沿山谷底部的推进依然缓慢。
\end{itemize}
\end{example}

例~\ref{ex:v4-ch1-zigzag} 揭示了一个极其深刻的普适规律：\textbf{在无加速的一阶梯度下降中，最陡峭的方向（最大曲率 $L$）强行锁死了全局允许的最大步长 $\eta \lesssim 1/L$，而最平坦的方向（最小曲率 $\mu$）又强行拖慢了全局收敛的进度。} 在多维逻辑回归或深度网络中，只要不同特征尺度的跨度达到 $15$ 倍，设计矩阵 $X^\top X$ 的特征值跨度就会达到 $15^2 = 225$ 倍乃至上千倍——这正是我们需要一套严格凸分析与收敛率理论的直接动因。

% ============================================================
\section{凸集与凸函数：无陷阱地形的数学刻画}
% ============================================================

对于一般光滑非凸函数，满足 $\nabla f(x) = 0$ 的驻点可能是局部极小点、局部极大点，甚至是鞍点（Saddle Point）；一阶算法一旦陷入次优驻点，仅凭局部梯度信息无法察觉远方还有更深的谷底。Rockafellar 在本章卷首语中一针见血地指出：优化问题的真正难易边界在于\textbf{凸性（Convexity）}。

\subsection{凸集与常见可行域}

\begin{defBox}
\begin{definition}[凸集]\label{def:v4-convex-set}
集合 $C \subseteq \R^d$ 称为\textbf{凸集（Convex Set）}，如果对任意两点 $x, y \in C$ 和任意 $\theta \in [0, 1]$，连接两点的线段完全落入 $C$ 中：
\[
  \theta x + (1-\theta) y \in C.
\]
\end{definition}
\end{defBox}

在经典机器学习与优化中，我们频繁遇到以下五类基本凸集（均可直接用定义~\ref{def:v4-convex-set} 一行验证）：
\begin{enumerate}
  \item \textbf{超平面与闭半空间}：$\{x \in \R^d : a^\top x = b\}$ 与 $\{x \in \R^d : a^\top x \le b\}$（第 4 章 SVM 线性分类间隔的基础几何构件）。
  \item \textbf{范数球}：对任意范数 $\norm{\cdot}$（如 $\ell_1, \ell_2, \ell_\infty$），球体 $\{x : \norm{x - x_0} \le r\}$ 都是凸集，因为由三角不等式与齐次性有 $\norm{\theta x + (1-\theta)y - x_0} \le \theta\norm{x-x_0} + (1-\theta)\norm{y-x_0} \le r$。
  \item \textbf{概率单纯形}：$\Delta^{d-1} = \{p \in \R^d : p_i \ge 0,\; \sum_{i=1}^d p_i = 1\}$（第 6 章 GMM 混合权重与多分类概率分布的自然定义域）。
  \item \textbf{半正定矩阵锥}：$\mathbb{S}_+^d = \{A \in \R^{d\times d} : A = A^\top,\; v^\top A v \ge 0\ \forall v\in\R^d\}$（协方差矩阵与核矩阵的天然栖息地）。
  \item \textbf{任意多个凸集的交集}：若 $\{C_\alpha\}_{\alpha \in \mathcal{I}}$ 是一族凸集，则 $\bigcap_{\alpha \in \mathcal{I}} C_\alpha$ 必为凸集。特别地，线性规划与 SVM 的多面体可行域 $\{x : Ax \le b\}$ 作为有限个闭半空间的交集，必然是凸集。
\end{enumerate}

\subsection{凸函数的零阶、一阶与二阶等价刻画}

\begin{defBox}
\begin{definition}[凸函数与严格凸函数]\label{def:v4-convex-fn}
设 $C \subseteq \R^d$ 为凸集，函数 $f: C \to \R$ 称为\textbf{凸函数（Convex Function）}，如果对任意 $x, y \in C$ 和任意 $\theta \in [0, 1]$，满足\textbf{零阶定义（弦在弧上方）}：
\begin{equation}\label{eq:v4-jensen2}
  f\big(\theta x + (1-\theta)y\big) \le \theta f(x) + (1-\theta)f(y).
\end{equation}
若对所有 $x \neq y$ 及 $\theta \in (0, 1)$ 上式严格成立，则称 $f$ 为\textbf{严格凸函数}。若 $-f$ 是凸函数，则称 $f$ 为\textbf{凹函数（Concave Function）}。
\end{definition}
\end{defBox}

通过数学归纳法，式~\eqref{eq:v4-jensen2} 立即可推广至有限凸组合 $\sum_{i=1}^k \theta_i = 1\;(\theta_i \ge 0)$ 以及一般随机变量的 \textbf{Jensen 不等式}：$f(\E[X]) \le \E[f(X)]$（我们在第 2 卷已用它证明 KL 散度非负，第 6 章还将再次用它推导 EM 算法的 ELBO 下界）。

凸函数与凸集之间存在一座完美的几何桥梁——\textbf{上境图（Epigraph）}：
\[
  \operatorname{epi}(f) \coloneqq \left\{(x, t) \in C \times \R : f(x) \le t\right\}.
\]
容易验证：\textbf{$f$ 是凸函数，当且仅当其上境图 $\operatorname{epi}(f)$ 是 $\R^{d+1}$ 中的凸集。} 这一几何视角在下一章处理不可微凸函数（如 $\ell_1$ 范数）与拉格朗日对偶时极具威力。

对于光滑函数，每次都去验证零阶弦不等式~\eqref{eq:v4-jensen2} 颇为繁琐。下面这条核心定理给出了一阶切平面与二阶海森矩阵的等价判据，它们分别是构建全局最优性与收敛率证明的基石。

\begin{thmBox}
\begin{theorem}[凸函数的一阶与二阶等价刻画]\label{thm:v4-convex-char}
设 $C \subseteq \R^d$ 为开凸集：
\begin{enumerate}
  \item \textbf{一阶条件（切超平面全局下界）}：若 $f: C \to \R$ 可微，则 $f$ 是凸函数当且仅当对任意 $x, y \in C$ 均有
  \begin{equation}\label{eq:v4-first-order-convex}
    f(y) \ge f(x) + \nabla f(x)^\top (y - x).
  \end{equation}
  \item \textbf{二阶条件（海森矩阵半正定）}：若 $f: C \to \R$ 二阶连续可微，则 $f$ 是凸函数当且仅当对任意 $x \in C$ 均有
  \begin{equation}\label{eq:v4-second-order-convex}
    \nabla^2 f(x) \succeq 0 \quad \text{（即对任意 } v \in \R^d,\; v^\top \nabla^2 f(x)\, v \ge 0\text{）}.
  \end{equation}
\end{enumerate}
\end{theorem}
\end{thmBox}

\begin{proof}
我们先证第一部分（一阶条件）。

\textbf{必要性（凸 $\implies$ 式~\eqref{eq:v4-first-order-convex}）}：任取 $x, y \in C$ 和 $\theta \in (0, 1]$，由凸性定义~\eqref{eq:v4-jensen2} 得
\[
  f\big(x + \theta(y - x)\big) = f\big((1-\theta)x + \theta y\big) \le (1-\theta)f(x) + \theta f(y).
\]
移项并两边同除以 $\theta > 0$：
\[
  f(y) - f(x) \ge \frac{f\big(x + \theta(y-x)\big) - f(x)}{\theta}.
\]
令 $\theta \to 0^+$，右侧由方向导数的定义趋于 $\nabla f(x)^\top (y - x)$，即得式~\eqref{eq:v4-first-order-convex}。

\textbf{充分性（式~\eqref{eq:v4-first-order-convex} $\implies$ 凸）}：任取 $x, y \in C$ 与 $\theta \in [0, 1]$，记 $z = \theta x + (1-\theta)y \in C$。分别在点 $z$ 处对 $x$ 与 $y$ 应用式~\eqref{eq:v4-first-order-convex}：
\[
  f(x) \ge f(z) + \nabla f(z)^\top (x - z), \qquad f(y) \ge f(z) + \nabla f(z)^\top (y - z).
\]
将第一式乘以 $\theta$、第二式乘以 $1-\theta$ 后相加，注意到 $\theta(x-z) + (1-\theta)(y-z) = \theta x + (1-\theta)y - z = 0$，梯度项精确抵消，立得 $\theta f(x) + (1-\theta)f(y) \ge f(z)$。

再证第二部分（二阶条件）。对任意固定 $x \in C$ 与方向 $v \in \R^d$，构造一元辅助函数 $\phi(t) = f(x + t v)$（定义在包含 $0$ 的开区间上）。由链式法则，$\phi'(t) = \nabla f(x+tv)^\top v$，$\phi''(t) = v^\top \nabla^2 f(x+tv)\, v$。
\begin{itemize}
  \item 若 $f$ 凸，则由一阶条件分别对 $(x, x+tv)$ 与 $(x+tv, x)$ 使用式~\eqref{eq:v4-first-order-convex} 并相加，得单调梯度不等式 $\big(\nabla f(x+tv)-\nabla f(x)\big)^\top (tv) \ge 0$。对 $t>0$ 两边同除以 $t^2$ 并令 $t \to 0^+$，即得 $\phi''(0) = v^\top \nabla^2 f(x)\, v \ge 0$。
  \item 反之，若 $\nabla^2 f(\cdot) \succeq 0$ 处处成立，则对任意 $x, y \in C$，记 $v = y-x$，由带积分余项（或 Lagrange 余项）的二阶泰勒展开，存在 $\xi \in (0, 1)$ 使得
  \[
    f(y) = \phi(1) = \phi(0) + \phi'(0) + \frac{1}{2}\phi''(\xi) = f(x) + \nabla f(x)^\top (y-x) + \frac{1}{2}v^\top \nabla^2 f(x+\xi v)\, v \ge f(x) + \nabla f(x)^\top (y-x).
  \]
  由此满足一阶条件~\eqref{eq:v4-first-order-convex}，故 $f$ 为凸函数。
\end{itemize}
\end{proof}

\subsection{保凸运算：不求海森矩阵也能一眼识别凸目标}

实际机器学习损失函数往往由多层运算复合而成，直接展开海森矩阵可能极其繁琐。利用以下三条\textbf{保凸运算（Convexity-Preserving Operations）}，我们可以像搭积木一样迅速判定凸性：

\begin{proposition}[三大常用保凸运算]\label{prop:v4-preserve}
\begin{enumerate}
  \item \textbf{非负加权和}：若 $f_1, \dots, f_m$ 为凸函数且 $\alpha_1, \dots, \alpha_m \ge 0$，则 $f(x) = \sum_{i=1}^m \alpha_i f_i(x)$ 为凸函数。
  \item \textbf{仿射复合}：若 $f: \R^k \to \R$ 为凸函数，且 $A \in \R^{k\times d},\, b \in \R^k$，则 $g(x) = f(Ax + b)$ 是关于 $x \in \R^d$ 的凸函数。
  \item \textbf{逐点上确界（逐点取最大）}：若对任意索引 $\alpha \in \mathcal{I}$ 函数 $f_\alpha(x)$ 均为凸函数，则 $h(x) = \sup_{\alpha \in \mathcal{I}} f_\alpha(x)$ 也是凸函数。
\end{enumerate}
\end{proposition}

\begin{proof}
第 1、2 条直接代入零阶定义~\eqref{eq:v4-jensen2} 即可验证。第 3 条从上境图立即可见：
\[
  \operatorname{epi}(h) = \{(x, t) : \sup_{\alpha} f_\alpha(x) \le t\} = \bigcap_{\alpha \in \mathcal{I}} \{(x, t) : f_\alpha(x) \le t\} = \bigcap_{\alpha \in \mathcal{I}} \operatorname{epi}(f_\alpha).
\]
凸集的任意交仍为凸集，故 $h$ 是凸函数。也可直接代数验证：$h(\theta x + (1-\theta)y) = \sup_\alpha f_\alpha(\theta x + (1-\theta)y) \le \sup_\alpha [\theta f_\alpha(x) + (1-\theta)f_\alpha(y)] \le \theta \sup_\alpha f_\alpha(x) + (1-\theta)\sup_\alpha f_\alpha(y) = \theta h(x) + (1-\theta)h(y)$。
\end{proof}

\begin{insightBox}{机器学习三大损失函数的凸性秒证与第 4 章伏笔}
\begin{itemize}
  \item \textbf{合页损失（Hinge Loss，第 4 章 SVM 核心）}：$\ell_{\text{hinge}}(w, b) = \max\big(0,\, 1 - y_i(w^\top x_i + b)\big)$ 是两个仿射函数（常数 $0$ 与 $1 - y_i(w^\top x_i + b)$）的逐点最大值，由命题~\ref{prop:v4-preserve} 第 3 条立知它是凸函数（尽管在 $y_i(w^\top x_i + b)=1$ 处不可微）！
  \item \textbf{指数与 Log-Sum-Exp 损失}：$\mathrm{LSE}(z) = \ln\sum_{k=1}^K e^{z_k}$ 是凸函数（留作习题~\ref{ex:4.1.2}），经仿射复合 $z = Xw$ 后即得逻辑回归与多分类 Softmax 交叉熵损失的全局凸性。
  \item \textbf{正则化目标}：经验风险（非负平均）与 $\lambda \norm{w}_1$ 或 $\frac{\lambda}{2}\norm{w}_2^2$（$\lambda \ge 0$）之和仍为凸函数。
\end{itemize}
\end{insightBox}

\subsection{凸函数的全局最优性定理}

一阶条件~\eqref{eq:v4-first-order-convex} 带来的最直接、也最珍贵的红利，就是局部驻点与全局极小点的完全统一。

\begin{thmBox}
\begin{theorem}[凸函数驻点的全局最优性]\label{thm:v4-global-opt}
设 $f: \R^d \to \R$ 为可微凸函数。则 $x^* \in \R^d$ 是 $f$ 的全局极小点，当且仅当
\[
  \nabla f(x^*) = 0.
\]
进一步地，若 $f$ 严格凸，则全局极小点（若存在）必然唯一。
\end{theorem}
\end{thmBox}

\begin{proof}
\textbf{必要性}：若 $x^*$ 为全局极小点，它必为局部极小点，由多元微积分 Fermat 引理立得 $\nabla f(x^*) = 0$。

\textbf{充分性}：若 $\nabla f(x^*) = 0$，对任意 $y \in \R^d$ 应用定理~\ref{thm:v4-convex-char} 的一阶切平面下界~\eqref{eq:v4-first-order-convex}：
\[
  f(y) \ge f(x^*) + \nabla f(x^*)^\top (y - x^*) = f(x^*) + 0 = f(x^*).
\]
因此 $x^*$ 为全局极小点。

\textbf{严格凸时的唯一性}：假设存在两个不同的全局极小点 $x_1^* \neq x_2^*$，且 $f(x_1^*) = f(x_2^*) = f^*$。取中点 $\bar{x} = \frac{1}{2}x_1^* + \frac{1}{2}x_2^*$，由严格凸性有 $f(\bar{x}) < \frac{1}{2}f(x_1^*) + \frac{1}{2}f(x_2^*) = f^*$，与 $f^*$ 是全局最小值矛盾！
\end{proof}

\begin{example}[$L_2$ 正则化逻辑回归的精确海森矩阵与曲率界]\label{ex:v4-ch1-logreg}
考虑第 1.6 节即将实测的目标函数（样本 $(x_i, y_i) \in \R^D \times \{0, 1\},\; i=1,\dots,N$，正则化系数 $\lambda > 0$）：
\[
  F(w) = \frac{1}{N}\sum_{i=1}^N \Big[\ln\big(1 + e^{x_i^\top w}\big) - y_i (x_i^\top w)\Big] + \frac{\lambda}{2}\norm{w}_2^2.
\]
记 $p_i(w) = \sigma(x_i^\top w) = \frac{1}{1+e^{-x_i^\top w}} \in (0, 1)$，由链式法则直接求一阶与二阶导数：
\[
  \nabla F(w) = \frac{1}{N}X^\top \big(p(w) - y\big) + \lambda w, \qquad
  \nabla^2 F(w) = \frac{1}{N}X^\top D(w) X + \lambda I_D,
\]
其中 $D(w) = \diag\big(p_1(1-p_1), \dots, p_N(1-p_N)\big)$。对任意 $p \in (0, 1)$，算术-几何平均不等式给出 $0 < p(1-p) \le \frac{1}{4}$。于是对任意非零向量 $v \in \R^D$，有双边谱估计：
\[
  \lambda \norm{v}_2^2 \le v^\top \nabla^2 F(w)\, v = \frac{1}{N}\sum_{i=1}^N p_i(1-p_i)(x_i^\top v)^2 + \lambda\norm{v}_2^2 \le \left(\frac{\lambda_{\max}(X^\top X)}{4N} + \lambda\right)\norm{v}_2^2.
\]
这表明 $F(w)$ 不仅是严格凸函数（具有唯一全局最优解 $w^*$），而且全局曲率被严丝合缝地夹在 $\mu = \lambda$ 与 $L = \frac{\sigma_{\max}(X)^2}{4N} + \lambda$ 之间！这个精确的常数 $L$ 将直接用来设定我们实验中的无调参步长 $\eta = 1/L$。
\end{example}

% ============================================================
\section{$L$-光滑性、下降引理与 $O(1/T)$ 收敛率证明}
% ============================================================

仅知道函数是凸的，还不足以保证固定步长梯度下降会收敛——例如 $f(x) = x^4$ 是凸的，但在远端梯度增长过快，任何固定步长从足够远处出发都会发散；而 $f(x) = |x|$ 也是凸的，但在原点不可微，固定步长只会在原点两侧来回弹跳。要保证一阶算法有可预期的下降幅度，必须限制梯度的变化速率，即引入\textbf{$L$-光滑性（$L$-Smoothness）}。

\begin{defBox}
\begin{definition}[$L$-光滑函数]\label{def:v4-lsmooth}
可微函数 $f: \R^d \to \R$ 称为\textbf{$L$-光滑（$L$-Smooth）}的（$L > 0$），如果其梯度满足 $L$-Lipschitz 连续条件：对任意 $x, y \in \R^d$，
\begin{equation}\label{eq:v4-lip-grad}
  \norm{\nabla f(y) - \nabla f(x)}_2 \le L \norm{y - x}_2.
\end{equation}
当 $f$ 二阶连续可微时，式~\eqref{eq:v4-lip-grad} 等价于海森矩阵谱范数处处有界：$-LI \preceq \nabla^2 f(x) \preceq LI$（对凸函数即 $0 \preceq \nabla^2 f(x) \preceq LI$）。
\end{definition}
\end{defBox}

\begin{thmBox}
\begin{lemma}[下降引理 / 二次上界引理，Descent Lemma]\label{lem:v4-descent}
若 $f: \R^d \to \R$ 是 $L$-光滑的（不要求 $f$ 为凸函数！），则对任意 $x, y \in \R^d$ 均有
\begin{equation}\label{eq:v4-descent-ub}
  f(y) \le f(x) + \nabla f(x)^\top (y - x) + \frac{L}{2}\norm{y - x}_2^2.
\end{equation}
特别地，令梯度下降迭代为 $x_{t+1} = x_t - \eta \nabla f(x_t)$，则
\begin{equation}\label{eq:v4-step-drop}
  f(x_{t+1}) \le f(x_t) - \eta\left(1 - \frac{L\eta}{2}\right)\norm{\nabla f(x_t)}_2^2.
\end{equation}
当取最优保守步长 $\eta = \frac{1}{L}$ 时，每步保证至少下降：
\begin{equation}\label{eq:v4-1overL-drop}
  f(x_{t+1}) \le f(x_t) - \frac{1}{2L}\norm{\nabla f(x_t)}_2^2.
\end{equation}
\end{lemma}
\end{thmBox}

\begin{proof}
构造沿线段的一元函数 $g(s) = f\big(x + s(y-x)\big),\; s \in [0, 1]$。由微积分基本定理：
\begin{align*}
  f(y) - f(x) &= g(1) - g(0) = \int_0^1 g'(s)\,ds = \int_0^1 \nabla f\big(x + s(y-x)\big)^\top (y - x)\,ds \\
  &= \nabla f(x)^\top (y-x) + \int_0^1 \Big(\nabla f\big(x + s(y-x)\big) - \nabla f(x)\Big)^\top (y - x)\,ds.
\end{align*}
对积分号内的内积依次应用 Cauchy-Schwarz 不等式与 $L$-光滑定义~\eqref{eq:v4-lip-grad}：
\begin{align*}
  f(y) - f(x) - \nabla f(x)^\top (y-x)
  &\le \int_0^1 \norm{\nabla f\big(x + s(y-x)\big) - \nabla f(x)}_2 \norm{y - x}_2\,ds \\
  &\le \int_0^1 \big(L \cdot s\norm{y - x}_2\big) \norm{y - x}_2\,ds
  = L\norm{y - x}_2^2 \int_0^1 s\,ds = \frac{L}{2}\norm{y - x}_2^2.
\end{align*}
移项即证得二次上界~\eqref{eq:v4-descent-ub}。将 $x = x_t,\; y = x_{t+1} = x_t - \eta \nabla f(x_t)$ 代入式~\eqref{eq:v4-descent-ub}，注意 $y - x = -\eta \nabla f(x_t)$：
\[
  f(x_{t+1}) \le f(x_t) - \eta\norm{\nabla f(x_t)}_2^2 + \frac{L\eta^2}{2}\norm{\nabla f(x_t)}_2^2
  = f(x_t) - \eta\left(1 - \frac{L\eta}{2}\right)\norm{\nabla f(x_t)}_2^2.
\]
最后，将右侧关于 $\eta$ 的二次函数 $\phi(\eta) = \eta - \frac{L}{2}\eta^2$ 求导并令其为零，得唯一极大值点 $\eta^* = 1/L$，此时系数 $\eta(1 - L\eta/2) = \frac{1}{2L}$，即得式~\eqref{eq:v4-1overL-drop}。
\end{proof}

\begin{figure}[htbp]
\centering
\begin{tikzpicture}[scale=1.06, >=Stealth]
  % Axes
  \draw[->, slate!75, thin] (-2.7, 0) -- (4.2, 0) node[right, font=\footnotesize\sffamily] {$x$};
  \draw[->, slate!75, thin] (-2.4, -0.3) -- (-2.4, 4.45) node[above, font=\footnotesize\sffamily] {$f(x)$};

  % Shaded sandwich region
  \fill[tealLight, opacity=0.85]
    plot[domain=-2.0:1.65, samples=60] (\x, {2.0052 - 1.232*(\x+0.8) + 0.625*(\x+0.8)^2})
    -- plot[domain=1.65:-2.0, samples=60] (\x, {max(0.05, 2.0052 - 1.232*(\x+0.8))}) -- cycle;

  % Upper L-smooth quadratic bound
  \draw[ochre, very thick, dashed, domain=-2.0:1.65, samples=80]
    plot (\x, {2.0052 - 1.232*(\x+0.8) + 0.625*(\x+0.8)^2});
  \node[ochre, font=\footnotesize\sffamily, align=left, anchor=west] at (0.35, 3.92)
    {$L$-光滑二次上界\\$q_t(y) = f(x_t) + \nabla f(x_t)^\top(y-x_t) + \frac{L}{2}\|y-x_t\|^2$};
  \draw[->, ochre!85, thin] (0.95, 3.48) -- (1.18, 2.95);

  % True convex function f(x)
  \draw[tealDark, line width=1.6pt, domain=-2.05:3.55, samples=80]
    plot (\x, {0.28*(\x-1.4)^2 + 0.65});
  \node[tealDark, font=\small\sffamily\bfseries] at (3.65, 2.25) {$f(x)$};

  % Lower convexity tangent line
  \draw[navy, thick, dash dot, domain=-2.05:0.82, samples=2]
    plot (\x, {2.0052 - 1.232*(\x+0.8)});
  \node[navy, font=\scriptsize\sffamily, align=left, fill=white, inner sep=1.5pt] (tanlbl) at (-1.55, 0.72)
    {凸性一阶切线下界\\$f(x_t)+\nabla f(x_t)^\top(y-x_t)$};
  \draw[->, navy!85, thin] (-0.82, 0.98) -- (-0.25, 1.32);

  % Key points: x_t, x_{t+1}, x^*
  \draw[dotted, slate, thick] (-0.8, 0) -- (-0.8, 2.0052);
  \fill[ink] (-0.8, 2.0052) circle (2.3pt);
  \node[below, font=\footnotesize] at (-0.8, -0.05) {$x_t$};

  \draw[dotted, ochre, thick] (0.186, 0) -- (0.186, 1.3981);
  \fill[ochre] (0.186, 1.3981) circle (2.3pt);
  \fill[teal] (0.186, 1.063) circle (2.3pt);
  \node[below, font=\footnotesize, color=ochre] at (0.186, -0.05) {$x_{t+1}$};

  % Drop bracket at x = 0.45 with clean non-overlapping callout
  \draw[dashed, slate!60] (-0.8, 2.0052) -- (0.50, 2.0052);
  \draw[dashed, slate!60] (0.186, 1.3981) -- (0.50, 1.3981);
  \draw[<->, teal, thick] (0.45, 2.0052) -- (0.45, 1.3981);
  \node[font=\scriptsize\sffamily, color=tealDark, fill=white, inner sep=1.8pt, anchor=west] (droplbl) at (1.58, 2.62)
    {保底下降 $\ge \frac{1}{2L}\|\nabla f(x_t)\|_2^2$};
  \draw[->, teal, semithick] (droplbl.west) -- (0.52, 1.72);

  % Global minimum x^*
  \draw[dotted, tealDark, thick] (1.4, 0) -- (1.4, 0.65);
  \fill[gold] (1.4, 0.65) circle (2.6pt);
  \node[below, font=\footnotesize, color=tealDark] at (1.4, -0.05) {$x^*$};
  \node[font=\footnotesize\sffamily, color=gold!80!ink, anchor=west] (optlbl) at (2.05, 0.40) {全局最优 $(x^*, f^*)$};
  \draw[->, gold!85!ink, thin] (optlbl.west) -- (1.50, 0.62);
\end{tikzpicture}
\caption{凸性与 $L$-光滑性构成的“夹心三明治”几何图景：凸性用一阶切线从下方托住 $f(x)$，$L$-光滑性用开口宽度由 $L$ 决定的二次抛物线 $q_t(y)$ 从上方罩住 $f(x)$。一步步长为 $1/L$ 的梯度下降，本质上就是精确跳到当前二次上界 $q_t(y)$ 的最低点 $x_{t+1}$，从而保证真实函数值至少下降 $\frac{1}{2L}\norm{\nabla f(x_t)}_2^2$。}\label{fig:v4-sandwich}
\end{figure}

图~\ref{fig:v4-sandwich} 揭示了梯度下降最优雅的优化-最小化（Majorization-Minimization, MM）解释：我们并不直接最小化复杂的 $f(y)$，而是在当前点 $x_t$ 构造一个处处高于 $f(y)$ 且在 $x_t$ 处相切的二次代理模型 $q_t(y) = f(x_t) + \nabla f(x_t)^\top(y-x_t) + \frac{L}{2}\norm{y-x_t}_2^2$，然后精确求解 $\argmin_y q_t(y) = x_t - \frac{1}{L}\nabla f(x_t)$！这一视角将在第 2 章推广为\textbf{近端梯度法（ISTA）}，并在第 6 章演化为 \textbf{EM 算法}。

\begin{failBox}{步长超过 $2/L$ 时为何必然面临发散风险？}
观察下降引理式~\eqref{eq:v4-step-drop} 中的系数 $\eta(1 - \frac{L\eta}{2})$：只要 $0 < \eta < \frac{2}{L}$，该系数严格为正，每步函数值单调不增；但一旦 $\eta > \frac{2}{L}$，上界右侧反而变成正项，二次代理模型不再保证下降。在例~\ref{ex:v4-ch1-zigzag} 中 $L=100$，临界步长正是 $2/L = 0.02$；实际工程中若无法精确知道 $L$，则必须配合\textbf{回溯线搜索（Armijo Backtracking）}或动态学习率衰减，切忌盲目放大固定步长。
\end{failBox}

现在，把“凸性切线下界”与“$L$-光滑二次上界”两块拼图扣合在一起，我们就能给出优化理论中最经典、最干净的\textbf{望远镜求和（Telescoping Sum）}收敛率证明。

\begin{thmBox}
\begin{theorem}[光滑凸函数梯度下降的 $O(1/T)$ 收敛率]\label{thm:v4-gd-convex}
设 $f: \R^d \to \R$ 为凸函数且 $L$-光滑，存在全局极小点 $x^*$，记 $f^* = f(x^*)$。从任意初值 $x_0 \in \R^d$ 出发，取固定步长 $\eta = \frac{1}{L}$ 执行梯度下降 $x_{t+1} = x_t - \frac{1}{L}\nabla f(x_t)$，则对任意步数 $T \ge 1$ 均有：
\begin{equation}\label{eq:v4-gd-o1t}
  f(x_T) - f^* \le \frac{L\norm{x_0 - x^*}_2^2}{2T}.
\end{equation}
\end{theorem}
\end{thmBox}

\begin{proof}
记 $g_t = \nabla f(x_t)$，迭代步为 $x_{t+1} = x_t - \frac{1}{L}g_t$。我们的目标是用相邻两步到最优解的距离平方差 $\norm{x_t - x^*}_2^2 - \norm{x_{t+1} - x^*}_2^2$ 去控制当前的函数次优度 $f(x_{t+1}) - f^*$。

\textbf{第一步：联合下降引理与凸性下界。} 由下降引理~\eqref{eq:v4-1overL-drop}，有
\[
  f(x_{t+1}) \le f(x_t) - \frac{1}{2L}\norm{g_t}_2^2.
\]
另一方面，由凸函数一阶条件~\eqref{eq:v4-first-order-convex}（取 $x=x_t,\, y=x^*$），有 $f(x^*) \ge f(x_t) + g_t^\top(x^* - x_t)$，移项得 $f(x_t) \le f^* + g_t^\top (x_t - x^*)$。将其代入上式：
\begin{equation}\label{eq:v4-tele-step}
  f(x_{t+1}) - f^* \le g_t^\top (x_t - x^*) - \frac{1}{2L}\norm{g_t}_2^2.
\end{equation}

\textbf{第二步：配方法化为距离平方差分。} 观察式~\eqref{eq:v4-tele-step} 右侧，凑出 $\frac{L}{2}$ 因子并直接配方：
\begin{align*}
  g_t^\top (x_t - x^*) - \frac{1}{2L}\norm{g_t}_2^2
  &= \frac{L}{2}\left[\frac{2}{L}g_t^\top (x_t - x^*) - \frac{1}{L^2}\norm{g_t}_2^2\right] \\
  &= \frac{L}{2}\left[\norm{x_t - x^*}_2^2 - \norm{x_t - \frac{1}{L}g_t - x^*}_2^2\right] \\
  &= \frac{L}{2}\left(\norm{x_t - x^*}_2^2 - \norm{x_{t+1} - x^*}_2^2\right).
\end{align*}
（展开第二行右侧的第二项即可逐项核对：$\norm{(x_t-x^*) - \frac{1}{L}g_t}_2^2 = \norm{x_t-x^*}_2^2 - \frac{2}{L}g_t^\top(x_t-x^*) + \frac{1}{L^2}\norm{g_t}_2^2$。）

\textbf{第三步：望远镜求和与单调性收尾。} 将上述不等式从 $t = 0$ 累加至 $T - 1$，中间项成对精确抵消：
\begin{align*}
  \sum_{t=0}^{T-1}\big(f(x_{t+1}) - f^*\big)
  &\le \frac{L}{2}\sum_{t=0}^{T-1}\left(\norm{x_t - x^*}_2^2 - \norm{x_{t+1} - x^*}_2^2\right) \\
  &= \frac{L}{2}\left(\norm{x_0 - x^*}_2^2 - \norm{x_T - x^*}_2^2\right)
  \le \frac{L}{2}\norm{x_0 - x^*}_2^2.
\end{align*}
再由下降引理~\eqref{eq:v4-1overL-drop} 知序列 $\{f(x_t)\}$ 单调不增，故对所有 $t \in \{0, \dots, T-1\}$ 均有 $f(x_T) \le f(x_{t+1})$。因此左侧满足
\[
  T\big(f(x_T) - f^*\big) \le \sum_{t=0}^{T-1}\big(f(x_{t+1}) - f^*\big) \le \frac{L}{2}\norm{x_0 - x^*}_2^2.
\]
两边同除以 $T$，即证得式~\eqref{eq:v4-gd-o1t}！
\end{proof}

这一证明堪称“天书”级别的典范：它告诉我们，对任意 $L$-光滑凸函数，要使误差 $f(x_T) - f^* \le \varepsilon$，梯度下降最多需要 $T = O(L/\varepsilon)$ 步。然而 $O(1/\varepsilon)$ 的代价依然不菲——要把精度从 $10^{-2}$ 提高到 $10^{-6}$，步数要暴涨 $10{,}000$ 倍！有没有办法把 $O(1/\varepsilon)$ 压缩到 $O(\log(1/\varepsilon))$？答案是肯定的：我们需要让山谷底部具有非零的最小曲率。

% ============================================================
\section{$\mu$-强凸性、PL 不等式与指数级线性收敛}
% ============================================================

一般凸函数允许谷底非常“平缓”（如 $f(x) = x^4$ 在 $x=0$ 处二阶导数为 $0$），导致靠近最优解时梯度迅速塌缩、动力不足。如果我们要求函数不仅在切平面上方，而且至少像一个开口向上、曲率为 $\mu > 0$ 的二次碗一样向上弯曲，情形就会发生质变。

\begin{defBox}
\begin{definition}[$\mu$-强凸函数]\label{def:v4-strong-convex}
设 $\mu > 0$。可微函数 $f: \R^d \to \R$ 称为\textbf{$\mu$-强凸（$\mu$-Strongly Convex）}的，如果对任意 $x, y \in \R^d$ 均有
\begin{equation}\label{eq:v4-sc-def}
  f(y) \ge f(x) + \nabla f(x)^\top (y - x) + \frac{\mu}{2}\norm{y - x}_2^2.
\end{equation}
等价地，$g(x) = f(x) - \frac{\mu}{2}\norm{x}_2^2$ 为凸函数；若 $f$ 二阶连续可微，则等价于 $\nabla^2 f(x) \succeq \mu I$ 处处成立。
\end{definition}
\end{defBox}

注意这个等价刻画的直接工程含义：\textbf{对任意凸损失函数 $f_0(w)$（例如逻辑回归交叉熵或最小二乘）加上 $L_2$ 正则项 $\frac{\lambda}{2}\norm{w}_2^2$（$\lambda > 0$），新目标 $F(w) = f_0(w) + \frac{\lambda}{2}\norm{w}_2^2$ 立刻变成 $\lambda$-强凸函数！} 正则化不仅像第 3 章将证明的那样能够压制方差、防止过拟合，还在优化层面直接赋予目标函数严格的最小曲率 $\mu \ge \lambda$。

在推导强凸收敛率之前，让我们先提炼一个比强凸更弱、却足以驱动指数级线性收敛的神奇不等式——\textbf{Polyak-\L{}ojasiewicz（PL）不等式}。

\begin{thmBox}
\begin{lemma}[$\mu$-强凸蕴含 PL 不等式]\label{lem:v4-sc-pl}
若 $f: \R^d \to \R$ 是 $\mu$-强凸函数且全局最小值为 $f^*$，则对任意 $x \in \R^d$ 均满足 \textbf{PL 不等式}：
\begin{equation}\label{eq:v4-pl}
  \frac{1}{2}\norm{\nabla f(x)}_2^2 \ge \mu\big(f(x) - f^*\big).
\end{equation}
\end{lemma}
\end{thmBox}

\begin{proof}
固定 $x \in \R^d$，在 $\mu$-强凸定义~\eqref{eq:v4-sc-def} 的左右两侧同时对 $y \in \R^d$ 取极小值：
\begin{itemize}
  \item 左侧 $\min_{y\in\R^d} f(y) = f^*$；
  \item 右侧是关于 $y$ 的开口向上的二次函数 $r(y) = f(x) + \nabla f(x)^\top(y-x) + \frac{\mu}{2}\norm{y-x}_2^2$。令 $\nabla_y r(y) = \nabla f(x) + \mu(y - x) = 0$，解得唯一极小点 $y^* = x - \frac{1}{\mu}\nabla f(x)$，代入得
  \[
    \min_{y\in\R^d} r(y) = f(x) - \frac{1}{\mu}\norm{\nabla f(x)}_2^2 + \frac{\mu}{2}\cdot\frac{1}{\mu^2}\norm{\nabla f(x)}_2^2 = f(x) - \frac{1}{2\mu}\norm{\nabla f(x)}_2^2.
  \]
\end{itemize}
因此 $f^* \ge f(x) - \frac{1}{2\mu}\norm{\nabla f(x)}_2^2$，整理即得式~\eqref{eq:v4-pl}。
\end{proof}

PL 不等式的物理意义一目了然：只要你离谷底还差 $f(x) - f^*$ 的高度，你脚下的坡度平方 $\norm{\nabla f(x)}_2^2$ 就绝不会太小（至少有 $2\mu(f(x)-f^*)$）。再把 PL 不等式代入下降引理~\eqref{eq:v4-1overL-drop}，只需三行代数，指数级收敛便跃然纸上！

\begin{thmBox}
\begin{theorem}[光滑 PL / $\mu$-强凸函数的线性收敛定理]\label{thm:v4-linear-conv}
设 $f: \R^d \to \R$ 是 $L$-光滑的且满足参数为 $\mu > 0$ 的 PL 不等式~\eqref{eq:v4-pl}（特别地，一切 $L$-光滑且 $\mu$-强凸的函数均满足条件，记条件数 $\kappa = L/\mu \ge 1$）。取步长 $\eta = \frac{1}{L}$ 执行梯度下降 $x_{t+1} = x_t - \frac{1}{L}\nabla f(x_t)$，则对任意 $T \ge 0$ 有：
\begin{equation}\label{eq:v4-linear-rate}
  f(x_T) - f^* \le \left(1 - \frac{\mu}{L}\right)^T \big(f(x_0) - f^*\big) = \left(1 - \frac{1}{\kappa}\right)^T \big(f(x_0) - f^*\big) \le e^{-T/\kappa}\big(f(x_0) - f^*\big).
\end{equation}
进一步地，若 $f$ 为 $\mu$-强凸函数（最优解 $x^*$ 唯一），则参数距离也以线性速率收敛：
\begin{equation}\label{eq:v4-param-dist}
  \norm{x_T - x^*}_2^2 \le \frac{2}{\mu}\left(1 - \frac{1}{\kappa}\right)^T \big(f(x_0) - f^*\big).
\end{equation}
\end{theorem}
\end{thmBox}

\begin{proof}
由下降引理~\eqref{eq:v4-1overL-drop} 与 PL 不等式~\eqref{eq:v4-pl} 直接串联：
\[
  f(x_{t+1}) - f(x_t) \le -\frac{1}{2L}\norm{\nabla f(x_t)}_2^2 \le -\frac{\mu}{L}\big(f(x_t) - f^*\big).
\]
两边同时加上 $f(x_t) - f^*$ 并整理：
\[
  f(x_{t+1}) - f^* \le \left(1 - \frac{\mu}{L}\right)\big(f(x_t) - f^*\big).
\]
从 $t=0$ 递推 $T$ 步，并利用初等不等式 $1 - u \le e^{-u}$，即证得式~\eqref{eq:v4-linear-rate}。

若 $f$ 还是 $\mu$-强凸的，在式~\eqref{eq:v4-sc-def} 中令 $x = x^*$（注意 $\nabla f(x^*)=0$）、$y = x_T$，得\textbf{二次增长性质}：
\[
  f(x_T) - f^* \ge \frac{\mu}{2}\norm{x_T - x^*}_2^2 \implies \norm{x_T - x^*}_2^2 \le \frac{2}{\mu}\big(f(x_T) - f^*\big),
\]
结合式~\eqref{eq:v4-linear-rate} 即得式~\eqref{eq:v4-param-dist}。
\end{proof}

\begin{insightBox}{为何称指数衰减为“线性收敛”？以及 PL 条件在非凸深度学习中的回响}
\begin{enumerate}
  \item \textbf{术语澄清}：在最优化文献中，若误差满足 $e_{t+1} \le \rho\, e_t\;(\rho < 1)$，在对数坐标 $\log e_t$ 下误差随步数 $t$ 呈\textbf{直线下降}（每固定步数多获得一位有效十进制精度），因此称为\textbf{线性收敛（Linear Convergence）}，其达到精度 $\varepsilon$ 所需的步数仅为 $T = O\big(\kappa \ln\frac{1}{\varepsilon}\big)$。
  \item \textbf{超越凸性的 PL 条件}：请注意，定理~\ref{thm:v4-linear-conv} 的证明\emph{根本没有用到凸性}，只用到了 $L$-光滑与 PL 不等式！存在许多非凸函数依然满足 PL 条件——例如一维波动函数 $f(x) = x^2 + 3\sin^2 x$（它并非凸函数，因为二阶导数 $2 + 6\cos(2x)$ 可为负，但读者可在习题~\ref{ex:4.1.6} 中验证它全局满足 PL 不等式且不存在非全局驻点！）。在现代深度学习理论（第 5 卷）中，宽过参数化神经网络在初始化邻域内常局部满足 PL 条件，这正是 SGD 竟能将非凸深度网络训练到零损失的重要数学解释之一。
\end{enumerate}
\end{insightBox}

% ============================================================
\section{突破条件数瓶颈：动量法与 Nesterov 加速梯度}
% ============================================================

定理~\ref{thm:v4-linear-conv} 虽然把对精度的依赖从 $1/\varepsilon$ 降到了 $\ln(1/\varepsilon)$，却留下了一座沉重的大山：步数仍正比于条件数 $\kappa = L/\mu$。在稍后的病态逻辑回归实验中，$\kappa = 1030.5$，意味着纯梯度下降每走一步，最坏情况下的误差倍率高达 $1 - 1/1030.5 \approx 0.99903$，走 150 步误差连一个数量级都下不去！

能否在不计算 $O(D^2)$ 海森矩阵的前提下，仅凭一阶梯度就大幅打破 $O(\kappa)$ 的魔咒？苏联数学家 Boris Polyak（1964）与 Yurii Nesterov（1983）给出了两套载入史册的方案。

\subsection{Polyak 重球动量法（Heavy-Ball Momentum）}

回到例~\ref{ex:v4-ch1-zigzag} 的狭长山谷。纯梯度下降像一个没有质量的纸片，每一步都完全忘掉刚才的方向，只能听从当前陡壁梯度的摆布来回弹跳。如果赋予迭代点一个“物理质量”，让它像一颗有惯性的重球沿山谷滚落，就会得到 \textbf{Heavy-Ball 动量法}：
\begin{equation}\label{eq:v4-momentum}
  v_{t+1} = \beta v_t + \nabla f(x_t), \qquad x_{t+1} = x_t - \eta v_{t+1},
\end{equation}
等价地消去速度向量 $v_t$ 后可写作二阶差分格式 $x_{t+1} = x_t - \eta \nabla f(x_t) + \beta(x_t - x_{t-1})$，其中 $\beta \in [0, 1)$ 为动量（阻尼）系数。其几何奇迹在于\textbf{方向选择性滤波}：
\begin{itemize}
  \item 在陡峭震荡方向上，相邻两步的梯度正负交替，历史动量项 $\beta v_t$ 与当前梯度相互抵消，剧烈震荡被平滑阻尼；
  \item 在平坦谷底方向上，每一步梯度指向一致，动量项不断同相累加，有效步长被放大至 $\frac{\eta}{1-\beta}$（例如 $\beta = 0.85$ 时有效步长放大约 $6.7$ 倍！）。
\end{itemize}
Polyak 证明，对强凸二次函数，取最优动量 $\beta^* = \left(\frac{\sqrt{\kappa}-1}{\sqrt{\kappa}+1}\right)^2$ 与步长 $\eta^* = \frac{4}{(\sqrt{L}+\sqrt{\mu})^2}$，重球法即可达到局部加速速率 $\frac{\sqrt{\kappa}-1}{\sqrt{\kappa}+1}$。但对于一般非二次强凸函数，Heavy-Ball 可能因为惯性过大“冲过头”而失去全局加速保证。

\subsection{Nesterov 加速梯度法（NAG）：带着“前瞻刹车”的动量}

1983 年，Nesterov 对动量法做了一个看似微小、实则化腐朽为神奇的改动：不再在当前位置 $x_t$ 求梯度，而是\textbf{先顺着惯性外推一步到达“前瞻点” $y_t$，再在 $y_t$ 处评估梯度并纠正方向}！对 $L$-光滑且 $\mu$-强凸函数，Nesterov 加速梯度（Nesterov Accelerated Gradient, NAG）的固定动量形式为：
\begin{equation}\label{eq:v4-nag}
  \left\{
  \begin{aligned}
    x_{t+1} &= y_t - \frac{1}{L}\nabla f(y_t), \\
    y_{t+1} &= x_{t+1} + \beta_{\text{nag}}\,(x_{t+1} - x_t), \qquad \beta_{\text{nag}} = \frac{\sqrt{\kappa} - 1}{\sqrt{\kappa} + 1},
  \end{aligned}
  \right.
\end{equation}
初值 $y_0 = x_0$。而对一般光滑凸函数（$\mu = 0$），只需采用时变动量 $\beta_t = \frac{\theta_t - 1}{\theta_{t+1}}$（其中 $\theta_0 = 1,\; \theta_{t+1} = \frac{1 + \sqrt{1+4\theta_t^2}}{2}$，近似于 $\beta_t \approx \frac{t-1}{t+2}$）。

为什么“先外推、再求梯度”就能驯服惯性？想象你正推着球冲向谷底：Heavy-Ball 是在原地看一眼坡度、然后闭着眼睛加上惯性往前冲，冲过谷底爬上对岸才发现坡度反向；而 NAG 是先预判自己靠惯性会滚到前方 $y_t$ 处，提前去感知 $y_t$ 的坡度——如果 $y_t$ 已经越过谷底开始上坡，$\nabla f(y_t)$ 就会立刻施加反向制动力，像一套自带“预判刹车”的控制系统！

\begin{thmBox}
\begin{theorem}[Nesterov 加速梯度的最优一阶收敛率，不证自引]\label{thm:v4-nag-rate}
\begin{enumerate}
  \item 若 $f$ 为 $L$-光滑凸函数，时变动量 NAG 满足
  \[
    f(x_T) - f^* \le \frac{2L\norm{x_0 - x^*}_2^2}{(T+1)^2} = O\!\left(\frac{1}{T^2}\right).
  \]
  \item 若 $f$ 为 $L$-光滑且 $\mu$-强凸函数，取式~\eqref{eq:v4-nag} 的固定动量 $\beta_{\text{nag}} = \frac{\sqrt{\kappa}-1}{\sqrt{\kappa}+1}$，则存在常数 $C$ 使得
  \[
    f(x_T) - f^* \le C \left(1 - \frac{1}{\sqrt{\kappa}}\right)^T \big(f(x_0) - f^*\big).
  \]
\end{enumerate}
\end{theorem}
\end{thmBox}

更重要的是，Nesterov 还通过构造高维三对角二次“最坏情况函数”证明了\textbf{一阶算法信息论下界}：任何只查询梯度的黑箱一阶算法，在凸光滑类上不可能快于 $\Omega(1/T^2)$，在强凸光滑类上不可能快于 $\Omega\big((1-1/\sqrt{\kappa})^T\big)$。也就是说，\textbf{NAG 已经撞上了光滑一阶凸优化的物理速度上限！} 表~\ref{tab:v4-ch1-rates} 总结了本章的核心收敛率阶梯。

\begin{table}[htbp]
\centering
\caption{光滑一阶优化算法收敛率与迭代复杂度对比（达到 $f(x_T)-f^* \le \varepsilon$，条件数 $\kappa = L/\mu$）}\label{tab:v4-ch1-rates}
\small
\begin{tabularx}{\linewidth}{@{}lcccX@{}}
\toprule
\textbf{算法与函数类} & \textbf{步长 / 参数} & \textbf{误差衰减率} & \textbf{迭代复杂度} & \textbf{核心特征} \\
\midrule
GD（光滑凸） & $\eta = 1/L$ & $O(L/T)$ & $O(L/\varepsilon)$ & 单调下降，但高精度极慢 \\
NAG（光滑凸） & $\eta = 1/L,\; \beta_t \approx \frac{t-1}{t+2}$ & $O(L/T^2)$ & $O(\sqrt{L/\varepsilon})$ & 达一阶信息下界，允许轻微震荡 \\
GD（光滑强凸 / PL） & $\eta = 1/L$ & $(1-1/\kappa)^T$ & $O(\kappa \ln\frac{1}{\varepsilon})$ & 指数收敛，但受制于大 $\kappa$ \\
NAG（光滑强凸） & $\eta = 1/L,\; \beta = \frac{\sqrt{\kappa}-1}{\sqrt{\kappa}+1}$ & $(1-1/\sqrt{\kappa})^T$ & $O(\sqrt{\kappa}\ln\frac{1}{\varepsilon})$ & $\kappa=1000$ 时提速约 $32\times$ \\
\bottomrule
\end{tabularx}
\end{table}

% ============================================================
\section{实测验证：病态 $L_2$ 逻辑回归上的三法对决}
% ============================================================

理论公式有多美，必须在真实数值实验里过招才算数。我们在配套脚本 \texttt{code/vol4\_experiments.py} 中构造了一个刻意制造量纲跨度的病态 $L_2$ 正则化逻辑回归问题：
\begin{itemize}
  \item 样本数 $N = 300$，特征维数 $D = 15$，各列特征的标准差按几何级数从 $1.0$ 跨越至 $15.0$，正则化系数 $\lambda = 0.05$。
  \item 利用例~\ref{ex:v4-ch1-logreg} 的精确公式，算得 $L = \frac{\sigma_{\max}(X)^2}{4N} + \lambda = 51.52$，$\mu = \lambda = 0.05$，因此条件数高达 $\kappa = L/\mu = 1030.5$，对应 Nesterov 理论最优动量 $\beta_{\text{nag}} = \frac{\sqrt{\kappa}-1}{\sqrt{\kappa}+1} = 0.940$。
  \item 为获得精确到机器精度的基准值 $f^*$，先运行 15 步牛顿法求解（最终梯度范数 $\norm{\nabla F(w^*)}_2 < 10^{-12}$），再从统一原点 $w_0 = 0$ 出发，使用完全相同的理论步长 $\eta = 1/L$ 对比 GD、Heavy-Ball 动量法（$\beta = 0.85$）与 NAG（$\beta_{\text{nag}} = 0.940$）在 150 步内的次优度 $F(w_t) - f^*$。
\end{itemize}

\begin{lstlisting}[caption={病态 $L_2$ 逻辑回归上 GD、Momentum 与 NAG 的核心实现（节选自 \texttt{vol4\_experiments.py}）},label={lst:v4-ch1-opt}]
smax = np.linalg.svd(X, compute_uv=False)[0]
L = 0.25 * (smax ** 2) / N + lam      # Exact smoothness bound: L = 51.52
mu, kappa, eta = lam, L / lam, 1.0 / L
beta_nag = (np.sqrt(kappa) - 1.0) / (np.sqrt(kappa) + 1.0)  # 0.940

w_gd, w_mom, v_mom = np.zeros(D), np.zeros(D), np.zeros(D)
w_nag, y_nag = np.zeros(D), np.zeros(D)
for t in range(151):
    _, g_gd = loss_grad(w_gd)
    _, g_mom = loss_grad(w_mom)
    w_gd -= eta * g_gd                       # Plain GD
    v_mom = 0.85 * v_mom + g_mom             # Heavy-Ball Momentum
    w_mom -= eta * v_mom
    _, g_y = loss_grad(y_nag)                # Nesterov look-ahead gradient
    w_next = y_nag - eta * g_y
    y_nag = w_next + beta_nag * (w_next - w_nag)
    w_nag = w_next
\end{lstlisting}

图~\ref{fig:v4-ch1-accel} 与表~\ref{tab:v4-ch1-exp} 直接读取自实测数据文件 \texttt{data/opt\_accel.csv}，震撼地呈现了条件数 $\kappa$ 与 $\sqrt{\kappa}$ 之间的鸿沟。

\begin{figure}[htbp]
\centering
\begin{tikzpicture}
\begin{semilogyaxis}[
  width=0.90\linewidth,
  height=6.7cm,
  xlabel={迭代步数 $t$},
  ylabel={次优度 $F(w_t) - f^*$（对数坐标）},
  xmin=-3, xmax=153,
  ymin=2e-6, ymax=6e-1,
  grid=both,
  major grid style={tealLine!45, line width=0.4pt},
  minor grid style={tealLine!20, line width=0.25pt},
  axis line style={slate!80},
  tick label style={font=\footnotesize\sffamily, color=ink},
  label style={font=\small\sffamily, color=ink},
  legend style={
    at={(0.035,0.055)}, anchor=south west,
    fill=white, draw=tealLine, rounded corners=1.2mm,
    font=\footnotesize\sffamily, cells={anchor=west}
  },
  every axis plot/.append style={line width=1.45pt}
]
  \addplot[color=ochre, mark=square*, mark size=1.4pt, mark repeat=3]
    table[x=t, y=gd, col sep=comma]{data/opt_accel.csv};
  \addlegendentry{GD（$\eta = 1/L$）}

  \addplot[color=navy, mark=triangle*, mark size=1.7pt, mark repeat=3]
    table[x=t, y=mom, col sep=comma]{data/opt_accel.csv};
  \addlegendentry{Heavy-Ball 动量（$\beta = 0.85$）}

  \addplot[color=teal, mark=*, mark size=1.4pt, mark repeat=3]
    table[x=t, y=nag, col sep=comma]{data/opt_accel.csv};
  \addlegendentry{Nesterov NAG（$\beta_{\text{nag}} = 0.940$）}
\end{semilogyaxis}
\end{tikzpicture}
\caption{病态 $L_2$ 逻辑回归（$N=300, D=15, \lambda=0.05, L=51.52, \kappa=1030.5$）实测收敛曲线（读取自 \texttt{data/opt\_accel.csv}）。对数纵轴上的斜率清晰印证了线性收敛速率：150 步后，NAG 与动量法将误差压低至 $10^{-6}$--$10^{-5}$ 量级，领先纯 GD 整整四个数量级！}\label{fig:v4-ch1-accel}
\end{figure}

\begin{table}[htbp]
\centering
\caption{实测次优度 $F(w_t) - f^*$ 关键节点对比（\texttt{data/opt\_accel.csv}）}\label{tab:v4-ch1-exp}
\small
\begin{tabular}{@{}rccc@{}}
\toprule
\textbf{迭代步 $t$} & \textbf{GD（$\eta=1/L$）} & \textbf{Momentum（$\beta=0.85$）} & \textbf{Nesterov NAG（$\beta=0.940$）} \\
\midrule
$0$   & $3.19 \times 10^{-1}$ & $3.19 \times 10^{-1}$ & $3.19 \times 10^{-1}$ \\
$50$  & $5.90 \times 10^{-2}$ & $3.93 \times 10^{-3}$ & $2.43 \times 10^{-3}$ \\
$100$ & $3.51 \times 10^{-2}$ & $2.26 \times 10^{-4}$ & $2.70 \times 10^{-5}$ \\
$150$ & $2.22 \times 10^{-2}$ & $1.56 \times 10^{-5}$ & $3.36 \times 10^{-6}$ \\
\bottomrule
\end{tabular}
\end{table}

细看图~\ref{fig:v4-ch1-accel} 中的绿色 NAG 曲线，你会发现一个极具教学价值的细节：在 $t=35 \to 40$（误差从 $2.23\times 10^{-3}$ 轻微回弹至 $2.70\times 10^{-3}$）以及 $t=105 \to 125$（从 $1.09\times 10^{-5}$ 轻微回弹至 $2.48\times 10^{-5}$ 后又迅速跌至 $3.36\times 10^{-6}$），NAG 并不是每一步都单调下降的！这正是\textbf{加速方法不是下降方法（Not a Descent Method）}的典型标志——为了换取全局 $\sqrt{\kappa}$ 级的惊人提速，NAG 允许像欠阻尼弹簧一样在谷底附近轻微过冲并迅速拉回。而在工程实践中，如果我们在前瞻步检测到函数值回弹，配合“自适应重启（Adaptive Restart，清零动量）”甚至能把这些小波纹彻底抹平。

\begin{extBox}{二阶牛顿法：局部二次收敛的“核动力”与 $O(D^3)$ 代价}
在刚才的实验中，我们只用了 \textbf{15 步牛顿法}就把基准解 $w^*$ 算到了梯度范数 $< 10^{-12}$ 的机器精度，为什么它比一阶法快如此之多？

牛顿法的核心思想是将图~\ref{fig:v4-sandwich} 中的各向同性二次上界 $\frac{L}{2}\norm{y-x_t}_2^2$，替换为目标函数在当前点 $x_t$ 的\textbf{精确二阶泰勒展开}：
\begin{align*}
  x_{t+1} &= \argmin_y \left\{f(x_t) + \nabla f(x_t)^\top(y-x_t) + \frac{1}{2}(y-x_t)^\top \nabla^2 f(x_t)(y-x_t)\right\} \\
  &= x_t - \big[\nabla^2 f(x_t)\big]^{-1}\nabla f(x_t).
\end{align*}
由于海森矩阵逆 $[\nabla^2 f(x_t)]^{-1}$ 自动把狭长椭球“拉圆”成了标准球体，牛顿法具备\textbf{仿射不变性}——无论特征尺度相差 $15$ 倍还是 $10^6$ 倍，条件数都被海森逆瞬间抵消！在最优解 $x^*$ 邻域内，只要 $\nabla^2 f$ 满足 Lipschitz 连续，牛顿法具备\textbf{局部二次收敛（Quadratic Convergence）}：
\[
  \norm{x_{t+1} - x^*}_2 \le C \norm{x_t - x^*}_2^2,
\]
这意味着一旦误差进入 $10^{-1}$，接下来的误差序列将按 $10^{-2} \to 10^{-4} \to 10^{-8} \to 10^{-16}$ 坍缩，每步有效数字\textbf{翻倍}！

那为何我们在大规模机器学习与深度学习里不直接全用牛顿法？瓶颈在于每一步都需要存储 $O(D^2)$ 的海森矩阵并花费 $O(D^3)$ 的算力求解线性方程组 $H \Delta x = -g$：当 $D = 15$（本节实验）时 $15^3 \approx 3000$ 瞬息即成；当 $D = 10^6$ 时 $D^3 = 10^{18}$ 完全不可承受。因此在大维场景下，人们转向只用 $O(D)$ 存储的\textbf{拟牛顿法（如 L-BFGS）}或本章的加速一阶法与自适应梯度法（Adam 等，第 5 卷）。
\end{extBox}

% ============================================================
\section*{本章分层习题}
\addcontentsline{toc}{section}{本章分层习题}
% ============================================================

\exTier{A 层}{概念核对与基础验算}
\begin{enumerate}[label=\textbf{4.1.\arabic*}, leftmargin=3.3em, itemsep=5pt]
  \item\label{ex:4.1.1} \textbf{凸性判别与反例}：判断下列函数在各自定义域上是否为凸函数、严格凸函数或非凸函数，并给出简明证明或反例：
  (a) $f_1(x) = e^{a^\top x + b}\;(x\in\R^d)$；
  (b) $f_2(x_1, x_2) = x_1 x_2\;(x \in \R_{++}^2)$；
  (c) $f_3(x) = -\sum_{i=1}^d \ln x_i\;(x \in \R_{++}^d)$；
  (d) $f_4(w) = \max_{1\le k \le K} (a_k^\top w + b_k)\;(w\in\R^d)$。
  \exHint{结合二阶海森矩阵判据与命题~\ref{prop:v4-preserve} 的保凸运算；对 $f_2$ 检查特征值符号。}

  \item\label{ex:4.1.2} \textbf{Log-Sum-Exp 函数的凸性与光滑常数}：证明 $\mathrm{LSE}(z) = \ln\left(\sum_{k=1}^K e^{z_k}\right)$ 是 $\R^K$ 上的凸函数，并证明其梯度是 $1$-Lipschitz 连续的（即 $\nabla^2 \mathrm{LSE}(z) \preceq I_K$）。
  \exHint{令 $p_k = e^{z_k}/\sum_j e^{z_j}$，导出 $\nabla^2 \mathrm{LSE}(z) = \diag(p) - pp^\top$，并用 Cauchy-Schwarz 验证 $0 \le v^\top(\diag(p)-pp^\top)v \le \norm{v}_2^2$。}
\end{enumerate}

\exTier{B 层}{定理推演与条件数分析}
\begin{enumerate}[label=\textbf{4.1.\arabic*}, leftmargin=3.3em, itemsep=5pt, start=3]
  \item\label{ex:4.1.3} \textbf{最小二乘与岭回归的精确条件数}：考虑 $F_\lambda(w) = \frac{1}{2N}\norm{Xw - y}_2^2 + \frac{\lambda}{2}\norm{w}_2^2$，设 $X \in \R^{N\times D}$ 的奇异值为 $\sigma_1 \ge \dots \ge \sigma_D \ge 0$。求出 $F_\lambda$ 的紧致光滑常数 $L$ 与强凸常数 $\mu$，并说明当 $X$ 秩亏（$\sigma_D = 0$）时为何加入 $\lambda > 0$ 会彻底改变梯度下降的收敛阶。
  \exHint{直接写出海森矩阵 $\nabla^2 F_\lambda(w) = \frac{1}{N}X^\top X + \lambda I_D$ 的最大与最小特征值。}

  \item\label{ex:4.1.4} \textbf{一般步长 $\eta \in (0, 2/L)$ 的线性收敛率}：修改定理~\ref{thm:v4-linear-conv} 的证明，在 $L$-光滑且满足 PL 不等式的假设下，求出任意固定步长 $\eta \in (0, 2/L)$ 对应的单步收缩系数 $\rho(\eta)$，并验证 $\eta = 1/L$ 恰好使 $\rho(\eta)$ 最小。
  \exHint{直接从下降引理~\eqref{eq:v4-step-drop} 出发，将 $\norm{\nabla f(x_t)}_2^2 \ge 2\mu(f(x_t)-f^*)$ 代入即可。}
\end{enumerate}

\exTier{C 层}{深度证明与非凸边界探索}
\begin{enumerate}[label=\textbf{4.1.\arabic*}, leftmargin=3.3em, itemsep=5pt, start=5]
  \item\label{ex:4.1.5} \textbf{共强制性（Co-coercivity）与强凸参数距离的直接收缩}：对 $L$-光滑凸函数，利用辅助函数 $g_x(z) = f(z) - \nabla f(x)^\top z$ 证明：
  \[
    \big(\nabla f(y) - \nabla f(x)\big)^\top (y - x) \ge \frac{1}{L}\norm{\nabla f(y) - \nabla f(x)}_2^2.
  \]
  再利用该不等式证明：若 $f$ 同时是 $\mu$-强凸的且 $\eta \le 1/L$，则有更紧的参数一步收缩
  \[
    \norm{x_{t+1} - x^*}_2^2 \le (1 - \mu\eta)\norm{x_t - x^*}_2^2.
  \]
  \exHint{注意 $g_x$ 仍 $L$-光滑且在 $z=x$ 处取全局最小，应用下降引理得 $g_x(x) \le g_x(y) - \frac{1}{2L}\norm{\nabla g_x(y)}_2^2$，交换 $x, y$ 后相加。}

  \item\label{ex:4.1.6} \textbf{满足 PL 条件的非凸函数}：考察一维函数 $f(x) = x^2 + 3\sin^2 x$。
  (a) 计算 $f''(x)$ 并说明 $f$ 在 $\R$ 上不是凸函数；
  (b) 证明 $x^* = 0$ 是唯一全局极小点且 $f^* = 0$；
  (c) 证明存在常数 $\mu > 0$（例如 $\mu = \frac{1}{32}$）使 $f$ 在整个 $\R$ 上满足 PL 不等式~\eqref{eq:v4-pl}，从而梯度下降对其依然具有全局线性收敛性！
  \exHint{注意到 $f'(x) = 2x + 3\sin(2x)$；可分 $|x| \le \frac{\pi}{2}$（利用 $\sin(2x)$ 与 $x$ 同号或界定比值）与 $|x| > \frac{\pi}{2}$ 两段放缩证明 $(f'(x))^2 / f(x)$ 有正下界。}
\end{enumerate}

\exTier{D 层}{数值实验与工程诊断}
\begin{enumerate}[label=\textbf{4.1.\arabic*}, leftmargin=3.3em, itemsep=5pt, start=7]
  \item\label{ex:4.1.7} \textbf{特征标准化对条件数 $\kappa$ 的震撼效应}：修改脚本 \texttt{vol4\_experiments.py} 中的实验 1 函数，在运行优化器之前仅对设计矩阵 $X$ 做零均值单位方差标准化（列归一化）。重新计算 $L, \mu, \kappa$ 并运行 150 步 GD 与 NAG，比较标准化前后 $t=150$ 时的次优度变化。这一结果对机器学习预处理有何启示？
  \exHint{标准化后 $X^\top X / N$ 的最大特征值将骤降两个数量级，纯 GD 也能迅速逼近机器精度。}

  \item\label{ex:4.1.8} \textbf{Nesterov 自适应重启（Adaptive Restart）实验}：在第 1.6 节实验中我们观察到 NAG 曲线在 $t \approx 40$ 与 $t \approx 115$ 出现轻微回弹。在代码清单~\ref{lst:v4-ch1-opt} 中加入梯度重启判据：每当 $\nabla F(y_t)^\top (w_{t+1} - w_t) > 0$（说明动量方向与当前负梯度形成钝角、已经冲过谷底）时，令 $y_{t+1} = w_{t+1}$ 强制清零惯性。绘制重启版 NAG 与原版 NAG 的对比曲线，验证 $t=150$ 时的精度提升。
  \exHint{参见 O'Donoghue \& Candès (2015) 经典启发式重启策略，仅需在循环末尾判断 \texttt{np.dot(g\_y, w\_next - w\_nag) > 0} 并重置 \texttt{y\_nag = w\_next.copy()}。}
\end{enumerate}

\chapter[拉格朗日对偶、KKT 与近端梯度]{拉格朗日对偶、KKT 条件与近端梯度}\label{ch:v4-duality}

\chapEpigraph{我所提出的方法既不需要几何作图，也不需要机械推理，只需要按统一而有规律的代数演算即可。}{Joseph-Louis Lagrange（《分析力学》，1788）}

\begin{chapIntro}{本章导读：当无约束的光滑世界撞上“围墙”与“尖角”}
在第 1 章中，我们所有的定理都隐含着两个理想假设：变量 $x$ 可以在整个 $\R^d$ 自由奔跑（无约束），且目标函数处处光滑可微（存在 $L$-光滑梯度）。然而走进经典机器学习的真实战场，这两条假设立刻同时被打破：
\begin{itemize}
  \item \textbf{第一堵墙：不等式约束}。第 4 章的支持向量机（SVM）要求每一个样本满足间隔不等式 $y_i(w^\top x_i + b) \ge 1$；概率建模要求权重落入单纯形 $p_i \ge 0,\; \sum_i p_i = 1$。在可行域边界上，$\nabla f(x^*) = 0$ 通常不再成立！
  \item \textbf{第二道坎：不可微尖角}。为了从高维冗余特征中自动筛选出真正重要的少数变量，Lasso 回归加入了 $\ell_1$ 惩罚项 $\lambda\norm{w}_1 = \lambda\sum_j |w_j|$。它在任何坐标 $w_j = 0$ 处都存在尖锐的不可微折角，梯度根本不存在！
\end{itemize}
本章将锻造两把专治“围墙”与“尖角”的利器：前半章建立\textbf{拉格朗日对偶与 KKT 互补松弛条件}（直接打通第 4 章 SVM 对偶与支持向量理论）；后半章建立\textbf{次梯度、近端算子与软阈值迭代（ISTA）}，并在 20 维稀疏回归数据上亲手跑出精确零解的 Lasso 正则化路径。
\end{chapIntro}

% ============================================================
\section{当最优解顶在围墙上：不等式约束的几何直觉}
% ============================================================

我们先通过一个二维极简例子，看清为什么无约束条件 $\nabla f(x^*) = 0$ 在不等式约束面前会失效。

\begin{example}[半空间投影：最优解处的“推力平衡”]\label{ex:v4-ch2-halfspace}
考虑在二维闭半空间 $C = \{x \in \R^2 : x_1 + x_2 \le 2\}$ 上寻找离目标点 $a = (3, 2)^\top$ 最近的点，即求解二次规划：
\[
  \min_{x \in \R^2}\; f_0(x) = \frac{1}{2}\norm{x - a}_2^2 = \frac{1}{2}\big((x_1-3)^2 + (x_2-2)^2\big)
  \quad\text{s.t.}\quad f_1(x) = x_1 + x_2 - 2 \le 0.
\]
无约束极小点 $a = (3, 2)^\top$ 满足 $3+2=5 > 2$，位于围墙之外。几何上，离 $(3, 2)$ 最近的半空间点必落在边界直线 $x_1 + x_2 = 2$ 上，且恰好是垂直投影点 $x^* = \left(\frac{3}{2}, \frac{1}{2}\right)^\top$。

检查 $x^*$ 处的目标梯度：$\nabla f_0(x^*) = x^* - a = \left(-\frac{3}{2}, -\frac{3}{2}\right)^\top \neq 0$！为什么球没有顺着负梯度 $-\nabla f_0(x^*) = \left(\frac{3}{2}, \frac{3}{2}\right)^\top$ 继续往右上角滚？因为边界围墙 $f_1(x) = 0$ 的外法向量恰好是 $\nabla f_1(x^*) = (1, 1)^\top$，围墙产生了一股向内推的\textbf{支撑反作用力} $-\lambda_1^* \nabla f_1(x^*)$（这里 $\lambda_1^* = \frac{3}{2} > 0$），与驱动球向外滚的拉力 $-\nabla f_0(x^*)$ 精确平衡：
\[
  -\nabla f_0(x^*) = \lambda_1^* \nabla f_1(x^*) \iff \nabla f_0(x^*) + \lambda_1^* \nabla f_1(x^*) = 0,\quad \lambda_1^* \ge 0.
\]
如果把目标点换成已经在墙内的 $a' = (0, 0)^\top$，则 $x^* = (0, 0)^\top$ 根本没碰到围墙（$f_1(x^*) = -2 < 0$），围墙不提供任何推力（$\lambda_1^* = 0$），恢复 $\nabla f_0(x^*) = 0$。注意两种情形竟然可以统一写成同一个乘积方程：
\[
  \lambda_1^* f_1(x^*) = 0!
\]
这正是名震优化史的\textbf{互补松弛条件（Complementary Slackness）}的初现端倪！在第 4 章的 SVM 中，每一个训练样本都是一堵这样的“间隔围墙”，绝大多数远离边界的样本都没被最优超平面碰到（推力 $\alpha_i = 0$），唯有死死顶住超平面的少数样本（$\alpha_i > 0$）成为\textbf{支持向量}。
\end{example}

% ============================================================
\section{拉格朗日函数、对偶问题与弱/强对偶定理}
% ============================================================

现在我们把这种“带反作用力的围墙”推广到一般带约束优化问题（暂不要求凸性！）：

\subsection{原始问题与拉格朗日对偶函数}

考虑如下\textbf{标准形式原始问题（Primal Problem）}：
\begin{equation}\label{eq:v4-primal}
  \begin{aligned}
    p^* \coloneqq \min_{x \in \mathcal{D}} \quad & f_0(x) \\
    \text{s.t.} \quad & f_i(x) \le 0, \quad i = 1, \dots, m, \\
                      & h_j(x) = 0, \quad j = 1, \dots, p,
  \end{aligned}
\end{equation}
其中定义域 $\mathcal{D} = \bigcap_{i=0}^m \operatorname{dom}(f_i) \cap \bigcap_{j=1}^p \operatorname{dom}(h_j) \subseteq \R^d$ 非空，$p^*$ 称为原始最优值。

\begin{defBox}
\begin{definition}[拉格朗日函数与对偶函数]\label{def:v4-lagrangian}
引入不等式乘子向量 $\lambda \in \R^m$（要求 $\lambda \succeq 0$，即 $\lambda_i \ge 0$）和等式乘子向量 $\nu \in \R^p$，定义\textbf{拉格朗日函数（Lagrangian）} $\mathcal{L}: \mathcal{D} \times \R^m \times \R^p \to \R$ 为
\begin{equation}\label{eq:v4-lagrangian}
  \mathcal{L}(x, \lambda, \nu) \coloneqq f_0(x) + \sum_{i=1}^m \lambda_i f_i(x) + \sum_{j=1}^p \nu_j h_j(x).
\end{equation}
进一步，定义\textbf{拉格朗日对偶函数（Lagrange Dual Function）} $g: \R^m \times \R^p \to \R \cup \{-\infty\}$ 为拉格朗日函数关于原始变量 $x$ 的下确界：
\begin{equation}\label{eq:v4-dual-fn}
  g(\lambda, \nu) \coloneqq \inf_{x \in \mathcal{D}} \mathcal{L}(x, \lambda, \nu)
  = \inf_{x \in \mathcal{D}} \left(f_0(x) + \sum_{i=1}^m \lambda_i f_i(x) + \sum_{j=1}^p \nu_j h_j(x)\right).
\end{equation}
\end{definition}
\end{defBox}

拉格朗日函数有一种极为传神的\textbf{博弈论（零和极大极小）解释}：把 $\lambda_i \ge 0$ 看作违反第 $i$ 条不等式约束 $f_i(x) \le 0$ 的单价罚款，把 $\nu_j \in \R$ 看作违反第 $j$ 条等式约束 $h_j(x) = 0$ 的双向罚款。如果对手（乘子玩家）可以任意哄抬罚款，则：
\[
  \sup_{\lambda \succeq 0,\, \nu} \mathcal{L}(x, \lambda, \nu) =
  \begin{cases}
    f_0(x), & \text{若 } x \text{ 满足所有约束（可行），}\\
    +\infty, & \text{若 } x \text{ 违反任意一条约束（不可行）。}
  \end{cases}
\]
因此，原始带约束问题~\eqref{eq:v4-primal} 完全等价于无约束极小极大博弈 $p^* = \inf_{x\in\mathcal{D}} \sup_{\lambda \succeq 0,\, \nu} \mathcal{L}(x, \lambda, \nu)$；而如果\textbf{交换 $\inf$ 与 $\sup$ 的次序}，就得到了\textbf{拉格朗日对偶问题（Dual Problem）}：
\begin{equation}\label{eq:v4-dual-prob}
  d^* \coloneqq \sup_{\lambda \succeq 0,\, \nu}\; g(\lambda, \nu) = \sup_{\lambda \succeq 0,\, \nu}\;\inf_{x \in \mathcal{D}} \mathcal{L}(x, \lambda, \nu).
\end{equation}

对偶函数有一个令所有优化学家赞叹的普遍性质——无论原始问题长得多么奇形怪状、多么非凸，对偶问题永远是一个凹最大化（等价于凸最小化）问题！

\begin{thmBox}
\begin{theorem}[对偶函数的普适凹性]\label{thm:v4-dual-concave}
即使原始问题~\eqref{eq:v4-primal} 的目标函数 $f_0$ 与约束函数 $f_i, h_j$ 均\textbf{非凸}，拉格朗日对偶函数 $g(\lambda, \nu)$ 关于 $(\lambda, \nu)$ 也\textbf{永远是凹函数}。
\end{theorem}
\end{thmBox}

\begin{proof}
固定任意 $x \in \mathcal{D}$，观察 $\mathcal{L}(x, \lambda, \nu) = f_0(x) + \sum_i f_i(x)\lambda_i + \sum_j h_j(x)\nu_j$：它是关于 $(\lambda, \nu)$ 的仿射函数（线性加常数），因此 $-\mathcal{L}(x, \lambda, \nu)$ 是关于 $(\lambda, \nu)$ 的凸函数。于是
\[
  -g(\lambda, \nu) = -\inf_{x \in \mathcal{D}} \mathcal{L}(x, \lambda, \nu) = \sup_{x \in \mathcal{D}} \big(-\mathcal{L}(x, \lambda, \nu)\big).
\]
由第 1 章命题~\ref{prop:v4-preserve} 第 3 条，一族凸函数的逐点上确界必为凸函数，故 $-g(\lambda, \nu)$ 为凸函数，即 $g(\lambda, \nu)$ 恒为凹函数。
\end{proof}

\subsection{弱对偶定理、对偶间隙与 Slater 强对偶条件}

\begin{thmBox}
\begin{theorem}[弱对偶定理，Weak Duality]\label{thm:v4-weak-duality}
对任意优化问题~\eqref{eq:v4-primal}（无需凸性！），只要 $\lambda \succeq 0$ 且 $\tilde{x}$ 为原始可行点，就有 $g(\lambda, \nu) \le f_0(\tilde{x})$。进而恒成立：
\begin{equation}\label{eq:v4-weak-dual}
  d^* \le p^*.
\end{equation}
非负差值 $p^* - d^* \ge 0$ 称为\textbf{最优对偶间隙（Optimal Duality Gap）}。
\end{theorem}
\end{thmBox}

\begin{proof}
任取原始可行点 $\tilde{x} \in \mathcal{D}$（满足 $f_i(\tilde{x}) \le 0$ 与 $h_j(\tilde{x}) = 0$）以及任意对偶可行乘子 $\lambda \succeq 0,\; \nu \in \R^p$。直接利用下确界定义与符号判定：
\[
  g(\lambda, \nu) = \inf_{x \in \mathcal{D}} \mathcal{L}(x, \lambda, \nu)
  \le \mathcal{L}(\tilde{x}, \lambda, \nu)
  = f_0(\tilde{x}) + \sum_{i=1}^m \underbrace{\lambda_i}_{\ge 0}\underbrace{f_i(\tilde{x})}_{\le 0} + \sum_{j=1}^p \nu_j \underbrace{h_j(\tilde{x})}_{= 0}
  \le f_0(\tilde{x}).
\]
由于上式对所有原始可行点 $\tilde{x}$ 成立，对右侧取下确界得 $g(\lambda, \nu) \le p^*$；再对左侧取 $\sup_{\lambda \succeq 0, \nu}$，立得 $d^* \le p^*$。
\end{proof}

弱对偶为我们提供了一张“免死金牌”式的全局下界证书：在求解任何难题（哪怕是 NP-难组合优化）时，随便算一个对偶可行点 $(\lambda, \nu)$ 对应的 $g(\lambda, \nu)$，就能严格保证最优值绝不可能低于 $g(\lambda, \nu)$。而当 $d^* = p^*$（对偶间隙为零）时，我们称\textbf{强对偶（Strong Duality）}成立——此时直接去解更简单、更低维或带有核结构的对偶问题，就能精准复原原始问题的最优值！

强对偶并非免费午餐。对于凸优化问题，只要满足温和的约束规范——\textbf{Slater 条件}，强对偶便必然成立：

\begin{thmBox}
\begin{theorem}[Slater 强对偶条件，不证自引]\label{thm:v4-slater}
若原始问题~\eqref{eq:v4-primal} 是\textbf{凸优化问题}（即 $f_0, f_1, \dots, f_m$ 均为凸函数，等式约束为仿射函数 $h_j(x) = a_j^\top x - b_j$），且存在一个\textbf{严格可行内点} $\tilde{x} \in \operatorname{relint}(\mathcal{D})$ 满足
\[
  f_i(\tilde{x}) < 0 \quad (i = 1, \dots, m), \qquad A\tilde{x} = b,
\]
则强对偶成立：$d^* = p^*$。进一步地，如果某些不等式约束 $f_i(x)$ 本身是仿射函数（例如线性不等式 $c_i^\top x \le d_i$），则对这些仿射约束无需严格小于，仅需 $f_i(\tilde{x}) \le 0$ 即可。
\end{theorem}
\end{thmBox}

特别地，如果所有约束都是仿射线性约束（例如线性规划 LP 与第 4 章的 SVM 二次规划 QP），只要原始可行域非空且 $p^* > -\infty$，强对偶自动成立！反之，当原始问题\textbf{非凸}时，即使存在严格内点，强对偶也可能彻底崩塌。

\begin{failBox}{非凸问题存在严格对偶间隙（$d^* < p^*$）的极简反例}
考察定义域为闭区间 $\mathcal{D} = [0, 2]$ 上的一维\textbf{非凸（凹）}优化问题：
\[
  \min_{x \in [0, 2]}\; f_0(x) = -x^2 \quad\text{s.t.}\quad f_1(x) = x - 1 \le 0.
\]
\begin{enumerate}
  \item \textbf{求原始最优值 $p^*$}：可行域为 $[0, 2] \cap (-\infty, 1] = [0, 1]$，在 $[0, 1]$ 上最小化 $-x^2$ 显然在 $x^* = 1$ 处取到，故 $p^* = -1$。
  \item \textbf{求对偶函数 $g(\lambda)$}：对 $\lambda \ge 0$，拉格朗日函数 $\mathcal{L}(x, \lambda) = -x^2 + \lambda(x - 1)$ 关于 $x$ 的二阶导数为 $-2 < 0$（严格凹函数），因此在闭区间 $[0, 2]$ 上的最小值\textbf{必然落在两个端点 $x = 0$ 或 $x = 2$ 之一}：
  \[
    g(\lambda) = \min_{x \in [0, 2]} \big(-x^2 + \lambda x - \lambda\big)
    = \min\big\{\mathcal{L}(0, \lambda),\; \mathcal{L}(2, \lambda)\big\}
    = \min\{-\lambda,\; \lambda - 4\}.
  \]
  注意 $g(\lambda)$ 确实是凹函数（两条直线的下包络，印证定理~\ref{thm:v4-dual-concave}）。
  \item \textbf{求对偶最优值 $d^*$}：两条射线 $-\lambda$ 与 $\lambda - 4$ 在 $-\lambda = \lambda - 4 \implies \lambda^* = 2 \ge 0$ 处相交达到最高峰，故 $d^* = g(2) = -2$。
\end{enumerate}
于是我们得到了精确的严格对偶间隙：$d^* = -2 < -1 = p^*$，$p^* - d^* = 1 > 0$！即便存在严格可行内点 $\tilde{x} = 0.5$（满足 $f_1(0.5) = -0.5 < 0$），一旦失去凸性，强对偶便不再有保障。
\end{failBox}

% ============================================================
\section{KKT 条件：带约束优化的终极判据}
% ============================================================

当强对偶成立时，原始最优解 $x^*$ 与对偶最优解 $(\lambda^*, \nu^*)$ 之间会锁死成一组极为紧致的代数方程与不等式——这便是 Karush（1939）与 Kuhn-Tucker（1951）各自独立发现的 \textbf{KKT 条件}。

\begin{defBox}
\begin{definition}[KKT 条件，Karush-Kuhn-Tucker Conditions]\label{def:v4-kkt}
假设原始问题~\eqref{eq:v4-primal} 中的所有函数 $f_0, f_i, h_j$ 均可微。称点组 $(x^*, \lambda^*, \nu^*)$ 满足 \textbf{KKT 条件}，如果以下四组条件同时成立：
\begin{enumerate}
  \item \textbf{原始可行性（Primal Feasibility）}：$f_i(x^*) \le 0\;(i=1,\dots,m)$，$h_j(x^*) = 0\;(j=1,\dots,p)$；
  \item \textbf{对偶可行性（Dual Feasibility）}：$\lambda_i^* \ge 0\;(i=1,\dots,m)$；
  \item \textbf{互补松弛性（Complementary Slackness）}：$\lambda_i^* f_i(x^*) = 0\;(i=1,\dots,m)$；
  \item \textbf{平稳性（Stationarity）}：
  \[
    \nabla_x \mathcal{L}(x^*, \lambda^*, \nu^*) = \nabla f_0(x^*) + \sum_{i=1}^m \lambda_i^* \nabla f_i(x^*) + \sum_{j=1}^p \nu_j^* \nabla h_j(x^*) = 0.
  \]
\end{enumerate}
\end{definition}
\end{defBox}

下面这条定理堪称带约束凸优化的“总纲”，它的证明极其精炼，请务必亲手跟着每一行不等式走一遍！

\begin{thmBox}
\begin{theorem}[KKT 条件的必要性与凸问题充分性]\label{thm:v4-kkt-thm}
设问题~\eqref{eq:v4-primal} 中所有函数在开集 $\mathcal{D} = \R^d$ 上可微：
\begin{enumerate}
  \item \textbf{必要性（强对偶 $\implies$ KKT）}：无论原问题是否凸，只要强对偶成立且 $x^*$ 与 $(\lambda^*, \nu^*)$ 分别是原始最优解与对偶最优解，则 $(x^*, \lambda^*, \nu^*)$ 必满足全部四组 KKT 条件。
  \item \textbf{充分性（凸问题 + KKT $\implies$ 全局最优且强对偶）}：若原问题为凸优化问题（$f_0, f_i$ 凸，$h_j$ 仿射），且存在 $(x^*, \lambda^*, \nu^*)$ 满足 KKT 条件，则 $x^*$ 必为原始全局最优解，$(\lambda^*, \nu^*)$ 必为对偶全局最优解，且强对偶成立（$d^* = p^*$）。
\end{enumerate}
\end{theorem}
\end{thmBox}

\begin{proof}
\textbf{第一部分（必要性）}：由于 $x^*$ 与 $(\lambda^*, \nu^*)$ 分别可行，原始可行性与对偶可行性自动满足。由强对偶条件 $f_0(x^*) = g(\lambda^*, \nu^*)$，写下贯穿全书的核心“夹逼链”：
\begin{align*}
  f_0(x^*) &= g(\lambda^*, \nu^*) = \inf_{x \in \R^d} \mathcal{L}(x, \lambda^*, \nu^*) \\
  &\le \mathcal{L}(x^*, \lambda^*, \nu^*) \\
  &= f_0(x^*) + \sum_{i=1}^m \underbrace{\lambda_i^* f_i(x^*)}_{\le 0} + \sum_{j=1}^p \nu_j^* \underbrace{h_j(x^*)}_{= 0}
  \le f_0(x^*).
\end{align*}
首尾都是 $f_0(x^*)$，因此链条中的\textbf{每一个“$\le$”都必须强制取“$=$”}！
\begin{itemize}
  \item 由第二个“$\le$”取等，得 $\sum_{i=1}^m \lambda_i^* f_i(x^*) = 0$。由于求和中的每一项都满足 $\lambda_i^* f_i(x^*) \le 0$，一群非正数之和为零当且仅当\textbf{每一项各自为零}：$\lambda_i^* f_i(x^*) = 0$（互补松弛性证毕！）。
  \item 由第一个“$\le$”取等，得 $\mathcal{L}(x^*, \lambda^*, \nu^*) = \inf_{x\in\R^d} \mathcal{L}(x, \lambda^*, \nu^*)$，说明 $x^*$ 是无约束函数 $x \mapsto \mathcal{L}(x, \lambda^*, \nu^*)$ 的全局极小点，由 Fermat 引理必有 $\nabla_x \mathcal{L}(x^*, \lambda^*, \nu^*) = 0$（平稳性证毕！）。
\end{itemize}

\textbf{第二部分（充分性）}：设原问题为凸问题且 $(x^*, \lambda^*, \nu^*)$ 满足四组 KKT 条件。因为 $\lambda_i^* \ge 0$、$f_0, f_i$ 为凸函数且 $h_j$ 为仿射函数，由第 1 章保凸运算（命题~\ref{prop:v4-preserve}），函数 $x \mapsto \mathcal{L}(x, \lambda^*, \nu^*)$ 是关于 $x$ 的可微凸函数。结合平稳性条件 $\nabla_x \mathcal{L}(x^*, \lambda^*, \nu^*) = 0$ 与定理~\ref{thm:v4-global-opt}，可知 $x^*$ 是 $\mathcal{L}(\cdot, \lambda^*, \nu^*)$ 的全局极小点！于是：
\[
  g(\lambda^*, \nu^*) = \inf_{x\in\R^d}\mathcal{L}(x, \lambda^*, \nu^*) = \mathcal{L}(x^*, \lambda^*, \nu^*)
  = f_0(x^*) + \sum_{i=1}^m \underbrace{\lambda_i^* f_i(x^*)}_{= 0} + \sum_{j=1}^p \nu_j^* \underbrace{h_j(x^*)}_{= 0} = f_0(x^*).
\]
最后联合弱对偶定理~\ref{thm:v4-weak-duality} 与可行性定义：$f_0(x^*) = g(\lambda^*, \nu^*) \le d^* \le p^* \le f_0(x^*)$。夹逼立得 $g(\lambda^*, \nu^*) = d^* = p^* = f_0(x^*)$，故 $x^*$ 与 $(\lambda^*, \nu^*)$ 分别为原始与对偶全局最优解！
\end{proof}

\subsection{手算演练：用 KKT 求解二维不等式约束二次规划}

\begin{example}[含两条不等式约束的二维 QP 完整 KKT 求解]\label{ex:v4-ch2-qp}
扩展例~\ref{ex:v4-ch2-halfspace}，增加第二条约束 $x_2 \ge 0$，求解凸二次规划：
\[
  \min_{x_1, x_2 \in \R}\; \frac{1}{2}\big((x_1 - 3)^2 + (x_2 - 2)^2\big)
  \quad\text{s.t.}\quad f_1(x) = x_1 + x_2 - 2 \le 0,\quad f_2(x) = -x_2 \le 0.
\]
\textbf{第 1 步：列出 KKT 方程组。} 拉格朗日函数为
\[
  \mathcal{L}(x, \lambda) = \frac{1}{2}(x_1 - 3)^2 + \frac{1}{2}(x_2 - 2)^2 + \lambda_1(x_1 + x_2 - 2) - \lambda_2 x_2.
\]
平稳性条件 $\nabla_x \mathcal{L} = 0$ 给出：
\[
  x_1 - 3 + \lambda_1 = 0 \implies x_1 = 3 - \lambda_1, \qquad
  x_2 - 2 + \lambda_1 - \lambda_2 = 0 \implies x_2 = 2 - \lambda_1 + \lambda_2.
\]
互补松弛与可行性条件为：$\lambda_1(x_1 + x_2 - 2) = 0,\; \lambda_2 x_2 = 0,\; \lambda_1, \lambda_2 \ge 0,\; x_1 + x_2 \le 2,\; x_2 \ge 0$。

\textbf{第 2 步：按积极约束（Active Set）分类讨论。}
\begin{itemize}
  \item \emph{假设 $\lambda_1 = 0$（第一条约束松弛）}：则 $x_1 = 3$。若 $\lambda_2 = 0$ 则 $x_2 = 2$，得 $x_1 + x_2 = 5 > 2$ 矛盾；若 $\lambda_2 > 0$ 则由互补松弛必有 $x_2 = 0$，得 $x_1 + x_2 = 3 > 2$ 仍矛盾！因此必有 $\lambda_1 > 0$，即约束 1 处于\textbf{积极（激活）状态}：$x_1 + x_2 = 2$。
  \item \emph{在 $\lambda_1 > 0$ 下测试 $\lambda_2 = 0$（约束 2 松弛）}：将 $x_1 = 3 - \lambda_1$ 与 $x_2 = 2 - \lambda_1$ 代入 $x_1 + x_2 = 2$：
  \[
    (3 - \lambda_1) + (2 - \lambda_1) = 2 \implies \lambda_1^* = \frac{3}{2} > 0.
  \]
  回代得 $x_1^* = 3 - \frac{3}{2} = \frac{3}{2}$，且 $x_2^* = 2 - \frac{3}{2} = \frac{1}{2} > 0$。
\end{itemize}
\textbf{第 3 步：核对全部四组 KKT 条件。} 原始可行（$\frac{3}{2}+\frac{1}{2}=2\le 2,\; -\frac{1}{2}\le 0$）、对偶可行（$\lambda_1^*=\frac{3}{2}\ge 0,\; \lambda_2^*=0\ge 0$）、互补松弛（$\frac{3}{2}\cdot 0 = 0,\; 0\cdot(-\frac{1}{2})=0$）、平稳性全部成立！由定理~\ref{thm:v4-kkt-thm} 的充分性，这组解就是唯一的全局最优解，最优值 $p^* = \frac{1}{2}\big(\frac{9}{4}+\frac{9}{4}\big) = \frac{9}{4}$。
\end{example}

\begin{insightBox}{直通第 4 章 SVM：互补松弛如何孕育“支持向量”}
把例~\ref{ex:v4-ch2-qp} 中的两条不等式换作 $N$ 个样本的分类间隔约束 $f_i(w, b) = 1 - y_i(w^\top x_i + b) \le 0$（乘子记为 $\alpha_i \ge 0$），拉格朗日平稳性将直接把法向量表示为样本的线性组合 $w^* = \sum_{i=1}^N \alpha_i^* y_i x_i$；而互补松弛条件
\[
  \alpha_i^* \Big(1 - y_i({w^*}^\top x_i + b^*)\Big) = 0
\]
立刻给出惊人的稀疏结构：任何严格落在安全间隔之外的样本（$y_i({w^*}^\top x_i + b^*) > 1$）必有 $\alpha_i^* = 0$，对决策边界毫无贡献；唯有乘子 $\alpha_i^* > 0$ 的点必满足 $y_i({w^*}^\top x_i + b^*) = 1$——它们正是紧紧顶住间隔边界的\textbf{支持向量（Support Vectors）}！
\end{insightBox}

% ============================================================
\section{当目标出现尖角：次梯度与近端梯度法（ISTA）}
% ============================================================

解决了“围墙”，现在我们迎战第二道坎：\textbf{不可微凸函数}。在机器学习中，最著名的不可微凸项当属 Lasso 回归中的 $\ell_1$ 正则项 $\lambda\norm{w}_1 = \lambda\sum_{j=1}^D |w_j|$ 与 SVM 的合页损失 $\max(0, 1 - z)$。在尖角处，切平面不再唯一，但“从下方托住函数图像”这一凸性本质依然完好无损！

\subsection{次梯度与次微分}

\begin{defBox}
\begin{definition}[次梯度与次微分]\label{def:v4-subgrad}
设 $f: \R^d \to \R$ 为凸函数（无需可微）。向量 $g \in \R^d$ 称为 $f$ 在点 $x$ 处的\textbf{次梯度（Subgradient）}，如果它定义的仿射超平面在全局范围内从下方托住 $f$：
\begin{equation}\label{eq:v4-subgrad-def}
  f(y) \ge f(x) + g^\top (y - x), \qquad \forall y \in \R^d.
\end{equation}
$f$ 在 $x$ 处所有次梯度的集合称为\textbf{次微分（Subdifferential）}，记作 $\partial f(x)$。
\end{definition}
\end{defBox}

当 $f$ 在 $x$ 处可微时，切平面唯一，故 $\partial f(x) = \{\nabla f(x)\}$。而在不可微尖角处，次微分是一个闭凸集。以最核心的绝对值函数 $f(x) = |x|$ 为例：
\begin{itemize}
  \item $x > 0$ 时可微，$\partial |x| = \{+1\}$；$x < 0$ 时可微，$\partial |x| = \{-1\}$；
  \item $x = 0$ 时，斜率 $g$ 满足 $|y| \ge 0 + g\cdot(y - 0) = g y$ 对一切 $y \in \R$ 成立，等价于 $|g| \le 1$！因此
  \begin{equation}\label{eq:v4-subdiff-abs}
    \partial |x| =
    \begin{cases}
      \{-1\}, & x < 0, \\
      [-1,\, 1], & x = 0, \\
      \{+1\}, & x > 0.
    \end{cases}
  \end{equation}
\end{itemize}
与可微凸函数定理~\ref{thm:v4-global-opt} 一脉相承，次梯度拥有一行即可证明的全局最优性判据：

\begin{proposition}[广义 Fermat 最优性条件]\label{prop:v4-subgrad-opt}
对任意凸函数 $f: \R^d \to \R$，$x^*$ 是 $f$ 的全局极小点，当且仅当 $0 \in \partial f(x^*)$。
\end{proposition}
\begin{proof}
由定义~\ref{def:v4-subgrad}，$0 \in \partial f(x^*) \iff f(y) \ge f(x^*) + 0^\top(y - x^*) = f(x^*)\ \forall y \in \R^d \iff x^*$ 是全局极小点。
\end{proof}

\begin{failBox}{为什么直接做“次梯度下降”既慢又得不到精确稀疏解？}
既然有了次梯度，很自然会想到把梯度下降改为 $x_{t+1} = x_t - \eta_t g_t$（$g_t \in \partial f(x_t)$）。但这套朴素方案有两大硬伤：
\begin{enumerate}
  \item \textbf{无法使用常数步长，收敛退化为 $O(1/\sqrt{T})$}：以 $f(x) = |x|$ 为例，即使无限接近最优解 $x^* = 0$，次梯度的模长始终为 $1$ 而不会趋于零！若取固定步长 $\eta$，迭代点将在 $[-\eta, \eta]$ 内永无止境地来回跳跃。迫不得已只能使用衰减步长 $\eta_t \propto 1/\sqrt{t}$，收敛率暴跌至 $O(1/\sqrt{T})$（误差降到 $10^{-3}$ 需要百万步！）。
  \item \textbf{几乎永远产不出精确零分量}：浮点数累加 $w_j - \eta_t (\nabla_j f + \lambda \sign(w_j))$ 很难恰好踩中 $0.0$，最终结果往往是一堆 $10^{-4}$ 量级的非零毛刺，彻底毁掉 Lasso 特征选择的稀疏美学。
\end{enumerate}
\end{failBox}

\subsection{复合凸优化、近端算子与软阈值闭式解}

怎么才能既保留 $L$-光滑的 $O(1/T)$ 高速，又干净利落地劈出精确零解？核心洞察在于\textbf{拆分结构}：Lasso 这类问题并非一团浆糊式的不可微函数，而是\textbf{“光滑凸损失 + 简单不可微凸正则”的复合结构}：
\begin{equation}\label{eq:v4-composite}
  \min_{x \in \R^d}\; F(x) = f(x) + g(x),
\end{equation}
其中 $f(x)$ 是 $L$-光滑凸函数（如最小二乘 $\frac{1}{2N}\norm{Xw-y}_2^2$），$g(x)$ 是凸但可能不可微的简单函数（如 $\lambda\norm{w}_1$）。

回想第 1 章图~\ref{fig:v4-sandwich} 的二次代理思想：\textbf{只对光滑部分 $f(x)$ 在 $x_t$ 处做二次上界展开，而将不可微部分 $g(x)$ 原封不动地保留下来！} 于是每一步我们求解代理子问题：
\begin{align}
  x_{t+1} &= \argmin_{x \in \R^d} \left\{ f(x_t) + \nabla f(x_t)^\top (x - x_t) + \frac{1}{2\eta}\norm{x - x_t}_2^2 + g(x) \right\} \nonumber \\
  &= \argmin_{x \in \R^d} \left\{ \frac{1}{2}\norm{x - \big(x_t - \eta \nabla f(x_t)\big)}_2^2 + \eta g(x) \right\}. \label{eq:v4-prox-gd-deriv}
\end{align}
（从第一行到第二行只需把常数项丢掉、乘以 $\eta$ 并对 $x$ 配方即可验证！）由此，我们隆重定义优化领域最优雅的算子之一——\textbf{近端算子（Proximal Operator）}。

\begin{defBox}
\begin{definition}[近端算子与近端梯度法]\label{def:v4-prox}
对任意闭真凸函数 $g: \R^d \to \R \cup \{+\infty\}$ 与步长参数 $\tau > 0$，定义其\textbf{近端算子}为
\begin{equation}\label{eq:v4-prox-def}
  \prox_{\tau g}(z) \coloneqq \argmin_{x \in \R^d} \left\{ \frac{1}{2}\norm{x - z}_2^2 + \tau g(x) \right\}.
\end{equation}
由于目标函数关于 $x$ 是 $1$-强凸的，上述极小点必然\textbf{存在且唯一}。利用它，复合优化~\eqref{eq:v4-composite} 的\textbf{近端梯度法（Proximal Gradient Descent）}简洁地写作两步曲：
\begin{equation}\label{eq:v4-pgd}
  x_{t+1} = \prox_{\eta g}\!\Big(x_t - \eta \nabla f(x_t)\Big).
\end{equation}
\end{definition}
\end{defBox}

近端算子具有惊人的统一力：
\begin{itemize}
  \item 若 $g(x) \equiv 0$，则 $\prox_{\eta g}(z) = z$，退化为第 1 章的普通梯度下降；
  \item 若 $g(x) = I_C(x)$（凸集 $C$ 的示性函数：在 $C$ 内为 $0$，在 $C$ 外为 $+\infty$），则近端算子化为欧氏投影 $\prox_{\eta I_C}(z) = \argmin_{x\in C}\frac{1}{2}\norm{x-z}_2^2 = \Pi_C(z)$，恰好就是\textbf{投影梯度下降（Projected Gradient Descent）}！
  \item 若 $g(x) = \lambda \norm{x}_1$，则近端算子拥有令人拍案叫绝的逐坐标解析解——\textbf{软阈值算子（Soft-Thresholding）}！
\end{itemize}

\begin{thmBox}
\begin{theorem}[$\ell_1$ 范数的近端算子与软阈值闭式解]\label{thm:v4-soft-thresh}
对于 $g(x) = \norm{x}_1$ 与阈值 $\tau > 0$，近端算子 $\prox_{\tau\norm{\cdot}_1}(z)$ 按坐标完全解耦，其第 $j$ 个分量等于\textbf{软阈值算子} $S_\tau(z_j)$：
\begin{equation}\label{eq:v4-soft-thresh}
  \big[\prox_{\tau\norm{\cdot}_1}(z)\big]_j = S_\tau(z_j) \coloneqq \sign(z_j)\max\big(|z_j| - \tau,\, 0\big) =
  \begin{cases}
    z_j - \tau, & z_j > \tau, \\
    0, & |z_j| \le \tau, \\
    z_j + \tau, & z_j < -\tau.
  \end{cases}
\end{equation}
\end{theorem}
\end{thmBox}

\begin{proof}
将定义~\eqref{eq:v4-prox-def} 按坐标展开：
\[
  \frac{1}{2}\norm{x - z}_2^2 + \tau \norm{x}_1 = \sum_{j=1}^d \left[\frac{1}{2}(x_j - z_j)^2 + \tau |x_j|\right].
\]
各坐标互不耦合，因此只需对任意标量 $z \in \R$ 独立最小化一维严格凸函数 $\phi(x) = \frac{1}{2}(x - z)^2 + \tau|x|$。由广义 Fermat 最优性条件（命题~\ref{prop:v4-subgrad-opt}），$x^*$ 是唯一极小点当且仅当
\[
  0 \in \partial \phi(x^*) = (x^* - z) + \tau\,\partial |x^*|.
\]
结合式~\eqref{eq:v4-subdiff-abs} 分三种情形直接求解：
\begin{enumerate}
  \item \textbf{若 $x^* > 0$}：$\partial |x^*| = \{+1\}$，方程化为 $x^* - z + \tau = 0 \implies x^* = z - \tau$。这与前提 $x^* > 0$ 相容当且仅当 $z > \tau$。
  \item \textbf{若 $x^* < 0$}：$\partial |x^*| = \{-1\}$，方程化为 $x^* - z - \tau = 0 \implies x^* = z + \tau$。这与前提 $x^* < 0$ 相容当且仅当 $z < -\tau$。
  \item \textbf{若 $x^* = 0$}：$\partial |0| = [-1, 1]$，条件化为 $0 \in -z + [-\tau, \tau] \iff z \in [-\tau, \tau] \iff |z| \le \tau$。
\end{enumerate}
三种情形恰好将整个实轴 $z \in \R$ 互斥且无遗漏地完全覆盖，综合即得软阈值公式~\eqref{eq:v4-soft-thresh}！
\end{proof}

将定理~\ref{thm:v4-soft-thresh} 代入近端梯度框架~\eqref{eq:v4-pgd}（令 $\tau = \eta\lambda$），即得到著名的\textbf{迭代收缩阈值算法（ISTA, Iterative Shrinkage-Thresholding Algorithm）}求解 Lasso 问题 $\min_w \frac{1}{2N}\norm{Xw - y}_2^2 + \lambda\norm{w}_1$：
\begin{equation}\label{eq:v4-ista}
  w^{(t+1)} = S_{\eta\lambda}\!\left(w^{(t)} - \frac{\eta}{N}X^\top\big(Xw^{(t)} - y\big)\right), \qquad \eta = \frac{1}{L},\quad L = \frac{\sigma_{\max}(X)^2}{N}.
\end{equation}
令人振奋的是：尽管 Lasso 目标不可微，可以严格证明（沿用第 1 章定理~\ref{thm:v4-gd-convex} 的望远镜证明思路将梯度替换为“广义梯度映射”），取固定步长 $\eta = 1/L$ 的 ISTA 完整恢复了光滑问题的 \textbf{$O(1/T)$ 收敛率}，同时软阈值的死区 $[-\eta\lambda, \eta\lambda]$ 会在有限步内将所有微弱噪声特征直接精准归零！

\begin{figure}[htbp]
\centering
\begin{tikzpicture}[scale=0.95, >=Stealth]
  % ===== Left Panel: L1 Diamond vs L2 Circle Contours =====
  \begin{scope}[shift={(0,0)}]
    \node[font=\small\sffamily\bfseries, color=tealDark] at (0.2, 3.35) {(a) $L_1$ 菱形尖角与 $L_2$ 圆球的稀疏几何对比};
    % Axes
    \draw[->, slate!80] (-2.2, 0) -- (2.7, 0) node[right, font=\footnotesize] {$w_1$};
    \draw[->, slate!80] (0, -1.75) -- (0, 2.60) node[above, font=\footnotesize] {$w_2$};

    % L2 ball (radius 1.15)
    \fill[navyLight, opacity=0.65] (0,0) circle (1.15);
    \draw[navy, thick, dashed] (0,0) circle (1.15);
    \node[navy, font=\scriptsize\sffamily, anchor=east] (l2lbl) at (-1.30, -1.28) {$L_2$ 球 $\norm{w}_2 \le C$};
    \draw[->, navy!80, thin] (l2lbl.east) -- (-0.82, -0.82);

    % L1 diamond (radius 1.15)
    \fill[tealLight, opacity=0.85] (1.15,0) -- (0,1.15) -- (-1.15,0) -- (0,-1.15) -- cycle;
    \draw[tealDark, very thick] (1.15,0) -- (0,1.15) -- (-1.15,0) -- (0,-1.15) -- cycle;
    \node[tealDark, font=\scriptsize\sffamily, anchor=east] (l1lbl) at (-1.30, 0.85) {$L_1$ 菱形 $\norm{w}_1 \le C$};
    \draw[->, tealDark!85, thin] (l1lbl.east) -- (-0.58, 0.58);

    % Elliptical loss contours centered at w_ols = (1.55, 1.80)
    \begin{scope}[shift={(1.55, 1.80)}, rotate=28]
      \draw[ochre!75, thin] (0,0) ellipse (0.52 and 0.30);
      \draw[ochre!85, semithick] (0,0) ellipse (1.05 and 0.60);
      \draw[ochre, thick] (0,0) ellipse (1.62 and 0.92);
    \end{scope}
    \fill[ochre] (1.55, 1.80) circle (2pt);
    \node[ochre, font=\scriptsize\sffamily, right] at (1.60, 1.88) {$\hat{w}_{\mathrm{OLS}}$};

    % Intersection point on w2 axis for L1 (0, 1.15) - label in open upper-left
    \fill[gold] (0, 1.15) circle (3pt);
    \draw[tealDark, semithick] (0, 1.15) circle (4.8pt);
    \node[font=\scriptsize\sffamily, color=tealDark, align=center, anchor=east] (cornerlbl) at (-0.58, 2.02)
      {$L_1$ 切于尖角\\$w_1^* = 0$（精确稀疏）};
    \draw[->, tealDark, semithick] (cornerlbl.south east) -- (-0.10, 1.24);

    % Intersection point on L2 circle approx (0.66, 0.94) - label in open lower-right
    \fill[navy] (0.66, 0.94) circle (2.5pt);
    \node[font=\scriptsize\sffamily, color=navy, align=left, anchor=west] (arclbl) at (1.28, -0.78)
      {$L_2$ 切于光滑弧\\($w_1, w_2 \neq 0$)};
    \draw[->, navy!85, thin] (1.35, -0.45) -- (0.72, 0.86);
  \end{scope}

  % ===== Right Panel: Soft-Thresholding Operator =====
  \begin{scope}[shift={(7.7, 0.25)}]
    \node[font=\small\sffamily\bfseries, color=tealDark] at (0, 3.10) {(b) 软阈值算子 $S_\tau(z)$ 与死区归零机理};
    % Axes
    \draw[->, slate!80] (-2.4, 0) -- (2.55, 0) node[right, font=\footnotesize] {$z$};
    \draw[->, slate!80] (0, -1.95) -- (0, 2.35) node[above, font=\footnotesize] {$S_\tau(z)$};

    % Identity line y = z
    \draw[slate!55, dashed, thick] (-2.05, -2.05) -- (2.05, 2.05);
    \node[slate, font=\scriptsize\sffamily, anchor=west] at (0.88, 2.15) {$y = z$（无惩罚）};

    % Dead zone shading [-tau, tau], tau = 0.85
    \fill[ochreLight, opacity=0.9] (-0.85, -0.22) rectangle (0.85, 0.22);
    \draw[ochre!70, dotted, thick] (-0.85, -1.55) -- (-0.85, 0.75);
    \draw[ochre!70, dotted, thick] (0.85, -0.85) -- (0.85, 1.55);

    % Soft thresholding piecewise linear curve
    \draw[tealDark, line width=1.8pt] (-2.20, -1.35) -- (-0.85, 0) -- (0.85, 0) -- (2.20, 1.35);
    \fill[tealDark] (-0.85, 0) circle (2.2pt);
    \fill[tealDark] (0.85, 0) circle (2.2pt);

    % Labels for -tau and tau
    \node[below, font=\footnotesize, color=ochre] at (-0.85, -0.24) {$-\tau$};
    \node[below, font=\footnotesize, color=ochre] at (0.85, -0.24) {$+\tau$};
    \node[font=\scriptsize\sffamily, color=ochre, align=center] (deadlbl) at (-1.45, 1.25)
      {死区 $[-\tau,\tau]$\\精确截断为 $0$};
    \draw[->, ochre, semithick] (deadlbl.south east) -- (-0.28, 0.08);

    % Shift annotation
    \draw[<->, navy, semithick] (1.92, 1.92) -- (1.92, 1.07) node[midway, right, font=\scriptsize\sffamily] {收缩 $\tau$};
  \end{scope}
\end{tikzpicture}
\caption{$L_1$ 稀疏性的双重视角：(a) 在约束几何上，椭圆损失等高线向外扩张时极易先撞击 $\ell_1$ 菱形位于坐标轴上的尖角，使对应权重精确归零，而光滑的 $\ell_2$ 圆球几乎总在非坐标轴处相切；(b) 在算法解析上，软阈值算子 $S_\tau(z)$ 将幅值落入死区 $[-\tau, \tau]$ 的梯度步输出严格压回零点，大权重则向零平移 $\tau$。}\label{fig:v4-l1-prox}
\end{figure}

% ============================================================
\section{实测验证：纯 NumPy ISTA 求解 20 维 Lasso 稀疏路径}
% ============================================================

{\rightskip=0pt plus 1fil\relax 下面我们用纯 NumPy 实现式~\eqref{eq:v4-ista} 的 ISTA 算法（见 \texttt{code/vol4\_experiments.py} 中的 \texttt{exp2\_lasso\_ista}），在真实含噪稀疏回归问题上检验软阈值的威力：\par}
\begin{itemize}
  \item 数据规模：$N = 120$ 个样本，$D = 20$ 维已标准化特征（零均值、单位标准差），高斯观测噪声 $\sigma = 0.5$。
  \item 真实参数向量 $w_{\text{true}} \in \R^{20}$ 中\textbf{只有前 4 维非零}：$w_{\text{true}, 0:4} = [3.5,\; -2.5,\; 1.8,\; -1.2]$，其余 16 维均为纯干扰特征（真值为 $0$）。
  \item 固定步长 $\eta = 1/L$（其中 $L = \sigma_{\max}(X)^2/N$），对 12 个递增的正则化强度 $\lambda \in [0.00, 3.80]$ 各运行 1200 步 ISTA，记录全部非零权重个数 \texttt{nnz}（阈值 $|w_j| > 10^{-4}$）及关键特征轨迹，导出至 \texttt{data/lasso\_path.csv}。
\end{itemize}

\begin{lstlisting}[caption={纯 NumPy 实现近端梯度 ISTA 追踪 Lasso 路径（节选自 \texttt{vol4\_experiments.py}）},label={lst:v4-ch2-ista}]
L = float(np.linalg.svd(X, compute_uv=False)[0] ** 2 / N)
eta = 1.0 / L

def soft_thresh(z, tau):
    return np.sign(z) * np.maximum(np.abs(z) - tau, 0.0)

for lam in [0.0, 0.05, 0.1, 0.2, 0.4, 0.6, 0.9, 1.2, 1.6, 2.2, 3.0, 3.8]:
    w = np.zeros(D)
    for _ in range(1200):
        grad = (X.T @ (X @ w - y)) / N
        w = soft_thresh(w - eta * grad, eta * lam)
    nnz = int(np.sum(np.abs(w) > 1e-4))
\end{lstlisting}

图~\ref{fig:v4-ch2-lasso} 与表~\ref{tab:v4-ch2-lasso} 忠实展示了从 \texttt{data/lasso\_path.csv} 读取的完整正则化路径（Regularization Path）。

\begin{figure}[htbp]
\centering
\begin{tikzpicture}
\begin{axis}[
  width=0.92\linewidth,
  height=6.8cm,
  xlabel={$\ell_1$ 正则化强度 $\lambda$},
  ylabel={估计权重系数 $w_j(\lambda)$},
  xmin=-0.08, xmax=3.88,
  ymin=-3.1, ymax=3.9,
  grid=both,
  major grid style={tealLine!45, line width=0.4pt},
  axis line style={slate!80},
  tick label style={font=\footnotesize\sffamily, color=ink},
  label style={font=\small\sffamily, color=ink},
  legend style={
    at={(0.98,0.97)}, anchor=north east,
    fill=white, draw=tealLine, rounded corners=1.2mm,
    font=\scriptsize\sffamily, cells={anchor=west}
  },
  every axis plot/.append style={line width=1.45pt}
]
  % Zero reference line
  \addplot[slate!65, thin, dashed, forget plot] coordinates {(0,0) (3.8,0)};
  % Vertical marker at exact recovery lam = 0.20
  \addplot[gold, thick, dotted, forget plot] coordinates {(0.20,-2.8) (0.20,3.9)};
  \node[font=\scriptsize\sffamily, color=gold!85!ink, anchor=south west, fill=white, inner sep=1.5pt]
    at (axis cs:0.24, -2.62) {$\lambda=0.20$：精确支撑恢复（$\mathrm{nnz}=4$）};

  \addplot[color=tealDark, mark=*, mark size=1.6pt]
    table[x=lam, y=w0, col sep=comma]{data/lasso_path.csv};
  \addlegendentry{$w_0$（真值 $+3.5$）}

  \addplot[color=ochre, mark=square*, mark size=1.6pt]
    table[x=lam, y=w1, col sep=comma]{data/lasso_path.csv};
  \addlegendentry{$w_1$（真值 $-2.5$）}

  \addplot[color=navy, mark=triangle*, mark size=1.9pt]
    table[x=lam, y=w2, col sep=comma]{data/lasso_path.csv};
  \addlegendentry{$w_2$（真值 $+1.8$）}

  \addplot[color=plum, mark=diamond*, mark size=1.9pt]
    table[x=lam, y=w3, col sep=comma]{data/lasso_path.csv};
  \addlegendentry{$w_3$（真值 $-1.2$）}

  \addplot[color=slate, dashed, mark=x, mark size=2pt]
    table[x=lam, y=w4_noise, col sep=comma]{data/lasso_path.csv};
  \addlegendentry{$w_4$（噪声特征，真值 $0$）}
\end{axis}
\end{tikzpicture}
\caption{20 维稀疏回归（4 个真实信号 + 16 个纯噪声特征）的实测 Lasso 正则化路径（读取自 \texttt{data/lasso\_path.csv}）。当 $\lambda = 0$ 时全部 20 个系数非零；刚增至 $\lambda = 0.20$ 时，16 个噪声特征全部被精确归零，4 个真实信号毫发无损地保留（$\mathrm{nnz}=4$）；继续增大 $\lambda$，真实信号按幅值从小到大（$w_3 \to w_2 \to w_1 \to w_0$）依次归零。}\label{fig:v4-ch2-lasso}
\end{figure}

\begin{table}[htbp]
\centering
\caption{ISTA 求解 20 维 Lasso 路径关键节点实测值（\texttt{data/lasso\_path.csv}）}\label{tab:v4-ch2-lasso}
\small
\begin{tabular}{@{}rrrrrrrc@{}}
\toprule
$\lambda$ & $w_0\,(+3.5)$ & $w_1\,(-2.5)$ & $w_2\,(+1.8)$ & $w_3\,(-1.2)$ & $w_4\,(0.0)$ & 非零个数 \texttt{nnz} & 阶段解读 \\
\midrule
$0.00$ & $3.5393$ & $-2.4444$ & $1.8266$ & $-1.1874$ & $-0.0489$ & $20$ & 无惩罚 OLS，全员非零 \\
$0.20$ & $3.3257$ & $-2.2861$ & $1.5831$ & $-0.9226$ & $0.0000$  & $4$  & \textbf{精确支撑恢复}（16 噪声归零） \\
$0.60$ & $2.8824$ & $-1.8848$ & $1.0838$ & $-0.3767$ & $0.0000$  & $4$  & 保持真支撑，系数整体收缩 \\
$1.20$ & $2.2588$ & $-1.3184$ & $0.3979$ & $0.0000$  & $0.0000$  & $3$  & 最弱信号 $w_3$ 率先归零 \\
$3.00$ & $0.5615$ & $0.0000$  & $0.0000$ & $0.0000$  & $0.0000$  & $1$  & 仅剩最强信号 $w_0$ \\
$3.80$ & $0.0000$ & $0.0000$  & $0.0000$ & $0.0000$  & $0.0000$  & $0$  & 超过 $\lambda_{\max}$，全零解 \\
\bottomrule
\end{tabular}
\end{table}

从表~\ref{tab:v4-ch2-lasso} 可以清楚地看到 Lasso 的一把“双刃剑”效应：在 $\lambda = 0.20$ 处，它干净利落地把 16 个纯噪声维度全部清零，实现了完美的模型稀疏化；但与此同时，真实信号 $w_0$ 也从 $3.5393$ 被略微向零拉低到了 $3.3257$（软阈值产生的\textbf{收缩偏差，Shrinkage Bias}）。这正好无缝引出第 3 章的中心命题——\textbf{偏差与方差的权衡}！

\begin{extBox}{加速近端梯度（FISTA）与交替方向乘子法（ADMM）}
当我们把第 1 章的 Nesterov 加速思想与本章的近端梯度法 ISTA 结合，便诞生了 Beck \& Teboulle（2009）著名的 \textbf{FISTA（Fast ISTA）}：
\[
  x_{t+1} = \prox_{\eta g}\!\big(y_t - \eta \nabla f(y_t)\big), \qquad
  \theta_{t+1} = \frac{1 + \sqrt{1 + 4\theta_t^2}}{2}, \qquad
  y_{t+1} = x_{t+1} + \frac{\theta_t - 1}{\theta_{t+1}}(x_{t+1} - x_t).
\]
仅仅增加两行外推运算，FISTA 就把复合非光滑问题的收敛率从 ISTA 的 $O(1/T)$ 直接提升至最优一阶速率 $O(1/T^2)$！

另一方面，如果目标函数包含\textbf{两个都很复杂}的凸项，或者带有线性耦合约束 $Ax + Bz = c$（例如分布式大规模机器学习或总变差图像去噪），单靠近端梯度就力不从心了。此时\textbf{交替方向乘子法（ADMM）}通过引入二次惩罚构成\textbf{增广拉格朗日函数}：
\[
  \mathcal{L}_\rho(x, z, u) = f(x) + g(z) + u^\top (Ax + Bz - c) + \frac{\rho}{2}\norm{Ax + Bz - c}_2^2,
\]
按“固定 $z$ 最小化 $x$ $\to$ 固定 $x$ 最小化 $z$ $\to$ 对偶变量 $u$ 做一步梯度上升”的节奏交替更新，将拉格朗日对偶与近端分裂的威力发挥到极致。
\end{extBox}

% ============================================================
\section*{本章分层习题}
\addcontentsline{toc}{section}{本章分层习题}
% ============================================================

\exTier{A 层}{对偶构建与软阈值手算}
\begin{enumerate}[label=\textbf{4.2.\arabic*}, leftmargin=3.3em, itemsep=5pt]
  \item\label{ex:4.2.1} \textbf{最小欧氏范数线性方程组的对偶问题}：求解凸优化问题 $\min_{x\in\R^d} \frac{1}{2}\norm{x}_2^2$，满足等式约束 $Ax = b$（其中 $A \in \R^{m\times d}$ 行满秩，$m \le d$）。请写出拉格朗日函数 $\mathcal{L}(x, \nu)$、对偶函数 $g(\nu)$、对偶最优解 $\nu^*$ 以及对应的原始最优解 $x^* = A^\top (AA^\top)^{-1} b$，并验证 $d^* = p^*$。
  \exHint{对 $\mathcal{L}(x,\nu) = \frac{1}{2}x^\top x + \nu^\top(Ax - b)$ 求梯度并令其为零得 $x = -A^\top \nu$，回代即得二次凹函数 $g(\nu)$。}

  \item\label{ex:4.2.2} \textbf{软阈值算子与 $\ell_1$ 次微分手算}：
  (a) 设向量 $z = (2.4,\; -0.6,\; 0.8,\; -3.1)^\top$，阈值 $\tau = 0.8$，手工计算 $\prox_{\tau\norm{\cdot}_1}(z)$；
  (b) 求三维函数 $g(x) = \norm{x}_1$ 在点 $x = (2, 0, -1)^\top$ 处的完整次微分集合 $\partial g(x)$；
  (c) 求合页损失 $h(z) = \max(0, 1 - z)$ 在 $z = 1$ 处的次微分 $\partial h(1)$。
  \exHint{利用公式~\eqref{eq:v4-soft-thresh} 与逐坐标次微分规则；对合页损失直接检验定义~\eqref{eq:v4-subgrad-def}。}
\end{enumerate}

\exTier{B 层}{KKT 完整推导与经典投影}
\begin{enumerate}[label=\textbf{4.2.\arabic*}, leftmargin=3.3em, itemsep=5pt, start=3]
  \item\label{ex:4.2.3} \textbf{最大熵分布与 Softmax 的 KKT 起源}：在概率单纯形上最大化 Shannon 熵：
  \[
    \min_{p \in \R^K}\; \sum_{k=1}^K p_k \ln p_k - \sum_{k=1}^K s_k p_k
    \quad\text{s.t.}\quad \sum_{k=1}^K p_k = 1,\quad p_k \ge 0.
  \]
  利用 KKT 条件证明唯一最优解恰好就是 \textbf{Softmax 分布} $p_k^* = \frac{e^{s_k}}{\sum_{j=1}^K e^{s_j}}$（注意为何不等式乘子自动为零！）。
  \exHint{因为 $\lim_{p_k\to 0^+} \ln p_k = -\infty$，最优解自动落在开集 $p_k > 0$ 内，由互补松弛得所有不等式乘子 $\lambda_k^* = 0$。}

  \item\label{ex:4.2.4} \textbf{注水算法（Water-Filling）的 KKT 求解}：在通信功率分配与矩阵特征值截断中常见如下凸问题（给定常数 $\alpha_i > 0$）：
  \[
    \min_{x \in \R^d}\; -\sum_{i=1}^d \ln(\alpha_i + x_i) \quad\text{s.t.}\quad x_i \ge 0,\quad \sum_{i=1}^d x_i = 1.
  \]
  用 KKT 条件证明最优解形如 $x_i^* = \max(0,\; \nu^{-1} - \alpha_i)$，其中水位线 $\nu^{-1}$ 由 $\sum_{i=1}^d \max(0, \nu^{-1}-\alpha_i) = 1$ 唯一确定。
  \exHint{列出平稳性 $-\frac{1}{\alpha_i + x_i^*} - \lambda_i^* + \nu^* = 0$，结合 $\lambda_i^* \ge 0$ 与 $\lambda_i^* x_i^* = 0$ 消去 $\lambda_i^*$。}
\end{enumerate}

\exTier{C 层}{对偶深入定理与弹性网近端算子}
\begin{enumerate}[label=\textbf{4.2.\arabic*}, leftmargin=3.3em, itemsep=5pt, start=5]
  \item\label{ex:4.2.5} \textbf{Lasso 全零解的精确临界阈值 $\lambda_{\max}$}：对于 Lasso 目标 $F_\lambda(w) = \frac{1}{2N}\norm{Xw - y}_2^2 + \lambda\norm{w}_1$，利用次梯度最优性条件 $0 \in \partial F_\lambda(0)$ 证明：
  \[
    w^* = 0 \iff \lambda \ge \lambda_{\max} \coloneqq \frac{1}{N}\norm{X^\top y}_\infty.
  \]
  结合第 2.4 节实验解释为何表~\ref{tab:v4-ch2-lasso} 在 $\lambda = 3.80$ 时全部系数已精准归零。
  \exHint{计算 $w=0$ 处光滑项的梯度 $-\frac{1}{N}X^\top y$，令其落入 $\lambda\,\partial\norm{0}_1 = [-\lambda, \lambda]^D$。}

  \item\label{ex:4.2.6} \textbf{弹性网（Elastic Net）的闭式近端算子}：弹性网结合了 $\ell_1$ 与 $\ell_2$ 正则项：$g(w) = \lambda_1\norm{w}_1 + \frac{\lambda_2}{2}\norm{w}_2^2$（$\lambda_1, \lambda_2 > 0$）。仿照定理~\ref{thm:v4-soft-thresh} 的证明，推导近端算子 $\prox_{\eta g}(z)$ 的精确逐坐标解析表达式。
  \exHint{先对 $\frac{1}{2}(x_j - z_j)^2 + \frac{\eta\lambda_2}{2}x_j^2 + \eta\lambda_1|x_j|$ 配方整理，将得到缩放后的软阈值 $\frac{1}{1+\eta\lambda_2} S_{\eta\lambda_1}(z_j)$。}
\end{enumerate}

\exTier{D 层}{代码实战与算法加速对照}
\begin{enumerate}[label=\textbf{4.2.\arabic*}, leftmargin=3.3em, itemsep=5pt, start=7]
  \item{\rightskip=0pt plus 1fil\relax \label{ex:4.2.7} \textbf{ISTA vs FISTA 与次梯度下降三法实测}：基于 \texttt{code/vol4\_experiments.py} 中的 \texttt{exp2\_lasso\_ista}（固定 $\lambda = 0.20$），用纯 NumPy 同时实现三种求解器：(1) 衰减步长次梯度下降（$\eta_t = \frac{0.5}{\sqrt{t+1}}$），(2) ISTA，(3) 拓展栏中的 FISTA。对比三者在前 200 步的目标次优度 $F(w^{(t)}) - F^*$ 与第 200 步的稀疏度 \texttt{nnz}，亲眼验证次梯度法为何无法产出精确零解、以及 FISTA 相比 ISTA 的提速幅度。\par}
  \exHint{可先运行 5000 步 FISTA 充当精确基准值 $F^*$，再绘制三者的半对数收敛对比曲线。}

  \item\label{ex:4.2.8} \textbf{单纯形投影算子（Simplex Projection）编程实现}：利用习题~\ref{ex:4.2.4} 类似的 KKT 推导，可以证明向量 $z \in \R^d$ 到概率单纯形 $\Delta^{d-1} = \{p \ge 0 : \sum_i p_i = 1\}$ 的欧氏投影为 $p_i^* = \max(z_i - \tau, 0)$，其中 $\tau$ 由将 $z$ 降序排列后在 $O(d\log d)$ 时间内解析求得。请推导该 $O(d\log d)$ 阈值公式，用纯 NumPy 实现 \texttt{proj\_simplex(z)}，并验证任意随机向量投影后均满足 KKT 全部条件。
  \exHint{记 $z$ 降序排列为 $u_1 \ge \dots \ge u_d$，寻找最大下标 $\rho$ 满足 $u_\rho > \frac{1}{\rho}\big(\sum_{j=1}^\rho u_j - 1\big)$，则 $\tau = \frac{1}{\rho}\big(\sum_{j=1}^\rho u_j - 1\big)$。}
\end{enumerate}

\chapter[统计学习：偏差-方差与泛化界]{统计学习：偏差-方差分解与泛化理论}\label{ch:v4-stat-learning}

\chapEpigraph{在解决给定问题时，尽量不要把解决一个更一般的问题作为中间步骤。}{Vladimir N. Vapnik（《统计学习理论的本质》）}

\begin{chapIntro}{本章导读：优化解得再准，为什么在测试集上仍可能一败涂地？}
前两章赋予了我们强大的最优化武器：给定任何一个凸目标函数，我们已经能精确分析它的收敛率并用加速算法或近端算子求解到机器精度。然而，机器学习与纯数值优化之间横亘着一条决定生死的分界线——\textbf{我们真正关心的不是在已见训练样本上把损失压到 $0$，而是在从未见过的未来样本上依然表现优异（泛化，Generalization）！}

本章将建立统领所有经典与现代机器学习的统计学习框架：
\begin{enumerate}
  \item 辨析\textbf{经验风险}与\textbf{期望泛化风险}，揭穿“零训练误差”的幻觉；
  \item 严格证明回归问题的\textbf{偏差-方差-噪声精确分解定理}，并用 400 次蒙特卡洛实验亲眼看到 U 形泛化曲线的每一笔拆解；
  \item 将第 2 章的拉格朗日对偶直接桥接到\textbf{正则化理论}，从 SVD 谱收缩与 $L_1/L_2$ 几何剖析岭回归与 Lasso 抑制方差的机理；
  \item 掌握\textbf{$K$ 折交叉验证、数据泄露排错与学习曲线诊断法}；
  \item 借助第 2 卷 Hoeffding 不等式严格证明\textbf{有限假设空间 PAC 泛化界与样本复杂度}，并建立无限假设空间的 \textbf{VC 维}直觉，最后远眺现代深度学习中的 Rademacher 复杂度与“双下降”现象。
\end{enumerate}
\end{chapIntro}

% ============================================================
\section{经验风险与期望风险：为什么“零训练误差”往往是灾难？}
% ============================================================

在统计学习理论的数学模型中，我们假设输入-输出对 $(x, y) \in \mathcal{X} \times \mathcal{Y}$ 服从某个\textbf{未知}的联合概率分布 $\mathcal{P}(X, Y)$。学习算法手头仅有一份从 $\mathcal{P}$ 中独立同分布（i.i.d.）抽取的有限训练集 $\mathcal{D} = \{(x_1, y_1), \dots, (x_N, y_N)\}$，并需在预先选定的\textbf{假设空间（Hypothesis Space）} $\mathcal{H}$ 中挑选一个预测函数 $h: \mathcal{X} \to \mathcal{Y}$。

给定非负损失函数 $\ell(h(x), y)$（如分类的 0-1 损失 $\mathbb{I}[h(x)\neq y]$ 或回归的平方损失 $(h(x)-y)^2$），我们必须严格区分两个截然不同的量：

\begin{defBox}
\begin{definition}[期望风险、经验风险与贝叶斯最优风险]\label{def:v4-risks}
\begin{enumerate}
  \item \textbf{期望风险 / 总体风险（Population Risk）}：假设 $h$ 在真实分布 $\mathcal{P}$ 上的期望损失，是我们\emph{真正想最小化却无法直接计算}的黄金标准：
  \[
    R(h) \coloneqq \E_{(x, y) \sim \mathcal{P}}\big[\ell(h(x), y)\big].
  \]
  \item \textbf{经验风险（Empirical Risk）}：假设 $h$ 在有限训练集 $\mathcal{D}$ 上的平均损失，是我们在算法中实际能算出来的代理目标：
  \[
    R_{\mathrm{emp}}(h) \coloneqq \frac{1}{N}\sum_{i=1}^N \ell\big(h(x_i), y_i\big).
  \]
  \item \textbf{贝叶斯最优风险（Bayes Risk）}：在一切可测函数上所能达到的理论最低误差 $R^* \coloneqq \inf_{h} R(h)$。由于真实世界存在固有随机噪声，通常有 $R^* > 0$。
\end{enumerate}
\end{definition}
\end{defBox}

如果我们不加节制地执行\textbf{经验风险最小化（ERM）} $\hat{h}_{\mathrm{ERM}} = \argmin_{h \in \mathcal{H}} R_{\mathrm{emp}}(h)$ 且放任假设空间 $\mathcal{H}$ 无限庞大，就会立刻掉进\textbf{死记硬背陷阱}：构造一个“查表记忆器”——只要输入 $x$ 等于某个训练样本 $x_i$，就原封不动输出 $y_i$，否则一律瞎猜 $0$。它在训练集上的经验风险精确等于 $R_{\mathrm{emp}} = 0$，但在连续空间的新测试点上几乎处处瞎猜，泛化风险接近灾难性的随机猜测！

类似地，在含噪一维回归中，用 $N-1$ 次多项式总能精确穿过全部 $N$ 个含噪数据点（$R_{\mathrm{emp}} = 0$），但在数据点之间剧烈振荡（Runge 现象），在测试集上误差暴涨。由此可见，一个模型的总泛化误差由两股彼此博弈的力量决定：\textbf{模型太简单学不会真规律（欠拟合），还是模型太灵活记住了随机噪声（过拟合）？} 在平方损失下，这一直觉可以升华为一条精确到小数点后任意一位的代数恒等式。

% ============================================================
\section{偏差-方差-噪声精确分解定理与实测验证}
% ============================================================

\subsection{核心定理与完整代数证明}

假设数据由确定性真实规律 $f(x)$ 叠加零均值独立观测噪声生成：
\[
  y = f(x) + \varepsilon, \qquad \E[\varepsilon] = 0, \qquad \Var(\varepsilon) = \E[\varepsilon^2] = \sigma^2.
\]
固定一个学习算法，当我们在随机抽取的训练集 $\mathcal{D}$ 上训练它时，学到的预测函数记为 $\hat{f}_{\mathcal{D}}(x)$——请注意，\textbf{由于 $\mathcal{D}$ 本身是随机抽取的，对任意固定的测试点 $x$，预测值 $\hat{f}_{\mathcal{D}}(x)$ 是一个随机变量！} 记它在所有可能训练集上的平均预测为
\[
  \bar{f}(x) \coloneqq \E_{\mathcal{D}}\big[\hat{f}_{\mathcal{D}}(x)\big],
\]
且测试点噪声 $\varepsilon$ 与训练集 $\mathcal{D}$ 相互独立。下面的定理把期望测试均方误差（MSE）天衣无缝地一分为三。

\begin{thmBox}
\begin{theorem}[偏差-方差-噪声精确分解定理]\label{thm:v4-bias-var}
在测试点 $x$ 处，预测模型 $\hat{f}_{\mathcal{D}}(x)$ 对新观测 $y = f(x) + \varepsilon$ 的期望均方误差精确分解为三项非负之和：
\begin{equation}\label{eq:v4-bv-decomp}
  \E_{\mathcal{D}, \varepsilon}\Big[\big(y - \hat{f}_{\mathcal{D}}(x)\big)^2\Big]
  = \underbrace{\big(\bar{f}(x) - f(x)\big)^2}_{\Bias^2(\hat{f}(x))}
  + \underbrace{\E_{\mathcal{D}}\Big[\big(\hat{f}_{\mathcal{D}}(x) - \bar{f}(x)\big)^2\Big]}_{\Var(\hat{f}(x))}
  + \underbrace{\sigma^2}_{\text{不可约噪声}}.
\end{equation}
\end{theorem}
\end{thmBox}

\begin{proof}
为书写简明，以下记 $f = f(x)$、$\hat{f} = \hat{f}_{\mathcal{D}}(x)$、$\bar{f} = \E_{\mathcal{D}}[\hat{f}]$（注意固定 $x$ 后 $f$ 与 $\bar{f}$ 均为确定常数，只有 $\hat{f}$ 依赖 $\mathcal{D}$、$\varepsilon$ 依赖测试噪声）。

\textbf{第一步：分离不可约噪声 $\varepsilon$。} 代入 $y = f + \varepsilon$ 并展开平方：
\begin{align*}
  \E_{\mathcal{D}, \varepsilon}\big[(y - \hat{f})^2\big]
  &= \E_{\mathcal{D}, \varepsilon}\big[(f - \hat{f} + \varepsilon)^2\big] \\
  &= \E_{\mathcal{D}}\big[(f - \hat{f})^2\big] + 2\,\E_{\mathcal{D}, \varepsilon}\big[(f - \hat{f})\,\varepsilon\big] + \E_{\varepsilon}[\varepsilon^2].
\end{align*}
因为测试噪声 $\varepsilon$ 与训练集 $\mathcal{D}$ 独立，且 $\E[\varepsilon] = 0$，第一个交叉项精确归零：
\[
  \E_{\mathcal{D}, \varepsilon}\big[(f - \hat{f})\,\varepsilon\big] = \E_{\mathcal{D}}[f - \hat{f}] \cdot \underbrace{\E_{\varepsilon}[\varepsilon]}_{= 0} = 0.
\]
而 $\E_{\varepsilon}[\varepsilon^2] = \Var(\varepsilon) + (\E[\varepsilon])^2 = \sigma^2$，因此
\begin{equation}\label{eq:v4-bv-step1}
  \E_{\mathcal{D}, \varepsilon}\big[(y - \hat{f})^2\big] = \E_{\mathcal{D}}\big[(f - \hat{f})^2\big] + \sigma^2.
\end{equation}

\textbf{第二步：插入平均预测 $\bar{f}$ 分离偏差与方差。} 在 $\E_{\mathcal{D}}[(f - \hat{f})^2]$ 中一加一减常数 $\bar{f}$ 并展开：
\begin{align*}
  \E_{\mathcal{D}}\big[(f - \hat{f})^2\big]
  &= \E_{\mathcal{D}}\Big[\big((f - \bar{f}) + (\bar{f} - \hat{f})\big)^2\Big] \\
  &= \E_{\mathcal{D}}\big[(f - \bar{f})^2\big] + 2\,\E_{\mathcal{D}}\big[(f - \bar{f})(\bar{f} - \hat{f})\big] + \E_{\mathcal{D}}\big[(\bar{f} - \hat{f})^2\big].
\end{align*}
观察三项：
\begin{itemize}
  \item 第一项中 $f - \bar{f}$ 是确定常数，直接走出期望得 $(\bar{f} - f)^2 = \Bias^2(\hat{f})$；
  \item 第二项中将常数 $f - \bar{f}$ 提到期望之外，利用 $\bar{f} = \E_{\mathcal{D}}[\hat{f}]$，第二个交叉项再次精确归零：
  \[
    2(f - \bar{f})\,\E_{\mathcal{D}}[\bar{f} - \hat{f}] = 2(f - \bar{f})\big(\bar{f} - \E_{\mathcal{D}}[\hat{f}]\big) = 0;
  \]
  \item 第三项正是随机变量 $\hat{f}$ 的方差定义 $\Var(\hat{f}) = \E_{\mathcal{D}}[(\hat{f} - \bar{f})^2]$。
\end{itemize}
将第二步代回式~\eqref{eq:v4-bv-step1}，即证得精确分解式~\eqref{eq:v4-bv-decomp}！
\end{proof}

定理~\ref{thm:v4-bias-var} 的三项各司其职，给出了打靶般的清晰直觉：
\begin{itemize}
  \item \textbf{平方偏差 $\Bias^2$}：平均预测 $\bar{f}(x)$ 距离真实目标 $f(x)$ 的系统性偏离（瞄准镜歪了）。模型越简单（如用直线拟合正弦波），偏差越大。
  \item \textbf{方差 $\Var$}：换一批训练数据，模型预测 $\hat{f}_{\mathcal{D}}(x)$ 上下剧烈摇摆的离散程度（手抖得厉害）。模型越复杂、参数越多或样本越少，方差越大。
  \item \textbf{不可约噪声 $\sigma^2$}：数据生成过程自带的固有随机性（即贝叶斯风险 $R^*$），任何模型都无法突破这道物理底线。
\end{itemize}

\subsection{实测验证：400 次多项式回归的 U 形曲线精确拆解}

为了让学生不再只看教科书里的“概念示意图”，我们在配套脚本 \texttt{vol4\_experiments.py} 的实验 3 中直接做了 400 轮真实蒙特卡洛仿真：
\begin{itemize}
  \item 真实函数 $f(x) = \sin(2\pi x)$，观测噪声标准差 $\sigma = 0.30$（因此理论不可约噪声严格等于 $\sigma^2 = 0.0900$）。
  \item 每一轮独立抽取 $n = 22$ 个含噪训练样本，用普通最小二乘法分别拟合阶数 $d = 1, 2, \dots, 9$ 的多项式模型，并在 $[0.05, 0.95]$ 的 50 个测试网格点上评估。
  \item 对全部 400 轮预测按定义计算经验 $\Bias^2$、$\Var$、$\sigma^2$ 与总误差，并将结果导出至 \texttt{data/bias\_var.csv}。
\end{itemize}

\begin{lstlisting}[caption={400 轮蒙特卡洛多项式回归偏差-方差分解（节选自 \texttt{vol4\_experiments.py}）},label={lst:v4-ch3-bv}]
for d in range(1, 10):
    preds = np.zeros((400, len(x_te)))
    A_te = np.column_stack([x_te ** k for k in range(d + 1)])
    for k in range(400):
        x_tr = np.clip(x_grid + rng.uniform(-0.02, 0.02, size=22), 0.0, 1.0)
        y_tr = np.sin(2.0 * np.pi * x_tr) + rng.normal(0.0, 0.30, size=22)
        A_tr = np.column_stack([x_tr ** j for j in range(d + 1)])
        w, *_ = np.linalg.lstsq(A_tr, y_tr, rcond=None)
        preds[k] = A_te @ w
    bias2 = float(np.mean((preds.mean(axis=0) - f_te) ** 2))
    var = float(np.mean(np.var(preds, axis=0)))
    total = bias2 + var + 0.30 ** 2
\end{lstlisting}

图~\ref{fig:v4-ch3-bv} 与表~\ref{tab:v4-ch3-bv} 直接读取 \texttt{data/bias\_var.csv}，呈现了极其纯粹的偏差-方差权衡 U 形曲线。

\begin{figure}[htbp]
\centering
\begin{tikzpicture}
\begin{axis}[
  width=0.92\linewidth,
  height=6.7cm,
  xlabel={多项式阶数 $d$（模型复杂度）},
  ylabel={测试均方误差（MSE）分解},
  xmin=0.8, xmax=9.2,
  xtick={1,2,3,4,5,6,7,8,9},
  ymin=-0.007, ymax=0.285,
  scaled y ticks=false,
  yticklabel style={/pgf/number format/fixed, /pgf/number format/precision=2},
  grid=both,
  major grid style={tealLine!45, line width=0.4pt},
  axis line style={slate!80},
  tick label style={font=\footnotesize\sffamily, color=ink},
  label style={font=\small\sffamily, color=ink},
  legend style={
    at={(0.98,0.97)}, anchor=north east,
    fill=white, draw=tealLine, rounded corners=1.2mm,
    font=\footnotesize\sffamily, cells={anchor=west}
  },
  every axis plot/.append style={line width=1.5pt}
]
  % Minimum annotation at deg=3
  \addplot[gold, thick, dotted, forget plot] coordinates {(3,0) (3,0.285)};
  \node[font=\scriptsize\sffamily, color=gold!85!ink, anchor=south west, fill=white, inner sep=1.5pt]
    at (axis cs:3.12, 0.145) {最优点 $d=3$（$\mathrm{Total}=0.1073$）};

  \addplot[color=tealDark, mark=*, mark size=1.8pt]
    table[x=deg, y=total, col sep=comma]{data/bias_var.csv};
  \addlegendentry{总测试 MSE（$\Bias^2 + \Var + \sigma^2$）}

  \addplot[color=ochre, mark=square*, mark size=1.6pt]
    table[x=deg, y=bias2, col sep=comma]{data/bias_var.csv};
  \addlegendentry{平方偏差 $\Bias^2$}

  \addplot[color=navy, mark=triangle*, mark size=1.9pt]
    table[x=deg, y=var, col sep=comma]{data/bias_var.csv};
  \addlegendentry{模型方差 $\Var$}

  \addplot[color=slate, dashed, line width=1.1pt]
    table[x=deg, y=noise, col sep=comma]{data/bias_var.csv};
  \addlegendentry{不可约噪声 $\sigma^2 = 0.0900$}
\end{axis}
\end{tikzpicture}
\caption{含噪 $\sin(2\pi x)$（$\sigma=0.30, n=22$）上 400 轮实测偏差-方差-噪声分解（读取自 \texttt{data/bias\_var.csv}）。当 $d=1, 2$ 时平方偏差高达 $0.159$（严重欠拟合）；在 $d=3$ 处三次多项式首次捕捉正弦波的反对称拐点，$\Bias^2$ 断崖跌至 $0.0035$，总误差触底 $0.1073$；继续增至 $d=9$，$\Bias^2$ 几乎为零但方差从 $0.0079$ 飙升近五倍至 $0.0383$，总误差重新抬头！}\label{fig:v4-ch3-bv}
\end{figure}

\begin{table}[htbp]
\centering
\caption{多项式阶数 $d$ 对应的实测 $\Bias^2$、$\Var$ 与总误差（\texttt{data/bias\_var.csv}）}\label{tab:v4-ch3-bv}
\small
\begin{tabular}{@{}rcccc@{}}
\toprule
\textbf{多项式阶数 $d$} & \textbf{平方偏差 $\Bias^2$} & \textbf{模型方差 $\Var$} & \textbf{不可约噪声 $\sigma^2$} & \textbf{总测试 MSE} \\
\midrule
$1$ & $0.1591$ & $0.0079$ & $0.0900$ & $0.2570$ \\
$2$ & $0.1592$ & $0.0113$ & $0.0900$ & $0.2605$ \\
\textbf{$3$（最优）} & $\mathbf{0.0035}$ & $\mathbf{0.0138}$ & $0.0900$ & $\mathbf{0.1073}$ \\
$5$ & $0.0001$ & $0.0193$ & $0.0900$ & $0.1094$ \\
$7$ & $0.0001$ & $0.0277$ & $0.0900$ & $0.1178$ \\
$9$ & $0.0001$ & $0.0383$ & $0.0900$ & $0.1284$ \\
\bottomrule
\end{tabular}
\end{table}

顺带欣赏一个令人拍案的数学彩蛋：为什么 $d=2$ 相比 $d=1$ 丝毫没有降低偏差（$\Bias^2$ 停在 $0.1592$），总误差反而因方差微增而变差了？因为真函数 $\sin(2\pi x)$ 关于区间中点 $x = 0.5$ 是\textbf{奇对称函数}，任何偶次抛物线 $(x-0.5)^2$ 与它在 $[0, 1]$ 上都正交，对逼近毫无帮助；直到 $d=3$ 引入奇次项 $(x-0.5)^3$，偏差才瞬间雪崩！这提醒我们：\textbf{盲目堆参数未必降低偏差，模型结构与问题对称性的匹配才是关键。}

% ============================================================
\section{正则化的对偶本质与岭回归 vs Lasso 的权衡}
% ============================================================

当模型特征维数较高（高方差）时，我们不必粗暴地删减特征或手动限制多项式阶数，而是通过\textbf{正则化（Regularization）}向优化目标注入先验平滑性，用极小的偏差增量换取方差的大幅跳水。

\subsection{从第 2 章拉格朗日对偶看“约束形式”与“惩罚形式”的等价性}

考虑带范数球约束的经验风险最小化问题（$p \in \{1, 2\}$，约束半径 $C > 0$）：
\begin{equation}\label{eq:v4-reg-constrained}
  \min_{w \in \R^D}\; R_{\mathrm{emp}}(w) \qquad\text{s.t.}\qquad \norm{w}_p^p \le C.
\end{equation}
若 $R_{\mathrm{emp}}(w)$ 为凸函数，且 $C > 0$（从而 $w = 0$ 满足 $\norm{0}_p^p = 0 < C$，严格满足第 2 章定理~\ref{thm:v4-slater} 的 Slater 条件！），则强对偶成立。写出其拉格朗日函数：
\[
  \mathcal{L}(w, \lambda) = R_{\mathrm{emp}}(w) + \lambda\big(\norm{w}_p^p - C\big), \qquad \lambda \ge 0.
\]
固定对偶最优乘子 $\lambda^* \ge 0$ 后，由 KKT 定理~\ref{thm:v4-kkt-thm}，原始最优解 $w^*$ 必同时最小化无约束的\textbf{惩罚形式目标（Penalized ERM）}：
\begin{equation}\label{eq:v4-reg-penalized}
  \min_{w \in \R^D}\; R_{\mathrm{emp}}(w) + \lambda^*\norm{w}_p^p.
\end{equation}
这在理论上彻底打通了第 2 章与第 3 章：\textbf{增大惩罚系数 $\lambda$，在几何上严格等价于把图~\ref{fig:v4-l1-prox}(a) 中的 $\ell_1$ 菱形或 $\ell_2$ 圆球半径 $C$ 向原点收紧！}

\subsection{岭回归（Ridge, $L_2$）的 SVD 谱收缩与显式降方差证明}

对最小二乘回归加上 $L_2$ 正则化，即得\textbf{岭回归（Ridge Regression）}：
\[
  \hat{w}_{\mathrm{ridge}}(\lambda) = \argmin_{w \in \R^D} \Big\{\norm{Xw - y}_2^2 + \lambda\norm{w}_2^2\Big\} = (X^\top X + \lambda I_D)^{-1} X^\top y.
\]
对设计矩阵 $X \in \R^{N\times D}$ 做瘦奇异值分解 $X = U \Sigma V^\top$（其中 $U^\top U = I_r,\; V^\top V = I_r,\; \Sigma = \diag(\sigma_1, \dots, \sigma_r)$，设 $X$ 列满秩 $r = D$）。代入 $X^\top X = V \Sigma^2 V^\top$，可将拟合值与估计器在奇异基 $V = [v_1, \dots, v_D]$ 下彻底对角化：
\begin{equation}\label{eq:v4-ridge-svd}
  \hat{w}_{\mathrm{ridge}}(\lambda) = \sum_{j=1}^D \underbrace{\frac{\sigma_j^2}{\sigma_j^2 + \lambda}}_{\text{谱收缩因子}\;\in (0, 1)} \frac{u_j^\top y}{\sigma_j}\, v_j.
\end{equation}
设真实数据满足 $y = X w_{\text{true}} + \varepsilon$（$\E[\varepsilon] = 0,\; \Cov(\varepsilon) = \sigma^2 I_N$）。直接由式~\eqref{eq:v4-ridge-svd} 算出估计误差的精确方差与偏差（对比 $\lambda = 0$ 时的 OLS）：
\begin{align}
  \text{OLS 总方差：}&\quad \tr\big(\Cov(\hat{w}_{\mathrm{OLS}})\big) = \sigma^2 \sum_{j=1}^D \frac{1}{\sigma_j^2}, \label{eq:v4-ols-var} \\
  \text{Ridge 总方差：}&\quad \tr\big(\Cov(\hat{w}_{\mathrm{ridge}})\big) = \sigma^2 \sum_{j=1}^D \frac{\sigma_j^2}{(\sigma_j^2 + \lambda)^2}, \label{eq:v4-ridge-var} \\
  \text{Ridge 平方偏差：}&\quad \norm{\E[\hat{w}_{\mathrm{ridge}}] - w_{\text{true}}}_2^2 = \sum_{j=1}^D \left(\frac{\lambda}{\sigma_j^2 + \lambda}\right)^2 (v_j^\top w_{\text{true}})^2. \label{eq:v4-ridge-bias}
\end{align}
看清这组公式的惊人威力了吗？当特征之间存在严重共线性时，某些奇异值 $\sigma_j \approx 0$：在 OLS 式~\eqref{eq:v4-ols-var} 中，$\frac{1}{\sigma_j^2}$ 会爆炸式放大噪声，导致方差失控；而岭回归在分母加上 $\lambda > 0$ 后，收缩因子 $\frac{\sigma_j^2}{\sigma_j^2+\lambda}$ 自动对微弱奇异方向施以重拳压制（方差大幅缩水），同时对信号强劲的主奇异方向（$\sigma_j^2 \gg \lambda$）几乎保持原样（偏差极小）！

\begin{table}[htbp]
\centering
\caption{岭回归（$L_2$）与 Lasso（$L_1$）正则化的多维对比}\label{tab:v4-ridge-lasso}
\small
\begin{tabularx}{\linewidth}{@{}lXX@{}}
\toprule
\textbf{维度} & \textbf{岭回归（Ridge, $\lambda\norm{w}_2^2$）} & \textbf{Lasso（$\lambda\norm{w}_1$）} \\
\midrule
约束域几何 & 光滑超球面（图~\ref{fig:v4-l1-prox}(a) 虚线圆） & 带坐标轴尖角的超菱形（图~\ref{fig:v4-l1-prox}(a) 实线菱形） \\
稀疏性 & 权重平滑收缩趋近零，但几乎不精确为零 & 自动产出精确零解，兼具\textbf{特征选择}功能 \\
最优解形式 & 闭式解 $(X^\top X + \lambda I)^{-1}X^\top y$，严格强凸 & 无一般闭式解，用第 2 章近端梯度 ISTA/FISTA 求解 \\
相关特征组 & 倾向于将权重均匀分摊给高度相关特征 & 倾向于挑出一个代表特征，把其余相关特征归零 \\
适用场景 & 多数特征均有贡献、存在多重共线性 & 高维稀疏信号、需精简特征提升可解释性 \\
\bottomrule
\end{tabularx}
\end{table}

% ============================================================
\section{模型选择、$K$ 折交叉验证与学习曲线诊断}
% ============================================================

正则化系数 $\lambda$、多项式阶数 $d$ 或第 4--5 章的核带宽 $\gamma$ 与树深度，统称为\textbf{超参数（Hyperparameters）}——它们无法靠最小化训练集经验风险来选（因为训练集永远偏爱 $\lambda = 0$ 或无限深的树！）。科学地选择模型与诊断问题，是工程实践中最值钱的基本功。

\subsection{三集合划分与 $K$ 折交叉验证}

\begin{enumerate}
  \item \textbf{严格三集合划分（Train / Validation / Test）}：
  \begin{itemize}
    \item \textbf{训练集（Training Set）}：仅用来拟合模型参数 $w$；
    \item \textbf{验证集（Validation Set）}：用来评估候选超参数（如不同 $\lambda$ 或不同模型族）并做最终模型选择；
    \item \textbf{测试集（Test Set）}：锁进“保险箱”，在所有调参和模型定型\textbf{彻底结束}之后，只评估一次，以获得对真实泛化误差 $R(h)$ 的无偏估计。
  \end{itemize}
  \item \textbf{$K$ 折交叉验证（$K$-Fold Cross-Validation）}：当总样本量不大（例如几百到几千）时，单次划出验证集既浪费训练数据，又可能因划分运气带来高方差。$K$ 折交叉验证将全部非测试数据随机均分为 $K$ 份（工程常取 $K = 5$ 或 $10$）：轮流取第 $k$ 份作验证、其余 $K-1$ 份作训练，最后取 $K$ 轮验证误差的平均值作为该超参数的评分。其极限情形 $K = N$ 称为\textbf{留一法（LOOCV）}——偏差极小，但计算代价为 $N$ 倍训练（线性回归有闭式捷径）且各折训练集高度重叠导致估计方差偏大，故实践中 $K = 5$ 或 $10$ 是偏差与方差的最佳折中。
\end{enumerate}

\begin{failBox}{机器学习最隐蔽的“自欺欺人”：数据泄露（Data Leakage）三大典型案例}
哪怕写了严格的 $K$ 折交叉验证，只要下面三件事中做错了一件，你的验证集分数就会虚高得离谱、上线后立刻崩盘：
\begin{enumerate}
  \item \textbf{划分之前先做全局标准化或缺失值填充}：用全数据集（含验证集）算出特征均值 $\mu_{\text{all}}$ 与标准差 $\sigma_{\text{all}}$ 去归一化，把未来测试分布的信息提前泄漏进了训练过程！正确做法：\textbf{只在当前折的训练集上 \texttt{fit} 均值方差，再 \texttt{transform} 验证集。}
  \item \textbf{交叉验证之外先做全量特征筛选}：在 10{,}000 个纯噪声特征里，先用全部样本挑出与标签 $y$ 相关性最高的 10 个特征，再跑 5 折 CV——哪怕 $y$ 完全是抛硬币随机生成的，CV 准确率也能骗出 $85\%$ 以上！特征选择、PCA 降维都必须严密封装在每一折的训练 Pipeline 内部。
  \item \textbf{时间序列随机打乱划分}：对股票、气温或用户行为日志随机做 $K$ 折划分，等于用“明天的数据”去预测“今天”，属于严重的时间穿越泄露，必须采用按时间向前的\textbf{滚动时间窗验证}。
\end{enumerate}
\end{failBox}

\subsection{学习曲线：一眼确诊“高偏差”还是“高方差”}

面对一个验证误差不理想的模型，新手常盲目尝试“多采集数据”或“换更大模型”，结果事倍功半。绘制\textbf{学习曲线（Learning Curves）}——即随着训练样本量 $N$ 增大，训练误差与验证误差的演变轨迹——能像医院的 X 光片一样精准开方（见图~\ref{fig:v4-learning-curves}）：

\begin{figure}[htbp]
\centering
\begin{tikzpicture}[scale=0.95, >=Stealth]
  % ===== Left Panel: High Bias (Underfitting) =====
  \begin{scope}[shift={(0,0)}]
    \node[font=\small\sffamily\bfseries, color=ochre] at (2.85, 4.35) {(a) 高偏差（欠拟合）：双双高位收敛};
    \draw[->, slate!80] (0,0) -- (5.8,0) node[right, font=\footnotesize\sffamily] {$N$};
    \draw[->, slate!80] (0,0) -- (0,3.85) node[left, font=\footnotesize\sffamily, yshift=-2pt] {误差};

    % Desired / Bayes baseline
    \draw[slate!70, dashed, thick] (0, 0.75) -- (5.5, 0.75);
    \node[right, font=\scriptsize\sffamily, color=slate] at (3.25, 0.48) {期望 / 贝叶斯底线 $R^*$};

    % Validation curve: starts high, plateaus high at ~2.55
    \draw[ochre, line width=1.6pt]
      (0.4, 3.45) .. controls (1.4, 2.75) and (2.6, 2.60) .. (5.4, 2.55);
    \node[ochre, font=\scriptsize\sffamily, above] at (4.35, 2.60) {验证误差};

    % Training curve: starts low, rises quickly to plateau near ~2.24
    \draw[tealDark, line width=1.6pt]
      (0.4, 0.95) .. controls (1.3, 1.95) and (2.4, 2.18) .. (5.4, 2.24);
    \node[tealDark, font=\scriptsize\sffamily, below] at (2.45, 2.10) {训练误差};

    % High bias gap arrow
    \draw[<->, plum, thick] (5.05, 0.75) -- (5.05, 2.24)
      node[midway, left, font=\scriptsize\sffamily, align=right, color=plum] {巨大偏差\\（加数据无用！）};
    \draw[<->, gold!90!ink, semithick] (5.38, 2.24) -- (5.38, 2.55);
  \end{scope}

  % ===== Right Panel: High Variance (Overfitting) =====
  \begin{scope}[shift={(7.45,0)}]
    \node[font=\small\sffamily\bfseries, color=navy] at (2.85, 4.35) {(b) 高方差（过拟合）：训练极低、间隙巨大};
    \draw[->, slate!80] (0,0) -- (5.8,0) node[right, font=\footnotesize\sffamily] {$N$};
    \draw[->, slate!80] (0,0) -- (0,3.85) node[left, font=\footnotesize\sffamily, yshift=-2pt] {误差};

    % Desired / Bayes baseline
    \draw[slate!70, dashed, thick] (0, 0.75) -- (5.5, 0.75);
    \node[right, font=\scriptsize\sffamily, color=slate] at (3.25, 0.48) {期望 / 贝叶斯底线 $R^*$};

    % Validation curve: starts very high, steadily drops with N
    \draw[ochre, line width=1.6pt]
      (0.4, 3.60) .. controls (1.5, 2.85) and (3.2, 2.25) .. (5.4, 1.85);
    \node[ochre, font=\scriptsize\sffamily, above] at (4.45, 2.02) {验证误差};

    % Training curve: stays very low near 0.90
    \draw[tealDark, line width=1.6pt]
      (0.4, 0.18) .. controls (1.6, 0.55) and (3.2, 0.78) .. (5.4, 0.90);
    \node[tealDark, font=\scriptsize\sffamily, above] at (2.05, 0.72) {训练误差};

    % Large generalization gap arrow
    \draw[<->, navy, thick] (5.05, 0.90) -- (5.05, 1.85)
      node[midway, left, font=\scriptsize\sffamily, align=right, color=navy] {大泛化间隙\\（加数据/加正则！）};
  \end{scope}
\end{tikzpicture}
\caption{学习曲线的临床诊断图谱：(a) 若训练误差与验证误差双双收敛在远离底线 $R^*$ 的高位且彼此很近，说明模型\textbf{高偏差（欠拟合）}，此时花大价钱增加训练样本几乎毫无帮助，应增大模型容量、构造非线性特征或减小正则化 $\lambda$；(b) 若训练误差极低而验证误差高悬、两者存在宽阔的泛化间隙，说明模型\textbf{高方差（过拟合）}，应增加样本量 $N$、加大正则化 $\lambda$、做特征选择或使用第 5 章的 Bagging 集成。}\label{fig:v4-learning-curves}
\end{figure}

% ============================================================
\section{PAC 学习理论：有限假设空间的泛化界与样本复杂度}
% ============================================================

偏差-方差分解优雅地解剖了固定样本量下的回归误差，但它没有回答一个更根本的概率问题：\textbf{如果我有 $N$ 个 i.i.d. 样本，如何以 $1-\delta$ 的高置信度，用数学定理严格保证训练误差 $R_{\mathrm{emp}}(\hat{h})$ 与未知的真实风险 $R(\hat{h})$ 之间的差距不超过 $\varepsilon$？}

这一框架由 Leslie Valiant（1984）与 Vladimir Vapnik、Alexey Chervonenkis 奠基，称为 \textbf{PAC（Probably Approximately Correct，概率近似正确）学习理论}。

\begin{recallBox}{第 2 卷浓度不等式回顾：Hoeffding 不等式}
设 $Z_1, \dots, Z_N$ 为独立同分布随机变量，且几乎必然取值在有界区间 $[0, 1]$ 内，均值 $\mu = \E[Z_i]$，样本均值 $\bar{Z} = \frac{1}{N}\sum_{i=1}^N Z_i$。则对任意 $\varepsilon > 0$，双边偏差概率呈指数级衰减：
\begin{equation}\label{eq:v4-hoeffding}
  \Prob\!\big(|\bar{Z} - \mu| \ge \varepsilon\big) \le 2\exp(-2N\varepsilon^2).
\end{equation}
\end{recallBox}

\textbf{警惕一个致命的初学者陷阱！} 设损失函数取值在 $[0, 1]$（例如分类的 0-1 损失）。对任意一个\emph{在看数据之前就预先固定好的}假设 $h \in \mathcal{H}$，令 $Z_i = \ell(h(x_i), y_i) \in [0, 1]$，则 $\bar{Z} = R_{\mathrm{emp}}(h),\; \mu = R(h)$，直接套用 Hoeffding 不等式立得 $\Prob(|R(h) - R_{\mathrm{emp}}(h)| \ge \varepsilon) \le 2e^{-2N\varepsilon^2}$。

\textbf{但是，我们绝不能把这个不等式直接套在 ERM 选出来的 $\hat{h} = \argmin_{h\in\mathcal{H}} R_{\mathrm{emp}}(h)$ 头上！} 为什么？因为 $\hat{h}$ 是“看过训练数据之后挑出来的”，它与样本 $\mathcal{D}$ 高度相关，导致 $\ell(\hat{h}(x_i), y_i)$ 不再是独立的！如果假设空间里塞了一百万个随机函数，即使每个函数单独看都很靠谱，总会有某个倒霉（或“幸运”）的假设恰好撞大运把这 $N$ 个样本全猜对了！

破解之道在于\textbf{一致收敛（Uniform Convergence）}：我们要求假设空间里的\emph{所有}假设 $h \in \mathcal{H}$ 同时不偏离真实风险，再借助最基础的\textbf{并集界（Union Bound）} $\Prob(\bigcup_k A_k) \le \sum_k \Prob(A_k)$ 一次锁死全部假设！

\begin{thmBox}
\begin{theorem}[有限假设空间的一致收敛与 PAC 泛化界]\label{thm:v4-pac-finite}
设假设空间 $\mathcal{H}$ 有限（基数 $|\mathcal{H}| < \infty$），损失函数取值于 $[0, 1]$，训练集 $\mathcal{D}$ 包含 $N$ 个 i.i.d. 样本。则对任意置信水平参数 $\delta \in (0, 1)$，至少以 $1 - \delta$ 的概率，\textbf{同时对所有 $h \in \mathcal{H}$} 成立：
\begin{equation}\label{eq:v4-pac-bound}
  \sup_{h \in \mathcal{H}} \big|R(h) - R_{\mathrm{emp}}(h)\big| \le \sqrt{\frac{\ln(2|\mathcal{H}| / \delta)}{2N}}.
\end{equation}
记 $\hat{h} = \argmin_{h\in\mathcal{H}} R_{\mathrm{emp}}(h)$ 为经验风险最小化器，$h^* = \argmin_{h\in\mathcal{H}} R(h)$ 为假设空间内的最优假设。则在相同的 $1-\delta$ 高概率事件上，超额风险满足：
\begin{equation}\label{eq:v4-erm-excess}
  R(\hat{h}) - R(h^*) \le 2\sqrt{\frac{\ln(2|\mathcal{H}| / \delta)}{2N}}.
\end{equation}
等价地，欲以至少 $1-\delta$ 的概率保证一致泛化间隙不超过 $\varepsilon$，所需\textbf{样本复杂度（Sample Complexity）}为：
\begin{equation}\label{eq:v4-sample-comp}
  N \ge \frac{1}{2\varepsilon^2}\ln\frac{2|\mathcal{H}|}{\delta}.
\end{equation}
\end{theorem}
\end{thmBox}

\begin{proof}
\textbf{第一步：对所有假设应用并集界。} 对每个固定 $h \in \mathcal{H}$，定义“坏事件” $E_h = \{|R(h) - R_{\mathrm{emp}}(h)| \ge \varepsilon\}$。由 Hoeffding 不等式~\eqref{eq:v4-hoeffding}，对每一个 $h$ 均有 $\Prob(E_h) \le 2\exp(-2N\varepsilon^2)$。考察“至少有一个假设发生大偏差”的并集事件：
\begin{align*}
  \Prob\!\left(\sup_{h \in \mathcal{H}} |R(h) - R_{\mathrm{emp}}(h)| \ge \varepsilon\right)
  &= \Prob\!\left(\bigcup_{h \in \mathcal{H}} E_h\right) \\
  &\le \sum_{h \in \mathcal{H}} \Prob(E_h) \quad\text{（并集界）} \\
  &\le \sum_{h \in \mathcal{H}} 2\exp(-2N\varepsilon^2) = 2|\mathcal{H}|\exp(-2N\varepsilon^2).
\end{align*}

\textbf{第二步：令失败概率上限等于 $\delta$ 并反解 $\varepsilon$ 与 $N$。} 令 $2|\mathcal{H}|\exp(-2N\varepsilon^2) = \delta$，两边取对数：
\[
  2N\varepsilon^2 = \ln\frac{2|\mathcal{H}|}{\delta}
  \iff \varepsilon = \sqrt{\frac{\ln(2|\mathcal{H}|/\delta)}{2N}}
  \iff N = \frac{1}{2\varepsilon^2}\ln\frac{2|\mathcal{H}|}{\delta}.
\]
这就同时证明了一致泛化界~\eqref{eq:v4-pac-bound} 与样本复杂度~\eqref{eq:v4-sample-comp}！

\textbf{第三步：导出 ERM 超额风险界~\eqref{eq:v4-erm-excess}。} 在上述概率 $\ge 1-\delta$ 的好事件上，对 $\hat{h}$ 与 $h^*$ 均有一致逼近 $|R(h) - R_{\mathrm{emp}}(h)| \le \varepsilon$，叠加 $\hat{h}$ 的经验最优性 $R_{\mathrm{emp}}(\hat{h}) \le R_{\mathrm{emp}}(h^*)$：
\[
  R(\hat{h}) \le R_{\mathrm{emp}}(\hat{h}) + \varepsilon \le R_{\mathrm{emp}}(h^*) + \varepsilon \le \big(R(h^*) + \varepsilon\big) + \varepsilon = R(h^*) + 2\varepsilon.
\]
代入 $\varepsilon$ 的表达式即证得式~\eqref{eq:v4-erm-excess}！
\end{proof}

\begin{example}[样本复杂度的数量级直觉]\label{ex:v4-ch3-pac-num}
假设我们在一个包含 $|\mathcal{H}| = 10^4$ 个候选分类器的有限空间里搜索，要求以 $95\%$ 的置信度（$\delta = 0.05$）保证一致泛化间隙不超过 $\varepsilon = 0.05$。代入式~\eqref{eq:v4-sample-comp}：
\[
  N \ge \frac{1}{2 \times 0.05^2}\ln\frac{2 \times 10^4}{0.05} = 200 \times \ln(4 \times 10^5) \approx 200 \times 12.90 = 2580.
\]
请注意公式中最美妙的地方：\textbf{样本复杂度对假设空间规模 $|\mathcal{H}|$ 仅是对数依赖 $\ln|\mathcal{H}|$！} 即使把假设数量暴增一万倍到 $10^8$，所需样本量也只从 $2580$ 增加到 $200 \times \ln(4\times 10^9) \approx 4422$——还不到两倍！另一方面，精度要求 $\varepsilon$ 是反平方依赖 $O(1/\varepsilon^2)$，这意味着泛化间隙随样本量按 $O(1/\sqrt{N})$ 衰减。
\end{example}

% ============================================================
\section{超越有限集合：VC 维的几何直觉与无限假设空间}
% ============================================================

定理~\ref{thm:v4-pac-finite} 有一个明显的短板：即便是最简单的一维阈值分类器 $\mathcal{H} = \{h_\theta(x) = \mathbb{I}[x \ge \theta] : \theta \in \R\}$，实数参数 $\theta$ 有不可数无穷多个，$|\mathcal{H}| = \infty$，式~\eqref{eq:v4-pac-bound} 的分子瞬间爆炸成无意义的 $+\infty$！

然而冷静思考一下：给定数轴上 $N$ 个固定样本点 $x_1 < x_2 < \dots < x_N$，虽然阈值 $\theta$ 有无穷多种取法，但无论怎么移动 $\theta$，它在这 $N$ 个点上所能切出来的不同二分类标记结果\textbf{最多只有 $N + 1$ 种}（而不是 $2^N$ 种，更不是 $\infty$ 种）！这正是 Vapnik 与 Chervonenkis 在 1971 年的天才洞察：\textbf{衡量无限假设空间“真实容量”的尺子，不是参数有多少个，而是它能在有限点集上实现多少种不同的二分类划分！}

\begin{defBox}
\begin{definition}[增长函数、打散与 VC 维]\label{def:v4-vc}
设 $\mathcal{H}$ 为二分类假设空间（$h: \mathcal{X} \to \{-1, +1\}$）：
\begin{enumerate}
  \item \textbf{增长函数（Growth Function）} $\Pi_{\mathcal{H}}(N)$：在所有大小为 $N$ 的点集 $\{x_1, \dots, x_N\} \subseteq \mathcal{X}$ 上，$\mathcal{H}$ 所能产生的不同标记向量 $\big(h(x_1), \dots, h(x_N)\big) \in \{-1, +1\}^N$ 的最大个数。显然 $\Pi_{\mathcal{H}}(N) \le 2^N$。
  \item \textbf{打散（Shatter）}：若存在 $N$ 个点使得 $\mathcal{H}$ 能实现它们全部 $2^N$ 种可能的 $\pm 1$ 标记组合，则称 $\mathcal{H}$ 能够\textbf{打散}这 $N$ 个点。
  \item \textbf{VC 维（Vapnik-Chervonenkis Dimension）} $d_{\mathrm{VC}} = \mathrm{VC}(\mathcal{H})$：$\mathcal{H}$ 能够打散的\textbf{最大点集规模}（即满足 $\Pi_{\mathcal{H}}(N) = 2^N$ 的最大整数 $N$）。
\end{enumerate}
\end{definition}
\end{defBox}

\begin{example}[二维平面线性分类器的 VC 维 $= 3$]\label{ex:v4-ch3-vc2d}
考虑二维平面上的线性分类器 $\mathcal{H}_{\mathrm{lin}}^{(2)} = \{x \mapsto \sign(w_1 x_1 + w_2 x_2 + b)\}$：
\begin{itemize}
  \item \textbf{能打散 3 个点（$\mathrm{VC} \ge 3$）}：任取不共线的 3 个点构成三角形（图~\ref{fig:v4-vc-shatter}(a)）。全部 $2^3 = 8$ 种标记可分为两类：3 点同色（2 种，把直线放三角形外侧即可）与 2 同 1 异（6 种，用直线切开单独那个顶点与另外两点即可）。8 种组合无一遗漏！
  \item \textbf{无法打散任何 4 个点（$\mathrm{VC} < 4$）}：任取平面上 4 个点，由平面几何 Radon 定理，其凸包只可能是两种情形：(i) 4 点构成凸四边形——此时把对角两点标为 $+$、另两对角点标为 $-$（著名的 \textbf{XOR 异或构型}，图~\ref{fig:v4-vc-shatter}(b)），两条对角线必然相交，任何一条直线都无法将两组对角点分开！(ii) 其中 1 点落在另 3 点构成的三角形内部或边上——此时把外围 3 顶点标为 $+$、内部点标为 $-$，直线同样无解。因此没有任何 4 个点能被平面直线打散！
\end{itemize}
故 $\mathrm{VC}(\mathcal{H}_{\mathrm{lin}}^{(2)}) = 3$。推广至一般 $d$ 维空间，线性超平面分类器的 VC 维精确等于其自由参数个数：$\mathrm{VC}(\mathcal{H}_{\mathrm{lin}}^{(d)}) = d + 1$。
\end{example}

\begin{figure}[htbp]
\centering
\begin{tikzpicture}[scale=1.0, >=Stealth]
  % ===== Left Panel: 3 Points Shattered =====
  \begin{scope}[shift={(0,0)}]
    \node[font=\small\sffamily\bfseries, color=tealDark] at (1.5, 3.35) {(a) 平面 3 点可被直线打散（全 8 种可行）};
    \draw[tealLine!80, rounded corners=2mm, fill=tealLight!45] (-0.4,-0.63) rectangle (3.4,2.85);

    % Triangle vertices: A=(0.6,0.5), B=(2.5,0.6), C=(1.5,2.2)
    \fill[teal] (0.6, 0.5) circle (4pt);
    \node[white, font=\tiny\bfseries] at (0.6, 0.5) {$+$};
    \fill[teal] (2.5, 0.6) circle (4pt);
    \node[white, font=\tiny\bfseries] at (2.5, 0.6) {$+$};
    \fill[ochre] (1.5, 2.2) circle (4pt);
    \node[white, font=\tiny\bfseries] at (1.5, 2.2) {$-$};

    % Separating line cleanly separating top point
    \draw[gold, very thick] (0.05, 1.45) -- (3.05, 1.55);
    \node[font=\scriptsize\sffamily, color=tealDark, align=center] at (1.5, -0.10)
      {任意 $2^3=8$ 种正负标记\\均可被某条直线一刀切开};
  \end{scope}

  % ===== Right Panel: 4 Points XOR Fails =====
  \begin{scope}[shift={(6.8,0)}]
    \node[font=\small\sffamily\bfseries, color=ochre] at (1.6, 3.35) {(b) 平面 4 点无法打散（XOR 对角不可分）};
    \draw[ochre!35, rounded corners=2mm, fill=ochreLight!55] (-0.4,-0.63) rectangle (3.6,2.85);

    % Diagonals crossing
    \draw[teal!50, dashed, thick] (0.55, 0.55) -- (2.65, 2.25);
    \draw[ochre!50, dashed, thick] (0.55, 2.25) -- (2.65, 0.55);

    % 4 corners in XOR configuration
    \fill[teal] (0.55, 0.55) circle (4.2pt);
    \node[white, font=\tiny\bfseries] at (0.55, 0.55) {$+$};
    \fill[teal] (2.65, 2.25) circle (4.2pt);
    \node[white, font=\tiny\bfseries] at (2.65, 2.25) {$+$};

    \fill[ochre] (0.55, 2.25) circle (4.2pt);
    \node[white, font=\tiny\bfseries] at (0.55, 2.25) {$-$};
    \fill[ochre] (2.65, 0.55) circle (4.2pt);
    \node[white, font=\tiny\bfseries] at (2.65, 0.55) {$-$};

    % Cross mark at intersection
    \node[font=\normalsize\bfseries, color=ochre] at (1.6, 1.40) {$\times$};
    \node[font=\scriptsize\sffamily, color=ochre, align=center] at (1.6, -0.10)
      {凸包对角线必相交\\不存在线性超平面可分离 XOR};
  \end{scope}
\end{tikzpicture}
\caption{二维线性分类器 VC 维的直观几何证明：(a) 三个不共线点的 8 种二分类标记都能被直线分离；(b) 四个点只要呈 XOR 交错标记（或一点落入三点三角形内部），便无法线性分离，故 $\mathrm{VC}(\text{平面直线}) = 3$。注意：第 4 章的核技巧（RBF 核）正是把这种低维不可分的构型映射到高维空间变成线性可分！}\label{fig:v4-vc-shatter}
\end{figure}

有了 VC 维，Sauer-Shelah 引理证明了一个令人震撼的相变定理：一旦样本数 $N$ 超过 $d_{\mathrm{VC}}$，增长函数 $\Pi_{\mathcal{H}}(N)$ 不再按指数 $2^N$ 狂飙，而是被多项式 $\left(\frac{eN}{d_{\mathrm{VC}}}\right)^{d_{\mathrm{VC}}}$ 死死压住！于是用 $\Pi_{\mathcal{H}}(2N)$ 替代定理~\ref{thm:v4-pac-finite} 中的 $|\mathcal{H}|$，就得到了统计学习理论的皇冠明珠——\textbf{VC 泛化界}：以至少 $1-\delta$ 的概率，对所有 $h \in \mathcal{H}$ 同时成立
\begin{equation}\label{eq:v4-vc-bound}
  R(h) \le R_{\mathrm{emp}}(h) + O\!\left(\sqrt{\frac{d_{\mathrm{VC}}\ln(N / d_{\mathrm{VC}}) + \ln(1/\delta)}{N}}\right).
\end{equation}
式~\eqref{eq:v4-vc-bound} 给出了\textbf{结构风险最小化（SRM）}的理论依据：当 $N \gg d_{\mathrm{VC}}$ 时，泛化间隙以 $\tilde{O}(\sqrt{d_{\mathrm{VC}}/N})$ 趋于零；反之，如果盲目扩大模型容量使 $d_{\mathrm{VC}}$ 逼近甚至超过样本数 $N$，泛化保证就会彻底崩塌。这也正是 Vapnik 卷首语的深意——面对特定判别任务时，直接控制判别边界的有效复杂度（例如第 4 章的大间隔 SVM），往往远胜于先去估计整个高维概率密度！

\begin{extBox}{从经典到现代：Rademacher 复杂度与深度学习的“双下降”现象}
VC 维完全只看假设空间 $\mathcal{H}$ 的最坏情况几何容量，完全不关心数据分布本身。但在实践中，即使参数维数高达百万（如第 4 章的无穷维 RBF 核 SVM 或第 5 卷的深度神经网络），只要数据本身结构良好、分类间隔够大，模型依然泛化极佳！这催生了两项通往现代理论的关键突破：
\begin{enumerate}
  \item \textbf{Rademacher 复杂度（Rademacher Complexity）}：给定样本 $S = \{x_1, \dots, x_N\}$，独立抛 $N$ 枚公平硬币生成随机噪声标签 $\sigma_i \in \{-1, +1\}$，定义经验 Rademacher 复杂度
  \[
    \widehat{\mathfrak{R}}_S(\mathcal{H}) \coloneqq \E_{\sigma}\!\left[\sup_{h \in \mathcal{H}} \frac{1}{N}\sum_{i=1}^N \sigma_i h(x_i)\right].
  \]
  它直观衡量假设空间 $\mathcal{H}$ “强行拟合纯随机硬币噪声”的能力，并且\textbf{依赖于具体数据分布}！利用它可直接推导出\textbf{与输入维度 $d$ 无关、仅与分类间隔 $\gamma$ 相关的 SVM 泛化界}，完美解释为何无穷维 RBF 核不会瞬间过拟合。
  \item \textbf{现代双下降现象（Double Descent，直通第 5 卷）}：经典偏差-方差 U 形曲线（图~\ref{fig:v4-ch3-bv}）在模型参数数 $P$ 逼近样本数 $N$（插值阈值 $P \approx N$）时确实会冲上尖峰；但 Belkin 等（2019）发现，当参数规模继续暴增进入\textbf{严重过参数化区间（$P \gg N$）}时，虽然训练误差保持为 $0$，测试误差竟然会\textbf{第二次下降}甚至低于经典 U 形的谷底！其背后的奥秘在于：在无数个都能零误差插值训练集的解中，梯度下降 / SGD 的\textbf{隐式正则化（Implicit Bias）}自动偏好于最小范数（最平滑）的插值解。我们将在第 5 卷深度学习中重新审视这一激动人心的前沿图景。
\end{enumerate}
\end{extBox}

% ============================================================
\section*{本章分层习题}
\addcontentsline{toc}{section}{本章分层习题}
% ============================================================

\exTier{A 层}{核心概念与数值验算}
\begin{enumerate}[label=\textbf{4.3.\arabic*}, leftmargin=3.3em, itemsep=5pt]
  \item\label{ex:4.3.1} \textbf{常值估计器的偏差与方差手算}：设真实分布满足 $y = f(x) + \varepsilon$（$\E[\varepsilon]=0,\; \Var(\varepsilon)=\sigma^2$），考虑极简模型“永远输出训练标签均值” $\hat{f}_{\mathcal{D}}(x) \equiv \bar{y} = \frac{1}{N}\sum_{i=1}^N y_i$（其中训练点 $x_i$ 固定）。求出任意测试点 $x$ 处的精确 $\Bias^2(\hat{f}(x))$、$\Var(\hat{f}(x))$ 及期望测试 MSE，并说明当 $N \to \infty$ 时三项各自趋向何值。
  \exHint{先算 $\bar{f}(x) = \E[\bar{y}] = \frac{1}{N}\sum_{i=1}^N f(x_i)$，再由独立性直接算 $\Var(\bar{y}) = \sigma^2 / N$。}

  \item\label{ex:4.3.2} \textbf{PAC 样本复杂度计算}：假设一个二分类任务包含 $|\mathcal{H}| = 2^{20}$（约一百万）个候选决策树桩组合。若希望以至少 $99\%$ 的置信度（$\delta = 0.01$）保证一致泛化间隙不超过 $\varepsilon = 0.02$，至少需要多少个 i.i.d. 样本？若将置信度放松到 $90\%$（$\delta = 0.10$），样本需求减少的百分比是多少？
  \exHint{直接代入公式~\eqref{eq:v4-sample-comp}，感受 $\ln(1/\delta)$ 对置信度的极低敏感度。}
\end{enumerate}

\exTier{B 层}{岭回归推导与 VC 维计算}
\begin{enumerate}[label=\textbf{4.3.\arabic*}, leftmargin=3.3em, itemsep=5pt, start=3]
  \item\label{ex:4.3.3} \textbf{岭回归的存在性优势：一定存在 $\lambda > 0$ 严格胜过 OLS}：基于式~\eqref{eq:v4-ridge-var} 与式~\eqref{eq:v4-ridge-bias}，记岭回归的总参数均方误差为 $\mathrm{MSE}(\lambda) = \tr(\Cov(\hat{w}_{\mathrm{ridge}})) + \norm{\E[\hat{w}_{\mathrm{ridge}}] - w_{\text{true}}}_2^2$。求导计算 $\left.\frac{d}{d\lambda}\mathrm{MSE}(\lambda)\right|_{\lambda = 0}$，并证明：只要 $\sigma^2 > 0$，总存在足够小的 $\lambda > 0$ 使得岭回归的期望总误差\textbf{严格低于}无偏的普通最小二乘 OLS！
  \exHint{检查 $\lambda = 0$ 处的一阶导数：平方偏差对 $\lambda$ 的导数在 $\lambda=0$ 处为 $0$，而方差的导数严格为负（$-2\sigma^2\sum \sigma_j^{-4} < 0$）！}

  \item\label{ex:4.3.4} \textbf{三类几何假设空间的精确 VC 维}：分别求出以下三个假设空间的精确 VC 维，并给出“能打散 $d_{\mathrm{VC}}$ 个点”的构造与“不能打散 $d_{\mathrm{VC}}+1$ 个点”的证明：
  (a) 一维实轴上的单侧阈值 $\mathcal{H}_1 = \{\mathbb{I}[x \ge \theta] : \theta \in \R\}$；
  (b) 一维实轴上的闭区间指示器 $\mathcal{H}_2 = \{\mathbb{I}[a \le x \le b] : a \le b\}$；
  (c) 二维平面上轴对齐矩形内部为正的分类器 $\mathcal{H}_3 = \{\mathbb{I}[a_1 \le x_1 \le b_1,\; a_2 \le x_2 \le b_2]\}$。
  \exHint{答案依次为 $1,\; 2,\; 4$；对 (c) 的 5 点反证：取最左、最右、最上、最下 4 个极值点标为 $+$，剩下的第 5 个点必被包含在它们张成的最小包围盒中，无法单独标为 $-$！}
\end{enumerate}

\exTier{C 层}{统计学习深度证明}
\begin{enumerate}[label=\textbf{4.3.\arabic*}, leftmargin=3.3em, itemsep=5pt, start=5]
  \item\label{ex:4.3.5} \textbf{留一法（LOOCV）线性回归的 $O(1)$ 快捷公式}：对线性最小二乘回归 $\hat{y} = H y$（帽子矩阵 $H = X(X^\top X)^{-1}X^\top$，对角元素 $h_{ii} \in (0, 1)$），利用第 1 卷 Sherman-Morrison 秩-1 逆矩阵公式证明：删去第 $i$ 个样本重新拟合得到的留一预测误差满足精确解析恒等式
  \[
    y_i - \hat{y}_i^{(-i)} = \frac{y_i - \hat{y}_i}{1 - h_{ii}},
  \]
  从而只需训练\textbf{一次}全量模型，即可精确算出全部 $N$ 折留一法交叉验证均方误差！
  \exHint{利用 $(X^\top X - x_i x_i^\top)^{-1} = (X^\top X)^{-1} + \frac{(X^\top X)^{-1}x_i x_i^\top (X^\top X)^{-1}}{1 - h_{ii}}$ 展开 $x_i^\top \hat{w}^{(-i)}$。}

  \item\label{ex:4.3.6} \textbf{单边可实现情形（Realizable PAC）的 $O(1/\varepsilon)$ 加速样本复杂度}：若假设空间 $\mathcal{H}$ 有限，且已知存在零真实误差的完美假设 $h^* \in \mathcal{H}$（即 $R(h^*) = 0$，无噪声可实现情形），证明对任意使 $R_{\mathrm{emp}}(\hat{h}) = 0$ 的一致学习器 $\hat{h}$，只需
  \[
    N \ge \frac{1}{\varepsilon}\ln\frac{|\mathcal{H}|}{\delta}
  \]
  个样本即可以至少 $1-\delta$ 的概率保证 $R(\hat{h}) \le \varepsilon$（依赖从 $1/\varepsilon^2$ 骤降至 $1/\varepsilon$！）。
  \exHint{若某个坏假设满足 $R(h) > \varepsilon$，它在 $N$ 个 i.i.d. 样本上侥幸全对的概率至多为 $(1 - \varepsilon)^N \le e^{-N\varepsilon}$；再对所有坏假设取并集界即可。}
\end{enumerate}

\exTier{D 层}{诊断实验与泄露复现}
\begin{enumerate}[label=\textbf{4.3.\arabic*}, leftmargin=3.3em, itemsep=5pt, start=7]
  \item\label{ex:4.3.7} \textbf{岭回归压制高阶多项式方差实验}：修改脚本 \texttt{vol4\_experiments.py} 中的实验 3，固定图~\ref{fig:v4-ch3-bv} 中的 $d = 9$ 阶多项式，将普通最小二乘替换为岭回归（注意不对常数截距项 $w_0$ 惩罚），扫描 $\lambda \in \{0, 10^{-8}, 10^{-6}, 10^{-4}, 10^{-2}, 10^{-1}\}$。打印并绘制 400 轮实测的 $\Bias^2$、$\Var$ 与总 MSE 随 $\lambda$ 的变化曲线，验证是否能把 $d=9$ 的总测试误差压到比无正则 $d=3$（$0.1073$）更低或持平！
  \exHint{先对高次项 $x^k$ 做中心化/标准化或采用正交多项式基，可显著改善范德蒙矩阵的数值条件数。}

  \item\label{ex:4.3.8} \textbf{亲手复现“特征选择数据泄露”假象}：用 NumPy 生成 $N = 100$ 个样本、$D = 5000$ 维纯独立标准正态噪声特征 $X$，以及与 $X$ \textbf{完全独立}的随机二分类标签 $y \in \{0, 1\}$（真实最优准确率严格等于 $50\%$）。比较以下两种 5 折交叉验证流程的验证准确率：
  (a) \textbf{错误流程（泄露）}：先在全部 100 个样本上计算每列特征与 $y$ 的相关系数，选出绝对值最大的 Top-20 特征，再跑 5 折交叉验证逻辑回归；
  (b) \textbf{正确流程（Pipeline 隔离）}：在每一折内仅用 80 个训练样本挑选 Top-20 特征并拟合，再去预测该折 20 个验证样本。
  报告两者的 5 折平均准确率差异，写下你的工程防错准则。
  \exHint{错误流程常会荒谬地报出 $80\%$--$90\%$ 的虚假高分，而正确流程会诚实地回归至 $50\%$ 附近！}
\end{enumerate}

\chapter[支持向量机与核技巧]{支持向量机（SVM）与核技巧}\label{ch:svm}
\chapEpigraph{在实数域中，两个真理之间的最短路径往往要穿过复数域。}{Jacques Hadamard（雅克·阿达马）}

\begin{chapIntro}{本章导读}
第 2 章我们系统建立了拉格朗日对偶与 KKT 互补松弛理论，第 3 章又剖析了正则化与泛化误差的深层联系。现在，这两条主线将在\textbf{支持向量机（Support Vector Machine, SVM）}中迎来最辉煌的合流。

SVM 的思想优雅得近乎艺术：它不满足于“把两类样本分开”，而是追问“在无穷多条分离超平面中，哪一条留出了最宽的安全缓冲带”；借助拉格朗日对偶，它把原始几何问题转化为仅依赖\textbf{样本内积}的凸二次规划，并通过 KKT 互补松弛自动筛选出极少数决定全局边界的\textbf{支持向量}；当数据在原空间线性不可分时，它顺理成章地祭出\textbf{核技巧（Kernel Trick）}——正如阿达马所言，借道高维甚至无穷维特征空间，一条在原空间盘根错节的非线性边界瞬间化为坦途。
\end{chapIntro}

% ============================================================
\section{最大间隔分类器与规范化凸二次规划}\label{sec:svm-margin}
% ============================================================

\subsection{为什么“分开”还不够？大间隔的几何与泛化直觉}

给定线性可分的二分类训练集 $\mathcal{D} = \{(x_i, y_i)\}_{i=1}^{N}$，其中 $x_i \in \R^d$，标签约定为 $y_i \in \{-1, +1\}$。寻找一个仿射超平面
\[
  \mathcal{H} = \{x \in \R^d \mid w^\top x + b = 0\}, \qquad w \in \R^d \setminus \{0\},\ b \in \R,
\]
使得所有正样本位于一侧（$w^\top x_i + b > 0$）、所有负样本位于另一侧（$w^\top x_i + b < 0$），这等价于要求
\[
  y_i (w^\top x_i + b) > 0, \qquad \forall\, i = 1, \dots, N.
\]
一旦数据线性可分，满足上述条件的超平面便有\textbf{无穷多条}。然而从第 3 章的泛化视角来看，这些超平面绝非生而平等：
\begin{itemize}
  \item 如果某条超平面紧紧擦着某几个训练样本的边儿划过，哪怕测试样本只受到极微小的测量噪声扰动，就可能越过边界被误判；
  \item 反之，若超平面稳居两类样本点云的正中央，为两侧都留出尽可能宽的\textbf{安全隔离带（Margin）}，它就能容忍更大范围的局部扰动，展现出更强的鲁棒性与更紧的泛化界。
\end{itemize}

\subsection{点到超平面的精确欧氏距离推导}

要精确刻画“隔离带宽度”，必须先严格算出空间中任意一点 $x \in \R^d$ 到超平面 $\mathcal{H}: w^\top z + b = 0$ 的欧氏距离。

\begin{thmBox}
\begin{proposition}[点到超平面的有符号与无符号距离]\label{prop:svm-dist}
设超平面 $\mathcal{H} = \{z \in \R^d \mid w^\top z + b = 0\}$，其中 $w \neq 0$。则任意点 $x \in \R^d$ 到 $\mathcal{H}$ 的最短欧氏距离为
\begin{equation}\label{eq:point-plane-dist}
  d(x, \mathcal{H}) = \frac{\abs{w^\top x + b}}{\norm{w}_2}.
\end{equation}
若 $(x_i, y_i)$ 被该超平面正确分类，其\textbf{几何间隔（Geometric Margin）}为
\[
  \gamma_i := \frac{y_i(w^\top x_i + b)}{\norm{w}_2} > 0.
\]
\end{proposition}
\end{thmBox}

\begin{proof}
首先确认 $w$ 是超平面 $\mathcal{H}$ 的法向量：对任意两点 $z_1, z_2 \in \mathcal{H}$，有 $w^\top z_1 + b = 0$ 且 $w^\top z_2 + b = 0$，两式相减得 $w^\top (z_1 - z_2) = 0$，故 $w$ 与 $\mathcal{H}$ 内一切切向量正交。

设 $x_\perp \in \mathcal{H}$ 为点 $x$ 在超平面 $\mathcal{H}$ 上的正交投影，则向量 $x - x_\perp$ 平行于单位法向量 $\hat{w} := w / \norm{w}_2$，即存在唯一标量 $r \in \R$（有符号距离）使得
\[
  x = x_\perp + r \frac{w}{\norm{w}_2}.
\]
两边左乘 $w^\top$ 并加上 $b$，利用 $w^\top x_\perp + b = 0$ 与 $w^\top w = \norm{w}_2^2$：
\[
  w^\top x + b = (w^\top x_\perp + b) + r \frac{w^\top w}{\norm{w}_2} = r \norm{w}_2 \implies r = \frac{w^\top x + b}{\norm{w}_2}.
\]
因此点到超平面的最短欧氏距离即为 $\norm{x - x_\perp}_2 = \abs{r} = \abs{w^\top x + b}/\norm{w}_2$。当样本被正确分类时，$y_i$ 与 $w^\top x_i + b$ 同号，故 $y_i(w^\top x_i + b) = \abs{w^\top x_i + b}$，几何间隔即等于该点到超平面的真实欧氏距离。
\end{proof}

\subsection{尺度不变性与规范化：从非凸分式到严格凸二次规划}

整个训练集上的\textbf{最小几何间隔}（简称超平面的间隔）定义为距离超平面最近的那个样本点的距离：
\[
  \gamma(w, b) := \min_{i=1,\dots,N} \frac{y_i(w^\top x_i + b)}{\norm{w}_2}.
\]
我们的目标自然是最大化 $\gamma(w, b)$：
\[
  \max_{w, b} \;\gamma \qquad \text{s.t.}\quad \frac{y_i(w^\top x_i + b)}{\norm{w}_2} \ge \gamma,\quad i=1,\dots,N.
\]
这一原始形式的分母里含 $\norm{w}_2$，约束并非凸集。但请留意一个关键的\textbf{尺度不变性（Scale Invariance）}：对任意正数 $c > 0$，将 $(w, b)$ 缩放为 $(cw, cb)$ 完全不改变超平面本身（$c(w^\top x + b) = 0 \iff w^\top x + b = 0$），也不改变任何样本的几何间隔 $\gamma_i$；但它会把\textbf{函数间隔（Functional Margin）} $\hat{\gamma}_i := y_i(w^\top x_i + b)$ 同步缩放 $c$ 倍！

既然 $(w, b)$ 多出一个无关的全局缩放自由度，我们可以主动\textbf{钉死这把尺子}：令离超平面最近的样本满足函数间隔恰好等于 $1$，即
\[
  \hat{\gamma} := \min_{i=1,\dots,N} y_i(w^\top x_i + b) = 1.
\]
在这种\textbf{规范化（Canonical Hyperplane）}约定下，最小几何间隔简化为 $\gamma = \hat{\gamma}/\norm{w}_2 = 1/\norm{w}_2$，两类规范边界超平面 $w^\top x + b = +1$ 与 $w^\top x + b = -1$ 之间的安全隔离带总宽度为 $2/\norm{w}_2$（如图~\ref{fig:svm-geometry} 所示）。最大化 $1/\norm{w}_2$ 等价于极小化 $\norm{w}_2$，为了求导干净利落，进一步等价地极小化 $\frac{1}{2}\norm{w}_2^2$。

\begin{defBox}
\begin{definition}[硬间隔支持向量机原始问题（Hard-Margin SVM Primal）]\label{def:svm-primal}
对于线性可分训练集 $\{(x_i, y_i)\}_{i=1}^{N}$（$y_i \in \{-1, +1\}$），硬间隔最大间隔分类器定义为如下凸二次规划（Quadratic Program, QP）：
\begin{equation}\label{eq:svm-primal}
  \min_{w \in \R^d,\, b \in \R}\; \frac{1}{2}\norm{w}_2^2
  \qquad \text{s.t.}\quad y_i(w^\top x_i + b) \ge 1, \quad i = 1, \dots, N.
\end{equation}
\end{definition}
\end{defBox}

检查式~\eqref{eq:svm-primal} 的凸性：目标函数 $f(w, b) = \frac{1}{2}\norm{w}_2^2$ 关于 $w$ 是 $1$-强凸的（Hessian 矩阵对 $w$ 块为单位阵 $I_d \succ 0$），全部 $N$ 条不等式约束 $g_i(w, b) := 1 - y_i(w^\top x_i + b) \le 0$ 均为仿射函数。只要数据线性可分（即存在严格分离超平面使可行域非空），该凸 QP 的最优解 $(w^*, b^*)$ 必然存在，且法向量 $w^*$ \textbf{严格唯一}（若存在两个不同最优解 $w_1 \neq w_2$，中点 $(w_1+w_2)/2$ 同样满足所有线性约束且使 $\frac{1}{2}\norm{w}_2^2$ 严格更小，矛盾！）。

\begin{figure}[htbp]
\centering
\begin{tikzpicture}[scale=0.95]
  % Grid / axes subtle frame
  \draw[->, >=Stealth, slate!75, thin] (-2.9, -2.0) -- (4.1, -2.0) node[right, font=\footnotesize] {$x_1$};
  \draw[->, >=Stealth, slate!75, thin] (-2.6, -2.2) -- (-2.6, 2.8) node[above, font=\footnotesize] {$x_2$};

  % Margin slab (rotated coordinate: line x2 = 0.55*x1 + 0.1)
  % Let w = (-0.55, 1) normalized direction -> slab half-height in x2 is 0.92
  \fill[tealLight, opacity=0.82] (-2.2, -2.03) -- (3.6, 1.16) -- (3.6, 2.74) -- (-2.2, -0.45) -- cycle;

  % Decision boundary & canonical margin boundaries
  \draw[tealDark, very thick] (-2.2, -1.24) -- (3.6, 1.95)
    node[pos=0.95, below right, font=\footnotesize\sffamily\bfseries, text=tealDark, sloped, xshift=-10mm, yshift=-1pt] {$w^\top x + b = 0$};
  \draw[teal, thick, dashed] (-2.2, -0.45) -- (3.6, 2.74)
    node[pos=0.44, above, font=\scriptsize, text=teal, sloped] {$w^\top x + b = +1$};
  \draw[ochre, thick, dashed] (-2.2, -2.03) -- (3.6, 1.16)
    node[pos=0.76, below, font=\scriptsize, text=ochre, sloped] {$w^\top x + b = -1$};

  % Width arrow perpendicular to w^\top x + b = 0
  \draw[<->, >=Stealth, gold, very thick] (1.334, -0.087) -- (0.666, 1.127);
  \node[font=\footnotesize\bfseries, text=gold, fill=tealLight, inner sep=1.2pt, rounded corners=1pt]
    at (1.60, 0.42) {$\dfrac{2}{\norm{w}_2}$};

  % Normal vector w from center line (placed cleanly on left away from labels & SVs)
  \draw[->, >=Stealth, navy, very thick] (-1.45, -0.83) -- (-1.91, 0.01)
    node[above left, font=\footnotesize\bfseries, text=navy] {$w$};

  % Positive samples (y = +1): y - 0.55x >= 0.76
  % Support vectors exactly on y = 0.55x + 0.76:
  % (-0.8, 0.32), (1.8, 1.75)
  \foreach \x/\y in {-0.8/0.32, 1.8/1.75} {
    \fill[teal] (\x, \y) circle (2.8pt);
    \draw[tealDark, thick] (\x, \y) circle (5.6pt);
  }
  % Non-SV positive points
  \foreach \x/\y in {-1.8/0.9, -1.2/1.6, -0.2/1.8, 0.4/2.2, 0.9/1.9, -1.7/2.1} {
    \fill[teal!70] (\x, \y) circle (2.6pt);
  }
  \node[font=\scriptsize\sffamily, text=tealDark] at (-1.55, 2.48) {正类 ($y_i = +1$)};

  % Negative samples (y = -1): y - 0.55x <= -0.82
  % Support vectors on y = 0.55x - 0.82:
  % (0.2, -0.71), (2.4, 0.50)
  \foreach \x/\y in {0.2/-0.71, 2.4/0.50} {
    \fill[ochre] (\x, \y) circle (2.8pt);
    \draw[ochre, thick] (\x, \y) circle (5.6pt);
  }
  % Non-SV negative points
  \foreach \x/\y in {-0.4/-1.5, 0.8/-1.4, 1.6/-1.1, 2.5/-0.5, 3.1/0.1, 2.1/-1.5, 3.4/-0.9} {
    \fill[ochre!70] (\x, \y) circle (2.6pt);
  }
  \node[font=\scriptsize\sffamily, text=ochre] at (3.11, -1.76) {负类 ($y_i = -1$)};

  % Callout for Support Vectors
  \node[font=\scriptsize\sffamily, text=ink, align=left, fill=white, draw=tealLine, rounded corners=1.5pt, inner sep=3pt]
    (svlbl) at (-0.88, -2.78) {圆圈标记为\textbf{支持向量}\\满足 $y_i(w^\top x_i + b) = 1$ 且 $\alpha_i^* > 0$};
  \draw[->, >=Stealth, slate, thin] (svlbl.north east) to[out=100,in=-70] (0.27, -0.90);
\end{tikzpicture}
\caption{硬间隔支持向量机的最大间隔几何结构。规范化超平面 $w^\top x + b = \pm 1$ 之间形成宽度为 $2/\norm{w}_2$ 的无样本安全隔离带，只有恰好压在两条虚线边界上的支持向量对决策边界起支撑作用。}
\label{fig:svm-geometry}
\end{figure}

% ============================================================
\section{拉格朗日对偶、KKT 互补松弛与支持向量}\label{sec:svm-dual}
% ============================================================

直接求解原始 QP~\eqref{eq:svm-primal} 在低维特征空间本不困难；但真正让 SVM 封神的，是把第 2 章的拉格朗日对偶理论工具作用上去之后发生的两个奇迹：\textbf{支持向量稀疏性}与\textbf{纯内积形式}。

\begin{recallBox}{拉格朗日强对偶与 KKT 条件（第 2 章回顾）}
对于凸优化问题 $\min_x f_0(x)\ \text{s.t.}\ g_i(x)\le 0$，若全部不等式约束均为仿射函数且可行域非空，则弱化的 Slater 条件自动满足，\textbf{强对偶性成立}（原始最优值与对偶最优值相等），且 $(x^*, \alpha^*)$ 为原始-对偶最优解的充要条件是同时满足：原始可行性 $g_i(x^*)\le 0$、对偶可行性 $\alpha_i^* \ge 0$、稳定性 $\nabla_x \mathcal{L}(x^*, \alpha^*) = 0$ 以及\textbf{互补松弛条件} $\alpha_i^* g_i(x^*) = 0$。
\end{recallBox}

\subsection{硬间隔 SVM 对偶问题的完整推导}

为每条不等式约束 $1 - y_i(w^\top x_i + b) \le 0$ 引入非负拉格朗日乘子 $\alpha_i \ge 0$，构造拉格朗日函数：
\begin{equation}\label{eq:svm-lagrangian}
  \mathcal{L}(w, b, \alpha) = \frac{1}{2}\norm{w}_2^2 + \sum_{i=1}^{N} \alpha_i \big(1 - y_i(w^\top x_i + b)\big).
\end{equation}
为了计算对偶函数 $g(\alpha) = \inf_{w, b} \mathcal{L}(w, b, \alpha)$，先对无约束变量 $w$ 与 $b$ 求偏导并令其为零：
\begin{align}
  \nabla_w \mathcal{L}(w, b, \alpha) &= w - \sum_{i=1}^{N} \alpha_i y_i x_i = 0 \implies w = \sum_{i=1}^{N} \alpha_i y_i x_i, \label{eq:svm-w-rep}\\
  \frac{\partial \mathcal{L}(w, b, \alpha)}{\partial b} &= -\sum_{i=1}^{N} \alpha_i y_i = 0 \implies \sum_{i=1}^{N} \alpha_i y_i = 0. \label{eq:svm-b-cond}
\end{align}
注意若 $\sum_{i=1}^N \alpha_i y_i \neq 0$，则 $\mathcal{L}$ 关于 $b$ 是斜率非零的线性函数，下确界为 $-\infty$。因此在 $\sum_i \alpha_i y_i = 0$ 成立的前提下，将式~\eqref{eq:svm-w-rep} 代回式~\eqref{eq:svm-lagrangian}：
\begin{align*}
  \mathcal{L}(w, b, \alpha)
  &= \frac{1}{2} w^\top w + \sum_{i=1}^{N} \alpha_i - w^\top \!\left(\sum_{i=1}^{N} \alpha_i y_i x_i\right) - b \underbrace{\sum_{i=1}^{N} \alpha_i y_i}_{=0} \\
  &= \sum_{i=1}^{N} \alpha_i - \frac{1}{2} w^\top w
   = \sum_{i=1}^{N} \alpha_i - \frac{1}{2}\sum_{i=1}^{N}\sum_{j=1}^{N} \alpha_i \alpha_j y_i y_j x_i^\top x_j.
\end{align*}
由此便得到 SVM 著名的\textbf{拉格朗日对偶问题}。

\begin{thmBox}
\begin{theorem}[硬间隔 SVM 的拉格朗日对偶与支持向量定理]\label{thm:svm-dual}
若训练集线性可分，则原始问题~\eqref{eq:svm-primal} 与如下对偶二次规划之间强对偶成立（零对偶间隙）：
\begin{equation}\label{eq:svm-dual}
  \max_{\alpha \in \R^N}\; \sum_{i=1}^{N} \alpha_i - \frac{1}{2}\sum_{i=1}^{N}\sum_{j=1}^{N} \alpha_i \alpha_j y_i y_j\, (x_i^\top x_j)
  \qquad \text{s.t.}\quad \alpha_i \ge 0\ (\forall i),\quad \sum_{i=1}^{N} \alpha_i y_i = 0.
\end{equation}
设 $\alpha^*$ 为对偶最优解，记支持向量下标集 $\mathcal{S} := \{i \mid \alpha_i^* > 0\}$。则：
\begin{enumerate}
  \item \textbf{权重表示}：$w^* = \sum_{i \in \mathcal{S}} \alpha_i^* y_i x_i$；
  \item \textbf{互补松弛与边界紧致性}：对一切 $i \in \mathcal{S}$，必有 $y_i((w^*)^\top x_i + b^*) = 1$（恰好落在间隔边界上）；对一切满足 $y_i((w^*)^\top x_i + b^*) > 1$ 的非边界样本，必有 $\alpha_i^* = 0$；
  \item \textbf{截距恢复}：任取 $s \in \mathcal{S}$，由 $y_s((w^*)^\top x_s + b^*) = 1$ 且 $y_s^2 = 1$ 得
  \[
    b^* = y_s - (w^*)^\top x_s = y_s - \sum_{i \in \mathcal{S}} \alpha_i^* y_i\, (x_i^\top x_s);
  \]
  \item \textbf{删除不变性}：将所有非支持向量（$i \notin \mathcal{S}$）从训练集中彻底删除后重新训练，最优超平面 $(w^*, b^*)$ \textbf{保持绝对不变}。
\end{enumerate}
\end{theorem}
\end{thmBox}

\begin{proof}
由于原始目标函数为凸函数、全部 $N$ 条约束均为仿射函数且线性可分保证可行域非空，第 2 章的仿射 Slater 条件满足，故强对偶成立且原始-对偶最优解 $(w^*, b^*, \alpha^*)$ 满足全套 KKT 条件。
(1) 直接来自稳定性条件~\eqref{eq:svm-w-rep} 并剔除 $\alpha_i^* = 0$ 的项；
(2) 直接来自 KKT 互补松弛条件 $\alpha_i^* \big(1 - y_i((w^*)^\top x_i + b^*)\big) = 0$：若 $\alpha_i^* > 0$，括号内因子必为零；反之若 $y_i((w^*)^\top x_i + b^*) > 1$，则括号内非零，迫使 $\alpha_i^* = 0$；
(3) 将 $y_s((w^*)^\top x_s + b^*) = 1$ 两边同乘 $y_s \in \{-1, +1\}$ 立即得 $b^* = y_s - (w^*)^\top x_s$（工程上常对全部 $s \in \mathcal{S}$ 取平均 $b^* = \frac{1}{|\mathcal{S}|}\sum_{s \in \mathcal{S}}(y_s - \sum_{i\in\mathcal{S}}\alpha_i^* y_i x_i^\top x_s)$ 以抵消浮点误差）；
(4) 删去所有 $i \notin \mathcal{S}$ 的样本后，其余样本对应的 $(w^*, b^*, \{\alpha_i^*\}_{i\in\mathcal{S}})$ 依然完整满足缩减问题的所有 KKT 条件（原始可行性、对偶可行性、稳定性与互补松弛丝毫不受影响），由凸问题 KKT 的充分性可知它仍是缩减问题的全局最优解！
\end{proof}

请停下来体会定理~\ref{thm:svm-dual} 第 (4) 条的震撼力：哪怕训练集有 10 万个样本，只要其中 99{,}950 个远离边界，把它们全部扔进垃圾桶，最大间隔超平面纹丝不动！整个分类器的全部记忆都被压缩进了极少数\textbf{支持向量（Support Vectors）}之中。

\subsection{手工完整推演：二维 3 点 SVM 与第 4 点的删除不变性}

\begin{example}[手工解对偶 QP 求 $\alpha^*, w^*, b^*$]\label{ex:svm-hand}
考虑二维平面上的 3 个训练样本：
\[
  x_1 = \begin{bmatrix}0 \\ 0\end{bmatrix},\ y_1 = -1; \qquad
  x_2 = \begin{bmatrix}2 \\ 0\end{bmatrix},\ y_2 = +1; \qquad
  x_3 = \begin{bmatrix}0 \\ 2\end{bmatrix},\ y_3 = +1.
\]
我们直接通过求解对偶问题~\eqref{eq:svm-dual} 算出最优解，再检验加入第 4 个样本 $x_4 = (3, 2)^\top, y_4 = +1$ 时的表现。
\end{example}

\paragraph{第 1 步：写出内积矩阵与对偶目标}
计算样本间的 Gram 内积矩阵 $G_{ij} = x_i^\top x_j$：
\[
  G = \begin{bmatrix}
    0 & 0 & 0 \\
    0 & 4 & 0 \\
    0 & 0 & 4
  \end{bmatrix}.
\]
对偶等式约束 $\sum_{i=1}^3 \alpha_i y_i = 0$ 给出 $-\alpha_1 + \alpha_2 + \alpha_3 = 0$，即 $\alpha_1 = \alpha_2 + \alpha_3$。由于 $x_1 = (0,0)^\top$ 与所有样本内积均为零，二次项中只剩下 $(2,2)$ 与 $(3,3)$ 两项：
\[
  W(\alpha) = (\alpha_1 + \alpha_2 + \alpha_3) - \frac{1}{2}\big(4\alpha_2^2 + 4\alpha_3^2\big)
  = 2\alpha_2 + 2\alpha_3 - 2\alpha_2^2 - 2\alpha_3^2.
\]

\paragraph{第 2 步：求对偶极大值}
$W$ 关于 $\alpha_2, \alpha_3$ 完全解耦且严格凹，令偏导为零：
\[
  \frac{\partial W}{\partial \alpha_2} = 2 - 4\alpha_2 = 0 \implies \alpha_2^* = \frac{1}{2}; \qquad
  \frac{\partial W}{\partial \alpha_3} = 2 - 4\alpha_3 = 0 \implies \alpha_3^* = \frac{1}{2}.
\]
进而 $\alpha_1^* = \alpha_2^* + \alpha_3^* = 1$。检查非负性：$\alpha^* = (1,\; \frac{1}{2},\; \frac{1}{2})^\top > 0$，三个点全部都是支持向量！

\paragraph{第 3 步：还原 $(w^*, b^*)$ 与强对偶核验}
由式~\eqref{eq:svm-w-rep}：
\[
  w^* = \alpha_1^* y_1 x_1 + \alpha_2^* y_2 x_2 + \alpha_3^* y_3 x_3
      = 1(-1)\begin{bmatrix}0\\0\end{bmatrix} + \frac{1}{2}(+1)\begin{bmatrix}2\\0\end{bmatrix} + \frac{1}{2}(+1)\begin{bmatrix}0\\2\end{bmatrix}
      = \begin{bmatrix}1 \\ 1\end{bmatrix}.
\]
任取支持向量 $x_1 = (0,0)^\top$ 计算截距：$b^* = y_1 - (w^*)^\top x_1 = -1 - 0 = -1$（用 $x_2$ 验算：$b^* = 1 - 2 = -1$；用 $x_3$ 验算：$b^* = 1 - 2 = -1$，三者严丝合缝）。因此最优超平面为
\[
  x_1 + x_2 - 1 = 0, \qquad \gamma^* = \frac{1}{\norm{w^*}_2} = \frac{1}{\sqrt{2}}, \qquad \text{隔离带总宽 } \frac{2}{\norm{w^*}_2} = \sqrt{2}.
\]
核验强对偶间隙：原始最优值 $\frac{1}{2}\norm{w^*}_2^2 = \frac{1}{2}(1^2 + 1^2) = 1$；对偶最优值 $W(\alpha^*) = 2(\frac{1}{2}) + 2(\frac{1}{2}) - 2(\frac{1}{4}) - 2(\frac{1}{4}) = 1$，二者精确相等！

\paragraph{第 4 步：加入第 4 个非支持向量 $x_4 = (3, 2)^\top, y_4 = +1$}
检验其在当前超平面下的函数间隔：
\[
  y_4 \big((w^*)^\top x_4 + b^*\big) = (+1)(3 + 2 - 1) = 4 > 1.
\]
该点严格落在正类规范边界之外。令 $\alpha_4^* = 0$，则 $(w^*, b^*, \alpha_1^*, \alpha_2^*, \alpha_3^*, 0)$ 满足 4 点问题的全部 KKT 条件，故最优超平面依然是 $x_1 + x_2 - 1 = 0$。

% ============================================================
\section{软间隔 SVM 与合页损失的严格等价性}\label{sec:svm-soft}
% ============================================================

\subsection{为什么需要软间隔？噪声与离群点的致命破坏}

硬间隔 SVM 的约束 $y_i(w^\top x_i + b) \ge 1$ 过于刚性：在真实数据中，哪怕只有一个正样本因标注错误掉进了负类腹地，硬间隔问题就会直接\textbf{无可行解}；即使勉强线性可分，单个紧贴对手阵营的离群点也会强行扭转超平面，将整个安全隔离带挤压得惨不忍睹（典型的过拟合与高方差）。

解决之道是允许每个样本付出可控的“违约金”：为每个样本 $i$ 引入非负\textbf{松弛变量（Slack Variable）} $\xi_i \ge 0$，允许其函数间隔向内退让 $\xi_i$ 的幅度，同时在目标函数中按惩罚系数 $C > 0$ 收取总违约罚金：
\begin{equation}\label{eq:svm-soft-primal}
  \min_{w,\, b,\, \xi}\; \frac{1}{2}\norm{w}_2^2 + C\sum_{i=1}^{N} \xi_i
  \qquad \text{s.t.}\quad y_i(w^\top x_i + b) \ge 1 - \xi_i,\quad \xi_i \ge 0,\quad i = 1,\dots,N.
\end{equation}
这里的 $\xi_i$ 具有清晰的几何分层意义：$\xi_i = 0$ 表示样本位于正确一侧的间隔边界上或边界外；$0 < \xi_i < 1$ 表示样本侵入了隔离带但仍被正确分类；$\xi_i \ge 1$ 则表示样本已越过决策边界（或压在边界上）遭到误分类。

\subsection{盒式约束对偶与三类乘子的完整刻画}

为两类不等式约束分别引入拉格朗日乘子 $\alpha_i \ge 0$ 与 $\mu_i \ge 0$，写出软间隔拉格朗日函数：
\[
  \mathcal{L}(w, b, \xi, \alpha, \mu)
  = \frac{1}{2}\norm{w}_2^2 + C\sum_{i=1}^{N}\xi_i
  + \sum_{i=1}^{N}\alpha_i\big(1 - \xi_i - y_i(w^\top x_i + b)\big) - \sum_{i=1}^{N}\mu_i \xi_i.
\]
对原始变量 $w, b, \xi_i$ 求偏导置零：
\begin{align*}
  \nabla_w \mathcal{L} = 0 &\implies w = \sum_{i=1}^{N}\alpha_i y_i x_i, \\
  \frac{\partial \mathcal{L}}{\partial b} = 0 &\implies \sum_{i=1}^{N}\alpha_i y_i = 0, \\
  \frac{\partial \mathcal{L}}{\partial \xi_i} = 0 &\implies C - \alpha_i - \mu_i = 0 \implies \mu_i = C - \alpha_i.
\end{align*}
请留意第三行导出的神奇消去：由于 $\mu_i \ge 0$ 且 $\alpha_i \ge 0$，关系式 $\mu_i = C - \alpha_i$ 直接把乘子 $\mu_i$ 从对偶问题中彻底消去，化作对 $\alpha_i$ 的\textbf{盒式约束（Box Constraint）} $0 \le \alpha_i \le C$！同时拉格朗日函数中关于 $\xi_i$ 的全部项 $(C - \alpha_i - \mu_i)\xi_i$ 恒等于零。代回后我们惊讶地发现，\textbf{软间隔 SVM 的对偶目标函数与硬间隔一字不差}：
\begin{equation}\label{eq:svm-soft-dual}
  \max_{\alpha \in \R^N}\; \sum_{i=1}^{N} \alpha_i - \frac{1}{2}\sum_{i=1}^{N}\sum_{j=1}^{N} \alpha_i \alpha_j y_i y_j\, (x_i^\top x_j)
  \qquad \text{s.t.}\quad 0 \le \alpha_i \le C\ (\forall i),\quad \sum_{i=1}^{N} \alpha_i y_i = 0.
\end{equation}
此时两套互补松弛条件为
\[
  \alpha_i^*\big(1 - \xi_i^* - y_i((w^*)^\top x_i + b^*)\big) = 0, \qquad
  (C - \alpha_i^*)\xi_i^* = 0.
\]
由此可以将全部训练样本精确划分为三种互斥状态（这也是 SMO 求解器判定 KKT 违背的核心依据）：
\begin{enumerate}
  \item $\alpha_i^* = 0$：由 $(C - 0)\xi_i^* = 0$ 得 $\xi_i^* = 0$，样本位于间隔外侧 $y_i f^*(x_i) \ge 1$，为\textbf{非支持向量}；
  \item $0 < \alpha_i^* < C$：两项互补松弛同时生效，必有 $\xi_i^* = 0$ 且 $y_i f^*(x_i) = 1$，为\textbf{自由支持向量（Free SV）}，恰好落在规范边界上，\textbf{只有此类样本才能用来求解截距 $b^*$}；
  \item $\alpha_i^* = C$：必有 $1 - \xi_i^* - y_i f^*(x_i) = 0$，即 $y_i f^*(x_i) = 1 - \xi_i^* \le 1$，为\textbf{受限支持向量（Bounded SV）}，位于隔离带内部或已被误分；其对 $w^*$ 的影响力被封顶在 $C$，从而天然防止单个离群点无限拉偏决策边界！
\end{enumerate}

\subsection{软间隔 SVM 与 $L_2$ 正则化合页损失的等价性}

\begin{thmBox}
\begin{theorem}[软间隔 SVM $\iff$ 正则化合页损失经验风险最小化]\label{thm:svm-hinge}
软间隔 SVM 原始问题~\eqref{eq:svm-soft-primal} 与如下无约束凸优化问题拥有完全相同的最优超平面 $(w^*, b^*)$：
\begin{equation}\label{eq:svm-hinge}
  \min_{w \in \R^d,\, b \in \R}\; \frac{1}{2}\norm{w}_2^2 + C \sum_{i=1}^{N} \max\!\big(0,\; 1 - y_i(w^\top x_i + b)\big).
\end{equation}
其中 $\ell_{\mathrm{hinge}}(z) := \max(0, 1 - z)$ 称为\textbf{合页损失（Hinge Loss）}，$z_i = y_i(w^\top x_i + b)$ 为样本函数间隔。
\end{theorem}
\end{thmBox}

\begin{proof}
对任意固定的一组超平面参数 $(w, b)$，软间隔问题~\eqref{eq:svm-soft-primal} 关于松弛变量 $\xi_1, \dots, \xi_N$ 完全分离；对每个 $i$，约束要求 $\xi_i \ge 1 - y_i(w^\top x_i + b)$ 且 $\xi_i \ge 0$，即 $\xi_i \ge \max(0,\, 1 - y_i(w^\top x_i + b))$。由于目标函数中 $\xi_i$ 的系数 $C > 0$ 严格为正，在任意给定 $(w, b)$ 下使目标最小的 $\xi_i$ 唯一取在下界处：
\[
  \xi_i^*(w, b) = \max\!\big(0,\; 1 - y_i(w^\top x_i + b)\big).
\]
将 $\xi_i^*(w, b)$ 代回原目标函数，即得无约束问题~\eqref{eq:svm-hinge}。反之，任何式~\eqref{eq:svm-hinge} 的解 $(w^*, b^*)$ 配上 $\xi_i^* := \max(0, 1 - y_i((w^*)^\top x_i + b^*))$ 均满足式~\eqref{eq:svm-soft-primal} 的全部约束且达到同一目标值，故两问题严格等价。
\end{proof}

定理~\ref{thm:svm-hinge} 将 SVM 彻底纳入了第 3 章的\textbf{正则化经验风险最小化（Regularized ERM）}大统一框架！将式~\eqref{eq:svm-hinge} 两边同除以 $C$，记 $\lambda = \frac{1}{2C}$：
\[
  \min_{w, b}\; \sum_{i=1}^{N} \ell_{\mathrm{hinge}}\!\big(y_i(w^\top x_i + b)\big) + \lambda \norm{w}_2^2.
\]
于是超参数 $C$ 的偏差-方差含义一目了然：
\begin{itemize}
  \item \textbf{$C$ 越大（$\lambda$ 越小）}：对误分类与间隔违规的惩罚越严厉，间隔越窄，模型逼近硬间隔 SVM——\textbf{低偏差、高方差}，容易过拟合噪声；
  \item \textbf{$C$ 越小（$\lambda$ 越大）}：正则化越强，允许更多样本落入隔离带换取更宽的全局间隔——\textbf{高偏差、低方差}，抗噪更强但过小时会欠拟合。
\end{itemize}

图~\ref{fig:loss-compare} 将 0-1 损失、合页损失与逻辑回归对数损失画在同一坐标系下，揭示了 SVM 与逻辑回归的根本分水岭。

\begin{figure}[htbp]
\centering
\begin{tikzpicture}
\begin{axis}[
  width=0.84\linewidth,
  height=5.9cm,
  xlabel={函数间隔 $z = y(w^\top x + b)$},
  ylabel={单样本损失 $\ell(z)$},
  xmin=-2.2, xmax=2.8,
  ymin=-0.12, ymax=3.45,
  xtick={-2,-1,0,1,2},
  ytick={0,1,2,3},
  label style={font=\footnotesize},
  tick label style={font=\scriptsize},
  grid=both,
  grid style={line width=0.25pt, draw=slate!18},
  legend pos=north east,
  legend style={font=\scriptsize, fill=white, fill opacity=0.92, draw=tealLine, rounded corners=1.5pt}
]
  % Reference axes z=0 and l=0
  \draw[slate!55, thin] (axis cs:0, -0.12) -- (axis cs:0, 3.45);
  \draw[slate!55, thin] (axis cs:-2.2, 0) -- (axis cs:2.8, 0);

  % 0-1 Loss (discontinuous step at z=0)
  \addplot[ochre, very thick, dashed, domain=-2.2:0, samples=2] {1};
  \addplot[ochre, very thick, dashed, domain=0:2.8, samples=2, forget plot] {0};
  \addplot[ochre, only marks, mark=*, mark size=1.8pt, forget plot] coordinates {(0,1)};
  \addplot[ochre, only marks, mark=o, mark size=1.8pt, fill=white, forget plot] coordinates {(0,0)};
  \addlegendentry{0-1 损失 $\mathbb{I}(z \le 0)$（非凸、零梯度）}

  % Hinge Loss max(0, 1-z)
  \addplot[tealDark, very thick, domain=-2.2:1, samples=2] {1 - x};
  \addplot[tealDark, very thick, domain=1:2.8, samples=2, forget plot] {0};
  \addlegendentry{合页损失 $\max(0, 1-z)$（$z\ge 1$ 严格为零 $\Rightarrow$ 稀疏）}

  % Logistic Loss ln(1+exp(-z)) / ln(2) so it passes through (0,1)
  \addplot[navy, thick, smooth, domain=-2.2:2.8, samples=120] {ln(1 + exp(-x))/ln(2)};
  \addlegendentry{逻辑损失 $\log_2(1 + e^{-z})$（处处光滑、永不为零）}

  % Highlight kink at z=1
  \draw[dotted, teal, thick] (axis cs:1, 0) -- (axis cs:1, 1.35);
  \node[font=\scriptsize, text=tealDark, align=left, fill=white, inner sep=1.5pt] at (axis cs:1.88, 0.82) {合页折点 $z=1$：\\超过安全间隔即零梯度};
  \draw[->, >=Stealth, teal, thin] (axis cs:1.45, 0.55) -- (axis cs:1.05, 0.08);
\end{axis}
\end{tikzpicture}
\caption{二分类三大损失函数关于函数间隔 $z = y(w^\top x + b)$ 的对比。0-1 损失非凸且几乎处处梯度为零；合页损失与逻辑损失均为其凸上界代理。合页损失在 $z \ge 1$ 区间彻底“躺平”为零，赋予 SVM 独有的支持向量稀疏性；逻辑损失处处光滑且尾部指数衰减，适合输出校准概率但所有样本都参与贡献梯度。}
\label{fig:loss-compare}
\end{figure}

\begin{failBox}{工程排错：为什么使用 SVM 前几乎必须做特征标准化？}
观察原始目标 $\frac{1}{2}\norm{w}_2^2 + C\sum_i \xi_i$：如果特征 $x^{(1)}$ 的量纲是“毫米”（数值 $10^3$ 量级），而特征 $x^{(2)}$ 是“米”（数值 $1$ 量级），为了产生相同的判别贡献，$w_1$ 只需 $10^{-3}$ 量级而 $w_2$ 需要 $1$ 量级。于是正则项 $\frac{1}{2}(w_1^2 + w_2^2)$ 会几乎只惩罚 $w_2$ 而放纵 $w_1$！进入下一节的 RBF 核 $\exp(-\gamma\norm{x-z}_2^2)$ 后灾难更剧：大尺度特征会彻底吞没欧氏距离，令其余特征形同虚设。因此，\textbf{在训练线性或核 SVM 前，务必在训练集上计算均值与标准差进行 Z-score 标准化}（并严格用训练集统计量变换测试集，避免数据泄漏）。
\end{failBox}

% ============================================================
\section{核技巧、Mercer 正定性与非线性升维实测}\label{sec:svm-kernel}
% ============================================================

\subsection{核技巧：跳过显式高维映射的内积捷径}

如果数据在原空间 $\R^d$ 中本质上线性不可分——例如内圈全是负类、外圈全是正类的\textbf{同心圆分布}——任何线性超平面都无能为力。一个自然想法是构造非线性特征映射 $\phi: \R^d \to \mathcal{H}$，把原数据抬升到更高维（甚至无穷维）的希尔伯特特征空间 $\mathcal{H}$ 中，再在新空间里做线性 SVM。

但高维映射带来可怕的维数灾难：如果 $d = 100$，仅做 4 次多项式展开就需要 $\binom{100+4}{4} = 4{,}598{,}126$ 维向量，显式计算与存储 $\phi(x)$ 根本不可行！

现在回头凝视对偶问题~\eqref{eq:svm-soft-dual} 与预测公式——一个令人拍案叫绝的奇迹就在眼前：
\begin{align}
  \max_{0 \le \alpha \le C}\;& \sum_{i=1}^{N}\alpha_i - \frac{1}{2}\sum_{i=1}^{N}\sum_{j=1}^{N}\alpha_i \alpha_j y_i y_j \inner{\phi(x_i)}{\phi(x_j)}, \qquad \text{s.t.}\quad \sum_{i=1}^{N}\alpha_i y_i = 0, \label{eq:svm-kernel-dual}\\
  f(x) &= (w^*)^\top \phi(x) + b^* = \sum_{i \in \mathcal{S}} \alpha_i^* y_i \inner{\phi(x_i)}{\phi(x)} + b^*. \label{eq:svm-kernel-pred}
\end{align}
\textbf{无论是训练阶段还是预测阶段，$\phi(x)$ 从未单独出现过！它永远只以两两内积 $\inner{\phi(x)}{\phi(z)}$ 的形式成对现身！}

因此，只要存在一个可在原空间直接高效计算的二元函数 $K(x, z) := \inner{\phi(x)}{\phi(z)}$（称为\textbf{核函数，Kernel Function}），我们就可以把所有 $\inner{\phi(x_i)}{\phi(x_j)}$ 替换为 $K(x_i, x_j)$，完全不必显式构造、存储或计算高维向量 $\phi(x)$——这就是享誉整个机器学习史的\textbf{核技巧（Kernel Trick）}。

\subsection{两大经典核函数与无穷维特征映射的严格证明}

\paragraph{1. 二次多项式核（显式验证）}
取 $x = (x_1, x_2)^\top, z = (z_1, z_2)^\top \in \R^2$ 与常数 $c > 0$，定义二次多项式核 $K_{\mathrm{poly}}(x, z) := (x^\top z + c)^2$。将其直接代数展开：
\begin{align*}
  (x_1 z_1 + x_2 z_2 + c)^2
  &= x_1^2 z_1^2 + 2 x_1 x_2 z_1 z_2 + x_2^2 z_2^2 + 2c x_1 z_1 + 2c x_2 z_2 + c^2 \\
  &= \inner{\phi(x)}{\phi(z)},
\end{align*}
其中显式 6 维特征映射为
\[
  \phi(x) = \Big(x_1^2,\; \sqrt{2}\, x_1 x_2,\; x_2^2,\; \sqrt{2c}\, x_1,\; \sqrt{2c}\, x_2,\; c\Big)^\top \in \R^6.
\]
推广到一般的 $p$ 次多项式核 $(x^\top z + c)^p$：原空间求值仅需 $O(d)$ 次浮点运算，却等价于在 $O(\binom{d+p}{p})$ 维多项式空间里做精确内积！

\paragraph{2. 高斯径向基核（RBF Kernel）与无穷维特征空间}
工程中最常用的非线性核当属\textbf{高斯径向基核（RBF / Gaussian Kernel）}：
\[
  K_{\mathrm{RBF}}(x, z) := \exp\!\big(-\gamma \norm{x - z}_2^2\big), \qquad \gamma = \frac{1}{2\sigma^2} > 0.
\]
它不仅合法，而且隐式对应着一个\textbf{无穷维}特征空间！下面给出零跳步证明。

\begin{thmBox}
\begin{theorem}[RBF 核的无穷维特征映射展开]\label{thm:rbf-inf}
对任意 $\gamma > 0$ 与 $x, z \in \R^d$，高斯核 $K_{\mathrm{RBF}}(x, z) = \exp(-\gamma\norm{x - z}_2^2)$ 可表示为平方可和序列空间 $\ell_2$ 中的内积 $\inner{\phi(x)}{\phi(z)}_{\ell_2}$。
\end{theorem}
\end{thmBox}

\begin{proof}
先展开指数中的欧氏距离平方 $\norm{x - z}_2^2 = \norm{x}_2^2 + \norm{z}_2^2 - 2 x^\top z$：
\[
  \exp\!\big(-\gamma\norm{x - z}_2^2\big)
  = \exp\!\big(-\gamma\norm{x}_2^2\big) \exp\!\big(-\gamma\norm{z}_2^2\big) \exp\!\big(2\gamma\, x^\top z\big).
\]
先看一维情形（$d = 1$）。对 $\exp(2\gamma x z)$ 作全域收敛的泰勒级数展开：
\[
  \exp\!\big(-\gamma(x - z)^2\big)
  = e^{-\gamma x^2} e^{-\gamma z^2} \sum_{k=0}^{\infty} \frac{(2\gamma)^k}{k!} x^k z^k
  = \sum_{k=0}^{\infty} \underbrace{\left(e^{-\gamma x^2}\sqrt{\frac{(2\gamma)^k}{k!}}\, x^k\right)}_{\phi_k(x)}
                        \underbrace{\left(e^{-\gamma z^2}\sqrt{\frac{(2\gamma)^k}{k!}}\, z^k\right)}_{\phi_k(z)}.
\]
且 $\norm{\phi(x)}_{\ell_2}^2 = \sum_{k=0}^{\infty} \phi_k(x)^2 = e^{-2\gamma x^2} e^{2\gamma x^2} = 1 < \infty$，故 $\phi(x) \in \ell_2$ 严格良定义！

对一般 $d$ 维情形，利用指数函数把和化为积的性质 $\exp(2\gamma \sum_{m=1}^d x_m z_m) = \prod_{m=1}^d \exp(2\gamma x_m z_m)$，对每一维展开泰勒级数并用多重指标 $\alpha = (\alpha_1, \dots, \alpha_d) \in \N^d$ 合并乘积：
\[
  K_{\mathrm{RBF}}(x, z)
  = \sum_{\alpha \in \N^d} \phi_\alpha(x)\,\phi_\alpha(z),
  \qquad
  \phi_\alpha(x) := e^{-\gamma \norm{x}_2^2} \prod_{m=1}^{d} \sqrt{\frac{(2\gamma)^{\alpha_m}}{\alpha_m!}}\, x_m^{\alpha_m}.
\]
同样易验 $\sum_{\alpha \in \N^d} \phi_\alpha(x)^2 = K_{\mathrm{RBF}}(x, x) = 1 < \infty$。因此 RBF 核等价于把样本映射到了包含\textbf{全部阶次多项式特征}的无穷维单位球面上！
\end{proof}

\subsection{什么样的二元函数是合法核？Gram 矩阵半正定性与闭包法则}

并非随便写一个对称函数 $K(x, z)$ 都能塞进对偶 QP~\eqref{eq:svm-kernel-dual}。如果核矩阵 $K_{ij} = K(x_i, x_j)$ 不是半正定的，二次项 $-\frac{1}{2}\alpha^\top (Y K Y) \alpha$ 就不再是凹函数，对偶问题会失去凸性甚至趋于 $+\infty$！

\begin{thmBox}
\begin{theorem}[Mercer 条件 / Gram 矩阵半正定判据]\label{thm:mercer}
对称函数 $K: \mathcal{X} \times \mathcal{X} \to \R$ 是某个希尔伯特空间中内积 $K(x, z) = \inner{\phi(x)}{\phi(z)}$ 的\textbf{充要条件}（在离散有限点集层面亦称正定核判据）为：对任意有限个样本点 $x_1, \dots, x_N \in \mathcal{X}$，其所诱导的 $N \times N$ \textbf{Gram 矩阵} $K_{ij} := K(x_i, x_j)$ 都是对称半正定矩阵（$K \succeq 0$），即对任意 $c \in \R^N$，
\begin{equation}\label{eq:gram-psd}
  c^\top K c = \sum_{i=1}^{N}\sum_{j=1}^{N} c_i c_j K(x_i, x_j) \ge 0.
\end{equation}
\end{theorem}
\end{thmBox}

必要性一行即证：若 $K(x_i, x_j) = \inner{\phi(x_i)}{\phi(x_j)}$，则
\[
  c^\top K c = \sum_{i,j} c_i c_j \inner{\phi(x_i)}{\phi(x_j)} = \norm{\sum_{i=1}^{N} c_i \phi(x_i)}_{\mathcal{H}}^2 \ge 0.
\]
充分性由 Moore--Aronszajn 再生核希尔伯特空间（RKHS）构造定理保证（见本节末拓展框）。在实践中，我们很少从头验证 Mercer 积分条件，而是直接利用\textbf{合法核的代数闭包法则}像搭积木一样组合复杂核：
\begin{itemize}
  \item \textbf{正数倍与加法}：若 $K_1, K_2$ 为合法核，$a, b \ge 0$，则 $a K_1 + b K_2$ 合法（因 $c^\top(a K_1 + b K_2)c = a c^\top K_1 c + b c^\top K_2 c \ge 0$，特征映射为加权直和 $(\sqrt{a}\phi_1, \sqrt{b}\phi_2)$）；
  \item \textbf{逐点乘积}：$K_1(x, z) K_2(x, z)$ 合法（对应 Gram 矩阵的 Hadamard 积 $K_1 \had K_2 \succeq 0$，即舒尔乘积定理，特征映射为张量积 $\phi_1(x) \otimes \phi_2(x)$）；
  \item \textbf{函数外积}：对任意实值函数 $h: \mathcal{X} \to \R$，$h(x) K_1(x, z) h(z)$ 合法（如 RBF 核前后的 $e^{-\gamma\norm{x}_2^2}$ 因子！）；
  \item \textbf{正系数幂级数}：若 $K$ 合法且 $a_k \ge 0$，则 $\sum_{k=0}^{\infty} a_k K(x, z)^k$（在收敛域内，特例包括 $\exp(K(x, z))$）仍为合法核！
\end{itemize}

\subsection{同心圆非线性升维几何与实测验证}

图~\ref{fig:circle-lift} 直观呈现了阿达马题记背后的升维魔术：二维平面上线性不可分的同心圆，只需增加第三维径向特征 $x_3 = x_1^2 + x_2^2$ 映射到三维抛物面上，原本纠缠的内外两圈立刻在高度轴上泾渭分明，一张水平平面 $x_3 = r_0^2$ 便一刀两断！

\begin{figure}[htbp]
\centering
\resizebox{0.98\linewidth}{!}{%
\begin{tikzpicture}
  % Left panel: 2D concentric circles (non-linearly separable)
  \begin{scope}[xshift=-4.7cm]
    \draw[rounded corners=2mm, draw=tealLine, fill=codebg!35] (-2.75,-2.35) rectangle (2.75,2.75);
    \node[font=\sffamily\bfseries\small, text=tealDark] at (0, 2.38) {原输入空间 $\R^2$：线性不可分};
    \draw[->, >=Stealth, slate!70, thin] (-2.1, -0.05) -- (2.15, -0.05) node[right, font=\scriptsize] {$x_1$};
    \draw[->, >=Stealth, slate!70, thin] (0, -1.80) -- (0, 1.85) node[right, font=\scriptsize] {$x_2$};
    % Dashed nonlinear boundary circle
    \draw[gold, thick, dashed] (0, -0.05) circle (1.02cm);
    % Inner circle points (ochre, y = -1)
    \foreach \ang/\r in {15/0.55, 55/0.68, 95/0.48, 138/0.72, 182/0.60, 225/0.66, 268/0.52, 310/0.70, 350/0.44, 75/0.35, 200/0.32} {
      \fill[ochre] ({\r*cos(\ang)}, {-0.05+\r*sin(\ang)}) circle (2.2pt);
    }
    % Outer ring points (teal, y = +1)
    \foreach \ang/\r in {10/1.45, 38/1.55, 68/1.38, 102/1.52, 132/1.42, 165/1.56, 198/1.40, 230/1.52, 260/1.46, 292/1.36, 322/1.54, 348/1.42} {
      \fill[teal] ({\r*cos(\ang)}, {-0.05+\r*sin(\ang)}) circle (2.2pt);
    }
    \node[font=\scriptsize, text=slate] at (0, -2.06) {任何直线都无法分开内圈与外圈};
  \end{scope}

  % Center lifting arrow
  \draw[->, >=Stealth, gold, very thick] (-1.65, 0.15) -- (1.55, 0.15)
    node[midway, above=2pt, font=\scriptsize\bfseries, text=navy, align=center] {特征升维 $\phi$\\$(x_1, x_2) \mapsto$\\$(x_1, x_2, x_1^2+x_2^2)$};

  % Right panel: 3D lifted paraboloid + separating horizontal plane
  \begin{scope}[xshift=4.75cm]
    \draw[rounded corners=2mm, draw=tealLine, fill=tealLight!35] (-2.85,-2.35) rectangle (2.85,2.75);
    \node[font=\sffamily\bfseries\small, text=tealDark] at (0, 2.38) {高维特征空间 $\R^3$：线性可分！};
    % Axes
    \draw[->, >=Stealth, slate!70, thin] (0, -1.65) -- (0, 1.88) node[right, font=\scriptsize] {$x_3 = x_1^2+x_2^2$};
    \draw[->, >=Stealth, slate!70, thin] (-2.1, -1.65) -- (2.1, -1.65) node[right, font=\scriptsize] {$x_1$};
    % Paraboloid silhouette
    \draw[slate!45, thick] plot[domain=-1.72:1.72, samples=60] (\x, {-1.55 + 0.96*\x*\x});
    % Separating horizontal hyperplane
    \fill[gold!25, opacity=0.78] (-2.25, -0.62) -- (1.45, -0.62) -- (2.25, -0.05) -- (-1.45, -0.05) -- cycle;
    \draw[gold, very thick] (-2.25, -0.62) -- (1.45, -0.62) -- (2.25, -0.05) -- (-1.45, -0.05) -- cycle;
    \node[font=\scriptsize\bfseries, text=ochre!90!black] at (1.72, -0.84) {分离超平面};
    % Inner points sitting at bottom of paraboloid (below plane)
    \foreach \x/\y in {-0.72/-1.05, -0.45/-1.34, -0.15/-1.50, 0.22/-1.48, 0.55/-1.24, 0.78/-0.98, 0.02/-1.36, -0.35/-1.18, 0.38/-1.15} {
      \fill[ochre] (\x, \y) circle (2.2pt);
    }
    % Outer points lifted high above the plane
    \foreach \x/\y in {-1.58/0.85, -1.35/0.35, -1.12/0.55, -0.85/0.95, -0.42/0.68, 0.15/0.82, 0.65/0.92, 1.08/0.52, 1.38/0.40, 1.56/0.82} {
      \fill[teal] (\x, \y) circle (2.2pt);
    }
    \node[font=\scriptsize, text=tealDark] at (0, -2.06) {水平超平面干净切开两类，投影回 $\R^2$ 即圆环边界};
  \end{scope}
\end{tikzpicture}}
\caption{核升维的几何本质。原平面上线性不可分的同心圆数据经 $\phi(x_1, x_2) = (x_1, x_2, x_1^2+x_2^2)$ 映射到三维抛物面后，被一张水平超平面线性分开；RBF 核则隐式地把数据抬升至无穷维空间自动雕刻任意平滑非线性边界。}
\label{fig:circle-lift}
\end{figure}

现在我们检阅配套脚本 \texttt{code/\allowbreak vol4\_experiments.py} 中 \texttt{exp4\_rbf\_svm\_circles} 的实测数据。实验在含径向高斯噪声的二维同心圆数据集上运行（训练集 $N_{\text{tr}} = 160$，测试集 $N_{\text{te}} = 80$；内圆半径 $r \sim \mathcal{N}(1.0, 0.28^2)$，外环半径 $r \sim \mathcal{N}(2.0, 0.30^2)$，两类在 $r \approx 1.5$ 处存在真实的径向重叠），对比线性 SVM 与固定 $C = 5.0$、不同核宽参数 $\gamma$ 的纯 NumPy 对偶 RBF-SVM（由 SMO 坐标上升求解），精确结果见表~\ref{tab:svm-circles}。

\begin{table}[htbp]
\centering
\caption{含噪同心圆数据集（$N_{\text{tr}}=160, N_{\text{te}}=80$）上线性 SVM 与对偶 RBF-SVM（$C=5.0$）的实测对比}
\label{tab:svm-circles}
\small
\begin{tabular}{@{}lcccc@{}}
\toprule
\textbf{模型与核参数} & \textbf{支持向量数 / 总数} & \textbf{支持向量占比} & \textbf{训练准确率} & \textbf{测试准确率} \\
\midrule
线性 SVM（原始次梯度） & --- & --- & --- & $46.25\%$ \\
RBF-SVM（$\gamma = 0.1$） & $57 / 160$ & $35.6\%$ & $93.75\%$ & $\mathbf{96.25\%}$ \\
RBF-SVM（$\gamma = 0.8$） & $47 / 160$ & $29.4\%$ & $96.25\%$ & $95.00\%$ \\
RBF-SVM（$\gamma = 12.0$） & $124 / 160$ & $77.5\%$ & $98.75\%$ & $92.50\%$ \\
\bottomrule
\end{tabular}
\end{table}

表~\ref{tab:svm-circles} 的每一行都是统计学习理论的活教科书：
\begin{enumerate}
  \item \textbf{线性模型彻底失效}：在对称同心圆上，任何直线都只能像切西瓜一样同时切穿内外圈，线性 SVM 测试准确率仅 \textbf{46.25\%}（接近随机猜测的 50\%）；
  \item \textbf{适中核宽大获全胜}：当 $\gamma = 0.1$ 或 $0.8$ 时，RBF 核仅靠约 $30\%\text{--}36\%$ 的边界支持向量便勾勒出圆润的环形决策边界，测试准确率一跃升至 \textbf{96.25\%}（$\gamma=0.1$）与 \textbf{95.00\%}（$\gamma=0.8$）；
  \item \textbf{过大 $\gamma$ 诱发经典过拟合}：回想 $K_{\mathrm{RBF}}(x, z) = \exp(-\gamma\norm{x-z}_2^2)$ 的物理意义——$\gamma$ 越大，每个样本周围的高斯“小山包”越尖锐、辐射半径越短。当 $\gamma = 12.0$ 时，支持向量暴增至 \textbf{124/160（77.5\%）}——模型几乎为每个训练样本都单独筑起一座孤岛，训练准确率飙升至 $98.75\%$，但测试准确率反降到 $92.50\%$！这就是“支持向量过多即过拟合警报”的直接实证。
\end{enumerate}

\subsection{可运行的纯 NumPy 对偶 RBF-SVM（SMO 算法）实现}

与普通梯度下降不同，对偶问题~\eqref{eq:svm-soft-dual} 带有等式约束 $\sum_i \alpha_i y_i = 0$，不能只更新单个 $\alpha_i$。Platt 于 1998 年提出的\textbf{序列最小优化（Sequential Minimal Optimization, SMO）}巧妙地\textbf{每次同时选取一对乘子 $(\alpha_i, \alpha_j)$}，在保持 $\alpha_i y_i + \alpha_j y_j = \text{const}$ 的前提下解析求出二元二次函数的闭式最优步长并裁剪到盒式区间 $[L, H]$。下面的纯 NumPy 实现直接复现表~\ref{tab:svm-circles} 的全部结果。

\begin{lstlisting}[caption={纯 NumPy 对偶 RBF-SVM 与简化 SMO 坐标上升求解器（复现表~\ref{tab:svm-circles}）},label={lst:svm-smo}]
import numpy as np

def rbf_kernel(A, B, gamma=0.8):
    d2 = np.sum(A**2, axis=1, keepdims=True) + np.sum(B**2, axis=1) - 2.0 * (A @ B.T)
    return np.exp(-gamma * np.maximum(d2, 0.0))

def fit_rbf_svm_smo(X, y, C=5.0, gamma=0.8, epochs=35, seed=2026):
    rng = np.random.default_rng(seed)
    n = len(y)
    K = rbf_kernel(X, X, gamma=gamma)
    alpha = np.zeros(n)
    b = 0.0
    for _ in range(epochs):
        for i in range(n):
            Ei = np.dot(alpha * y, K[:, i]) + b - y[i]
            # Check KKT violation for alpha_i
            if (y[i] * Ei < -1e-3 and alpha[i] < C) or (y[i] * Ei > 1e-3 and alpha[i] > 0):
                j = (i + 1 + int(rng.integers(0, n - 1))) % n
                Ej = np.dot(alpha * y, K[:, j]) + b - y[j]
                ai_old, aj_old = alpha[i], alpha[j]
                if y[i] != y[j]:
                    L_b, H_b = max(0.0, aj_old - ai_old), min(C, C + aj_old - ai_old)
                else:
                    L_b, H_b = max(0.0, ai_old + aj_old - C), min(C, ai_old + aj_old)
                if L_b == H_b:
                    continue
                eta = 2.0 * K[i, j] - K[i, i] - K[j, j]
                if eta >= 0.0:
                    continue
                alpha[j] = np.clip(aj_old - y[j] * (Ei - Ej) / eta, L_b, H_b)
                if abs(alpha[j] - aj_old) < 1e-5:
                    continue
                alpha[i] = ai_old + y[i] * y[j] * (aj_old - alpha[j])
                b1 = b - Ei - y[i]*(alpha[i]-ai_old)*K[i,i] - y[j]*(alpha[j]-aj_old)*K[i,j]
                b2 = b - Ej - y[i]*(alpha[i]-ai_old)*K[i,j] - y[j]*(alpha[j]-aj_old)*K[j,j]
                if 0.0 < alpha[i] < C:   b = b1
                elif 0.0 < alpha[j] < C: b = b2
                else:                    b = 0.5 * (b1 + b2)
    return alpha, b

def predict_rbf_svm(X_new, X_tr, y_tr, alpha, b, gamma=0.8):
    K_new = rbf_kernel(X_new, X_tr, gamma=gamma)
    return np.sign(K_new @ (alpha * y_tr) + b)
\end{lstlisting}

\begin{extBox}{表示定理（Representer Theorem）、RKHS 直觉与 SMO 启发式}
为什么核 SVM 的最优预测函数必定形如有限线性组合 $f^*(x) = \sum_{i=1}^N c_i K(x_i, x) + b$，即使特征空间 $\mathcal{H}$ 是无穷维的？这背后是泛函分析中极具威力的\textbf{表示定理（Representer Theorem）}：设 $\mathcal{H}_K$ 是由正定核 $K$ 唯一诱导的\textbf{再生核希尔伯特空间（RKHS）}，满足再生性质 $\inner{f}{K(x, \cdot)}_{\mathcal{H}_K} = f(x)$。将任意 $f \in \mathcal{H}_K$ 正交分解为 $f = f_\parallel + f_\perp$，其中 $f_\parallel \in \operatorname{span}\{K(x_1,\cdot), \dots, K(x_N,\cdot)\}$ 而 $f_\perp$ 与该子空间正交。由再生性质，$f_\perp(x_i) = \inner{f_\perp}{K(x_i,\cdot)} = 0$，故 $f_\perp$ 对所有样本的经验损失毫无影响，却严格增大了正则项 $\norm{f}_{\mathcal{H}_K}^2 = \norm{f_\parallel}^2 + \norm{f_\perp}^2$！因此任何单调递增正则化的经验风险极小化解必有 $f_\perp = 0$，即无穷维优化问题自动坍缩到了 $N$ 维有限子空间中。

此外，工业级 SVM 求解器（如 \texttt{LIBSVM}）在 SMO 中并非随机挑选第二个变量 $j$，而是优先挑选使 $|E_i - E_j|$ 最大的“最大步长配对”并缓存核矩阵高频列，将计算耗时再压缩一个数量级；不过当样本量 $N > 10^5$ 时，$O(N^2)$ 核矩阵仍会成为瓶颈，此时可采用随机傅里叶特征（Random Fourier Features）显式近似 RBF 核或改用线性 SVM/树模型。
\end{extBox}

\begin{insightBox}{核方法在深度学习时代的当代回声}
核方法从未过时，它以全新的面貌渗透进了现代深度学习：
\begin{itemize}
  \item \textbf{Transformer 注意力即“可学习的数据相关核”}：第 5 卷将揭示，自注意力权重 $\operatorname{softmax}(QK^\top / \sqrt{d_k})$ 本质上是以 $\exp(q_i^\top k_j / \sqrt{d_k})$ 为非归一化核函数、动态学习投影矩阵 $W_Q, W_K$ 的核平滑器；
  \item \textbf{神经切线核（Neural Tangent Kernel, NTK）}：在无限宽极限下，梯度下降训练深度神经网络的动力学精确等价于使用固定核函数 $K_{\mathrm{NTK}}(x, z) = \inner{\nabla_\theta f_{\theta_0}(x)}{\nabla_\theta f_{\theta_0}(z)}$ 做核岭回归！理解了本章的核技巧，便拥有了透视现代神经网络泛化理论的 X 光机。
\end{itemize}
\end{insightBox}

% ============================================================
\section{本章习题}\label{sec:ch4-ex}
% ============================================================

\exTier{A 层：基础巩固}{掌握几何间隔、规范超平面与基本核函数判别}

\begin{enumerate}[label=\textbf{4.\arabic*}]
  \item \textbf{规范超平面与间隔计算}：给定二维二分类超平面 $3x_1 - 4x_2 + 5 = 0$，以及四个已正确分类的样本：$x_1=(1, 1)^\top\ (y_1=+1)$，$x_2=(0, 2)^\top\ (y_2=-1)$，$x_3=(3, 1)^\top\ (y_3=+1)$，$x_4=(-1, 3)^\top\ (y_4=-1)$。（1）分别计算这四个样本的函数间隔与几何间隔；（2）将 $(w, b)$ 等比例缩放使其成为规范超平面（最小函数间隔恰为 $1$），写出缩放后的 $(w_{\mathrm{can}}, b_{\mathrm{can}})$ 及对应的隔离带总宽度。
  \exHint{先算出 $\hat{\gamma} = \min_i y_i(w^\top x_i + b)$，再将 $(w, b)$ 同除以 $\hat{\gamma}$ 即可，检查几何间隔是否保持不变。}

  \item \textbf{合法核函数的快速甄别}：已知 $K_1(x, z)$ 与 $K_2(x, z)$ 是 $\R^d$ 上的合法正定核函数，逐一判断下列函数是否必为合法核，若成立请简述依据，若不成立请给出显式反例：
  （1）$K(x, z) = 3 K_1(x, z) + 2 K_2(x, z)$；（2）$K(x, z) = K_1(x, z) - K_2(x, z)$；（3）$K(x, z) = \big(K_1(x, z) + 1\big)^3$；（4）$K(x, z) = \cos(x - z)$（$x, z \in \R$）。
  \exHint{第 (4) 问可将 $\cos(x - z)$ 用三角和差公式展开为两项乘积之和，一眼看出显式二维特征映射 $\phi(x)$！}
\end{enumerate}

\exTier{B 层：方法与推导}{熟练手解小型对偶 QP 与分析软间隔 KKT 状态}

\begin{enumerate}[label=\textbf{4.\arabic*}, resume]
  \item \textbf{一维非线性可分数据的二次核手算}：一维轴上有 3 个样本：$x_1 = -1\ (y_1 = +1)$，$x_2 = 0\ (y_2 = -1)$，$x_3 = 1\ (y_3 = +1)$，显然在原空间 $\R^1$ 线性不可分。使用特征映射 $\phi(x) = (x, x^2)^\top$（对应核函数 $K(x, z) = xz + x^2 z^2$）：（1）写出 $3\times 3$ 核矩阵 $K$；（2）手解硬间隔对偶 QP~\eqref{eq:svm-kernel-dual}，求出精确的 $\alpha_1^*, \alpha_2^*, \alpha_3^*$ 与 $b^*$；（3）写出最终判别函数 $f(x)$ 并求出原一维轴上的两个决策边界点。
  \exHint{利用对称性 $\alpha_1 = \alpha_3$ 以及等式约束 $\alpha_1 + \alpha_3 = \alpha_2$，整个对偶问题会退化为一元二次函数极值！}

  \item \textbf{软间隔支持向量状态诊断}：在某个 $C = 2.0$ 的软间隔线性 SVM 求解完毕后，得到 $w^* = (1, 0)^\top, b^* = -1$。现有四个训练样本：$A:(3, 1)^\top, y=+1$；$B:(2, -2)^\top, y=+1$；$C:(1.5, 0)^\top, y=+1$；$D:(0.5, 2)^\top, y=+1$。（1）计算各样本的函数间隔与最优松弛变量 $\xi_i^*$；（2）由 KKT 互补松弛条件判断各样本对应的对偶乘子 $\alpha_i^*$ 处于哪一区间（$=0$、$(0, C)$ 还是 $=C$）？哪几个样本可用于求解 $b^*$？
  \exHint{对照第~\ref{sec:svm-soft} 节的三类乘子刻画，分别检验 $y_i f^*(x_i) > 1$、$=1$ 与 $<1$ 三种情形。}
\end{enumerate}

\exTier{C 层：理论证明与反例构造}{深化对偶目标几何含义、核性质与正则化联系}

\begin{enumerate}[label=\textbf{4.\arabic*}, resume]
  \item \textbf{支持向量乘子之和与最大间隔的精确恒等式}：对于线性可分的硬间隔 SVM，证明最优对偶变量满足如下优雅的几何恒等式：
  \[
    \sum_{i=1}^{N} \alpha_i^* = \norm{w^*}_2^2 = \frac{1}{(\gamma^*)^2},
  \]
  其中 $\gamma^* = 1/\norm{w^*}_2$ 为最大几何间隔。进而说明为何数据类别越难以分开（间隔 $\gamma^*$ 越小），对偶变量总和必然越大。
  \exHint{同时利用强对偶等式（原始最优值 $\frac{1}{2}\norm{w^*}_2^2$ 等于对偶最优值 $\sum_i \alpha_i^* - \frac{1}{2}\norm{w^*}_2^2$），一行即可证毕！}

  \item \textbf{多项式核的特征维度与舒尔乘积合法性证明}：（1）证明任意两个对称半正定矩阵 $A, B \succeq 0$ 的逐元素乘积（Hadamard 积）$A \had B$ 依然是半正定矩阵（提示：将 $A$ 作谱分解 $A = \sum_r \lambda_r u_r u_r^\top$，计算 $c^\top (A \had B) c$）；（2）据此证明：若 $K_1, K_2$ 为合法正定核，则 $K_1(x, z) K_2(x, z)$ 必为合法正定核，并推导出 $K(x, z) = (x^\top z + c)^p$（$c \ge 0, p \in \N^+$）对任意维度 $d$ 均合法。
  \exHint{将 $c^\top ((u_r u_r^\top) \had B) c$ 改写为 $(c \had u_r)^\top B (c \had u_r)$，再利用 $B \succeq 0$ 与 $\lambda_r \ge 0$。}
\end{enumerate}

\exTier{D 层：编程实现与实验验证}{亲手验证 SMO 求解器、次梯度与核参数敏感性}

\begin{enumerate}[label=\textbf{4.\arabic*}, resume]
  \item \textbf{原始次梯度求解线性软间隔 SVM 与对偶解一致性校验}：基于式~\eqref{eq:svm-hinge} 的合页损失目标函数，写出 $L_2$ 正则化合页损失对 $(w, b)$ 的次梯度公式，用纯 NumPy 实现批量次梯度下降（步长 $\eta_t = \eta_0 / \sqrt{t+1}$）；在一个二维线性可分加少量噪声的数据集上（$N=100$），分别运行你的原始次梯度算法与代码~\ref{lst:svm-smo} 中线性核版 SMO（$K(x, z) = x^\top z$），验证两者得到的 $(w^*, b^*)$ 与目标函数值相对误差小于 $1\%$。
  \exHint{对合页项 $\max(0, 1 - y_i(w^\top x_i + b))$，当 $y_i(w^\top x_i + b) < 1$ 时关于 $w$ 的次梯度为 $-y_i x_i$，$>1$ 时为 $0$，恰等于 $1$ 时可取 $0$。}

  \item \textbf{RBF-SVM 的 $(C, \gamma)$ 二维网格搜索与支持向量热力图}：复用 \texttt{code/\allowbreak vol4\_experiments.py} 中的含噪同心圆数据生成代码与代码~\ref{lst:svm-smo}，在 $C \in \{0.5, 5.0, 50.0\}$ 与 $\gamma \in \{0.05, 0.2, 0.8, 4.0, 12.0\}$ 的 $3 \times 5$ 网格上系统记录：（1）测试集准确率；（2）支持向量数量占比；（3）训练与测试准确率之差（过拟合间隙）。结合偏差-方差权衡分析为何“高 $C$ + 高 $\gamma$”是最危险的过拟合角落。
  \exHint{注意固定随机种子并在每次实验前重置 \texttt{alpha} 与 \texttt{b}，可将结果整理为 $3\times 5$ 对照矩阵并与表~\ref{tab:svm-circles} 交叉核验。}
\end{enumerate}

\chapter[决策树、随机森林与 GBDT]{决策树、随机森林与梯度提升树（GBDT）}\label{ch:trees}
\chapEpigraph{整体大于部分之和。}{Aristotle（亚里士多德，《形而上学》）}

\begin{chapIntro}{本章导读}
第 4 章的核支持向量机通过全局平滑的相似性度量雕刻非线性边界，但它对特征量纲极度敏感，且面对混合类型、含大量噪声特征的表格数据时往往捉襟见肘。本章转向一套截然不同的建模哲学——\textbf{轴对齐递归空间划分（决策树）}，并在此基础上展开机器学习史上最富成果的\textbf{集成学习（Ensemble Learning）}双峰并峙：
\begin{itemize}
  \item \textbf{Bagging 与随机森林}：并行训练大批高方差深树并平均，精准打击第 3 章偏差-方差分解中的\textbf{方差项}；
  \item \textbf{梯度提升树（GBDT / XGBoost）}：串行拟合前序模型的函数空间负梯度（伪残差），逐步消融弱学习器的\textbf{偏差项}。
\end{itemize}
即便在深度学习席卷图像与语言的今天，GBDT 家族依然是表格数据预测当之无愧的王者。
\end{chapIntro}

% ============================================================
\section{CART 决策树：轴对齐递归划分与不纯度度量}\label{sec:trees-cart}
% ============================================================

\subsection{为什么轴对齐划分在表格数据上如此锋利？}

一棵\textbf{分类与回归树（Classification and Regression Tree, CART）}将输入空间 $\R^d$ 递归地用一系列垂直于坐标轴的阈值切面 $x^{(j)} \le s$ 切割为 $M$ 个互不相交的超矩形叶区域 $R_1, R_2, \dots, R_M$，在每个叶区域上输出一个常数值 $c_m$：
\begin{equation}\label{eq:cart-model}
  T(x) = \sum_{m=1}^{M} c_m \,\mathbb{I}(x \in R_m).
\end{equation}
与线性模型或 RBF-SVM 相比，这种看似朴素的“逐维问答机”在现实表格数据（金融风控、医疗体检、推荐点击率）上拥有四项先天优势：
\begin{enumerate}
  \item \textbf{单调变换不变性与零标准化需求}：决策树只关心每个特征的\textbf{排序位置}（$x^{(j)} \le s$），对任何严格单调递增变换（如取对数、乘以 $1000$、Z-score 缩放）完全免疫！
  \item \textbf{天然捕捉高阶特征交互}：沿着树根到叶子的路径“$(\text{年龄}>50) \land (\text{收缩压}>150) \land (\text{血糖}>7.5)$”，本质上自动构造了局部高阶逻辑乘积，无需人工手写交叉特征；
  \item \textbf{对冗余噪声特征具有鲁棒性}：每次分裂只挑选信息收益最大的单一特征，完全不相关的噪声列很少被选中（第~\ref{sec:trees-exp} 节实验将直观验证这一点）；
  \item \textbf{自然兼容异构特征与缺失值}：连续值、有序离散值与类别特征可统一处理，且单棵浅树具备“白盒规则”级可解释性。
\end{enumerate}

\subsection{回归树的最优叶预测与方差缩减准则}

在回归任务中，给定落入某一叶区域 $R_m$ 的训练样本下标集 $I_m = \{i \mid x_i \in R_m\}$（记样本数 $n_m = |I_m|$），采用均方误差（平方损失）衡量叶内拟合误差。

\begin{thmBox}
\begin{proposition}[回归树叶节点的最优常数预测与贪心分裂目标]\label{prop:cart-reg}
在平方损失下：
\begin{enumerate}
  \item 对任意固定的叶区域 $R_m$，使叶内残差平方和 $\operatorname{RSS}(R_m, c) := \sum_{i \in I_m}(y_i - c)^2$ 最小的唯一常数预测为该区域内标签的\textbf{样本均值}：
  \[
    c_m^* = \bar{y}_{R_m} := \frac{1}{n_m}\sum_{i \in I_m} y_i,
  \]
  此时该叶节点的最小残差平方和恰等于 $n_m \Var(y_{R_m}) := \sum_{i \in I_m}(y_i - \bar{y}_{R_m})^2$。
  \item 当考虑用特征 $j$ 与阈值 $s$ 将当前节点 $R$ 一分为二为 $R_L(j, s) = \{x \in R \mid x^{(j)} \le s\}$ 与 $R_R(j, s) = \{x \in R \mid x^{(j)} > s\}$（样本数分别为 $n_L, n_R$）时，最优贪心分裂等价于最小化\textbf{加权叶内方差}：
  \begin{equation}\label{eq:cart-reg-split}
    (j^*, s^*) = \argmin_{j,\, s}\; \Big[ n_L \Var(y_L) + n_R \Var(y_R) \Big].
  \end{equation}
\end{enumerate}
\end{proposition}
\end{thmBox}

\begin{proof}
将 $g(c) := \sum_{i \in I_m}(y_i - c)^2$ 关于 $c$ 求一阶导与二阶导：
\[
  g'(c) = -2\sum_{i \in I_m}(y_i - c) = 0 \implies c_m^* = \frac{1}{n_m}\sum_{i \in I_m} y_i = \bar{y}_{R_m}, \qquad
  g''(c) = 2 n_m > 0.
\]
二阶导严格为正，故样本均值是全局唯一极小值点。将其代回即得最小残差平方和 $\sum_{i \in I_m}(y_i - \bar{y}_{R_m})^2 = n_m \Var(y_{R_m})$（这里 $\Var(y_R)$ 指经验总体方差 $\frac{1}{n_R}\sum (y_i - \bar{y}_R)^2$）。对分裂后的两个子节点分别代入各自最优均值预测，分裂后的总残差平方和即为 $n_L \Var(y_L) + n_R \Var(y_R)$，最小化它等价于最大化总方差缩减量 $n\Var(y_R) - [n_L \Var(y_L) + n_R \Var(y_R)]$。
\end{proof}

由于寻找全局最优树划分是 NP-难问题，CART 采用自顶向下的\textbf{贪心递归二分法}：在每个节点遍历全部候选特征 $j$ 与候选切分点 $s$，按式~\eqref{eq:cart-reg-split} 选出当前一步最优切分，然后分别对左右子节点递归建树，直至满足停止条件。

\subsection{分类树：信息增益、Gini 不纯度与二者的解析联系}

对于 $K$ 分类任务（标签 $y_i \in \{1, \dots, K\}$），设某节点 $R$ 中各类样本的经验占比为 $p = (p_1, \dots, p_K)$（满足 $\sum_{k=1}^K p_k = 1$）。若所有样本同属一类（某 $p_k = 1$），则节点完全“纯净”；若各类均匀混杂（$p_k = 1/K$），则节点最“浑浊”。衡量节点浑浊度的两大主流指标是\textbf{香农熵}与\textbf{Gini 不纯度}。

\begin{recallBox}{香农熵与信息增益（第 2 卷回顾）}
离散分布 $p = (p_1, \dots, p_K)$ 的香农熵定义为 $H(p) := -\sum_{k=1}^K p_k \log_2 p_k$（约定 $0\log_2 0 = 0$）。在决策树（ID3/C4.5）中，以切分 $(j, s)$ 将样本集 $D$ 分为 $D_L, D_R$ 后的\textbf{信息增益（Information Gain）}等于分裂前后条件熵的减少量：
\[
  \operatorname{IG}(D;\, j, s) := H(D) - \frac{n_L}{n} H(D_L) - \frac{n_R}{n} H(D_R).
\]
\end{recallBox}

CART 分类树则默认采用不含对数运算、计算更迅捷的\textbf{Gini 不纯度（Gini Impurity）}：
\begin{equation}\label{eq:gini-def}
  G(p) := \sum_{k=1}^{K} p_k (1 - p_k) = 1 - \sum_{k=1}^{K} p_k^2.
\end{equation}
相应地，分裂 $(j, s)$ 的 Gini 增益（不纯度下降量）为 $\Delta G(D;\, j, s) := G(D) - \frac{n_L}{n} G(D_L) - \frac{n_R}{n} G(D_R)$。

\begin{thmBox}
\begin{proposition}[Gini 不纯度的概率意义及其与香农熵的泰勒近似关系]\label{prop:gini-entropy}
对任意概率分布 $p = (p_1, \dots, p_K)$：
\begin{enumerate}
  \item \textbf{随机标签分类器误差解释}：从该节点按分布 $p$ 独立抽取一个样本，并按同一分布 $p$ 随机赋予预测标签，误分类概率恰好等于 $\sum_{k=1}^K p_k(1 - p_k) = G(p)$；
  \item \textbf{与自然对数熵的泰勒联系}：记自然对数香农熵 $H_e(p) := -\sum_{k=1}^K p_k \ln p_k$。将 $-\ln t$ 在 $t = 1$ 处作一阶展开 $-\ln t = (1 - t) + \frac{1}{2}(1-t)^2 + \cdots$，截取首项即精准还原 Gini 不纯度，且处处满足 $G(p) \le H_e(p)$；在二分类 $p = (\frac{1}{2}+\delta, \frac{1}{2}-\delta)$ 对称点附近作二阶展开，有
  \[
    H_e(p) = \ln 2 - 2\delta^2 + O(\delta^4), \qquad G(p) = \frac{1}{2} - 2\delta^2,
  \]
  两者的二次曲率完全一致！
  \item \textbf{严格凹性与非负增益}：$G(p)$ 与 $H_e(p)$ 均是概率单纯形上的严格凹函数；由 Jensen 不等式，只要分裂后的左右子节点类别分布不完全相同（$p_L \neq p_R$），就必然严格有 $\Delta G > 0$ 与 $\operatorname{IG} > 0$。
\end{enumerate}
\end{proposition}
\end{thmBox}

\begin{proof}
(1) 抽中真实类别 $k$ 的概率为 $p_k$，独立随机赋标签却未抽中 $k$ 的概率为 $1 - p_k$，全概率求和即 $\sum_k p_k(1 - p_k) = G(p)$。
(2) 由基本不等式 $\ln t \le t - 1$（$t > 0$，当且仅当 $t = 1$ 取等）得 $-\ln p_k \ge 1 - p_k$；两边同乘 $p_k \ge 0$ 并对 $k$ 求和即得 $H_e(p) = \sum_k p_k(-\ln p_k) \ge \sum_k p_k(1 - p_k) = G(p)$。对于二分类 $p_1 = \frac{1}{2} + \delta,\ p_2 = \frac{1}{2} - \delta$，直接计算得 $G(p) = 2(\frac{1}{2}+\delta)(\frac{1}{2}-\delta) = \frac{1}{2} - 2\delta^2$；而对 $H_e(p)$ 用 $\ln(1 \pm 2\delta) = \pm 2\delta - 2\delta^2 + O(\delta^3)$ 展开：
\begin{align*}
  H_e(p) &= -\Big(\frac{1}{2}+\delta\Big)\Big(-\ln 2 + 2\delta - 2\delta^2\Big) - \Big(\frac{1}{2}-\delta\Big)\Big(-\ln 2 - 2\delta - 2\delta^2\Big) + O(\delta^4) \\
         &= \ln 2 - 2\delta^2 + O(\delta^4).
\end{align*}
(3) 因 $G(p) = 1 - \norm{p}_2^2$ 的 Hessian 为 $-2 I_K \prec 0$，故严格凹。又母节点分布恰好是子节点分布的凸组合 $p = \frac{n_L}{n} p_L + \frac{n_R}{n} p_R$，当 $p_L \neq p_R$ 时严格 Jensen 不等式给出 $G(p) > \frac{n_L}{n} G(p_L) + \frac{n_R}{n} G(p_R)$，即 $\Delta G > 0$。
\end{proof}

\subsection{手工完整推演：6 样本分类树的 Gini 与信息增益比较}

\begin{example}[手算 CART 分类分裂]\label{ex:cart-hand}
考虑包含 $N = 6$ 个样本的一维二分类数据集：特征 $x_1 = (1, 2, 3, 4, 5, 6)$，对应标签 $y = (0, 0, 1, 0, 1, 1)$。请分别用 Gini 不纯度与香农信息增益评估两个候选切分点 $s_1 = 2.5$ 与 $s_2 = 3.5$，确定最优根节点分裂。
\end{example}

\paragraph{第 1 步：计算根节点不纯度}
根节点含 3 个负例与 3 个正例（$p = (0.5, 0.5)$）：
\[
  G(D) = 1 - (0.5^2 + 0.5^2) = 0.5000, \qquad
  H(D) = -\big(0.5\log_2 0.5 + 0.5\log_2 0.5\big) = 1.0000\text{ bit}.
\]

\paragraph{第 2 步：评估候选切分 $s_1 = 2.5$（$x_1 \le 2.5$）}
左子节点 $D_L$ 含样本 $\{1, 2\}$，标签为 $(0, 0)$（$n_L = 2$，完全纯净！）；右子节点 $D_R$ 含样本 $\{3, 4, 5, 6\}$，标签为 $(1, 0, 1, 1)$（$n_R = 4$，比例 $1/4$ 与 $3/4$）：
\begin{align*}
  G(D_L) &= 0, & G(D_R) &= 1 - \Big(\frac{1}{4}\Big)^2 - \Big(\frac{3}{4}\Big)^2 = \frac{6}{16} = 0.3750, \\
  G_{\mathrm{split}}(2.5) &= \frac{2}{6}\cdot 0 + \frac{4}{6}\cdot 0.3750 = \mathbf{0.2500}, &
  \Delta G(2.5) &= 0.5000 - 0.2500 = \mathbf{0.2500}.
\end{align*}
再算香农熵：
\begin{align*}
  H(D_L) &= 0, & H(D_R) &= -\frac{1}{4}\log_2\frac{1}{4} - \frac{3}{4}\log_2\frac{3}{4} \approx 0.8113\text{ bit}, \\
  \operatorname{IG}(2.5) &= 1.0000 - \left(\frac{2}{6}\cdot 0 + \frac{4}{6}\cdot 0.8113\right) \approx \mathbf{0.4591\text{ bit}}.
\end{align*}

\paragraph{第 3 步：评估候选切分 $s_2 = 3.5$（$x_1 \le 3.5$）}
左子节点标签为 $(0, 0, 1)$（$n_L = 3$，比例 $2/3$ 与 $1/3$）；右子节点标签为 $(0, 1, 1)$（$n_R = 3$，比例 $1/3$ 与 $2/3$）：
\begin{align*}
  G(D_L) = G(D_R) &= 1 - \Big(\frac{2}{3}\Big)^2 - \Big(\frac{1}{3}\Big)^2 = \frac{4}{9} \approx 0.4444, &
  \Delta G(3.5) &= 0.5000 - 0.4444 = \mathbf{0.0556}, \\
  H(D_L) = H(D_R) &= -\frac{2}{3}\log_2\frac{2}{3} - \frac{1}{3}\log_2\frac{1}{3} \approx 0.9183\text{ bit}, &
  \operatorname{IG}(3.5) &= 1.0000 - 0.9183 = \mathbf{0.0817\text{ bit}}.
\end{align*}
可见无论按 Gini 还是按信息增益，切分 $x_1 \le 2.5$（以及对称的 $x_1 \le 4.5$）都以压倒性优势击败 $x_1 \le 3.5$，原因恰在于它立刻剥离出了一个零不纯度的纯净叶节点！

\subsection{过拟合与剪枝：预剪枝与代价复杂度剪枝}

如果不加限制地持续分裂，一棵决策树可以把每个训练样本都单独关进一间叶节点牢房——训练误差归零，但方差爆表，对新数据的预测完全崩溃。抑制单棵树过拟合有两条路线：
\begin{itemize}
  \item \textbf{预剪枝（Pre-Pruning）}：在生长过程中提前刹车，常用超参数包括最大树深 \texttt{max\_depth}、叶节点最小样本数 \texttt{min\_leaf} 以及最小不纯度下降阈值 $\Delta_{\min}$；实现简单但有“视野短视”风险（有时当前一步增益很小，但下一步分裂能带来巨额收益）；
  \item \textbf{后剪枝 / 代价复杂度剪枝（Cost-Complexity Pruning, CCP）}：先长出一棵充分生长的大树 $T_0$，再引入带叶节点数正则惩罚的目标函数
  \begin{equation}\label{eq:cart-ccp}
    R_\alpha(T) := R(T) + \alpha \abs{T}, \qquad \alpha \ge 0,
  \end{equation}
  其中 $R(T)$ 为树 $T$ 的总经验误差，$\abs{T}$ 为叶节点总数。对树中任意内部节点 $t$（其所领子树为 $T_t$），把整棵子树 $T_t$ 剪成单叶节点 $t$ 带来的单叶代价阈值为
  \[
    \alpha(t) = \frac{R(t) - R(T_t)}{\abs{T_t} - 1}.
  \]
  不断剪去 $\alpha(t)$ 最小的“最薄弱环节（Weakest Link）”，可生成一族严格嵌套的最优剪枝子树序列 $T_0 \supset T_1 \supset \cdots \supset T_{\text{root}}$，最后通过第 3 章的 $K$ 折交叉验证选出泛化误差最小的 $\alpha^*$。
\end{itemize}

\begin{failBox}{轴对齐决策树的几何盲区：45 度对角边界与高基数分裂偏置}
\begin{enumerate}
  \item \textbf{锯齿阶梯困境}：若真实分类边界是简单的斜直线 $x_1 + x_2 = 0$，线性分类器或 SVM 只需 2 个参数即可精确刻画；而轴对齐决策树只能用大量水平与垂直线段拼出锯齿状阶梯去艰难逼近，树深迅速膨胀且泛化急剧变差；
  \item \textbf{高基数特征偏好}：若直接用信息增益切分无序离散特征，取值极多的 ID 类特征（如“用户学号”）会把每个样本分进单独子节点，使条件熵骤降为零从而骗取满分信息增益！（C4.5 通过除以分裂信息得到增益率，或工程上在 CART 中控制树深与编码方式加以规避。）
\end{enumerate}
\end{failBox}

% ============================================================
\section{Bagging 与随机森林：降方差的数学证明}\label{sec:trees-rf}
% ============================================================

单棵深树的致命病根不是偏差（它几乎能拟合任何复杂边界），而是\textbf{极高的方差}——训练集中稍微增删几个样本，根节点的第一次分裂就可能面目全非，整棵树随之剧烈抖动。如果把许多棵略有差异的高方差深树并行平均起来，是否就能在几乎不增加偏差的前提下把方差压扁？这就是 Leo Breiman 在 1996 年提出的 \textbf{Bagging（Bootstrap Aggregating）}与 2001 年提出的\textbf{随机森林（Random Forest）}。

\subsection{自助采样（Bootstrap）与 $36.8\%$ 袋外（OOB）法则}

手头只有一份大小为 $N$ 的训练集 $\mathcal{D}$，如何制造出 $M$ 份略有差异的“平行世界训练集”？答案是\textbf{自助采样（Bootstrap）}：每次从 $\mathcal{D}$ 中\textbf{有放回地均匀随机抽取 $N$ 个样本}构成第 $m$ 个自助集 $\mathcal{D}^{(m)}$，并在其上独立训练第 $m$ 棵决策树 $T_m$。

\begin{thmBox}
\begin{proposition}[自助采样的袋外（Out-of-Bag, OOB）概率极限]\label{prop:oob-limit}
对任意给定训练样本 $x_i \in \mathcal{D}$，它在一次大小为 $N$ 的自助采样中\textbf{一次都未被抽中}（成为该树的袋外样本）的概率为
\[
  \Prob(x_i \notin \mathcal{D}^{(m)}) = \left(1 - \frac{1}{N}\right)^{N} \xrightarrow{\;N \to \infty\;} e^{-1} \approx \mathbf{36.79\%}.
\]
\end{proposition}
\end{thmBox}

\begin{proof}
每次有放回抽取不命中 $x_i$ 的概率为 $1 - \frac{1}{N}$；独立重复 $N$ 次均不命中的概率为 $(1 - \frac{1}{N})^N$。取自然对数并利用 $\ln(1 - u) = -u + O(u^2)$：
\[
  N \ln\!\left(1 - \frac{1}{N}\right) = N\left(-\frac{1}{N} + O\!\left(\frac{1}{N^2}\right)\right) = -1 + O\!\left(\frac{1}{N}\right) \xrightarrow{\;N\to\infty\;} -1.
\]
两边取指数即得极限 $e^{-1} \approx 0.367879$。
\end{proof}

这一简短的极限带来一份美妙的“免费午餐”——\textbf{袋外评估（OOB Validation）}：对于每个训练样本 $x_i$，大约有 $36.8\%$ 的树在训练时压根没见过它！只用这批“没见过 $x_i$ 的树”对 $x_i$ 投票或平均求预测，再汇总所有 $N$ 个样本的误差，就能在不切分独立验证集的情况下，得到与留一交叉验证高度接近的无偏泛化误差估计。

\subsection{核心定理：相关集成方差公式与特征随机子采样的数学必然}

现在我们严格证明集成平均究竟能把方差降到多低，以及为什么单纯 Bagging 还不够、必须升级为随机森林。

\begin{thmBox}
\begin{theorem}[相关同分布集成的精确方差公式]\label{thm:rf-variance}
设在给定测试点 $x$ 处，集成中的 $M$ 个基学习器预测值 $T_1(x), \dots, T_M(x)$ 具有相同的方差 $\Var(T_m(x)) = \sigma^2$，且任意两棵不同树之间的成对相关系数均为 $\operatorname{Corr}(T_m(x), T_{m'}(x)) = \rho$（$m \neq m'$）。则等权平均集成 $\bar{T}_M(x) := \frac{1}{M}\sum_{m=1}^{M} T_m(x)$ 的方差满足精确恒等式：
\begin{equation}\label{eq:rf-var-formula}
  \Var\!\left(\frac{1}{M}\sum_{m=1}^{M} T_m(x)\right) = \rho\sigma^2 + \frac{1 - \rho}{M}\sigma^2.
\end{equation}
\end{theorem}
\end{thmBox}

\begin{proof}
利用方差的双线性展开公式 $\Var(\sum_{m=1}^M a_m Z_m) = \sum_{m=1}^M a_m^2 \Var(Z_m) + \sum_{m \neq m'} a_m a_{m'} \Cov(Z_m, Z_{m'})$：
\begin{align*}
  \Var\!\left(\frac{1}{M}\sum_{m=1}^{M} T_m(x)\right)
  &= \frac{1}{M^2}\left[\sum_{m=1}^{M}\Var\big(T_m(x)\big) + \sum_{m \neq m'}\Cov\big(T_m(x), T_{m'}(x)\big)\right].
\end{align*}
对角项共有 $M$ 个，每项等于 $\sigma^2$；非对角项共有 $M(M - 1)$ 个，由相关系数定义，每项等于 $\Cov(T_m, T_{m'}) = \rho\sqrt{\sigma^2 \cdot \sigma^2} = \rho\sigma^2$。代入化简：
\begin{align*}
  \Var\!\left(\frac{1}{M}\sum_{m=1}^{M} T_m(x)\right)
  &= \frac{1}{M^2}\Big(M\sigma^2 + M(M - 1)\rho\sigma^2\Big)
   = \frac{\sigma^2}{M} + \frac{M - 1}{M}\rho\sigma^2
   = \rho\sigma^2 + \frac{1 - \rho}{M}\sigma^2. \qedhere
\end{align*}
\end{proof}

请仔细端详式~\eqref{eq:rf-var-formula}，它是整个随机森林算法设计的灯塔：
\begin{itemize}
  \item \textbf{第二项 $\frac{1-\rho}{M}\sigma^2$} 随着树的数量 $M \to \infty$ 迅速衰减趋近于 $0$——这就是增大森林规模能平滑预测、且\textbf{几乎不会因树太多而过拟合}的数学保证；
  \item \textbf{第一项 $\rho\sigma^2$} 与 $M$ 无关！它构成了集成方差\textbf{永远跨不过去的下界地板}。在纯 Bagging 中，尽管每棵树用了不同自助集，但只要数据中存在一两个极强的核心预测特征，几乎每棵树的根节点都会争先恐后地选同一个特征切分，导致所有树长得大同小异（$\rho$ 高达 $0.5\text{--}0.8$），再多种 1000 棵树也无济于事！
\end{itemize}

\textbf{随机森林的天才一击}直指 $\rho$：在每棵树的\textbf{每一次节点分裂}时，不再遍历全部 $d$ 个特征，而是先随机抽取一个规模为 $m_{\text{try}} < d$ 的特征子集（分类经验默认 $m_{\text{try}} = \lfloor \sqrt{d} \rfloor$，回归经验默认 $m_{\text{try}} = \lfloor d/3 \rfloor$），仅在这 $m_{\text{try}}$ 个候选中寻找最佳切分！这一举措强行剥夺了强势特征对根节点的垄断权，迫使不同树探索多样化的特征组合，大幅拉低树间相关性 $\rho$，从而以微小的单树偏差代价换来集成总均方误差的断崖式下降！同时，正因平均操作专职降方差，随机森林通常使用\textbf{充分生长的深树}（低偏差、高方差）作为基学习器。

% ============================================================
\section{Boosting 与 GBDT：函数空间中的梯度下降}\label{sec:trees-gbdt}
% ============================================================

Bagging 把高方差的“天才乱猜手”拉在一起民主投票以求稳健；而 \textbf{Boosting（提升法）}走了截然相反的另一条大道：它雇佣一队极浅的弱学习器（例如深度仅 $d = 2\text{--}4$ 的“决策树桩/浅树”，本身偏差极大但方差很小），让它们\textbf{串行接力}——每一棵新树专门盯着前面所有树合起来仍未解决的错误猛攻，步步为营地将整体偏差磨平！图~\ref{fig:bag-vs-boost} 并置对比了这两大集成范式的计算拓扑。

\begin{figure}[htbp]
\centering
\resizebox{0.98\linewidth}{!}{%
\begin{tikzpicture}[
  box/.style={rectangle, rounded corners=1.5mm, draw=navy, thick, fill=navyLight, minimum width=18mm, minimum height=7.5mm, align=center, font=\sffamily\scriptsize, text=ink},
  gbox/.style={rectangle, rounded corners=1.5mm, draw=tealDark, thick, fill=tealLight, minimum width=18mm, minimum height=7.5mm, align=center, font=\sffamily\scriptsize, text=ink},
  rbox/.style={rectangle, rounded corners=1.5mm, draw=ochre, thick, fill=ochreLight, minimum width=16mm, minimum height=6.5mm, align=center, font=\sffamily\scriptsize, text=ochre},
  arr/.style={-{Stealth[length=2mm]}, thick, slate}
]
  % Top panel: Parallel Bagging / Random Forest
  \draw[rounded corners=2mm, draw=navy!45, fill=navyLight!25] (-0.3, 1.55) rectangle (14.45, 4.75);
  \node[font=\sffamily\bfseries\small, text=navy, anchor=west] at (0.0, 4.42) {范式一：Bagging / 随机森林（并行降方差，基树深、高方差）};

  \node[box, minimum width=19mm, minimum height=10mm] (D1) at (1.1, 2.95) {训练集 $\mathcal{D}$};
  \node[box] (B1) at (4.1, 3.85) {Bootstrap $\mathcal{D}^{(1)}$};
  \node[box] (B2) at (4.1, 2.95) {Bootstrap $\mathcal{D}^{(2)}$};
  \node[box] (BM) at (4.1, 2.05) {Bootstrap $\mathcal{D}^{(M)}$};

  \node[box, fill=white] (T1) at (7.3, 3.85) {深树 $T_1$ ($m_{\text{try}}$)};
  \node[box, fill=white] (T2) at (7.3, 2.95) {深树 $T_2$ ($m_{\text{try}}$)};
  \node[box, fill=white] (TM) at (7.3, 2.05) {深树 $T_M$ ($m_{\text{try}}$)};

  \node[box, draw=navy, very thick, minimum width=27mm, minimum height=11mm] (AVG) at (11.35, 2.95) {等权平均 / 投票\\$\hat{f}(x) = \frac{1}{M}\sum_{m=1}^M T_m(x)$};

  \draw[arr] (D1.east) -- (B1.west);
  \draw[arr] (D1.east) -- (B2.west);
  \draw[arr] (D1.east) -- (BM.west);
  \draw[arr] (B1) -- (T1);
  \draw[arr] (B2) -- (T2);
  \draw[arr] (BM) -- (TM);
  \draw[arr] (T1.east) -- (AVG.west);
  \draw[arr] (T2.east) -- (AVG.west);
  \draw[arr] (TM.east) -- (AVG.west);

  % Bottom panel: Sequential GBDT Boosting
  \draw[rounded corners=2mm, draw=teal!50, fill=tealLight!25] (-0.3, -2.45) rectangle (14.45, 1.25);
  \node[font=\sffamily\bfseries\small, text=tealDark, anchor=west] at (0.0, 0.92) {范式二：梯度提升树 GBDT（串行降偏差，基树浅、拟合负梯度伪残差）};

  \node[gbox, minimum width=17mm] (F0) at (0.95, -0.58) {常值初始化\\$F_0(x) = c_0$};
  \node[rbox] (R1) at (3.25, 0.28) {伪残差 $r^{(1)}$};
  \node[gbox, fill=white] (H1) at (3.25, -1.52) {浅树 $h_1(x)$};

  \node[gbox, minimum width=21mm] (F1) at (5.95, -0.58) {$F_1 = F_0 + \eta h_1$};
  \node[rbox] (R2) at (8.45, 0.28) {伪残差 $r^{(2)}$};
  \node[gbox, fill=white] (H2) at (8.45, -1.52) {浅树 $h_2(x)$};

  \node[font=\bfseries\color{tealDark}] (DOTS) at (10.35, -0.58) {$\cdots$};
  \node[gbox, draw=tealDark, very thick, minimum width=21mm, minimum height=11mm] (FM) at (12.65, -0.58) {最终累加模型\\$F_M = F_0 + \eta\sum h_m$};

  \draw[arr, tealDark] (F0.east) -- (R1.west);
  \draw[arr, ochre] (R1.south) -- node[pos=0.28, right=1pt, font=\tiny, text=ochre, fill=tealLight!25, inner sep=0.8pt] {拟合} (H1.north);
  \draw[arr, tealDark] (F0.east) -- (F1.west);
  \draw[arr, tealDark] (H1.east) -- node[below right, font=\tiny, text=tealDark] {$+\eta h_1$} (F1.south west);

  \draw[arr, tealDark] (F1.east) -- (R2.west);
  \draw[arr, ochre] (R2.south) -- node[pos=0.28, right=1pt, font=\tiny, text=ochre, fill=tealLight!25, inner sep=0.8pt] {拟合} (H2.north);
  \draw[arr, tealDark] (F1.east) -- (DOTS.west);
  \draw[arr, tealDark] (H2.east) -- node[below right, font=\tiny, text=tealDark] {$+\eta h_2$} (DOTS.south west);
  \draw[arr, tealDark] (DOTS.east) -- (FM.west);
\end{tikzpicture}}
\caption{Bagging/随机森林与 GBDT 梯度提升树的架构对照。随机森林通过并行独立自助集与随机选特征压低方差地板 $\rho\sigma^2$；GBDT 通过串行追踪函数空间负梯度（伪残差）逐步消减系统偏差。}
\label{fig:bag-vs-boost}
\end{figure}

\subsection{函数空间中的梯度下降（Friedman 2001）}

普通梯度下降在第 1 章是在\textbf{参数空间} $\theta \in \R^d$ 里沿负梯度走步：$\theta_{m} = \theta_{m-1} - \eta \nabla_\theta \mathcal{L}(\theta_{m-1})$。Jerome Friedman 的惊世洞察在于：\textbf{为什么不把模型本身 $F(\cdot)$ 在训练集上的预测向量 $F = (F(x_1), \dots, F(x_N))^\top \in \R^N$ 直接视为优化变量，在“函数空间”里做梯度下降？}

给定任意可微损失函数 $\ell(y_i, F(x_i))$ 与经验总损失 $\mathcal{L}(F) = \sum_{i=1}^N \ell(y_i, F(x_i))$，第 $m$ 轮迭代时，我们在当前累计模型 $F_{m-1}$ 处计算总损失对每个样本预测值 $F(x_i)$ 的\textbf{负梯度}——称为第 $m$ 轮的\textbf{伪残差（Pseudo-Residual）}：
\begin{equation}\label{eq:gbdt-pseudores}
  r_{im} := -\left.\frac{\partial \ell\big(y_i, F(x_i)\big)}{\partial F(x_i)}\right|_{F(x_i) = F_{m-1}(x_i)}, \qquad i = 1, \dots, N.
\end{equation}
直接把 $F(x_i)$ 沿 $r_{im}$ 移动固然能让训练损失骤降，但那样只能更新这 $N$ 个训练点的取值，无法泛化到未见过的测试样本 $x_{\text{new}}$！于是 GBDT 做了最关键的一步\textbf{投影}：用一棵浅层 CART 回归树 $h_m(x)$ 去最小二乘拟合伪残差训练集 $\{(x_i, r_{im})\}_{i=1}^N$，使得 $h_m$ 成为浅树假设空间里与负梯度方向最贴合的函数方向，然后乘以收缩学习率（Shrinkage）$\eta \in (0, 1]$ 进行加法更新：
\begin{equation}\label{eq:gbdt-update}
  F_m(x) = F_{m-1}(x) + \eta\, h_m(x), \qquad F_0(x) \equiv \argmin_{c \in \R} \sum_{i=1}^{N} \ell(y_i, c).
\end{equation}

让我们立即检验两个核心损失函数下的伪残差公式，见证其浑然天成的统一性：
\begin{enumerate}
  \item \textbf{平方损失回归}：取 $\ell(y, F) = \frac{1}{2}(y - F)^2$，求导得
  \[
    r_{im} = -\frac{\partial}{\partial F}\left[\frac{1}{2}(y_i - F)^2\right]_{F=F_{m-1}(x_i)} = y_i - F_{m-1}(x_i).
  \]
  此时伪残差就是字面意义上的\textbf{真实残差}！每棵新树都在拟合前一轮还没预测完的剩余误差。
  \item \textbf{二分类对数损失（Logistic Loss）}：设 $y_i \in \{0, 1\}$，$F(x_i)$ 为对数几率（logit），预测概率 $p_i = \sigma(F(x_i)) = \frac{1}{1+e^{-F(x_i)}}$，损失为 $\ell(y_i, F) = \ln(1 + e^F) - y_i F$。求导得
  \[
    r_{im} = -\left(\frac{e^{F_{m-1}(x_i)}}{1 + e^{F_{m-1}(x_i)}} - y_i\right) = y_i - p_{i, m-1} \in (-1, 1).
  \]
  惊不惊喜？二分类 GBDT 的每棵基树竟然仍是\textbf{回归树}，拟合的正是真实 $0/1$ 标签与当前预测概率之差！（同理，绝对值损失 $|y-F|$ 对应的伪残差为符号函数 $\sign(y_i - F_{m-1}(x_i))$，天然抗离群点！）
\end{enumerate}

\textbf{学习率 $\eta$ 的正则化作用}：为什么不干脆取 $\eta = 1$ 一步到位？经验与理论均表明，把每棵树的步子收小（如 $\eta = 0.05\text{--}0.10$）并相应增加树的数量 $M$，能防止单棵树对早期残差过度拟合，让大量弱树以“细水长流”的慢速推进方式共同雕琢目标函数——这本质上等价于函数空间中的隐式 $L_1$ 正则化（与第 2 章 Lasso 的前向分步路径高度呼应）。不过请注意：与随机森林不同，若 GBDT 的树数 $M$ 无限增大，它最终会开始拟合训练噪声而\textbf{出现过拟合}，因此必须配合验证集早停（Early Stopping）与树深限制使用。

\subsection{从一阶 GBDT 到二阶牛顿提升：XGBoost 与 LightGBM 精要}

经典 GBDT 只用了一阶负梯度；陈天奇于 2016 年推出的 \textbf{XGBoost} 将函数空间梯度下降直接升级为\textbf{带显式正则化的二阶牛顿步}，推导出极其漂亮的闭式叶权重与分裂增益公式。

在第 $m$ 轮加入新树 $f_m(x) = w_{q(x)}$（将样本 $x$ 映射到第 $q(x) \in \{1,\dots,T\}$ 个叶子并输出权重 $w_j$），定义带树复杂度惩罚的正则化目标：
\[
  \mathcal{L}^{(m)} = \sum_{i=1}^{N} \ell\!\big(y_i,\, \hat{y}_i^{(m-1)} + f_m(x_i)\big) + \Omega(f_m),
  \qquad \Omega(f_m) := \gamma T + \frac{1}{2}\lambda \sum_{j=1}^{T} w_j^2.
\]
对损失项在 $\hat{y}_i^{(m-1)}$ 处作二阶泰勒展开，记一阶导 $g_i := \partial_{\hat{y}}\ell(y_i, \hat{y}_i^{(m-1)})$ 与二阶导 $h_i := \partial_{\hat{y}}^2\ell(y_i, \hat{y}_i^{(m-1)}) > 0$，略去与 $f_m$ 无关的常数项，并将样本按所属叶节点 $I_j := \{i \mid q(x_i) = j\}$ 分组：
\[
  \widetilde{\mathcal{L}}^{(m)}
  = \sum_{i=1}^{N}\!\left[g_i f_m(x_i) + \frac{1}{2} h_i f_m(x_i)^2\right] + \gamma T + \frac{1}{2}\lambda \sum_{j=1}^{T} w_j^2
  = \sum_{j=1}^{T}\!\left[G_j w_j + \frac{1}{2}(H_j + \lambda) w_j^2\right] + \gamma T,
\]
其中 $G_j := \sum_{i \in I_j} g_i$、$H_j := \sum_{i \in I_j} h_i$。各个叶节点 $w_j$ 彻底解耦！对每个 $w_j$ 求导置零，立得\textbf{闭式最优叶权重}与\textbf{最优树结构评分}：
\begin{equation}\label{eq:xgb-leaf-gain}
  w_j^* = -\frac{G_j}{H_j + \lambda}, \qquad
  \widetilde{\mathcal{L}}^{(m)*} = -\frac{1}{2}\sum_{j=1}^{T} \frac{G_j^2}{H_j + \lambda} + \gamma T.
\end{equation}
从而将任一叶节点分裂为左右两叶 $(I_L, I_R)$ 的精确\textbf{分裂收益（Gain）}一气呵成：
\begin{equation}\label{eq:xgb-split-gain}
  \operatorname{Gain} = \frac{1}{2}\left[\frac{G_L^2}{H_L + \lambda} + \frac{G_R^2}{H_R + \lambda} - \frac{(G_L + G_R)^2}{H_L + H_R + \lambda}\right] - \gamma.
\end{equation}
式~\eqref{eq:xgb-leaf-gain}--\eqref{eq:xgb-split-gain} 堪称“天书”级公式：不仅自动把二阶曲率信息（牛顿步长）和 $L_2$ 叶权正则 $\lambda$ 融为一体，而且当 $\operatorname{Gain} \le 0$ 时直接拒绝分裂（将叶节点数惩罚 $\gamma$ 内化为自动剪枝！）。微软后续推出的 \textbf{LightGBM} 在同一二阶数学框架上引入了特征直方图分桶（将切分搜索复杂度从排序 $O(N\log N)$ 降至桶数 $O(B)$）与按叶生长（Leaf-wise）策略，使千万级样本训练提速一个数量级。

\begin{extBox}{AdaBoost 与指数损失的前向分步推导}
Boosting 家族的鼻祖 \textbf{AdaBoost}（Freund \& Schapire, 1997）表面上通过“提高错分样本权重”来迭代训练二分类弱分类器 $h_m(x) \in \{-1, +1\}$。其实它正是\textbf{指数损失} $\ell_{\exp}(y, F) = \exp(-y F)$ 上的前向分步加性建模！证明极简：设第 $m$ 轮要最小化
\[
  \sum_{i=1}^{N} \exp\!\big(-y_i(F_{m-1}(x_i) + \alpha_m h_m(x_i))\big)
  = \sum_{i=1}^{N} \underbrace{e^{-y_i F_{m-1}(x_i)}}_{w_i^{(m)}} \exp\!\big(-\alpha_m y_i h_m(x_i)\big).
\]
将求和按“预测正确（$y_i h_m(x_i)=+1$）”与“预测错误（$y_i h_m(x_i)=-1$）”拆成两部分，记加权错误率 $\varepsilon_m := \frac{\sum_{y_i \neq h_m(x_i)} w_i^{(m)}}{\sum_i w_i^{(m)}}$，则目标正比于 $(1 - \varepsilon_m) e^{-\alpha_m} + \varepsilon_m e^{\alpha_m}$。
\begin{itemize}
  \item 固定 $\alpha_m > 0$，最小化目标等价于最小化加权错分率 $\varepsilon_m$（即挑选加权误差最小的弱分类器 $h_m$）；
  \item 固定 $h_m$，对 $\alpha_m$ 求导置零得 $-(1-\varepsilon_m)e^{-\alpha_m} + \varepsilon_m e^{\alpha_m} = 0 \implies \alpha_m^* = \frac{1}{2}\ln\frac{1-\varepsilon_m}{\varepsilon_m}$；
  \item 下一轮未归一化权重自动更新为 $w_i^{(m+1)} = w_i^{(m)}\exp(-\alpha_m^* y_i h_m(x_i))$——错分样本权重乘以 $e^{\alpha_m^*} > 1$，答对样本权重乘以 $e^{-\alpha_m^*} < 1$！
\end{itemize}
由于指数损失对极端错分样本的惩罚呈指数爆炸，AdaBoost 对标签噪声较为敏感；现代实践中通常优先使用对数损失或 Huber 损失的 GBDT。
\end{extBox}

% ============================================================
\section{三模型实测对比：单棵树、随机森林与 GBDT}\label{sec:trees-exp}
% ============================================================

现在我们直接检验配套脚本 \texttt{code/\allowbreak vol4\_experiments.py} 中 \texttt{exp5\_trees\_rf\_gbdt} 的实测曲线与数值（数据存于 \texttt{data/\allowbreak ensemble\_curve.csv}）。

\paragraph{实验设置}
在 5 维超立方体 $[-3, 3]^5$ 上生成非线性回归数据（$N_{\text{tr}} = 220$，$N_{\text{te}} = 220$），真实函数为
\[
  f_{\mathrm{true}}(x) = \sin(x_0) + 0.5 x_1^2 - 0.8 x_2 + 0.4 x_0 x_1,
\]
注意最后两维 $x_3, x_4$ 是\textbf{纯无关噪声特征}，且训练标签叠加了标准差 $\sigma = 0.70$ 的高斯噪声（测试集与无噪真实函数 $f_{\mathrm{true}}$ 比对以纯粹考察泛化估计误差）。我们用同一套纯 NumPy 回归树引擎对比三种模型：
\begin{enumerate}
  \item \textbf{单棵深 CART 树}（$\text{max\_depth}=8,\ \text{min\_leaf}=2$）：训练 MSE 低至 \textbf{0.2603}，但测试 MSE 高达 \textbf{1.6017}——严重过拟合噪声！
  \item \textbf{随机森林（RF）}（$M=60,\ \text{max\_depth}=8,\ \text{min\_leaf}=2,\ m_{\text{try}}=3$）：测试 MSE 从单棵自助随机树的 \textbf{2.9815}（$M=1$）在仅 10 棵树时就暴跌至 \textbf{1.0414}（$M=10$），随后在 \textbf{1.0217}（$M=60$）附近稳定进入平台期——精准印证定理~\ref{thm:rf-variance} 中 $\rho\sigma^2$ 的方差地板效应！
  \item \textbf{梯度提升树（GBDT）}（$M=60,\ \text{max\_depth}=3,\ \text{min\_leaf}=4,\ \eta=0.10$）：起步时单棵浅树欠拟合严重，测试 MSE 高达 \textbf{4.7747}（$M=1$）；但随着残差接力推进，它一路势如破竹：\textbf{2.2161}（$M=10$）$\to$ \textbf{0.9053}（$M=30$，反超随机森林！）$\to$ \textbf{0.5517}（$M=60$，将单棵深树误差削去近三分之二！）。
\end{enumerate}

图~\ref{fig:ensemble-curve} 直接读取实测数据文件 \texttt{data/\allowbreak ensemble\_curve.csv} 绘制了完整的测试 MSE 演化轨迹，表~\ref{tab:ensemble-metrics} 汇总了关键里程碑数值。

\begin{figure}[htbp]
\centering
\begin{tikzpicture}
\begin{axis}[
  width=0.84\linewidth,
  height=6.2cm,
  xlabel={集成基树数量 $M$},
  ylabel={测试集均方误差（Test MSE）},
  xmin=0, xmax=62,
  ymin=0.25, ymax=5.15,
  xtick={1,10,20,30,40,50,60},
  tick label style={font=\scriptsize},
  label style={font=\footnotesize},
  grid=both,
  grid style={line width=0.25pt, draw=slate!18},
  legend pos=north east,
  legend style={font=\scriptsize, fill=white, fill opacity=0.92, draw=tealLine, rounded corners=1.5pt}
]
  \addplot[ochre, very thick, dashed] table[x=m, y=single_te, col sep=comma] {data/ensemble_curve.csv};
  \addlegendentry{单棵深 CART（$d=8$，测试 MSE $= 1.6017$）}

  \addplot[navy, very thick, mark=square*, mark size=1.8pt] table[x=m, y=rf_te, col sep=comma] {data/ensemble_curve.csv};
  \addlegendentry{随机森林（$d=8, m_{\text{try}}=3$，快速触及方差地板）}

  \addplot[tealDark, very thick, mark=*, mark size=1.9pt] table[x=m, y=gbdt_te, col sep=comma] {data/ensemble_curve.csv};
  \addlegendentry{GBDT（$d=3, \eta=0.10$，持续降偏差至 $0.5517$）}

  % Annotations
  \node[font=\scriptsize, text=navy, align=left] at (axis cs:44, 1.38) {RF 平台期 $\approx 1.02$};
  \node[font=\scriptsize, text=tealDark, align=left] at (axis cs:46, 0.42) {GBDT 持续逼近真值};
  \draw[->, >=Stealth, gold, thick] (axis cs:24, 2.05) -- (axis cs:24.5, 1.18);
  \node[font=\scriptsize\bfseries, text=gold, fill=white, inner sep=1.2pt] at (axis cs:25, 2.28) {$M \approx 24$ 处 GBDT 反超 RF};
\end{axis}
\end{tikzpicture}
\caption{5 维非线性含噪回归（$N_{\text{tr}}=220, \sigma=0.70$）上单棵深 CART、随机森林与 GBDT 的测试 MSE 演化曲线（实测数据直接读自 \texttt{data/\allowbreak ensemble\_curve.csv}）。随机森林在前 10 棵树内通过降方差迅速击败单棵深树并稳居平台；GBDT 虽然单棵浅树起步偏差很高，却凭借串行负梯度拟合持续压低误差，在 $M \approx 24$ 后全面反超。}
\label{fig:ensemble-curve}
\end{figure}

\begin{table}[htbp]
\centering
\caption{5 维非线性回归实测关键节点测试 MSE（单棵深 CART 训练 MSE $=0.2603$，测试 MSE $=1.6017$）}
\label{tab:ensemble-metrics}
\small
\begin{tabular}{@{}lccccccc@{}}
\toprule
\textbf{集成模型} & $M=1$ & $M=5$ & $M=10$ & $M=20$ & $M=30$ & $M=50$ & $M=60$ \\
\midrule
随机森林（$d=8, m_{\text{try}}=3$） & $2.9815$ & $1.5006$ & $1.0414$ & $1.0753$ & $1.0799$ & $1.0504$ & $1.0217$ \\
GBDT（$d=3, \eta=0.10$） & $4.7747$ & $3.3744$ & $2.2161$ & $1.2893$ & $0.9053$ & $0.6303$ & $\mathbf{0.5517}$ \\
\bottomrule
\end{tabular}
\end{table}

\subsection{可运行的纯 NumPy 回归树、随机森林与 GBDT 引擎}

下面的紧凑代码用不到 50 行纯 NumPy 实现了 CART 回归树递归生长器，并直接驱动随机森林与 GBDT 两种集成训练循环（完整实验脚本见 \texttt{code/\allowbreak vol4\_experiments.py}）。

\begin{lstlisting}[caption={纯 NumPy 实现 CART 回归树、随机森林与 GBDT 核心循环},label={lst:trees-ensemble}]
import numpy as np

def fit_reg_tree(X, y, max_depth=3, min_leaf=4, max_features=None, rng=None):
    n, d = X.shape
    if max_depth == 0 or n <= min_leaf or np.var(y) < 1e-10:
        return float(np.mean(y))
    feats = range(d) if max_features is None else rng.choice(d, size=max_features, replace=False)
    best_score, best_split = float("inf"), None
    for j in feats:
        vals = np.unique(X[:, j])
        if len(vals) <= 1:
            continue
        qs = vals if len(vals) <= 16 else np.quantile(vals, np.linspace(0.1, 0.9, 15))
        for thr in qs:
            left = X[:, j] <= thr
            nl = int(np.sum(left))
            if nl < min_leaf or (n - nl) < min_leaf:
                continue
            score = nl * np.var(y[left]) + (n - nl) * np.var(y[~left])
            if score < best_score:
                best_score, best_split = score, (j, float(thr))
    if best_split is None:
        return float(np.mean(y))
    j, thr = best_split
    left = X[:, j] <= thr
    return (j, thr,
            fit_reg_tree(X[left], y[left], max_depth - 1, min_leaf, max_features, rng),
            fit_reg_tree(X[~left], y[~left], max_depth - 1, min_leaf, max_features, rng))

def pred_tree(tree, X):
    if isinstance(tree, float):
        return np.full(len(X), tree)
    j, thr, tl, tr = tree
    out = np.empty(len(X))
    left = X[:, j] <= thr
    out[left], out[~left] = pred_tree(tl, X[left]), pred_tree(tr, X[~left])
    return out

# Core ensemble loops (given X_tr, y_tr, X_te):
# 1) Random Forest (Bagging + feature subsampling):
#    idx = rng.integers(0, len(X_tr), size=len(X_tr))
#    t_rf = fit_reg_tree(X_tr[idx], y_tr[idx], max_depth=8, min_leaf=2, max_features=3, rng=rng)
#    rf_pred_sum += pred_tree(t_rf, X_te)
# 2) GBDT (Sequential residual fitting with shrinkage lr=0.10):
#    resid = y_tr - gb_tr
#    t_gb = fit_reg_tree(X_tr, resid, max_depth=3, min_leaf=4)
#    gb_tr += 0.10 * pred_tree(t_gb, X_tr); gb_te += 0.10 * pred_tree(t_gb, X_te)
\end{lstlisting}

\begin{insightBox}{工业界黄金分界线：表格选 GBDT，非结构化模态选深度学习}
为什么在大量中等规模（$10^3\text{--}10^6$ 行）的真实表格基准评测中，XGBoost / LightGBM / CatBoost 至今依然常常击败多层感知机与表格 Transformer？核心原因有三：（1）表格特征的真实响应面常含\textbf{非光滑的轴对齐阈值跳变}（例如收入过某一档位税率突变），而神经网络天然偏好平滑旋转不变函数；（2）现实表格充斥无关噪声列（正如本节实验中的 $x_3, x_4$），决策树在每次分裂时自动忽略无关维度，而全连接层容易被噪声特征的梯度稀释；（3）树集成的训练、调参与缺失值处理极其省心。反之，一旦数据具有强时空局部性或组合语义结构（图像像素、语音波形、自然语言词元），第 5 卷的卷积网络与 Transformer 便会以压倒性表征优势接管主场。
\end{insightBox}

% ============================================================
\section{本章习题}\label{sec:ch5-ex}
% ============================================================

\exTier{A 层：基础巩固}{掌握不纯度计算、Bagging/Boosting 核心差异与超参直觉}

\begin{enumerate}[label=\textbf{5.\arabic*}]
  \item \textbf{三分类节点不纯度计算}：某决策树节点包含 10 个三分类样本，各类别频数分别为 $(5, 3, 2)$。（1）计算该节点的 Gini 不纯度 $G(p)$ 与以 2 为底的香农熵 $H(p)$；（2）若某次分裂将其切分为左子节点 $(5, 0, 0)$ 与右子节点 $(0, 3, 2)$，计算此次分裂的 Gini 增益 $\Delta G$ 与信息增益 $\operatorname{IG}$。
  \exHint{左子节点完全纯净，不纯度与熵均为 $0$；只需算右子节点比例 $(0.6, 0.4)$ 的不纯度再按权重 $5/10$ 加权。}

  \item \textbf{随机森林 vs GBDT 诊断四连问}：结合本章理论简明判断下列陈述正误并说明核心数学依据：
  （1）把随机森林的树数 $M$ 从 100 增至 1000，通常会引起严重过拟合；
  （2）把 GBDT 的树数 $M$ 从 100 增至 5000（固定较大学习率 $\eta=0.2$），通常会引起严重过拟合；
  （3）对输入特征做单调对数变换 $x \mapsto \ln(1 + x)$（$x \ge 0$）不会改变 CART 分类树的分裂结构；
  （4）随机森林通常偏好浅树桩（$d=1$），而 GBDT 通常偏好不剪枝的极深树（$d=20$）。
  \exHint{紧扣方差公式~\eqref{eq:rf-var-formula} 与 Bagging 降方差、Boosting 降偏差的分工定位作答。}
\end{enumerate}

\exTier{B 层：方法与推导}{熟练手算回归树分裂、XGBoost 叶权重与分裂增益}

\begin{enumerate}[label=\textbf{5.\arabic*}, resume]
  \item \textbf{一维回归树手算与代价复杂度剪枝}：给定 4 个一维回归样本 $(x_i, y_i)$：$(1, 1),\, (2, 2),\, (3, 9),\, (4, 10)$。（1）评估全部 3 个中点切分 $s \in \{1.5, 2.5, 3.5\}$ 的加权残差平方和，找出最优第一次分裂；（2）在最优第一次分裂后继续完全生长至 4 个叶节点（令总误差为 $0$），计算完全树 $T_0$、1 次分裂树 $T_1$ 与单叶根节点 $T_{\text{root}}$ 在式~\eqref{eq:cart-ccp} 下的代价复杂度 $R_\alpha(T)$，求出使 $T_1$ 严格优于 $T_0$ 与 $T_{\text{root}}$ 的正则参数 $\alpha$ 区间！
  \exHint{记单节点残差平方和 $R(t) = \sum_{i\in t}(y_i - \bar{y}_t)^2$：根节点的 $R(T_{\text{root}}) = 65$，$T_1$ 的 $R(T_1) = 1$，$T_0$ 的 $R(T_0) = 0$。}

  \item \textbf{XGBoost 二阶叶权重与分裂增益手算}：在二分类对数损失下，当前某叶节点包含 4 个样本，其真实标签为 $y = (1, 1, 0, 0)$，上一轮预测概率分别为 $p = (0.8, 0.4, 0.6, 0.2)$。设正则超参数 $\lambda = 1.0,\ \gamma = 0.05$。（1）计算每个样本的一阶导 $g_i = p_i - y_i$ 与二阶导 $h_i = p_i(1 - p_i)$；（2）由式~\eqref{eq:xgb-leaf-gain} 算出分裂前该叶节点的最优权重 $w^*$；（3）若某候选切分将样本分为 $I_L = \{1, 2\}$ 与 $I_R = \{3, 4\}$，利用式~\eqref{eq:xgb-split-gain} 计算分裂增益 $\operatorname{Gain}$，判断是否应当执行该分裂？
  \exHint{注意先汇总 $G_L, H_L, G_R, H_R$，检查分裂前是否满足 $G = G_L + G_R$ 与 $H = H_L + H_R$。}
\end{enumerate}

\exTier{C 层：理论证明与反例构造}{证明带相关性的不等权集成界与 Huber 损失伪残差}

\begin{enumerate}[label=\textbf{5.\arabic*}, resume]
  \item \textbf{回归中位数叶预测与 $L_1$ 绝对值损失的 GBDT}：假设在回归任务中将平方损失替换为绝对值损失（MAE）$\ell_1(y, F) = \abs{y - F}$。（1）证明对任意叶节点样本集 $\{y_1, \dots, y_n\}$，最小化 $\sum_{i=1}^n \abs{y_i - c}$ 的最优常数预测 $c^*$ 是这组数的\textbf{样本中位数}（而不是均值）；（2）写出 $L_1$-GBDT 第 $m$ 轮的伪残差 $r_{im}$ 表达式，并解释为何在拟合回归树 $h_m(x)$ 确定叶区域 $R_{jm}$ 之后，标准 $L_1$-GBDT 会把每个叶子的输出值重新替换为 $\operatorname{median}_{i \in R_{jm}}(y_i - F_{m-1}(x_i))$！
  \exHint{第 (1) 问可利用第 2 章的次梯度最优性条件 $0 \in \partial_c \sum_i \abs{y_i - c}$，即满足 $y_i > c$ 的个数与 $y_i < c$ 的个数平衡。}

  \item \textbf{随机森林特征子采样相关性极限分析}：在定理~\ref{thm:rf-variance} 的模型假设下，假设把候选特征数从 $m_{\text{try}} = d$ 缩减到 $m_{\text{try}} = \sqrt{d}$ 时，单棵树的方差从 $\sigma^2$ 略升至 $1.25\sigma^2$，但任意两棵树之间的成对相关系数从 $\rho_0 = 0.60$ 降至 $\rho_1 = 0.25$。（1）计算在 $M \to \infty$ 极限下，特征子采样使集成方差降低了百分之多少？（2）求出随机森林（$m_{\text{try}}=\sqrt{d}$）开始严格优于纯 Bagging（$m_{\text{try}}=d$）的最小树数阈值 $M_{\min}$。
  \exHint{直接代入式~\eqref{eq:rf-var-formula}，解不等式 $1.25(\rho_1 + (1-\rho_1)/M) < \rho_0 + (1-\rho_0)/M$。}
\end{enumerate}

\exTier{D 层：编程实现与实验验证}{复现 OOB 评估与 GBDT 深度/学习率调参相图}

\begin{enumerate}[label=\textbf{5.\arabic*}, resume]
  \item \textbf{手写随机森林袋外（OOB）验证器}：在代码~\ref{lst:trees-ensemble} 与 \texttt{code/\allowbreak vol4\_experiments.py} 的 5 维非线性回归数据集上，扩展随机森林训练循环：对每棵自助树记录未被抽中的袋外样本掩码 \texttt{oob\_mask}，累计每个训练样本仅由其袋外树给出的预测均值，并计算 $M = 10, 20, 40, 60$ 时的 OOB MSE。验证：（1）每个样本被用作 OOB 的树数占比收敛在 $36.8\%$ 附近；（2）OOB MSE（注意加回真实标签噪声方差 $\sigma^2 = 0.49$）与独立测试集含噪 MSE 的相对偏差在 $10\%$ 以内。
  \exHint{可用 \texttt{in\_bag = np.bincount(b\_idx, minlength=N\_tr) > 0} 一行得到 \texttt{oob\_mask = \textasciitilde in\_bag}，只对至少有 1 棵 OOB 树的样本计算均方误差。}

  \item \textbf{GBDT 过拟合警报与早停（Early Stopping）实验}：复用 \texttt{exp5\_trees\_rf\_gbdt} 的数据，将 GBDT 的迭代轮数拉长至 $M = 250$，比较两组超参数：（A）浅树慢步长（$\text{max\_depth}=3,\ \eta=0.05$）；（B）深树快步长（$\text{max\_depth}=6,\ \eta=0.30$）。分别打印两者的训练 MSE 与测试 MSE 曲线，观察（B）组在多少轮之后测试误差掉头上升（过拟合），并实现一个基于 20\% 留出验证集的早停策略（连续 15 轮验证 MSE 不降即停）。
  \exHint{注意观察（B）组训练 MSE 迅速逼近 $0$ 而测试 MSE 反弹超过 $1.0$ 的拐点，体会“浅树 + 小学习率”为何是 GBDT 的黄金法则。}
\end{enumerate}

\chapter[K-Means、高斯混合与 EM 算法]{K-Means、高斯混合模型与 EM 算法}\label{ch:em}
\chapEpigraph{真正本质的东西，用眼睛是看不见的。}{Antoine de Saint-Exupéry（圣埃克苏佩里，《小王子》）}

\begin{chapIntro}{本章导读}
前两章的 SVM 与集成树都在监督信号的指引下工作：每个样本 $x_i$ 背后都贴着标准答案 $y_i$。现在请想象把全部标签一笔抹去——你面前只剩下一团散落的未标注点云 $\{x_1, \dots, x_N\}$，必须仅凭数据自身的结构推断出隐藏在幕后的类别归属（\textbf{隐变量，$z_i$}）。

本章是第 4 卷的收官压轴之章，也是通向第 5 卷现代深度生成模型的直接桥梁。我们将从最朴素的 \textbf{K-Means 交替极小化}出发，亲眼目睹其球形几何先验在倾斜数据上的惨烈崩塌；接着引入\textbf{高斯混合模型（GMM）}，借第 2 卷的 KL 散度与 Jensen 不等式，一气呵成地推导出贯穿整个机器学习与变分推断史的明珠——\textbf{EM 算法的 ELBO 单调上升定理}。当你读完本章末页时，第 5 卷变分自编码器（VAE）的损失函数已经在你掌中成形。
\end{chapIntro}

% ============================================================
\section{K-Means 聚类与 Lloyd 交替极小化}\label{sec:em-kmeans}
% ============================================================

\subsection{无监督聚类的本质：归纳偏置与隐式几何假设}

在监督学习中，好坏有客观标尺（测试误差）；而在没有标签的\textbf{聚类（Clustering）}任务中，并不存在脱离问题背景的“唯一正确划分”——同一堆水果图片，既可以按颜色分为红黄绿，也可以按形状分为圆与长，还可以按产地分类。从数学上讲，\textbf{无监督聚类本质上是一个欠定（Ill-Posed）问题}：任何聚类算法之所以能产出确定结果，都是因为它在目标函数里悄悄写进了一套关于“什么叫一簇”的\textbf{归纳偏置（Inductive Bias）}。看清这一偏置，才能一眼识破每种算法的适用疆界与失效模式。

\subsection{K-Means 目标函数与 Lloyd 算法的有限步单调收敛证明}

最经典的聚类算法 \textbf{K-Means} 将“一簇”定义为紧致围绕某一个中心点（质心）的欧氏球形点团。给定样本集 $\{x_1, \dots, x_N\} \subset \R^D$ 与预设簇数 $K$，引入二值指示变量 $r_{ik} \in \{0, 1\}$（满足 $\sum_{k=1}^K r_{ik} = 1$，当样本 $i$ 归属第 $k$ 簇时 $r_{ik} = 1$）与 $K$ 个质心 $\mu_1, \dots, \mu_K \in \R^D$，定义\textbf{簇内平方和（Within-Cluster Sum of Squares / Distortion）}目标函数：
\begin{equation}\label{eq:kmeans-obj}
  J\big(\{r_{ik}\},\, \{\mu_k\}\big) := \sum_{i=1}^{N}\sum_{k=1}^{K} r_{ik}\norm{x_i - \mu_k}_2^2.
\end{equation}
由于离散指派 $\{r_{ik}\}$ 与连续质心 $\{\mu_k\}$ 相互耦合，联合最小化式~\eqref{eq:kmeans-obj} 是 NP-难问题。\textbf{Lloyd 算法}（通常直接简称 K-Means 算法）采用极简的\textbf{两步交替坐标极小化}：
\begin{enumerate}
  \item \textbf{指派步（Assignment Step）}：固定当前质心 $\{\mu_k\}$，将每个样本 $x_i$ 划归欧氏距离最近的质心：
  \[
    r_{ik} = \begin{cases}
      1, & \text{若 } k = \argmin_{j \in \{1,\dots,K\}} \norm{x_i - \mu_j}_2^2 \text{（平局时按最小下标打破）}, \\
      0, & \text{否则};
    \end{cases}
  \]
  \item \textbf{更新步（Centroid Update Step）}：固定当前指派 $\{r_{ik}\}$，将每个非空簇的新质心更新为该簇内样本的算术均值：
  \[
    \mu_k = \frac{\sum_{i=1}^{N} r_{ik} x_i}{\sum_{i=1}^{N} r_{ik}}, \qquad k = 1, \dots, K.
  \]
\end{enumerate}

\begin{thmBox}
\begin{theorem}[Lloyd 算法的单调下降与有限步终止性]\label{thm:kmeans-conv}
从任意初始化出发，Lloyd 算法的每一轮指派步与更新步均使目标函数 $J$ 单调不增，且算法必在\textbf{有限步内严格终止}于一个局部驻点。
\end{theorem}
\end{thmBox}

\begin{proof}
\textbf{单调性}：首先固定 $\{\mu_k^{(t)}\}$。由于 $J$ 关于样本下标 $i$ 完全可加，对每个 $i$ 独立挑选使 $\norm{x_i - \mu_k^{(t)}}_2^2$ 最小的唯一 $k$ 置 $r_{ik}^{(t+1)} = 1$，必有
\[
  J\big(\{r_{ik}^{(t+1)}\}, \{\mu_k^{(t)}\}\big) \le J\big(\{r_{ik}^{(t)}\}, \{\mu_k^{(t)}\}\big).
\]
接着固定 $\{r_{ik}^{(t+1)}\}$。此时 $J$ 关于各质心 $\mu_k$ 是严格凸的二次函数且相互解耦，求梯度置零：
\[
  \nabla_{\mu_k} J = -2 \sum_{i=1}^{N} r_{ik}^{(t+1)}(x_i - \mu_k) = 0 \implies \mu_k^{(t+1)} = \frac{\sum_{i=1}^{N} r_{ik}^{(t+1)} x_i}{\sum_{i=1}^{N} r_{ik}^{(t+1)}}.
\]
因此更新步同样达到精确全局最小，有 $J(\{r_{ik}^{(t+1)}\}, \{\mu_k^{(t+1)}\}) \le J(\{r_{ik}^{(t+1)}\}, \{\mu_k^{(t)}\})$。两步合起来，$J^{(t+1)} \le J^{(t)}$ 恒成立。

\textbf{有限步终止性}：注意到若某轮指派 $\{r_{ik}^{(t+1)}\}$ 与上一轮完全相同，更新步所得质心不变，算法立刻停止；反之只要指派发生实质改变且采用确定性平局规则，必有至少一个样本切换到了严格更近的质心（或由前一步严格更新所致），从而 $J$ 严格下降。于是任何一种离散划分配置 $\{r_{ik}\}$ \textbf{至多只能被访问一次}。而 $N$ 个样本分入 $K$ 个簇的全部可能划分总数至多为有限数 $K^N$，因此算法绝不会无限循环，必在有限步内终止！
\end{proof}

\subsection{局部最优陷阱反例与 K-Means++ 初始化}

定理~\ref{thm:kmeans-conv} 只保证收敛到\textbf{局部极小值}，绝不保证全局最优！看一个仅 4 个点的极简反例：

\begin{example}[$3\times 2$ 矩形上的 K-Means 严格次优不动点]\label{ex:kmeans-localmin}
平面上 4 个点构成宽 3、高 2 的矩形顶点：$x_1=(0,0),\ x_2=(0,2),\ x_3=(3,0),\ x_4=(3,2)$，取 $K = 2$。
\begin{itemize}
  \item \textbf{全局最优划分（竖切）}：$\{x_1, x_2\}$ 与 $\{x_3, x_4\}$，质心 $\mu_1=(0,1),\ \mu_2=(3,1)$，每点到自身质心距离平方均为 $1$，总目标 $J^* = 4 \times 1 = \mathbf{4.0}$；
  \item \textbf{次优不动点（横切）}：若初始质心不幸落在上下两边中点附近，形成划分 $\{x_1, x_3\}$ 与 $\{x_2, x_4\}$，质心为 $\mu_1=(1.5, 0),\ \mu_2=(1.5, 2)$。检验 $x_1=(0,0)$ 到两个质心的距离平方：到 $\mu_1$ 为 $1.5^2 = 2.25$，到 $\mu_2$ 为 $1.5^2 + 2^2 = 6.25 > 2.25$！其余三点同理完全对称——\textbf{每个点都严格更靠近当前质心}，Lloyd 算法立刻死锁在此处，目标值高达 $J_{\mathrm{sub}} = 4 \times 2.25 = \mathbf{9.0} = 2.25\, J^*$！
\end{itemize}
\end{example}

为避免把初始质心挤在同一簇或卡死在糟糕的局部极小值，Arthur 与 Vassilvitskii（2007）提出了如今成为全行业标配的 \textbf{K-Means++} 初始化算法：
\begin{enumerate}
  \item 从数据集中均匀随机抽取第 1 个质心 $\mu_1$；
  \item 对每个样本 $x_i$，计算它到\textbf{已选质心集合}的最短欧氏距离 $D(x_i) := \min_{1 \le j \le m} \norm{x_i - \mu_j}_2$；
  \item 以正比于距离平方的概率 $P(x_i) = \frac{D(x_i)^2}{\sum_{l=1}^N D(x_l)^2}$ 抽取下一个质心 $\mu_{m+1}$（越远离已有质心的区域越容易被选中！）；
  \item 重复步骤 2--3 直至选满 $K$ 个质心，再运行标准 Lloyd 迭代。
\end{enumerate}
令人惊叹的是，仅靠这四行 $D(x)^2$ 加权采样，K-Means++ 就能在理论上保证期望目标值满足 $\E[J] \le 8(\ln K + 2) J^*$——即 $O(\log K)$ 竞争比近似保证，并在实践中大幅加速收敛（工程上再辅以多次重启取最小 $J$）。

\begin{failBox}{K-Means 的致命几何硬伤：各向同性球形 Voronoi 剖分}
即使拥有神谕帮你找到 K-Means 的\textbf{全局最优解}，它依然会在非球形数据上惨败！原因藏在指派步 $\argmin_k \norm{x_i - \mu_k}_2^2$ 中：任意两个质心之间的决策边界必是其连线的\textbf{垂直平分超平面}（构成 Voronoi 图）。这等价于强行假设所有簇都是\textbf{等半径、各向同性的正圆球}！当真实数据呈细长倾斜椭圆（例如两个平行的斜雪茄形点云）或两簇方差悬殊时，K-Means 的垂直平分线会毫不留情地拦腰斩断两个真实簇，把每根雪茄的左半段拼成一簇、右半段拼成另一簇——第~\ref{sec:em-gmm-derivation} 节的实测将展示它如何把纯度砸到仅有 \textbf{55.00\%}！
\end{failBox}

% ============================================================
\section{高斯混合模型（GMM）与直接极大似然的困境}\label{sec:em-gmm}
% ============================================================

要让聚类算法自动识别椭圆的伸展方向（协方差）、大小比例与软归属概率，我们必须从硬距离几何升级为\textbf{概率生成模型（Generative Model）}——\textbf{高斯混合模型（Gaussian Mixture Model, GMM）}。

\subsection{隐变量生成故事与边际对数似然}

GMM 设想每个观测样本 $x_i \in \R^D$ 都是通过如下两步随机过程生成的：
\begin{enumerate}
  \item 先按混合权重 $\pi = (\pi_1, \dots, \pi_K)$（满足 $\pi_k \ge 0,\ \sum_{k=1}^K \pi_k = 1$）掷一枚 $K$ 面有偏骰子，抽出样本所属的\textbf{隐藏簇标签（隐变量）} $z_i \sim \operatorname{Cat}(\pi)$，即 $P(z_i = k) = \pi_k$；
  \item 给定 $z_i = k$，再从第 $k$ 个多元高斯分布中采样观测向量：
  \[
    x_i \mid (z_i = k) \;\sim\; \mathcal{N}(x_i \mid \mu_k, \Sigma_k)
    := \frac{1}{(2\pi)^{D/2} \abs{\Sigma_k}^{1/2}} \exp\!\left(-\frac{1}{2}(x_i - \mu_k)^\top \Sigma_k^{-1}(x_i - \mu_k)\right).
  \]
\end{enumerate}
记全部未知参数为 $\theta = \{\pi_k, \mu_k, \Sigma_k\}_{k=1}^K$。由于观测阶段看不到隐变量 $z_i$，根据全概率公式对 $z_i$ 求和消元，得到单样本边际密度 $p_\theta(x_i) = \sum_{k=1}^K \pi_k \mathcal{N}(x_i \mid \mu_k, \Sigma_k)$，进而独立同分布数据集 $X = \{x_1, \dots, x_N\}$ 的\textbf{边际对数似然（Marginal Log-Likelihood）}为：
\begin{equation}\label{eq:gmm-loglik}
  \ell(\theta) := \ln p_\theta(X) = \sum_{i=1}^{N} \ln\!\left(\sum_{k=1}^{K} \pi_k\,\mathcal{N}(x_i \mid \mu_k, \Sigma_k)\right).
\end{equation}

\subsection{为什么直接求导失败？“和的对数”与协方差奇异性陷阱}

如果隐变量 $z_i$ 已知（即变成了监督高斯判别分析），完全数据对数似然 $\sum_i \ln(\pi_{z_i} \mathcal{N}(x_i \mid \mu_{z_i}, \Sigma_{z_i}))$ 的对数会直接把乘积拆成加法，所有参数一秒求得闭式样本均值与样本协方差。但在无监督 GMM 式~\eqref{eq:gmm-loglik} 中，\textbf{对数符号外面包着一层求和号（Log of a Sum）}！一旦对 $\mu_k$ 或 $\Sigma_k$ 求偏导，分母里会冒出纠缠了全部 $K$ 组参数的后验分数，令梯度方程组高度非线性且不存在闭式解析解。

\begin{failBox}{GMM 极大似然的奇异性崩溃（Singularity Collapse）与正则修正}
直接最大化式~\eqref{eq:gmm-loglik} 还暗藏一个极具欺骗性的数学悬崖：假设把第 1 个高斯分量的均值精确对准某一个数据点 $\mu_1 = x_1$，令其协方差矩阵为各向同性收缩阵 $\Sigma_1 = \sigma_1^2 I_D$，其余分量保持固定。则 $x_1$ 对似然的贡献项为
\[
  \pi_1 \mathcal{N}(x_1 \mid x_1, \sigma_1^2 I_D) = \frac{\pi_1}{(2\pi)^{D/2} \sigma_1^D} \xrightarrow{\;\sigma_1 \to 0^+\;} +\infty!
\]
而其余 $N-1$ 个样本仍由其他分量提供有限正概率密度，于是总对数似然 $\ell(\theta) \to +\infty$！这意味着 GMM 的无约束 MLE 根本没有全局有限极大值——只要某个高斯退化成一根插在单点上的狄拉克“细针”，对数似然就会瞬间爆表！

\textbf{工程必做修复}：在每次更新协方差矩阵 $\Sigma_k$ 时加上微小对角岭正则 $\Sigma_k \leftarrow \Sigma_k + \epsilon I_D$（例如配套代码中采用 $\epsilon = 10^{-6}$），既杜绝了奇异性崩溃，也保证了 $\Sigma_k$ 永远严格正定可逆；在贝叶斯框架下，这等价于为协方差矩阵施加逆 Wishart 先验做 MAP 估计。
\end{failBox}

% ============================================================
\section{EM 算法：ELBO 分解与单调上升定理}\label{sec:em-elbo}
% ============================================================

既然直接最大化“和的对数” $\ln \sum_z p_\theta(x, z)$ 如此棘手，能否在每一步构造一个\textbf{把对数推到求和号里面}的紧致下界，通过交替抬升下界来逼迫真实的边际对数似然步步攀升？这就是 1977 年 Dempster、Laird 与 Rubin 集大成的\textbf{期望最大化算法（Expectation-Maximization, EM）}——其核心恒等式堪称整个统计学与深度学习最美的公式之一！

\begin{recallBox}{KL 散度非负性与 Jensen 不等式（第 2 卷回顾）}
对任意两个定义在同一离散隐变量空间上的概率分布 $q(z)$ 与 $p(z)$，Kullback--Leibler（KL）散度定义为
\[
  \DKL(q \,\|\, p) := \sum_{z} q(z) \ln \frac{q(z)}{p(z)} = \E_{q(z)}\!\left[\ln \frac{q(z)}{p(z)}\right] \ge 0,
\]
且等号成立的\textbf{充要条件}是 $q(z) = p(z)$ 处处成立（该结论直接由严格凹函数 $\ln(\cdot)$ 的 Jensen 不等式证得）。
\end{recallBox}

\subsection{皇冠上的明珠：精确 ELBO--KL 分解恒等式}

考虑任意包含观测变量 $x$ 与隐变量 $z$ 的概率模型 $p_\theta(x, z)$（这里 $x, z$ 既可指单个样本，也可指整个数据集 $X, Z$）。引入任意一个关于隐变量 $z$ 的辅助概率分布 $q(z)$（要求在 $p_\theta(z \mid x) > 0$ 处 $q(z) > 0$）。

\begin{thmBox}
\begin{theorem}[精确 ELBO--KL 分解恒等式]\label{thm:elbo-decomp}
对任意参数 $\theta$ 与任意合法变分分布 $q(z)$，观测数据的边际对数似然满足如下精确分解：
\begin{equation}\label{eq:elbo-identity}
  \ln p_\theta(x) = \underbrace{\E_{q(z)}\!\left[\ln \frac{p_\theta(x, z)}{q(z)}\right]}_{\mathcal{L}(q, \theta)\text{（证据下界，ELBO）}} + \DKL\!\big(q(z) \,\|\, p_\theta(z \mid x)\big).
\end{equation}
由于 $\DKL(q \,\|\, p) \ge 0$，泛函 $\mathcal{L}(q, \theta)$ 恒为边际对数似然（亦称模型证据，Evidence）的下界：$\ln p_\theta(x) \ge \mathcal{L}(q, \theta)$，且等号成立\textbf{当且仅当} $q(z) = p_\theta(z \mid x)$（精确后验分布）。
\end{theorem}
\end{thmBox}

\begin{proof}
利用条件概率恒等式 $p_\theta(x, z) = p_\theta(x)\, p_\theta(z \mid x)$ 即 $\ln p_\theta(x) = \ln p_\theta(x, z) - \ln p_\theta(z \mid x)$，对两边关于 $q(z)$ 取期望（左边不依赖 $z$，期望即其自身），再在分子分母同时插入 $q(z)$：
\begin{align*}
  \ln p_\theta(x)
  &= \sum_{z} q(z) \ln p_\theta(x)
   = \sum_{z} q(z) \ln \frac{p_\theta(x, z)}{p_\theta(z \mid x)} \\
  &= \sum_{z} q(z) \ln \left(\frac{p_\theta(x, z)}{q(z)} \cdot \frac{q(z)}{p_\theta(z \mid x)}\right) \\
  &= \sum_{z} q(z) \ln \frac{p_\theta(x, z)}{q(z)} + \sum_{z} q(z) \ln \frac{q(z)}{p_\theta(z \mid x)} \\
  &= \mathcal{L}(q, \theta) + \DKL\!\big(q(z) \,\|\, p_\theta(z \mid x)\big).
\end{align*}
最后由 KL 散度非负性立得 $\ln p_\theta(x) \ge \mathcal{L}(q, \theta)$，且取等当且仅当 $\DKL = 0 \iff q(z) = p_\theta(z \mid x)$。
\end{proof}

请把 ELBO $\mathcal{L}(q, \theta)$ 拆写成两种等价展开形式，它们分别照亮本章的 EM 与第 5 卷的 VAE：
\begin{align}
  \mathcal{L}(q, \theta) &= \E_{q(z)}\big[\ln p_\theta(x, z)\big] + H(q), \label{eq:elbo-em-form}\\
  \mathcal{L}(q, \theta) &= \E_{q(z)}\big[\ln p_\theta(x \mid z)\big] - \DKL\!\big(q(z) \,\|\, p_\theta(z)\big). \label{eq:elbo-vae-form}
\end{align}
其中 $H(q) = -\E_q[\ln q(z)]$ 是 $q$ 的熵。注意式~\eqref{eq:elbo-em-form} 中：\textbf{对数符号已经彻底钻进了期望（求和）内部}，只作用于结构极简的完全数据联合分布 $p_\theta(x, z)$！

\subsection{交替坐标上升与单调不降定理}

EM 算法的全部精髓，就是在双变量泛函 $\mathcal{L}(q, \theta)$ 上对 $q$ 和 $\theta$ 轮流做\textbf{精确坐标上升}：
\begin{itemize}
  \item \textbf{E-步（Expectation Step）}：固定当前参数 $\theta^{(t)}$，最大化 $\mathcal{L}(q, \theta^{(t)})$ 关于分布 $q$。由恒等式~\eqref{eq:elbo-identity}，左边 $\ln p_{\theta^{(t)}}(X)$ 与 $q$ 无关，因此最大化下界 $\mathcal{L}$ 精确等价于把 KL 间隙压为零！唯一全局最优解就是当前参数下的\textbf{精确后验分布}：
  \[
    q^{(t+1)}(Z) := p_{\theta^{(t)}}(Z \mid X), \qquad \text{使得 } \mathcal{L}\big(q^{(t+1)}, \theta^{(t)}\big) = \ln p_{\theta^{(t)}}(X).
  \]
  几何上，这一步构造出了一条在 $\theta = \theta^{(t)}$ 处与真实对数似然曲线\textbf{紧密相切}的替代下界曲面！
  \item \textbf{M-步（Maximization Step）}：固定刚刚算好的 $q^{(t+1)}$，最大化 $\mathcal{L}(q^{(t+1)}, \theta)$ 关于模型参数 $\theta$。由式~\eqref{eq:elbo-em-form}，熵项 $H(q^{(t+1)})$ 与 $\theta$ 无关，故 M-步等价于最大化著名的\textbf{期望完全数据对数似然（$Q$-函数）}：
  \begin{equation}\label{eq:em-qfunc}
    \theta^{(t+1)} := \argmax_{\theta}\; Q(\theta \mid \theta^{(t)}), \qquad
    Q(\theta \mid \theta^{(t)}) := \E_{Z \sim p_{\theta^{(t)}}(Z \mid X)}\big[\ln p_\theta(X, Z)\big].
  \end{equation}
\end{itemize}

\begin{thmBox}
\begin{theorem}[EM 算法的边际对数似然单调上升定理]\label{thm:em-mono}
EM 算法产生的参数序列 $\{\theta^{(t)}\}_{t \ge 0}$ 恒保证真实观测数据的边际对数似然单调不降：
\begin{equation}\label{eq:em-mono}
  \ln p_{\theta^{(t+1)}}(X) \;\ge\; \ln p_{\theta^{(t)}}(X), \qquad \forall\, t \ge 0.
\end{equation}
\end{theorem}
\end{thmBox}

\begin{proof}
只需短短三步不等链，每一环都无懈可击：
\begin{align*}
  \ln p_{\theta^{(t+1)}}(X)
  &\;\ge\; \mathcal{L}\big(q^{(t+1)}, \theta^{(t+1)}\big) && \text{（定理~\ref{thm:elbo-decomp}：ELBO 恒为下界，差值为非负 KL）} \\
  &\;\ge\; \mathcal{L}\big(q^{(t+1)}, \theta^{(t)}\big) && \text{（M-步定义：$\theta^{(t+1)}$ 最大化 $\mathcal{L}(q^{(t+1)}, \theta)$）} \\
  &\;=\; \ln p_{\theta^{(t)}}(X). && \text{（E-步定义：$q^{(t+1)} = p_{\theta^{(t)}}(Z\mid X)$ 使下界紧贴）}
\end{align*}
更精细地看，总增益甚至包含两份非负红利：
\[
  \ln p_{\theta^{(t+1)}}(X) - \ln p_{\theta^{(t)}}(X)
  = \underbrace{\Big[\mathcal{L}\big(q^{(t+1)}, \theta^{(t+1)}\big) - \mathcal{L}\big(q^{(t+1)}, \theta^{(t)}\big)\Big]}_{\ge\, 0\text{（M-步抬升下界）}}
  + \underbrace{\DKL\!\big(p_{\theta^{(t)}}(Z\mid X) \,\|\, p_{\theta^{(t+1)}}(Z\mid X)\big)}_{\ge\, 0\text{（后验漂移带来额外抬升）}}. \qedhere
\]
\end{proof}

图~\ref{fig:em-tangency} 生动描绘了这种“构造紧贴下界 $\to$ 爬上下界顶点 $\to$ 重新紧贴”的优雅台阶式攀登几何（也是主化-极小化 MM 算法的典范）。需注意：真实似然函数 $\ln p_\theta(X)$ 通常是非凹的，因此 EM 与 K-Means 一样只保证收敛到局部极值或驻点，实践中仍需配合 K-Means++ 热启动或多次随机初始化。

\begin{figure}[htbp]
\centering
\begin{tikzpicture}[scale=1.02]
  % Axes
  \draw[->, >=Stealth, slate, thick] (0.0, 0.0) -- (10.2, 0.0) node[right, font=\footnotesize] {参数 $\theta$};
  \draw[->, >=Stealth, slate, thick] (0.0, 0.0) -- (0.0, 4.9) node[above, font=\footnotesize] {对数似然 / 下界};

  % True non-concave log-likelihood curve ln p_theta(X): smooth bump peaking around x=7.3
  % Let's use an explicit smooth formula so tangency points are EXACT!
  % f(x) = 1.2 + 0.85*(x-0.5) - 0.072*(x-0.5)^2 + 0.45*sin(0.85*(x-0.5))?
  % Simpler: f(x) = 4.15 - 0.11*(x - 7.0)^2 + 0.32*sin(0.9*(x - 2.0))
  % Wait: let's define an exact parabola-like concave envelope minus small wave, or exact tangency curves!
  % Even cleaner: draw f(x) = 4.25 - 0.095*(x - 7.1)^2 - 0.18*cos(1.1*(x - 1.5)) -> exact analytical lower bounds!
  % Better yet: define f(x) = 4.20 - 0.085*(x - 7.2)^2 (strictly below max 4.20), wait: make it visibly non-concave on left:
  % Let's pick exact smooth curve:
  % f(x) = 1.15 + 3.05 / (1 + exp(-0.85*(x - 3.6))) - 0.07*(x - 7.4)^2 * (x > 7.4) ... simpler:
  % Use exact f(x) = 0.8 + 0.92*x - 0.068*x^2 + 0.35*sin(0.85*x) over [0.6, 9.5]
  % Let's evaluate tangent parabola g_k(x) = f(x_k) + f'(x_k)*(x - x_k) - c_k*(x - x_k)^2!
  % At x0 = 1.8:
  %   f(1.8) = 0.8 + 1.656 - 0.2203 + 0.35*sin(1.53 rad=87.66 deg -> 0.999) = 2.585
  %   f'(x) = 0.92 - 0.136*x + 0.2975*cos(0.85*x)
  %   f'(1.8) = 0.92 - 0.2448 + 0.2975*0.041 = 0.687
  %   Choose c0 = 0.24 -> peak of g0 at x1 = 1.8 + 0.687/(2*0.24) = 3.23!
  % At x1 = 3.23:
  %   f(3.23) = 0.8 + 2.9716 - 0.7094 + 0.35*sin(2.745 rad -> 0.386) = 3.197
  %   g0(3.23) = 2.585 + 0.687*1.43 - 0.24*(1.43)^2 = 3.077 < 3.197 (strictly below!)
  %   f'(3.23) = 0.92 - 0.4393 + 0.2975*(-0.922) = 0.206 ... wait make slope nicer!
  % Even cleaner: use f(x) = 4.25 - 0.085*(x - 7.0)^2 on [1.0, 9.6]:
  % Wait, add small wave only outside? No, a smooth asymmetric arch:
  % f(x) = 0.65 + 1.08*x - 0.082*x^2 on [0.8, 9.6]:
  %   peak at x* = 1.08/0.164 = 6.585, f(x*) = 4.206!
  %   f(1.8) = 0.65 + 1.944 - 0.2657 = 2.328, f'(1.8) = 1.08 - 0.2952 = 0.7848.
  %   Lower bound 1: g1(x) = f(1.8) + f'(1.8)*(x-1.8) - 0.20*(x-1.8)^2
  %     Since f'' = -0.164 > -0.40 = g1'', g1(x) <= f(x) EVERYWHERE with exact tangency at x=1.8!
  %     Peak of g1: x1 = 1.8 + 0.7848/0.40 = 3.762!
  %     g1(3.762) = 2.328 + 0.7848*1.962 - 0.20*1.962^2 = 3.098;
  %     f(3.762)  = 0.65 + 4.063 - 1.1605 = 3.552! (Nice visible KL gap = 0.454!)
  %   f'(3.762) = 1.08 - 0.164*3.762 = 0.4630.
  %   Lower bound 2: g2(x) = f(3.762) + 0.4630*(x-3.762) - 0.16*(x-3.762)^2
  %     Since f'' = -0.164 > -0.32 = g2'', g2(x) <= f(x) EVERYWHERE with exact tangency at x=3.762!
  %     Peak of g2: x2 = 3.762 + 0.4630/0.32 = 5.209!
  %     g2(5.209) = 3.552 + 0.4630*1.447 - 0.16*1.447^2 = 3.887;
  %     f(5.209)  = 0.65 + 5.626 - 2.225 = 4.051!

  % Plot true marginal log-likelihood f(x)
  \draw[tealDark, very thick, domain=0.8:9.6, samples=80]
    plot (\x, {0.65 + 1.08*\x - 0.082*\x*\x});
  \node[font=\footnotesize\sffamily\bfseries, text=tealDark, anchor=west] at (7.55, 4.28) {真实边际对数似然 $\ln p_\theta(X)$};

  % Plot ELBO surrogate 1: L(q^(1), theta) tangent at x0 = 1.8
  \draw[ochre, thick, dashed, domain=0.65:6.45, samples=60]
    plot (\x, {2.3283 + 0.7848*(\x-1.8) - 0.20*(\x-1.8)*(\x-1.8)});
  \node[font=\scriptsize, text=ochre, anchor=north east] at (6.05, 1.92) {$\mathcal{L}(q^{(1)}, \theta)$};

  % Plot ELBO surrogate 2: L(q^(2), theta) tangent at x1 = 3.762
  \draw[navy, thick, dashed, domain=1.75:8.25, samples=60]
    plot (\x, {3.5525 + 0.4630*(\x-3.762) - 0.16*(\x-3.762)*(\x-3.762)});
  \node[font=\scriptsize, text=navy, anchor=north west] at (7.95, 2.78) {$\mathcal{L}(q^{(2)}, \theta)$};

  % Vertical guides & parameter ticks
  \draw[dotted, slate!70] (1.8, 0) -- (1.8, 2.328);
  \draw[dotted, slate!70] (3.762, 0) -- (3.762, 3.552);
  \draw[dotted, slate!70] (5.209, 0) -- (5.209, 4.051);
  \node[font=\scriptsize, below] at (1.8, -0.05) {$\theta^{(0)}$};
  \node[font=\scriptsize, below] at (3.762, -0.05) {$\theta^{(1)}$};
  \node[font=\scriptsize, below] at (5.209, -0.05) {$\theta^{(2)}$};

  % Highlight points
  \fill[ochre] (1.8, 2.328) circle (2.4pt);
  \fill[ochre] (3.762, 3.098) circle (2.2pt);
  \fill[navy] (3.762, 3.552) circle (2.4pt);
  \fill[navy] (5.209, 3.887) circle (2.2pt);
  \fill[tealDark] (5.209, 4.051) circle (2.4pt);

  % Step arrows showing M-step & E-step KL gap
  \draw[->, >=Stealth, gold, very thick] (1.8, 0.45) -- (3.72, 0.45)
    node[midway, above, font=\tiny\sffamily\bfseries, text=ochre] {M-步 1};
  \draw[->, >=Stealth, gold, very thick] (3.80, 0.45) -- (5.17, 0.45)
    node[midway, above, font=\tiny\sffamily\bfseries, text=navy] {M-步 2};

  % KL vertical gap brace/arrow at theta^(1)
  \draw[<->, >=Stealth, plum, thick] (3.88, 3.10) -- (3.88, 3.55);
  \node[font=\tiny\sffamily, text=plum, anchor=west, align=left] at (3.96, 3.32) {E-步消除 KL 间隙\\令新下界紧贴相切};

  % Callout at tangency theta^(0)
  \node[font=\tiny\sffamily, text=ochre, align=left] at (0.98, 3.10) {紧贴相切点\\$\DKL = 0$};
  \draw[->, >=Stealth, ochre, thin] (1.35, 2.86) -- (1.74, 2.40);
\end{tikzpicture}
\caption{EM 算法的紧致下界攀升几何。每一轮 E-步在当前点 $\theta^{(t)}$ 构造一条严格位于真实对数似然下方且在 $\theta^{(t)}$ 处紧贴相切的代理曲线 $\mathcal{L}(q^{(t+1)}, \theta)$；M-步直接跳到该易优化下界的顶点 $\theta^{(t+1)}$，由此保证真实似然 $\ln p_\theta(X)$ 稳步单调攀升。}
\label{fig:em-tangency}
\end{figure}

% ============================================================
\section{GMM 的闭式 E-步、M-步推导与倾斜椭圆实测}\label{sec:em-gmm-derivation}
% ============================================================

现在将通用 EM 框架精准套用到第~\ref{sec:em-gmm} 节的高斯混合模型上，体验每一步都化为闭式解析公式的丝滑快感。

\subsection{GMM 的 E-步（响应度）与 M-步闭式推导}

\paragraph{E-步：计算后验响应度（Responsibilities）}
给定当前参数 $\theta^{(t)} = \{\pi_k^{(t)}, \mu_k^{(t)}, \Sigma_k^{(t)}\}_{k=1}^K$，各样本条件独立，由贝叶斯定理直接算出样本 $x_i$ 由第 $k$ 个高斯分量生成的后验概率 $\gamma_{ik} := p_{\theta^{(t)}}(z_i = k \mid x_i)$（称为分量 $k$ 对样本 $i$ 的\textbf{响应度}）：
\begin{equation}\label{eq:gmm-estep}
  \gamma_{ik} = \frac{\pi_k^{(t)}\,\mathcal{N}\big(x_i \mid \mu_k^{(t)}, \Sigma_k^{(t)}\big)}{\sum_{j=1}^{K} \pi_j^{(t)}\,\mathcal{N}\big(x_i \mid \mu_j^{(t)}, \Sigma_j^{(t)}\big)}, \qquad i=1,\dots,N,\; k=1,\dots,K.
\end{equation}
\textbf{数值稳定性提示}：在高维空间，高斯密度值 $\mathcal{N}(x_i \mid \mu_k, \Sigma_k)$ 极易下溢为浮点零 $0.0$ 导致 $0/0 = \texttt{NaN}$！工程实现必须全程在对数域计算 $a_{ik} := \ln \pi_k + \ln \mathcal{N}(x_i \mid \mu_k, \Sigma_k)$，并借 \textbf{Log-Sum-Exp 技巧}（减去每行最大值 $m_i = \max_j a_{ij}$）：
\[
  \ln p_\theta(x_i) = m_i + \ln\!\left(\sum_{j=1}^{K} e^{a_{ij} - m_i}\right), \qquad
  \gamma_{ik} = \exp\!\big(a_{ik} - \ln p_\theta(x_i)\big).
\]

\paragraph{M-步：最大化 $Q(\theta \mid \theta^{(t)})$ 求闭式参数更新}
写出 GMM 的期望完全数据对数似然：
\begin{align}
  Q(\theta \mid \theta^{(t)})
  &= \sum_{i=1}^{N}\sum_{k=1}^{K} \gamma_{ik}\ln\!\Big(\pi_k\,\mathcal{N}(x_i \mid \mu_k, \Sigma_k)\Big) \nonumber\\
  &= \sum_{i=1}^{N}\sum_{k=1}^{K} \gamma_{ik}\!\left[\ln \pi_k - \frac{D}{2}\ln(2\pi) - \frac{1}{2}\ln\abs{\Sigma_k} - \frac{1}{2}(x_i - \mu_k)^\top \Sigma_k^{-1}(x_i - \mu_k)\right]. \label{eq:gmm-qfunc}
\end{align}
记第 $k$ 个分量的\textbf{有效样本数（Effective Count）} $N_k := \sum_{i=1}^{N} \gamma_{ik}$（显然 $\sum_{k=1}^K N_k = N$）。下面逐一求出三组参数的闭式最优解：

\begin{enumerate}
  \item {\rightskip=0pt plus 1em\relax \textbf{更新混合权重 $\pi_k$}（直接动用第 2 章拉格朗日乘子法！）：在等式约束 $\sum_{k=1}^K \pi_k = 1$ 下最大化 $\sum_{k=1}^K N_k \ln \pi_k$，构造拉格朗日函数 $\sum_k N_k \ln \pi_k + \lambda(1 - \sum_k \pi_k)$，求偏导置零得 $\frac{N_k}{\pi_k} - \lambda = 0 \implies \pi_k = N_k / \lambda$。对 $k$ 求和利用 $\sum_k \pi_k = 1$ 得 $\lambda = \sum_k N_k = N$，故\par}
  \begin{equation}\label{eq:gmm-m-pi}
    \pi_k^{(t+1)} = \frac{N_k}{N}.
  \end{equation}

  \item \textbf{更新均值向量 $\mu_k$}：对式~\eqref{eq:gmm-qfunc} 关于 $\mu_k$ 求梯度置零：
  \[
    \nabla_{\mu_k} Q = \sum_{i=1}^{N} \gamma_{ik} \Sigma_k^{-1}(x_i - \mu_k) = 0.
  \]
  左乘非奇异矩阵 $\Sigma_k$，立得加权样本均值：
  \begin{equation}\label{eq:gmm-m-mu}
    \mu_k^{(t+1)} = \frac{1}{N_k} \sum_{i=1}^{N} \gamma_{ik} x_i.
  \end{equation}

  \item \textbf{更新协方差矩阵 $\Sigma_k$}：为简化矩阵求导，令精度矩阵 $\Lambda_k := \Sigma_k^{-1} \succ 0$，利用 $\ln\abs{\Sigma_k} = -\ln\abs{\Lambda_k}$ 与矩阵微积分恒等式（第 1 卷回顾）$\nabla_{\Lambda} \ln\abs{\Lambda} = \Lambda^{-1}$、$\nabla_{\Lambda}(v^\top \Lambda v) = v v^\top$，将式~\eqref{eq:gmm-qfunc} 关于 $\Lambda_k$ 求梯度置零：
  \[
    \nabla_{\Lambda_k} Q = \frac{1}{2} N_k \Lambda_k^{-1} - \frac{1}{2}\sum_{i=1}^{N} \gamma_{ik}(x_i - \mu_k)(x_i - \mu_k)^\top = 0.
  \]
  将 $\Lambda_k^{-1} = \Sigma_k$ 代回并整理，立得加权样本协方差矩阵（工程中附加 $\epsilon I_D$ 防奇异）：
  \begin{equation}\label{eq:gmm-m-sigma}
    \Sigma_k^{(t+1)} = \frac{1}{N_k}\sum_{i=1}^{N} \gamma_{ik}\big(x_i - \mu_k^{(t+1)}\big)\big(x_i - \mu_k^{(t+1)}\big)^\top.
  \end{equation}
\end{enumerate}

\subsection{大统一桥梁：硬 K-Means 正是 GMM-EM 的零方差极限！}

\begin{thmBox}
\begin{proposition}[K-Means 作为 GMM-EM 的各向同性退化极限]\label{prop:kmeans-gmm-limit}
在 GMM 中固定等混合权重 $\pi_k \equiv 1/K$，并强制所有分量的协方差矩阵共享同一各向同性标量方差 $\Sigma_k \equiv \sigma^2 I_D$（不更新 $\pi_k, \Sigma_k$）。当 $\sigma^2 \to 0^+$ 时，GMM-EM 的 E-步与 M-步逐行精确退化为 K-Means（Lloyd 算法）的硬指派步与质心均值更新步！
\end{proposition}
\end{thmBox}

\begin{proof}
将 $\pi_k = 1/K$ 与 $\Sigma_k = \sigma^2 I_D$ 代入 E-步响应度公式~\eqref{eq:gmm-estep}，归一化常数 $(2\pi\sigma^2)^{-D/2} / K$ 上下约去：
\[
  \gamma_{ik} = \frac{\exp\!\left(-\dfrac{\norm{x_i - \mu_k}_2^2}{2\sigma^2}\right)}{\sum_{j=1}^{K} \exp\!\left(-\dfrac{\norm{x_i - \mu_j}_2^2}{2\sigma^2}\right)}.
\]
记 $k^*(i) := \argmin_j \norm{x_i - \mu_j}_2^2$（假设最近质心唯一），分子分母同乘 $\exp(\norm{x_i - \mu_{k^*(i)}}_2^2 / (2\sigma^2))$。对所有 $j \neq k^*(i)$，指数上出现严格负差值 $-(\norm{x_i - \mu_j}_2^2 - \norm{x_i - \mu_{k^*(i)}}_2^2)/(2\sigma^2) \to -\infty$（当 $\sigma^2 \to 0^+$），故对应指数项趋于 $0$！于是
\[
  \lim_{\sigma^2 \to 0^+} \gamma_{ik} = \begin{cases} 1, & k = k^*(i), \\ 0, & k \neq k^*(i), \end{cases}
\]
软响应度 $\gamma_{ik}$ 精确坍缩为 K-Means 的 $0/1$ 硬指派 $r_{ik}$！再代入 M-步均值公式~\eqref{eq:gmm-m-mu}，立刻化为 $\mu_k = \frac{\sum_i r_{ik} x_i}{\sum_i r_{ik}}$——这正是 Lloyd 算法的质心更新公式！
\end{proof}

这一命题彻底揭示了 K-Means 与 GMM 的血缘关系：K-Means 绝不是与概率模型无关的孤立启发式，它只是把 GMM 的全部自由度阉割为“等权重 + 等方差球面 + 零温度硬指派”后的极简特例！

\subsection{手工完整推演：一维 4 点二分量 GMM 的一轮 EM 迭代}

\begin{example}[手算 1D GMM 一轮完整 EM 与似然上升验证]\label{ex:gmm-hand}
一维轴上有 4 个观测点：$x = (-2,\, -1,\, 2,\, 3)$，拟合 $K = 2$ 分量一维 GMM。设初始参数为 $\pi^{(0)} = (0.5,\, 0.5)$，$\mu^{(0)} = (-1,\, 1)$，$(\sigma^2)^{(0)} = (1,\, 1)$。请手工完成第 1 轮 E-步与 M-步，并验算边际对数似然严格上升。
\end{example}

\paragraph{第 1 步：初始对数似然与 E-步响应度}
由于 $\pi_1^{(0)} = \pi_2^{(0)} = 0.5$ 且 $(\sigma_1^2)^{(0)} = (\sigma_2^2)^{(0)} = 1$，对数后验几率之比戏剧性地化简为线性函数：
\[
  \ln\frac{\pi_1^{(0)}\mathcal{N}(x_i\mid -1, 1)}{\pi_2^{(0)}\mathcal{N}(x_i\mid 1, 1)}
  = -\frac{1}{2}\big[(x_i + 1)^2 - (x_i - 1)^2\big] = -2 x_i \implies \gamma_{i1} = \frac{1}{1 + e^{2x_i}} = \sigma(-2 x_i)!
\]
逐点代入 $x = (-2, -1, 2, 3)$（同时记 $\gamma_{i2} = 1 - \gamma_{i1}$）：
\begin{align*}
  \gamma_{11} &= \sigma(4) = 0.9820, & \gamma_{21} &= \sigma(2) = 0.8808, & \gamma_{31} &= \sigma(-4) = 0.0180, & \gamma_{41} &= \sigma(-6) = 0.0025; \\
  \gamma_{12} &= 0.0180, & \gamma_{22} &= 0.1192, & \gamma_{32} &= 0.9820, & \gamma_{42} &= 0.9975.
\end{align*}
初始边际对数似然直接求和得 $\ell(\theta^{(0)}) = \mathbf{-9.2826}$。

\paragraph{第 2 步：M-步闭式更新}
先算有效样本数：$N_1 = 0.9820 + 0.8808 + 0.0180 + 0.0025 = 1.8833$，$N_2 = 4 - N_1 = 2.1167$。代入式~\eqref{eq:gmm-m-pi}--\eqref{eq:gmm-m-sigma}：
\begin{align*}
  \pi_1^{(1)} &= \frac{1.8833}{4} = 0.4708, \qquad \pi_2^{(1)} = \frac{2.1167}{4} = 0.5292; \\
  \mu_1^{(1)} &= \frac{0.9820(-2) + 0.8808(-1) + 0.0180(2) + 0.0025(3)}{1.8833} = \mathbf{-1.4875}, \\
  \mu_2^{(1)} &= \frac{0.0180(-2) + 0.1192(-1) + 0.9820(2) + 0.9975(3)}{2.1167} = \mathbf{2.2683}; \\
  (\sigma_1^2)^{(1)} &= \frac{1}{N_1}\sum_{i=1}^4 \gamma_{i1}(x_i + 1.4875)^2 = \mathbf{0.3907}, \\
  (\sigma_2^2)^{(1)} &= \frac{1}{N_2}\sum_{i=1}^4 \gamma_{i2}(x_i - 2.2683)^2 = \mathbf{1.0420}.
\end{align*}
仅经过\textbf{一轮}迭代，两个均值已从 $(-1, 1)$ 迅猛逼近真实簇中心 $(-1.5, 2.5)$，且新边际对数似然暴涨至 $\ell(\theta^{(1)}) = \mathbf{-6.4826}$（单步净增 $+2.8000$ 奈特！），完美兑现定理~\ref{thm:em-mono}。

\subsection{倾斜椭圆双簇实测：K-Means 的崩塌与 GMM-EM 的全胜}

现在检阅配套脚本 \texttt{code/\allowbreak vol4\_experiments.py} 中 \texttt{exp6\_kmeans\_vs\_gmm\_em} 的实测数据（日志存于 \texttt{data/\allowbreak gmm\_em.csv}）。

\paragraph{实验设置}
在二维平面上生成两个平行倾斜的细长高斯点团（每簇 $180$ 点，共 $N = 360$ 点），真实均值分别为 $\mu_0 = (-1.0, 1.2)^\top$ 与 $\mu_1 = (1.2, -1.0)^\top$，共享强正相关倾斜协方差矩阵
\[
  \Sigma_{\mathrm{tilt}} = \begin{bmatrix} 4.0 & 2.8 \\ 2.8 & 2.2 \end{bmatrix}
  \qquad (\text{相关系数 } \rho \approx 0.944\text{，呈细长 } 45^\circ \text{ 斜椭圆！}).
\]
对比 5 次随机重启取最优惯量的 K-Means 与带 $\epsilon = 10^{-6}$ 协方差正则的 GMM-EM：
\begin{itemize}
  \item \textbf{K-Means（5 次重启最佳）}：聚类纯度仅 \textbf{55.00\%}（接近随机乱分！因为两个细长雪茄沿主轴长度远超它们之间的垂直间距，K-Means 毫不客气地用一条垂直于主轴的平分线把两根雪茄同时拦腰腰斩）；
  \item \textbf{GMM-EM（30 轮迭代）}：聚类纯度高达 \textbf{100.00\%}（全部 360 个点无一错分！）。表~\ref{tab:gmm-em-iters} 与图~\ref{fig:gmm-em-curve}（直接读自 \texttt{data/\allowbreak gmm\_em.csv}）展示了边际对数似然的严格单调攀升轨迹：从第 1 轮的 $-1578.57$ 起步，在第 5--7 轮发生戏剧性“相变”（两个高斯椭圆成功旋转对齐到倾斜主轴！），对数似然飙升逾 $158$ 奈特，在第 8 轮精准收敛至 \textbf{-1286.45} 并纹丝不动。
\end{itemize}

\begin{figure}[htbp]
\centering
\begin{tikzpicture}
\begin{axis}[
  width=0.84\linewidth,
  height=5.9cm,
  xlabel={EM 迭代轮数 $t$},
  ylabel={边际对数似然 $\ln p_{\theta^{(t)}}(X)$},
  xmin=0, xmax=31,
  ymin=-1610, ymax=-1245,
  xtick={1,5,10,15,20,25,30},
  tick label style={font=\scriptsize},
  label style={font=\footnotesize},
  grid=both,
  grid style={line width=0.25pt, draw=slate!18},
  legend pos=south east,
  legend style={font=\scriptsize, fill=white, fill opacity=0.92, draw=tealLine, rounded corners=1.5pt}
]
  \addplot[tealDark, very thick, mark=*, mark size=1.6pt, mark repeat=1]
    table[x=iter, y=loglik, col sep=comma] {data/gmm_em.csv};
  \addlegendentry{GMM-EM 边际对数似然（严格单调不降）}

  % Converged plateau line
  \addplot[gold, thick, dashed, domain=1:30, samples=2] {-1286.4522};
  \addlegendentry{收敛极大值 $-1286.45$（聚类纯度 $100.00\%$）}

  % Callout for phase transition (iters 5->7)
  \draw[->, >=Stealth, ochre, thick] (axis cs:10.5, -1415) -- (axis cs:6.3, -1350);
  \node[font=\scriptsize, text=ochre, align=left, fill=white, inner sep=1.5pt]
    at (axis cs:15.8, -1432) {第 5--7 轮椭圆协方差旋转对齐\\似然跃升逾 $158$ 奈特！};
\end{axis}
\end{tikzpicture}
\caption{双倾斜椭圆簇（$N=360$）上 GMM-EM 的边际对数似然实测收敛曲线（直接读自 \texttt{data/\allowbreak gmm\_em.csv}）。整条轨迹严格满足 $\ln p_{\theta^{(t+1)}}(X) \ge \ln p_{\theta^{(t)}}(X)$（单步零回退），仅 8 轮即彻底对齐倾斜协方差并实现 100.00\% 聚类纯度，而 K-Means 纯度仅 55.00\%。}
\label{fig:gmm-em-curve}
\end{figure}

\begin{table}[htbp]
\centering
\caption{双倾斜椭圆簇（$N=360$）实测对比与 GMM-EM 关键轮次对数似然快照}
\label{tab:gmm-em-iters}
\small
\begin{tabular}{@{}lcccccc@{}}
\toprule
\textbf{GMM-EM 迭代轮次} & $t=1$ & $t=2$ & $t=5$ & $t=6$ & $t=7$ & $t=8\text{--}30$ \\
\midrule
边际对数似然 $\ell(\theta^{(t)})$ & $-1578.57$ & $-1514.27$ & $-1445.31$ & $-1322.18$ & $-1286.50$ & $\mathbf{-1286.45}$ \\
\midrule
\multicolumn{7}{@{}l@{}}{\footnotesize 最终聚类纯度对比：\quad K-Means（5 次重启最佳）$=\mathbf{55.00\%}$ \quad vs \quad GMM-EM $=\mathbf{100.00\%}$。} \\
\bottomrule
\end{tabular}
\end{table}

\subsection{可运行的数值稳定纯 NumPy GMM-EM 引擎}

下面的代码片段展示了采用 Log-Sum-Exp 技巧、\texttt{slogdet} 稳定对数行列式与 $\epsilon I$ 协方差正则的完整 GMM-EM 核心循环，并内置单调上升断言自动排错。

\begin{lstlisting}[caption={数值稳定的纯 NumPy GMM-EM 实现（含 Log-Sum-Exp 与单调性断言）},label={lst:gmm-em}]
import numpy as np

def log_gaussian_pdf(X, mu, Sigma):
    D = X.shape[1]
    diff = X - mu
    _, logdet = np.linalg.slogdet(Sigma)
    invS = np.linalg.inv(Sigma)
    quad = np.sum((diff @ invS) * diff, axis=1)
    return -0.5 * (D * np.log(2.0 * np.pi) + logdet + quad)

def fit_gmm_em(X, K=2, n_iters=30, reg=1e-6, seed=2026):
    rng = np.random.default_rng(seed)
    N, D = X.shape
    pis = np.full(K, 1.0 / K)
    mus = X[rng.choice(N, size=K, replace=False)].copy()
    Sigmas = np.array([np.cov(X, rowvar=False) + reg * np.eye(D) for _ in range(K)])
    history, prev_ll = [], -float("inf")

    for _ in range(n_iters):
        # E-step in log-space via Log-Sum-Exp
        log_prob = np.column_stack(
            [np.log(pis[k] + 1e-12) + log_gaussian_pdf(X, mus[k], Sigmas[k]) for k in range(K)]
        )
        m = np.max(log_prob, axis=1, keepdims=True)
        log_norm = m + np.log(np.sum(np.exp(log_prob - m), axis=1, keepdims=True))
        ll = float(np.sum(log_norm))
        assert ll >= prev_ll - 1e-7, f"ELBO theorem violated: {prev_ll} -> {ll}"
        prev_ll = ll; history.append(ll)
        gamma = np.exp(log_prob - log_norm)

        # M-step closed-form updates + diagonal ridge regularization
        Nk = gamma.sum(axis=0)
        pis = Nk / N
        for k in range(K):
            mus[k] = (gamma[:, k] @ X) / Nk[k]
            diff = X - mus[k]
            Sigmas[k] = (diff.T * gamma[:, k]) @ diff / Nk[k] + reg * np.eye(D)
    return pis, mus, Sigmas, gamma, history
\end{lstlisting}

\begin{insightBox}{直通第 5 卷：从本章 EM 的 ELBO 一步跨入变分自编码器（VAE）}
请把式~\eqref{eq:elbo-vae-form} 再读一遍：
\[
  \ln p_\theta(x) \;\ge\; \mathcal{L}(q, \theta) = \E_{q(z)}\big[\ln p_\theta(x \mid z)\big] - \DKL\!\big(q(z) \,\|\, p_\theta(z)\big).
\]
在 GMM 中，隐变量 $z \in \{1,\dots,K\}$ 只是离散有限集合，因此 E-步可以直接用贝叶斯公式算出\textbf{精确后验} $q(z) = p_\theta(z \mid x)$。但进入第 5 卷的深度生成模型后，隐变量变成了高维连续语义向量 $z \in \R^d$（例如描述一张人脸的姿态、光照与表情），解码器 $p_\theta(x \mid z)$ 是一个多层深度神经网络——此时分母积分 $p_\theta(x) = \int p_\theta(x \mid z) p(z)\,\mathrm{d}z$ 完全不可积，精确后验 $p_\theta(z \mid x)$ 根本算不出来！

\textbf{变分自编码器（Variational Autoencoder, VAE）}（Kingma \& Welling, 2013）的破局之道极其干脆：既然算不出精确后验，就训练另一个神经网络\textbf{编码器} $q_\phi(z \mid x) = \mathcal{N}(\mu_\phi(x), \diag(\sigma_\phi^2(x)))$（称为\textbf{摊销变分推断，Amortized Variational Inference}）来近似它，然后用随机梯度上升（配重参数化技巧）对编码器参数 $\phi$ 和解码器参数 $\theta$ \textbf{联合最大化同一个 ELBO 目标}：
\[
  \max_{\theta,\, \phi}\; \underbrace{\E_{z \sim q_\phi(z \mid x)}\big[\ln p_\theta(x \mid z)\big]}_{\text{图像/信号重建项}} - \underbrace{\DKL\!\big(q_\phi(z \mid x) \,\|\, \mathcal{N}(0, I)\big)}_{\text{潜空间正则对齐项}}.
\]
你看，跨越三十六年的经典 EM 算法与现代深度生成模型 VAE，骨子里跳动的竟是\textbf{同一颗 ELBO 心脏}！
\end{insightBox}

\begin{extBox}{广义 EM（GEM）、变分推断与如何用 BIC 选择簇数 $K$}
\begin{enumerate}
  \item \textbf{广义 EM（Generalized EM, GEM）}：回头检视定理~\ref{thm:em-mono} 的三行证明——你其实并不要求 M-步必须找到 $Q(\theta \mid \theta^{(t)})$ 的全局极大值！只要找到一个使 $Q(\theta^{(t+1)} \mid \theta^{(t)}) \ge Q(\theta^{(t)} \mid \theta^{(t)})$ 的新参数（例如对非闭式 $Q$ 只走几步梯度上升），边际对数似然单调上升定理\textbf{照样严格成立}！
  \item \textbf{如何为 GMM 自动选择分量个数 $K$？}不能直接挑边际对数似然 $\ell(\hat{\theta}_K)$ 最大的 $K$——因为多加一个高斯分量只会把似然越撑越高（直至每个点自成一簇过拟合）。实践中通常采用带模型参数个数 $p_K$ 惩罚的\textbf{贝叶斯信息准则（BIC）}或\textbf{赤池信息准则（AIC）}：
  \[
    \operatorname{BIC}(K) = -2\ell(\hat{\theta}_K) + p_K \ln N, \qquad
    \operatorname{AIC}(K) = -2\ell(\hat{\theta}_K) + 2 p_K,
  \]
  其中 $D$ 维满协方差 $K$ 分量 GMM 的独立参数个数为 $p_K = (K - 1) + K D + K \frac{D(D+1)}{2}$，挑选使 $\operatorname{BIC}(K)$ 最小的 $K$ 即可在拟合优度与模型复杂度之间取得优雅平衡。
\end{enumerate}
\end{extBox}

% ============================================================
\section{本章习题}\label{sec:ch6-ex}
% ============================================================

\exTier{A 层：基础巩固}{掌握 K-Means/GMM 参数计数、初始化与奇异性诊断}

\begin{enumerate}[label=\textbf{6.\arabic*}]
  \item \textbf{GMM 独立参数量计数与自由度对比}：在特征维度 $D = 10$、簇数 $K = 4$ 的数据上，分别计算以下三种聚类/混合模型所需的独立连续参数总数：（1）K-Means；（2）对角协方差 GMM（每个 $\Sigma_k = \diag(\sigma_{k1}^2, \dots, \sigma_{kD}^2)$）；（3）满协方差 GMM（每个 $\Sigma_k$ 为任意对称正定阵）。由此解释为何在高维小样本（$N$ 较小、$D$ 较大）场景下常优先选用对角协方差或球形协方差 GMM。
  \exHint{别忘了混合权重 $\pi_1, \dots, \pi_K$ 受 $\sum_k \pi_k = 1$ 约束只有 $K-1$ 个自由度，对称 $D\times D$ 矩阵有 $D(D+1)/2$ 个独立元素。}

  \item \textbf{K-Means++ 抽样概率手算}：一维轴上有 4 个样本点 $x = (0,\, 1,\, 4,\, 6)$。假设 K-Means++ 的第一步随机抽中了 $x_1 = 0$ 作为第一个质心 $\mu_1$。（1）计算其余三个样本被选为第二个质心 $\mu_2$ 的精确概率；（2）若第二步恰好抽中 $x_4 = 6$，继续计算剩余两个点被选为第三个质心 $\mu_3$ 的精确概率。
  \exHint{第二步计算到已有集合 $\{0, 6\}$ 的\textbf{最短距离} $D(x_i) = \min(|x_i - 0|, |x_i - 6|)$，再按 $D(x_i)^2$ 归一化。}
\end{enumerate}

\exTier{B 层：方法与推导}{熟练手算 K-Means 迭代、伯努利混合 EM 与 ELBO 拆分}

\begin{enumerate}[label=\textbf{6.\arabic*}, resume]
  \item \textbf{硬币投掷（伯努利混合）模型的 EM 推导}：设有两枚质地不均的硬币 $A$ 与 $B$，正面朝上概率分别为 $\theta_A, \theta_B$，每次实验等概率选一枚硬币（$\pi_A = \pi_B = 0.5$ 已知）连续抛掷 $L$ 次，但观测者只记录到第 $i$ 轮出现的正面次数 $h_i$（$i = 1, \dots, N$），不知道选的是哪枚硬币（隐变量 $z_i \in \{A, B\}$）。请完整推导该模型的：（1）E-步后验响应度 $\gamma_{iA} = P(z_i = A \mid h_i; \theta_A^{(t)}, \theta_B^{(t)})$；（2）M-步闭式更新公式 $\theta_A^{(t+1)}, \theta_B^{(t+1)}$。
  \exHint{给定硬币 $k$，$P(h_i \mid z_i = k) = \binom{L}{h_i}\theta_k^{h_i}(1-\theta_k)^{L-h_i}$；在 E-步分子分母中组合数 $\binom{L}{h_i}$ 会直接约去。}

  \item \textbf{例~\ref{ex:gmm-hand} 的对称检验与硬 K-Means 对比}：在例~\ref{ex:gmm-hand} 的一维数据 $x = (-2, -1, 2, 3)$ 上，以相同初始中心 $\mu^{(0)} = (-1, 1)$ 运行标准 K-Means（Lloyd 算法）。（1）写出全部迭代轮次的指派结果、质心变化与最终目标值 $J$；（2）解释为何在此例中 K-Means 1 轮即收敛到最终簇中心 $(-1.5, 2.5)$，而 GMM-EM 第 1 轮给出的是 $(-1.4875, 2.2683)$——两者差异的来源是什么？
  \exHint{注意样本 $-1$ 在软响应度下仍有 $11.92\%$ 概率归属右簇，从而把右簇均值略微向左拉了一点。}
\end{enumerate}

\exTier{C 层：理论证明与反例构造}{深化 ELBO 变分性质、EM 驻点与非球形 K-Means 边界}

\begin{enumerate}[label=\textbf{6.\arabic*}, resume]
  \item \textbf{ELBO 在切点处的一阶梯度一致性}：设对数似然 $\ell(\theta) = \ln p_\theta(x)$ 与下界 $g(\theta) := \mathcal{L}(q^{(t+1)}, \theta)$（其中 $q^{(t+1)}(z) = p_{\theta^{(t)}}(z \mid x)$）均在 $\theta^{(t)}$ 处可微。证明在切点 $\theta = \theta^{(t)}$ 处，两者不仅函数值相等，而且\textbf{一阶梯度也完全相等}：
  \[
    \nabla_\theta \mathcal{L}\big(q^{(t+1)}, \theta\big)\Big|_{\theta = \theta^{(t)}} = \nabla_\theta \ln p_\theta(x)\Big|_{\theta = \theta^{(t)}}.
  \]
  由此推论：若 $\theta^*$ 是 EM 算法的不动点（在无约束内部达到 M-步极值），则 $\theta^*$ 必为真实对数似然的驻点 $\nabla_\theta \ell(\theta^*) = 0$。
  \exHint{定义差值函数 $d(\theta) := \ln p_\theta(x) - \mathcal{L}(q^{(t+1)}, \theta) = \DKL(q^{(t+1)} \,\|\, p_\theta(z\mid x))$。由于 $d(\theta) \ge 0$ 恒成立且 $d(\theta^{(t)}) = 0$，故 $\theta^{(t)}$ 是可微非负函数 $d(\theta)$ 的无约束全局极小值点，必有 $\nabla_\theta d(\theta^{(t)}) = 0$！}

  \item \textbf{高斯混合马氏距离决策边界}：对于 $K = 2$ 个等权重（$\pi_1 = \pi_2 = 0.5$）二维高斯分量 $\mathcal{N}(\mu_1, \Sigma_1)$ 与 $\mathcal{N}(\mu_2, \Sigma_2)$，推导后验边界 $\gamma_{i1} = \gamma_{i2} = 0.5$ 的方程。证明：（1）当两簇共享同一非球形协方差阵 $\Sigma_1 = \Sigma_2 = \Sigma$ 时，决策边界仍是\textbf{直线}（仿射超平面），但法向量变为 $\Sigma^{-1}(\mu_1 - \mu_2)$（不再垂直于质心连线 $\mu_1 - \mu_2$！）；（2）当 $\Sigma_1 \neq \Sigma_2$ 时，决策边界一般为二次曲线（椭圆、抛物线或双曲线）。
  \exHint{化简对数似然比 $\ln \mathcal{N}(x\mid \mu_1, \Sigma_1) = \ln \mathcal{N}(x\mid \mu_2, \Sigma_2)$，观察二次项 $x^\top(\Sigma_2^{-1} - \Sigma_1^{-1})x$ 何时相消。}
\end{enumerate}

\exTier{D 层：编程实现与实验验证}{亲手实现 K-Means++、BIC 选簇与协方差奇异复现}

\begin{enumerate}[label=\textbf{6.\arabic*}, resume]
  \item \textbf{K-Means++ 实现与多重启稳定性对比}：用纯 NumPy 实现 K-Means++ 初始化函数与 Lloyd 算法。构造 5 个分离的二维高斯团簇（每簇 60 点，$N = 300, K = 5$），分别运行 50 次“完全随机均匀初始化 K-Means”与 50 次“K-Means++ 初始化”，统计并比较两者的：（1）最终惯量 $J$ 的均值、中位数与最差值；（2）收敛所需平均迭代轮数。验证 K-Means++ 显著减少落入劣质局部极小值的频率。
  \exHint{注意 K-Means++ 抽样概率 $p = D^2 / \sum D^2$ 可直接传入 \texttt{rng.choice(N, p=p)}，已选点距离为 $0$ 故绝不会被重复抽中。}

  \item \textbf{GMM-EM 协方差类型对比与 BIC 自动选 $K$ 实验}：基于代码~\ref{lst:gmm-em} 与 \texttt{code/\allowbreak vol4\_experiments.py} 中的倾斜双椭圆数据（$N = 360$）：（1）将代码~\ref{lst:gmm-em} 修改为对角协方差版（只保留 $\Sigma_k$ 对角线），观察聚类纯度相比满协方差版的 $100.00\%$ 下降了多少，并从几何上解释原因；（2）在满协方差 GMM 上遍历 $K \in \{1, 2, 3, 4, 5\}$，计算每组的最终对数似然与 $\operatorname{BIC}(K)$，验证虽然 $\ell(\hat{\theta}_K)$ 随 $K$ 增大持续上升，但 $\operatorname{BIC}(K)$ 精准在真实簇数 $K = 2$ 处取得全局最小值！
  \exHint{二维满协方差 GMM 独立参数个数为 $p_K = (K-1) + 2K + 3K = 6K - 1$；每个 $K$ 建议跑 3 次随机种子取最大似然以避免偶然局部极值。}
\end{enumerate}


% ==================== Back Matter / Appendix ====================
\appendix
% !TeX root = ../main.tex
\renewcommand{\chapWatermarkNum}{\thechapter}
\ctexset{chapter={name={附录\ ,}, number=\Alph{chapter}}}
\renewcommand{\chaptermark}[1]{\markboth{附录 \thechapter\quad #1}{}}
\chapter[全卷 48 道分层习题完整精解]{全卷 48 道分层习题完整精解}
\label{app:solutions}
\chapEpigraph{计算的目的不在于数字，而在于洞察。}{Richard W. Hamming（理查德·汉明，《科学计算的艺术》）}

\begin{chapIntro}{本附录使用说明}
本附录收录第 1--6 章全部 48 道分层习题（第 1--3 章编号为 4.1.1--4.3.8，第 4--6 章编号为 4.1--6.8）的完整、自洽精解，所有手算数值、反例验算、不等式常数与代码实测基准均经独立 Python/NumPy 脚本交叉验证：
\begin{itemize}
  \item \textbf{A 层（基础概念与快速验算）}：给出明确正误/凸性/合法性结论、闭式手算步骤与显式反例；
  \item \textbf{B 层（定理推演与手工求解）}：保留精确分数与四位小数，完整呈现 KKT 消元、对偶 QP、回归树剪枝与 EM 迭代细节；
  \item \textbf{C 层（深度理论证明）}：给出零跳步的严格推导与构造性证明；
  \item \textbf{D 层（代码实现与实验基准）}：提供核心 NumPy 实现逻辑、排错要点以及固定随机种子下的实测数值自检表。
\end{itemize}
\end{chapIntro}

% =========================================================================
\section*{第 1 章习题解答（共 8 题）}
\addcontentsline{toc}{section}{第 1 章习题解答（共 8 题）}
% =========================================================================

\subsection*{习题 4.1.1\quad 凸性判别与反例}
\begin{enumerate}[label=(\alph*), leftmargin=1.8em]
  \item \textbf{$f_1(x) = e^{a^\top x + b}\;(x\in\R^d)$}：\textbf{凸函数}；当且仅当 $d = 1$ 且 $a \neq 0$ 时为\textbf{严格凸函数}，否则（$d \ge 2$ 或 $a = 0$）非严格凸。
  \emph{证明}：计算海森矩阵得
  \[
    \nabla f_1(x) = e^{a^\top x + b} a, \qquad
    \nabla^2 f_1(x) = e^{a^\top x + b} a a^\top.
  \]
  对任意 $v \in \R^d$，有 $v^\top \nabla^2 f_1(x) v = e^{a^\top x + b}(a^\top v)^2 \ge 0$，故 $\nabla^2 f_1(x) \succeq 0$，$f_1$ 为凸函数。当 $d \ge 2$ 时，可取非零向量 $v \perp a$，沿直线 $x + tv$ 函数保持常数 $e^{a^\top x + b}$，违反严格凸性（例如 $a=(1,0)^\top, b=0$，取 $x=(0,0)^\top, y=(0,2)^\top$，有 $f_1(\frac{x+y}{2}) = 1 = \frac{1}{2}(f_1(x)+f_1(y))$）；当 $d = 1$ 且 $a \neq 0$ 时，$f_1''(x) = a^2 e^{ax+b} > 0$，严格凸。

  \item \textbf{$f_2(x_1, x_2) = x_1 x_2\;(x \in \R_{++}^2)$}：\textbf{非凸函数}。
  \emph{证明与反例}：其海森矩阵为常阵
  \[
    \nabla^2 f_2(x) = \begin{pmatrix} 0 & 1 \\ 1 & 0 \end{pmatrix},
  \]
  特征值为 $\lambda_1 = +1$ 与 $\lambda_2 = -1$，矩阵不定。显式反例：取 $x = (1, 2)^\top,\; y = (2, 1)^\top \in \R_{++}^2$，中点为 $\frac{x+y}{2} = (1.5, 1.5)^\top$，则
  \[
    f_2\!\left(\frac{x+y}{2}\right) = 1.5 \times 1.5 = 2.25 > \frac{1}{2}\big(f_2(x) + f_2(y)\big) = \frac{2 + 2}{2} = 2.0,
  \]
  违反 Jensen 不等式，故非凸。

  \item \textbf{$f_3(x) = -\sum_{i=1}^d \ln x_i\;(x \in \R_{++}^d)$}：\textbf{严格凸函数}（当然也是凸函数）。
  \emph{证明}：海森矩阵为对角阵 $\nabla^2 f_3(x) = \diag(x_1^{-2}, \dots, x_d^{-2})$。对任意 $x \in \R_{++}^d$ 与任意非零向量 $v \in \R^d$，有 $v^\top \nabla^2 f_3(x) v = \sum_{i=1}^d (v_i / x_i)^2 > 0$，即 $\nabla^2 f_3(x) \succ 0$，因此 $f_3$ 在凸域 $\R_{++}^d$ 上严格凸。

  \item \textbf{$f_4(w) = \max_{1\le k\le K}(a_k^\top w + b_k)\;(w\in\R^d)$}：\textbf{凸函数，但非严格凸}。
  \emph{证明}：每个仿射函数 $g_k(w) = a_k^\top w + b_k$ 均为凸函数，由命题~\ref{prop:v4-preserve}（逐点取最大值保凸），$f_4$ 为凸函数。但 $f_4$ 是分段仿射函数，在任意一段激活区域内部均呈线性，无法严格凸；显式反例：令 $K=2, d=1, f_4(w)=\max(w, -w) = |w|$，取 $w_1 = 1, w_2 = 3$，则 $f_4(\frac{1+3}{2}) = 2 = \frac{1}{2}(|1| + |3|)$。
\end{enumerate}

\subsection*{习题 4.1.2\quad Log-Sum-Exp 函数的凸性与光滑常数}
\begin{proof}
记 $S(z) = \sum_{j=1}^K e^{z_j} > 0$ 及 Softmax 概率向量 $p_k(z) = e^{z_k} / S(z)$（满足 $p_k > 0$ 且 $\sum_{k=1}^K p_k = 1$）。
\begin{enumerate}[label=(\arabic*), leftmargin=1.8em]
  \item \textbf{一阶与二阶导数}：直接求偏导得 $\frac{\partial \mathrm{LSE}(z)}{\partial z_k} = \frac{e^{z_k}}{S(z)} = p_k$，故 $\nabla \mathrm{LSE}(z) = p$。再次求导：
  \[
    \frac{\partial^2 \mathrm{LSE}(z)}{\partial z_k \partial z_j}
    = \frac{\partial}{\partial z_j}\!\left(\frac{e^{z_k}}{S(z)}\right)
    = \frac{\delta_{kj} e^{z_k} S(z) - e^{z_k} e^{z_j}}{S(z)^2}
    = p_k \delta_{kj} - p_k p_j,
  \]
  写成矩阵形式即为 $\nabla^2 \mathrm{LSE}(z) = \diag(p) - p p^\top$。

  \item \textbf{凸性证明（$\nabla^2 \mathrm{LSE}(z) \succeq 0$）}：对任意 $v \in \R^K$，记 $V$ 为以概率分布 $p$ 取值 $v_k$ 的离散随机变量，则
  \[
    v^\top \big(\diag(p) - p p^\top\big) v
    = \sum_{k=1}^K p_k v_k^2 - \left(\sum_{k=1}^K p_k v_k\right)^2
    = \Var_{p}(V) \ge 0.
  \]
  （等价地，由 Cauchy-Schwarz 不等式，对向量 $(\sqrt{p_k} v_k)$ 与 $(\sqrt{p_k})$ 有 $(\sum_k p_k v_k)^2 \le (\sum_k p_k v_k^2)(\sum_k p_k) = \sum_k p_k v_k^2$。）故 $\nabla^2 \mathrm{LSE}(z) \succeq 0$，从而 $\mathrm{LSE}$ 是 $\R^K$ 上的凸函数。

  \item \textbf{$1$-光滑性证明（$\nabla^2 \mathrm{LSE}(z) \preceq I_K$）}：因 $p_k \in (0, 1)$ 且 $(\sum_k p_k v_k)^2 \ge 0$，对任意 $v \in \R^K$ 有
  \[
    v^\top \big(\diag(p) - p p^\top\big) v
    \le \sum_{k=1}^K p_k v_k^2
    \le \left(\max_{1\le k\le K} p_k\right)\sum_{k=1}^K v_k^2
    \le \norm{v}_2^2.
  \]
  故对一切 $z \in \R^K$ 有 $0 \preceq \nabla^2 \mathrm{LSE}(z) \preceq I_K$，即其梯度是 $1$-Lipschitz 连续的。（注：实际上还存在更紧的谱上界 $\nabla^2 \mathrm{LSE}(z) \preceq \frac{1}{2}I_K$，因 $v^\top(\diag(p)-pp^\top)v = \frac{1}{2}\sum_{j,k} p_j p_k(v_j-v_k)^2 \le \frac{1}{2}\norm{v}_2^2$；此处 $L=1$ 亦严格成立。）\qedhere
\end{enumerate}
\end{proof}

\subsection*{习题 4.1.3\quad 最小二乘与岭回归的精确条件数}
\begin{enumerate}[label=(\arabic*), leftmargin=1.8em]
  \item \textbf{海森矩阵与精确特征值}：目标函数 $F_\lambda(w) = \frac{1}{2N}\norm{Xw - y}_2^2 + \frac{\lambda}{2}\norm{w}_2^2$ 是二次函数，其海森矩阵为常阵
  \[
    H_\lambda = \nabla^2 F_\lambda(w) = \frac{1}{N}X^\top X + \lambda I_D.
  \]
  由于 $X \in \R^{N\times D}$ 的奇异值为 $\sigma_1 \ge \dots \ge \sigma_D \ge 0$，矩阵 $X^\top X$ 的特征值恰好为 $\sigma_1^2 \ge \dots \ge \sigma_D^2 \ge 0$。加上 $\lambda I_D$ 使所有特征值统一平移 $\lambda$，故 $H_\lambda$ 的最大与最小特征值分别为
  \[
    L = \lambda_{\max}(H_\lambda) = \frac{\sigma_1^2}{N} + \lambda, \qquad
    \mu = \lambda_{\min}(H_\lambda) = \frac{\sigma_D^2}{N} + \lambda.
  \]
  因海森阵恒定，这两个常数正是 $F_\lambda$ 的\textbf{紧致}光滑常数与强凸常数，对应精确条件数为
  \[
    \kappa(F_\lambda) = \frac{L}{\mu} = \frac{\sigma_1^2 + N\lambda}{\sigma_D^2 + N\lambda}.
  \]

  \item \textbf{秩亏时 $\lambda > 0$ 对收敛阶的质变}：当 $X$ 秩亏（如高维情形 $D > N$ 或存在完全共线性特征）时，$\sigma_D = 0$：
  \begin{itemize}
    \item 若 $\lambda = 0$（无正则最小二乘），则 $\mu = 0$，目标仅光滑凸而不强凸（条件数 $\kappa = \infty$），梯度下降只能保证次线性收敛速率 $O(1/T)$，且最小二乘解不唯一；
    \item 一旦加入任意 $\lambda > 0$（岭回归），立刻将最小特征值抬升至 $\mu = \lambda > 0$，使目标变为全局 $\lambda$-强凸，条件数变为有限值 $\kappa = 1 + \frac{\sigma_1^2}{N\lambda}$，梯度下降的收敛率瞬间从次线性 $O(1/T)$ 跃升为\textbf{指数级线性收敛} $O\!\big((1 - 1/\kappa)^T\big)$，同时保证了全局最优解的唯一性与数值稳定性。
  \end{itemize}
\end{enumerate}

\subsection*{习题 4.1.4\quad 一般步长 $\eta \in (0, 2/L)$ 的线性收敛率}
\begin{proof}
由 $L$-光滑下降引理~\eqref{eq:v4-step-drop}，对一步梯度下降 $x_{t+1} = x_t - \eta\nabla f(x_t)$ 有
\[
  f(x_{t+1}) \le f(x_t) - \eta\!\left(1 - \frac{L\eta}{2}\right)\norm{\nabla f(x_t)}_2^2.
\]
当 $\eta \in (0, 2/L)$ 时，系数 $\eta(1 - L\eta/2) > 0$。将 PL 不等式~\eqref{eq:v4-pl} 即 $\norm{\nabla f(x_t)}_2^2 \ge 2\mu(f(x_t) - f^*)$ 代入上式，并在两边同时减去 $f^*$：
\[
  f(x_{t+1}) - f^*
  \le \big(f(x_t) - f^*\big) - 2\mu\eta\!\left(1 - \frac{L\eta}{2}\right)\big(f(x_t) - f^*\big)
  = \rho(\eta)\big(f(x_t) - f^*\big),
\]
其中单步收缩系数为
\[
  \rho(\eta) = 1 - 2\mu\eta + \mu L \eta^2 = 1 - \frac{\mu}{L} + \mu L\!\left(\eta - \frac{1}{L}\right)^2.
\]
验证区间与最优步长：
\begin{itemize}
  \item 因 $0 < \mu \le L$ 且 $\eta \in (0, 2/L)$，有 $0 < 2\mu\eta(1 - L\eta/2) \le \frac{\mu}{L} \le 1$，故 $0 \le \rho(\eta) < 1$，迭代对区间内任意固定步长均保持线性收敛：$f(x_T) - f^* \le \rho(\eta)^T (f(x_0) - f^*)$。
  \item 由配方形式（或求导 $\rho'(\eta) = -2\mu + 2\mu L\eta = 0$ 且 $\rho''(\eta) = 2\mu L > 0$）立得：当且仅当 $\eta^* = \frac{1}{L}$ 时，$\rho(\eta)$ 取得全局最小值 $\rho(1/L) = 1 - \frac{\mu}{L} = 1 - \frac{1}{\kappa}$。\qedhere
\end{itemize}
\end{proof}

\subsection*{习题 4.1.5\quad 共强制性（Co-coercivity）与强凸参数距离收缩}
\begin{proof}
分两步证明：
\begin{enumerate}[label=(\arabic*), leftmargin=1.8em]
  \item \textbf{证明共强制性}：固定 $x \in \R^d$，构造辅助函数 $g_x(z) := f(z) - \nabla f(x)^\top z$。
  \begin{itemize}
    \item $g_x$ 与 $f$ 相差一个线性项，因此 $g_x$ 依然是凸函数且 $L$-光滑，其梯度为 $\nabla g_x(z) = \nabla f(z) - \nabla f(x)$。
    \item 在 $z = x$ 处 $\nabla g_x(x) = 0$，由凸函数驻点的全局最优性（定理~\ref{thm:v4-global-opt}），$z = x$ 是 $g_x$ 的全局极小点。
    \item 从任意点 $y \in \R^d$ 出发对 $g_x$ 做一步步长为 $1/L$ 的梯度下降，由下降引理~\eqref{eq:v4-1overL-drop} 得
    \[
      g_x(x) \le g_x\!\left(y - \frac{1}{L}\nabla g_x(y)\right) \le g_x(y) - \frac{1}{2L}\norm{\nabla g_x(y)}_2^2.
    \]
    展开 $g_x$ 的定义并移项，即得
    \[
      f(y) - f(x) - \nabla f(x)^\top(y - x) \ge \frac{1}{2L}\norm{\nabla f(y) - \nabla f(x)}_2^2.
    \]
    \item 交换 $x$ 与 $y$ 的角色，同理有
    \[
      f(x) - f(y) - \nabla f(y)^\top(x - y) \ge \frac{1}{2L}\norm{\nabla f(x) - \nabla f(y)}_2^2.
    \]
    两式相加，$f(x)$ 与 $f(y)$ 精确相消，立得共强制性不等式：
    \[
      \big(\nabla f(y) - \nabla f(x)\big)^\top(y - x) \ge \frac{1}{L}\norm{\nabla f(y) - \nabla f(x)}_2^2.
    \]
  \end{itemize}

  \item \textbf{证明 $\mu$-强凸下的参数一步收缩}：因 $x^*$ 是全局极小点，有 $\nabla f(x^*) = 0$。一方面，由上述共强制性（取 $y = x_t, x = x^*$）得
  \[
    \nabla f(x_t)^\top(x_t - x^*) \ge \frac{1}{L}\norm{\nabla f(x_t)}_2^2.
  \]
  另一方面，由 $\mu$-强凸性的一阶下界 $f(x^*) \ge f(x_t) + \nabla f(x_t)^\top(x^* - x_t) + \frac{\mu}{2}\norm{x_t - x^*}_2^2$ 与 $f(x_t) \ge f(x^*) + \frac{\mu}{2}\norm{x_t - x^*}_2^2$ 相加，得强单调性
  \[
    \nabla f(x_t)^\top(x_t - x^*) \ge \mu\norm{x_t - x^*}_2^2.
  \]
  现在展开一步梯度下降 $x_{t+1} = x_t - \eta\nabla f(x_t)$ 距最优解的平方距离，并将交叉项 $-2\eta\nabla f(x_t)^\top(x_t - x^*)$ 拆为两半分别代入上述两个下界：
  \begin{align*}
    \norm{x_{t+1} - x^*}_2^2
    &= \norm{x_t - x^*}_2^2 - 2\eta\nabla f(x_t)^\top(x_t - x^*) + \eta^2\norm{\nabla f(x_t)}_2^2 \\
    &\le \norm{x_t - x^*}_2^2 - \eta\mu\norm{x_t - x^*}_2^2 - \frac{\eta}{L}\norm{\nabla f(x_t)}_2^2 + \eta^2\norm{\nabla f(x_t)}_2^2 \\
    &= (1 - \mu\eta)\norm{x_t - x^*}_2^2 + \eta\!\left(\eta - \frac{1}{L}\right)\norm{\nabla f(x_t)}_2^2.
  \end{align*}
  由于步长满足 $0 < \eta \le \frac{1}{L}$，末项 $\eta(\eta - 1/L)\norm{\nabla f(x_t)}_2^2 \le 0$，舍去即证
  \[
    \norm{x_{t+1} - x^*}_2^2 \le (1 - \mu\eta)\norm{x_t - x^*}_2^2. \qedhere
  \]
\end{enumerate}
\end{proof}

\subsection*{习题 4.1.6\quad 满足 PL 条件的非凸函数}
考察 $f(x) = x^2 + 3\sin^2 x$（偶函数）。
\begin{enumerate}[label=(\alph*), leftmargin=1.8em]
  \item \textbf{非凸性验证}：利用 $\sin^2 x = \frac{1 - \cos(2x)}{2}$ 求导得
  \[
    f'(x) = 2x + 3\sin(2x), \qquad f''(x) = 2 + 6\cos(2x).
  \]
  在 $x = \frac{\pi}{2}$ 处，$f''(\pi/2) = 2 + 6\cos\pi = -4 < 0$，故 $f$ 在 $\R$ 上不是凸函数。

  \item \textbf{唯一全局极小点}：对任意 $x \in \R$，$x^2 \ge 0$ 且 $3\sin^2 x \ge 0$，故 $f(x) \ge 0 = f(0)$。等号成立当且仅当 $x^2 = 0$ 且 $\sin^2 x = 0$，即 $x = 0$。因此 $x^* = 0$ 是唯一全局极小点且 $f^* = 0$。

  \item \textbf{全局 PL 不等式证明}：因 $f(-x) = f(x)$ 且 $(f'(-x))^2 = (f'(x))^2$，只需对 $x > 0$ 证明 $\frac{(f'(x))^2}{2f(x)} \ge \mu$（$x = 0$ 时两边均为 $0$）。将 $x > 0$ 分为三段：
  \begin{itemize}
    \item \textbf{区间 1：$x \in (0, \pi/2]$}。此时 $2x \in (0, \pi]$，故 $\sin(2x) \ge 0$，有 $f'(x) = 2x + 3\sin(2x) \ge 2x$；同时 $|\sin x| \le x$，故 $f(x) = x^2 + 3\sin^2 x \le 4x^2$。于是
    \[
      \frac{(f'(x))^2}{2f(x)} \ge \frac{4x^2}{8x^2} = \frac{1}{2}.
    \]
    \item \textbf{区间 2：$x \in (\pi/2, 3]$}。先求 $f'(x) = 2x + 3\sin(2x)$ 在此区间内的严格正下界：$f''(x) = 2 + 6\cos(2x) = 0 \iff \cos(2x) = -\frac{1}{3}$。在 $2x \in (\pi, 6]$ 内仅有一个驻点 $2x_0 = \pi + \arccos(1/3) \approx 4.3726$（此处 $f''$ 由负变正，为 $f'$ 的极小值点），对应 $\sin(2x_0) = -\frac{2\sqrt{2}}{3}$，故在此区间恒有
    \[
      f'(x) \ge f'(x_0) = \pi + \arccos\!\left(\frac{1}{3}\right) - 2\sqrt{2} \approx 1.5441 > \frac{3}{2}.
    \]
    另一方面，在 $(\pi/2, 3]$ 上 $f(x) = x^2 + 3\sin^2 x \le 3^2 + 3 = 12$，故
    \[
      \frac{(f'(x))^2}{2f(x)} > \frac{(3/2)^2}{2 \times 12} = \frac{9}{96} = \frac{3}{32} > \frac{1}{32}.
    \]
    \item \textbf{区间 3：$x > 3$}。因 $|\sin(2x)| \le 1$，有 $f'(x) \ge 2x - 3 > x$（当 $x > 3$ 时 $2x - 3 > x$）；同时 $f(x) \le x^2 + 3 < x^2 + \frac{1}{3}x^2 = \frac{4}{3}x^2$，故
    \[
      \frac{(f'(x))^2}{2f(x)} > \frac{x^2}{2 \cdot \frac{4}{3}x^2} = \frac{3}{8}.
    \]
  \end{itemize}
  综合三段，对一切 $x \in \R$ 均满足 $\frac{1}{2}(f'(x))^2 \ge \frac{3}{32} f(x) \ge \frac{1}{32}(f(x) - f^*)$（Python 密集扫描测得全局最紧 PL 常数为 $\mu^* \approx 0.1755$，在 $x \approx 2.2017$ 处达到），因此取 $\mu = \frac{1}{32}$ 严格满足全局 PL 条件，梯度下降对其仍享指数级线性收敛。
\end{enumerate}

\subsection*{习题 4.1.7\quad 特征标准化对条件数 $\kappa$ 的震撼效应}
\begin{enumerate}[label=(\arabic*), leftmargin=1.8em]
  \item \textbf{核心代码改动}：在 \texttt{exp1\_opt\_acceleration} 生成原始矩阵 $X$ 后，加入一行逐列零均值单位方差标准化：\texttt{X = (X - X.mean(axis=0)) / X.std(axis=0)}，其余优化器代码不变。
  \item \textbf{实测对比基准}（\texttt{seed=2026}, $N=300, D=15, \lambda=0.05$）：
  \begin{center}\footnotesize\setlength{\tabcolsep}{4pt}
  \begin{tabular}{@{}lcccccc@{}}
    \toprule
    \textbf{预处理状态} & $L$ & $\mu$ & 条件数 $\kappa$ & $t=50$ GD & $t=150$ GD & $t=150$ NAG \\
    \midrule
    未标准化（原实验） & $51.5233$ & $0.0500$ & $1030.47$ & $1.01\times 10^{-1}$ & $2.222\times 10^{-2}$ & $3.361\times 10^{-6}$ \\
    列标准化后 & $0.3975$ & $0.0500$ & $\mathbf{7.95}$ & $5.551\times 10^{-17}$ & $0.000\times 10^{0}$ & $0.000\times 10^{0}$ \\
    \bottomrule
  \end{tabular}
  \end{center}
  \item \textbf{工程启示}：标准化将设计矩阵各列尺度拉齐后，最大光滑常数 $L$ 暴跌两个数量级（$51.52 \to 0.40$），条件数从极度病态的 $\kappa = 1030.5$ 骤降至几乎各向同性的 $\kappa = 7.95$。此时连最朴素的纯 GD 仅用 50 步就直达双精度机器零（$5.55\times 10^{-17}$）。这揭示了机器学习中“\textbf{数据预处理即隐式预条件（Preconditioning）}”的本质——改变地形几何往往比更换更复杂的优化器收益高出一个数量级。
\end{enumerate}

\subsection*{习题 4.1.8\quad Nesterov 自适应重启（Adaptive Restart）实验}
\begin{enumerate}[label=(\arabic*), leftmargin=1.8em]
  \item \textbf{核心实现逻辑}：在 NAG 迭代循环中计算 \texttt{w\_next = y\_nag - eta * g\_y} 后，检测梯度方向与位移方向内积：
  \begin{lstlisting}
if np.dot(g_y, w_next - w_nag) > 0.0:
    y_nag = w_next.copy()   # 冲过谷底：清零惯性重启
else:
    y_nag = w_next + beta_nag * (w_next - w_nag)
w_nag = w_next
  \end{lstlisting}
  \item \textbf{实测数值基准}（未标准化病态逻辑回归，$\kappa = 1030.47$）：
  \begin{itemize}
    \item \textbf{原版 NAG}：$t=50$ 次优度 $2.76\times 10^{-3}$，$t=100$ 次优度 $5.92\times 10^{-5}$，$t=150$ 次优度 $3.361\times 10^{-6}$，且出现两次明显回弹波纹；
    \item \textbf{梯度重启版 NAG}：在前 150 步中共触发 \textbf{4 次重启}（位于 $t = 34, 65, 78, 121$），所有回弹被彻底抹平；实测 $t=50$ 降至 $4.399\times 10^{-4}$，$t=100$ 降至 $4.531\times 10^{-7}$，$t=150$ 降至 $\mathbf{3.198\times 10^{-10}}$，相比原版 NAG 精度足足提升超过 \textbf{4 个数量级}（约 $10{,}500\times$）！
  \end{itemize}
\end{enumerate}

% =========================================================================
\section*{第 2 章习题解答（共 8 题）}
\addcontentsline{toc}{section}{第 2 章习题解答（共 8 题）}
% =========================================================================

\subsection*{习题 4.2.1\quad 最小欧氏范数线性方程组的对偶问题}
\begin{enumerate}[label=(\arabic*), leftmargin=1.8em]
  \item \textbf{拉格朗日函数}：引入等式对偶乘子 $\nu \in \R^m$，写出
  \[
    \mathcal{L}(x, \nu) = \frac{1}{2}x^\top x + \nu^\top (Ax - b).
  \]
  \item \textbf{对偶函数 $g(\nu)$}：对 $x$ 求梯度并置零得 $\nabla_x \mathcal{L}(x, \nu) = x + A^\top \nu = 0 \implies x(\nu) = -A^\top \nu$。回代 $\mathcal{L}$ 得
  \[
    g(\nu) = \inf_{x\in\R^d} \mathcal{L}(x, \nu)
    = \frac{1}{2}\nu^\top A A^\top \nu - \nu^\top A A^\top \nu - \nu^\top b
    = -\frac{1}{2}\nu^\top (A A^\top)\nu - b^\top \nu.
  \]
  \item \textbf{对偶最优解与原始最优解}：因 $A \in \R^{m\times d}$ 行满秩，对称阵 $A A^\top \succ 0$ 可逆，故 $g(\nu)$ 为严格凹二次函数。令 $\nabla g(\nu) = -(AA^\top)\nu - b = 0$，得唯一对偶最优解
  \[
    \nu^* = -(A A^\top)^{-1} b, \qquad
    x^* = -A^\top \nu^* = A^\top (A A^\top)^{-1} b.
  \]
  直接验证可行性：$A x^* = A A^\top (A A^\top)^{-1} b = b$。
  \item \textbf{强对偶验证（$d^* = p^*$）}：计算原始最优值与对偶最优值：
  \begin{align*}
    p^* &= \frac{1}{2}(x^*)^\top x^* = \frac{1}{2} b^\top (AA^\top)^{-1} A A^\top (AA^\top)^{-1} b = \frac{1}{2} b^\top (A A^\top)^{-1} b, \\
    d^* &= g(\nu^*) = -\frac{1}{2} b^\top (AA^\top)^{-1} b + b^\top (AA^\top)^{-1} b = \frac{1}{2} b^\top (A A^\top)^{-1} b = p^*.
  \end{align*}
  对偶间隙精确为零（与无不等式约束的线性等式凸问题满足 Slater 条件一致）。
\end{enumerate}

\subsection*{习题 4.2.2\quad 软阈值算子与 $\ell_1$ 次微分手算}
\begin{enumerate}[label=(\alph*), leftmargin=1.8em]
  \item \textbf{软阈值算子计算}：逐坐标应用 $S_\tau(z_j) = \sign(z_j)\max(|z_j| - \tau, 0)$（$\tau = 0.8$）：
  \begin{itemize}
    \item $z_1 = 2.4$：$|2.4| > 0.8 \implies +(2.4 - 0.8) = 1.6$；
    \item $z_2 = -0.6$：$|-0.6| \le 0.8 \implies 0$；
    \item $z_3 = 0.8$：$|0.8| \le 0.8 \implies 0$；
    \item $z_4 = -3.1$：$|-3.1| > 0.8 \implies -(3.1 - 0.8) = -2.3$。
  \end{itemize}
  因此 $\prox_{0.8\norm{\cdot}_1}(z) = \mathbf{(1.6,\; 0,\; 0,\; -2.3)^\top}$。

  \item \textbf{$\ell_1$ 范数的三维次微分}：由可分函数的逐坐标次微分规则，对 $x = (2, 0, -1)^\top$：
  $x_1 = 2 > 0 \implies \partial |x_1| = \{+1\}$；$x_2 = 0 \implies \partial |x_2| = [-1, 1]$；$x_3 = -1 < 0 \implies \partial |x_3| = \{-1\}$。因此完整次微分集合为三维空间中的一条闭线段：
  \[
    \partial g(2, 0, -1) = \big\{(1,\, s,\, -1)^\top : s \in [-1, 1]\big\} = \{1\} \times [-1, 1] \times \{-1\}.
  \]

  \item \textbf{合页损失在折点 $z = 1$ 处的次微分}：$h(z) = \max(0, 1 - z)$ 在 $z=1$ 左侧斜率为 $-1$、右侧斜率为 $0$。检验次梯度定义 $h(y) \ge h(1) + g(y - 1) = g(y - 1)$：
  \begin{itemize}
    \item 当 $y > 1$ 时要求 $0 \ge g(y - 1) \iff g \le 0$；
    \item 当 $y < 1$ 时要求 $1 - y \ge g(y - 1) = -g(1 - y) \iff g \ge -1$。
  \end{itemize}
  故完整次微分集合为 $\partial h(1) = \mathbf{[-1, 0]}$。
\end{enumerate}

\subsection*{习题 4.2.3\quad 最大熵分布与 Softmax 的 KKT 起源}
\begin{proof}
目标函数 $f_0(p) = \sum_{k=1}^K (p_k \ln p_k - s_k p_k)$ 的海森矩阵为 $\diag(1/p_1, \dots, 1/p_K) \succ 0$，在单纯形上严格凸，且均匀分布 $p = (\frac{1}{K}, \dots, \frac{1}{K})^\top$ 严格满足 $p_k > 0$（Slater 条件成立），故 KKT 条件是全局唯一最优解的充要条件。

构造拉格朗日函数（不等式写为 $-p_k \le 0$，乘子 $\lambda_k \ge 0$；等式乘子 $\nu \in \R$）：
\[
  \mathcal{L}(p, \lambda, \nu) = \sum_{k=1}^K p_k \ln p_k - \sum_{k=1}^K s_k p_k - \sum_{k=1}^K \lambda_k p_k + \nu\!\left(\sum_{k=1}^K p_k - 1\right).
\]
\begin{enumerate}[label=(\arabic*), leftmargin=1.8em]
  \item \textbf{为何所有 $p_k^* > 0$ 且 $\lambda_k^* = 0$}：若存在某个坐标 $p_k^* = 0$，则必有另一坐标 $p_j^* > 0$。沿可行方向把极小质量 $\varepsilon > 0$ 从 $p_j^*$ 移至 $p_k^*$，目标函数一阶变化率为 $\lim_{\varepsilon\to 0^+}\big(\ln\varepsilon - \ln p_j^* - s_k + s_j\big) = -\infty < 0$，说明边界点绝不可能是极小值点。因此必有 $p_k^* > 0\;(\forall k)$。再由互补松弛 $\lambda_k^* p_k^* = 0$，立得所有不等式乘子自动为零：$\lambda_k^* = 0\;(\forall k)$。
  \item \textbf{平稳性与归一化}：令 $\frac{\partial \mathcal{L}}{\partial p_k} = \ln p_k^* + 1 - s_k + \nu^* = 0$，解得
  \[
    p_k^* = e^{-(1+\nu^*)} e^{s_k}.
  \]
  代入原始等式约束 $\sum_{j=1}^K p_j^* = 1$ 得 $e^{-(1+\nu^*)} = \left(\sum_{j=1}^K e^{s_j}\right)^{-1}$，回代即得唯一最优解为 Softmax 分布：
  \[
    p_k^* = \frac{e^{s_k}}{\sum_{j=1}^K e^{s_j}}, \qquad k = 1, \dots, K. \qedhere
  \]
\end{enumerate}
\end{proof}

\subsection*{习题 4.2.4\quad 注水算法（Water-Filling）的 KKT 求解}
\begin{proof}
目标 $f_0(x) = -\sum_{i=1}^d \ln(\alpha_i + x_i)$ 在可行域 $\{x \ge 0 : \sum_i x_i = 1\}$ 上严格凸，且内点 $x = (\frac{1}{d}, \dots, \frac{1}{d})^\top$ 满足 Slater 条件，故 KKT 条件充要。构造拉格朗日函数：
\[
  \mathcal{L}(x, \lambda, \nu) = -\sum_{i=1}^d \ln(\alpha_i + x_i) - \sum_{i=1}^d \lambda_i x_i + \nu\!\left(\sum_{i=1}^d x_i - 1\right).
\]
KKT 全部条件为（对 $i = 1, \dots, d$）：
\[
  -\frac{1}{\alpha_i + x_i^*} - \lambda_i^* + \nu^* = 0, \qquad
  \lambda_i^* \ge 0, \qquad x_i^* \ge 0, \qquad
  \lambda_i^* x_i^* = 0, \qquad \sum_{i=1}^d x_i^* = 1.
\]
先由平稳性得 $\nu^* = \lambda_i^* + \frac{1}{\alpha_i + x_i^*} > 0$（故水位线 $w := (\nu^*)^{-1} > 0$ 有定义），且 $\lambda_i^* = \nu^* - \frac{1}{\alpha_i + x_i^*} \ge 0$。分类讨论每个分量 $i$：
\begin{itemize}
  \item 若 $\alpha_i < (\nu^*)^{-1}$：假设 $x_i^* = 0$，则 $\lambda_i^* = \nu^* - \frac{1}{\alpha_i} < 0$，与对偶可行性 $\lambda_i^* \ge 0$ 矛盾！故必有 $x_i^* > 0$，由互补松弛得 $\lambda_i^* = 0$，代入平稳性解得 $x_i^* = (\nu^*)^{-1} - \alpha_i > 0$。
  \item 若 $\alpha_i \ge (\nu^*)^{-1}$：假设 $x_i^* > 0$，则由互补松弛 $\lambda_i^* = 0$，平稳性给出 $\alpha_i + x_i^* = (\nu^*)^{-1} \le \alpha_i \implies x_i^* \le 0$，矛盾！故必有 $x_i^* = 0$（此时 $\lambda_i^* = \nu^* - \alpha_i^{-1} \ge 0$ 自动满足）。
\end{itemize}
合并两种情形即得 $x_i^* = \max\!\big(0,\; (\nu^*)^{-1} - \alpha_i\big)$。最后令 $\phi(w) = \sum_{i=1}^d \max(0, w - \alpha_i)$：该函数在 $w \ge \min_i \alpha_i$ 上连续、严格单调递增，且 $\phi(\min_i \alpha_i) = 0$、$\lim_{w\to+\infty}\phi(w) = +\infty$，由介值定理方程 $\phi(w) = 1$ 存在唯一解 $w^* = (\nu^*)^{-1}$。
\end{proof}

\subsection*{习题 4.2.5\quad Lasso 全零解的精确临界阈值 $\lambda_{\max}$}
\begin{proof}
Lasso 目标函数 $F_\lambda(w) = \underbrace{\frac{1}{2N}\norm{Xw - y}_2^2}_{f(w)} + \lambda\norm{w}_1$ 为无约束凸函数，因此 $w^* = 0$ 为全局极小点的充要条件是次梯度最优性条件 $0 \in \partial F_\lambda(0)$。
\begin{enumerate}[label=(\arabic*), leftmargin=1.8em]
  \item 光滑项在 $w = 0$ 处的梯度为 $\nabla f(0) = \frac{1}{N}X^\top(X\cdot 0 - y) = -\frac{1}{N}X^\top y$。
  \item 非光滑项 $\lambda\norm{w}_1$ 在 $w = 0$ 处的次微分为超立方体 $\lambda\,\partial\norm{0}_1 = [-\lambda, \lambda]^D$。
  \item 由次微分加法规则，$\partial F_\lambda(0) = \left\{-\frac{1}{N}X^\top y + s : s \in [-\lambda, \lambda]^D\right\}$，故
  \[
    0 \in \partial F_\lambda(0)
    \iff \frac{1}{N}X^\top y \in [-\lambda, \lambda]^D
    \iff \max_{1\le j\le D}\abs{\frac{1}{N}x_j^\top y} \le \lambda
    \iff \lambda \ge \lambda_{\max} \coloneqq \frac{1}{N}\norm{X^\top y}_\infty.
  \]
\end{enumerate}
\textbf{结合第 2.4 节实测核验}：对 \texttt{exp2\_lasso\_ista}（$N=120, D=20$, \texttt{seed=42}）的标准化特征矩阵 $X$ 与目标 $y$ 直接计算得精确临界值
\[
  \lambda_{\max} = \frac{1}{120}\norm{X^\top y}_\infty = \mathbf{3.5615}.
\]
表~\ref{tab:v4-ch2-lasso} 中最后一个正则点 $\lambda = 3.80 > 3.5615 = \lambda_{\max}$，满足全零解充要条件，因此 ISTA 迭代将全部 20 个系数精准归零（\texttt{nnz = 0}）。
\end{proof}

\subsection*{习题 4.2.6\quad 弹性网（Elastic Net）的闭式近端算子}
\begin{proof}
根据近端算子定义，对 $g(w) = \lambda_1\norm{w}_1 + \frac{\lambda_2}{2}\norm{w}_2^2$（$\lambda_1, \lambda_2 > 0$）：
\[
  \prox_{\eta g}(z) = \argmin_{x \in \R^D}\left\{\frac{1}{2}\norm{x - z}_2^2 + \frac{\eta\lambda_2}{2}\norm{x}_2^2 + \eta\lambda_1\norm{x}_1\right\}.
\]
目标函数关于各个坐标 $x_j$ 完全可分。对第 $j$ 个坐标配方整理（忽略与 $x_j$ 无关的常数项）：
\begin{align*}
  \frac{1}{2}(x_j - z_j)^2 + \frac{\eta\lambda_2}{2}x_j^2 + \eta\lambda_1 |x_j|
  &= \frac{1 + \eta\lambda_2}{2}x_j^2 - z_j x_j + \eta\lambda_1 |x_j| + \frac{1}{2}z_j^2 \\
  &= (1 + \eta\lambda_2)\left[\frac{1}{2}\!\left(x_j - \frac{z_j}{1 + \eta\lambda_2}\right)^2 + \frac{\eta\lambda_1}{1 + \eta\lambda_2}|x_j|\right] + \text{const}.
\end{align*}
因 $1 + \eta\lambda_2 > 0$，极小化方括号内函数等价于对点 $\tilde{z}_j = \frac{z_j}{1+\eta\lambda_2}$ 施加阈值为 $\tilde{\tau} = \frac{\eta\lambda_1}{1+\eta\lambda_2}$ 的软阈值算子（定理~\ref{thm:v4-soft-thresh}）。再提取公共正因子 $\frac{1}{1+\eta\lambda_2}$ 即得闭式解析解：
\[
  \big[\prox_{\eta g}(z)\big]_j
  = S_{\frac{\eta\lambda_1}{1+\eta\lambda_2}}\!\left(\frac{z_j}{1+\eta\lambda_2}\right)
  = \frac{1}{1 + \eta\lambda_2}\, S_{\eta\lambda_1}(z_j)
  = \frac{\sign(z_j)\max(|z_j| - \eta\lambda_1,\, 0)}{1 + \eta\lambda_2}.
\]
几何意义极为清晰：先用 $\ell_1$ 软阈值把弱特征精确归零，再统一乘以 $\ell_2$ 阻尼因子 $\frac{1}{1+\eta\lambda_2}$ 收缩非零幅值。
\end{proof}

\subsection*{习题 4.2.7\quad ISTA vs FISTA 与次梯度下降三法实测}
\begin{enumerate}[label=(\arabic*), leftmargin=1.8em]
  \item \textbf{核心算法对比实现}（基于 \texttt{exp2\_lasso\_ista} 数据，$\lambda = 0.20, L = 1.7474, \eta = 1/L = 0.5723$）：
  \begin{itemize}
    \item \textbf{次梯度下降}：每步更新 $w^{(t+1)} = w^{(t)} - \frac{0.5}{\sqrt{t+1}}\big(\frac{1}{N}X^\top(Xw^{(t)}-y) + \lambda\sign(w^{(t)})\big)$；
    \item \textbf{ISTA}：每步更新 $w^{(t+1)} = S_{\eta\lambda}\big(w^{(t)} - \eta\nabla f(w^{(t)})\big)$；
    \item \textbf{FISTA}：维护辅助序列 $y^{(t)}$ 与动量标量 $t_0 = 1$：
    \begin{lstlisting}
w_next = soft_thresh(y_f - eta * grad(y_f), eta * lam)
t_next = 0.5 * (1.0 + np.sqrt(1.0 + 4.0 * tk * tk))
y_f = w_next + ((tk - 1.0) / t_next) * (w_next - w_f)
    \end{lstlisting}
  \end{itemize}
  \item \textbf{实测自检基准}（以 5000 步 FISTA 求得精确基准 $F^* = 1.848404$，真稀疏度 $\texttt{nnz}^* = 4$）：
  \begin{center}\small
  \begin{tabular}{@{}lcccc@{}}
    \toprule
    \textbf{算法} & $t=10$ 次优度 / \texttt{nnz} & $t=20$ 次优度 / \texttt{nnz} & $t=40$ 次优度 / \texttt{nnz} & $t=200$ 次优度 / \texttt{nnz} \\
    \midrule
    次梯度下降 & $8.76\times 10^{-2}$ / \textbf{20} & $3.28\times 10^{-2}$ / \textbf{20} & $2.34\times 10^{-2}$ / \textbf{20} & $1.01\times 10^{-2}$ / \textbf{20} \\
    ISTA & $2.39\times 10^{-5}$ / \textbf{4} & $1.11\times 10^{-10}$ / \textbf{4} & $<10^{-15}$ / \textbf{4} & $<10^{-15}$ / \textbf{4} \\
    FISTA & $3.30\times 10^{-5}$ / \textbf{4} & $1.48\times 10^{-8}$ / \textbf{4} & $1.02\times 10^{-14}$ / \textbf{4} & $0.00$ / \textbf{4} \\
    \bottomrule
  \end{tabular}
  \end{center}
  \item \textbf{关键洞察}：(i) 次梯度下降在 200 步后误差仍高达 $1.01\times 10^{-2}$，且全部 20 维系数在零点附近来回锯齿震荡（\texttt{nnz = 20}），\textbf{永远无法产出精确零解}；(ii) ISTA 与 FISTA 均凭借软阈值算子在第 10 步内就锁死真实稀疏支撑集（\texttt{nnz = 4}），40 步内达到双精度机器精度；(iii) 注意在本例标准化弱相关矩阵上（局部限制条件数 $\kappa_S$ 很小），ISTA 本身已呈极快线性收敛，FISTA 因无重启的惯性过冲在前 40 步略慢于 ISTA；一旦换成高相关病态特征矩阵（$\kappa_S \gg 100$），FISTA 的 $O(1/t^2)$ 加速优势将迅速甩开 ISTA 两个数量级。
\end{enumerate}

\subsection*{习题 4.2.8\quad 单纯形投影算子（Simplex Projection）推导与实现}
\begin{enumerate}[label=(\arabic*), leftmargin=1.8em]
  \item \textbf{$O(d\log d)$ 阈值公式推导}：求解 $\min_p \frac{1}{2}\norm{p - z}_2^2$ s.t. $p_i \ge 0,\; \sum_{i=1}^d p_i = 1$。拉格朗日函数为 $\mathcal{L}(p, \lambda, \tau) = \frac{1}{2}\norm{p - z}_2^2 - \lambda^\top p + \tau(\mathbf{1}^\top p - 1)$。由平稳性 $p_i^* - z_i - \lambda_i^* + \tau = 0$、$\lambda_i^* \ge 0$ 与互补松弛 $\lambda_i^* p_i^* = 0$，与习题 4.2.4 完全同理可得：
  \[
    p_i^* = \max(z_i - \tau,\, 0), \qquad \sum_{i=1}^d \max(z_i - \tau,\, 0) = 1.
  \]
  将 $z$ 降序排列为 $u_1 \ge u_2 \ge \dots \ge u_d$。因 $p_i^* > 0 \iff z_i > \tau$，正分量必对应排序后的前 $\rho$ 个最大元素 $u_1, \dots, u_\rho$，且满足归一化方程 $\sum_{j=1}^\rho (u_j - \tau) = 1 \implies \tau = \frac{1}{\rho}\big(\sum_{j=1}^\rho u_j - 1\big)$。同时必须满足截断边界条件 $u_\rho > \tau \ge u_{\rho+1}$。可以验证序列 $\Delta_k := u_k - \frac{1}{k}(\sum_{j=1}^k u_j - 1)$ 满足单调变号性质，因而满足 $u_k > \frac{1}{k}(\sum_{j=1}^k u_j - 1)$ 的\textbf{最大下标} $\rho$ 即为唯一精确支撑大小（排序耗时 $O(d\log d)$，前缀和扫描耗时 $O(d)$）。

  \item \textbf{纯 NumPy 实现与自检基准}：
  \begin{lstlisting}
def proj_simplex(z):
    u = np.sort(z)[::-1]
    cssv = np.cumsum(u) - 1.0
    ind = np.arange(1, len(z) + 1)
    rho = ind[u - cssv / ind > 0][-1]
    tau = cssv[rho - 1] / rho
    return np.maximum(z - tau, 0.0)
  \end{lstlisting}
  \textbf{手算与随机自检}：取 $z = (1.5, 0.2, -0.5, 0.8)^\top$，排序得 $u = (1.5, 0.8, 0.2, -0.5)$。检验：$k=1$ 时 $1.5 > 0.5$；$k=2$ 时 $0.8 > \frac{2.3-1}{2} = 0.65$；$k=3$ 时 $0.2 \not> \frac{2.5-1}{3} = 0.5$。故 $\rho = 2, \tau = 0.65$，投影结果为 $p^* = \mathbf{(0.85,\; 0,\; 0,\; 0.15)^\top}$。在 100 组随机 15 维高斯向量上检验 KKT 四项条件（$p_i \ge 0$、$\sum p_i = 1$、$\lambda_i = p_i - z_i + \tau \ge 0$、$\max_i |\lambda_i p_i| < 10^{-10}$），全部以机器精度通过。
\end{enumerate}

% =========================================================================
\section*{第 3 章习题解答（共 8 题）}
\addcontentsline{toc}{section}{第 3 章习题解答（共 8 题）}
% =========================================================================

\subsection*{习题 4.3.1\quad 常值估计器的偏差与方差手算}
训练输入 $x_1, \dots, x_N$ 固定，$y_i = f(x_i) + \varepsilon_i$，其中 $\varepsilon_i$ 独立且均值 $0$、方差 $\sigma^2$。模型预测为 $\hat{f}_{\mathcal{D}}(x) = \bar{y} = \frac{1}{N}\sum_{i=1}^N y_i$。
\begin{enumerate}[label=(\arabic*), leftmargin=1.8em]
  \item \textbf{期望预测与平方偏差}：
  \[
    \bar{f}(x) = \E_{\mathcal{D}}[\bar{y}] = \frac{1}{N}\sum_{i=1}^N \E[f(x_i) + \varepsilon_i] = \frac{1}{N}\sum_{i=1}^N f(x_i) \eqqcolon \bar{f}_{\mathrm{tr}}.
  \]
  故在任意测试点 $x$ 处的平方偏差为 $\Bias^2(\hat{f}(x)) = \big(\bar{f}_{\mathrm{tr}} - f(x)\big)^2$。
  \item \textbf{方差}：由训练噪声 $\varepsilon_1, \dots, \varepsilon_N$ 相互独立，
  \[
    \Var(\hat{f}(x)) = \Var\!\left(\frac{1}{N}\sum_{i=1}^N \varepsilon_i\right) = \frac{1}{N^2}\sum_{i=1}^N \sigma^2 = \frac{\sigma^2}{N}.
  \]
  \item \textbf{期望测试 MSE 与 $N \to \infty$ 极限}：由偏差-方差-噪声分解定理，
  \[
    \E\big[(y - \hat{f}_{\mathcal{D}}(x))^2\big]
    = \underbrace{\big(\bar{f}_{\mathrm{tr}} - f(x)\big)^2}_{\Bias^2}
    + \underbrace{\frac{\sigma^2}{N}}_{\Var}
    + \underbrace{\sigma^2}_{\text{不可约噪声}}.
  \]
  当 $N \to \infty$ 时，方差项 $\sigma^2 / N \to 0$，不可约噪声恒为 $\sigma^2$，而平方偏差收敛至 $\big(\E_X[f(X)] - f(x)\big)^2$（除非真实函数 $f$ 本身是常数，否则该项永不消失）。这清晰揭示了高偏差欠拟合模型的致命局限：单纯增加样本量只能消灭方差，无法降低系统性偏差。
\end{enumerate}

\subsection*{习题 4.3.2\quad PAC 样本复杂度计算}
代入有限假设空间样本复杂度公式~\eqref{eq:v4-sample-comp}：$N \ge \frac{1}{2\varepsilon^2}\ln\frac{2|\mathcal{H}|}{\delta}$，其中 $|\mathcal{H}| = 2^{20}$、$\varepsilon = 0.02$，故 $\frac{1}{2\varepsilon^2} = \frac{1}{2 \times 0.0004} = 1250$。
\begin{enumerate}[label=(\arabic*), leftmargin=1.8em]
  \item \textbf{置信度 $99\%$（$\delta = 0.01$）}：
  \[
    \ln\frac{2 \times 2^{20}}{0.01} = 21\ln 2 + \ln 100 \approx 21 \times 0.693147 + 4.605170 = 19.16126,
  \]
  因此 $N_{0.01} \ge 1250 \times 19.16126 = 23{,}951.58$，至少需要 $\mathbf{23{,}952}$ 个 i.i.d. 样本。
  \item \textbf{置信度 $90\%$（$\delta = 0.10$）}：
  \[
    \ln\frac{2 \times 2^{20}}{0.10} = 21\ln 2 + \ln 10 \approx 14.55609 + 2.302585 = 16.85868,
  \]
  因此 $N_{0.10} \ge 1250 \times 16.85868 = 21{,}073.34$，至少需要 $\mathbf{21{,}074}$ 个样本。
  \item \textbf{样本需求减少比例}：
  \[
    \frac{23951.58 - 21073.34}{23951.58} = \frac{\ln 10}{21\ln 2 + \ln 100} \approx \mathbf{12.02\%}
  \]
  （按向上取整后的整数算为 $(23952 - 21074)/23952 \approx 12.02\%$）。失败概率放大整整 10 倍，样本量却只省了 $12\%$——因为样本复杂度对 $\delta$ 仅为极温和的对数依赖 $\ln(1/\delta)$，真正昂贵的是精度项 $1/\varepsilon^2$。
\end{enumerate}

\subsection*{习题 4.3.3\quad 岭回归的存在性优势：一定存在 $\lambda > 0$ 严格胜过 OLS}
\begin{proof}
由式~\eqref{eq:v4-ridge-var} 与式~\eqref{eq:v4-ridge-bias}（设设计矩阵 $X$ 列满秩，即 $\sigma_1 \ge \dots \ge \sigma_D > 0$），岭回归总参数均方误差为
\[
  \mathrm{MSE}(\lambda)
  = \underbrace{\sigma^2 \sum_{j=1}^D \frac{\sigma_j^2}{(\sigma_j^2 + \lambda)^2}}_{V(\lambda)}
  + \underbrace{\sum_{j=1}^D \frac{\lambda^2}{(\sigma_j^2 + \lambda)^2}\,(v_j^\top w_{\mathrm{true}})^2}_{B(\lambda)}.
\]
对 $\lambda \ge 0$ 逐项求一阶导数：
\begin{align*}
  V'(\lambda) &= -2\sigma^2 \sum_{j=1}^D \frac{\sigma_j^2}{(\sigma_j^2 + \lambda)^3}, \\
  B'(\lambda) &= \sum_{j=1}^D \frac{2\lambda(\sigma_j^2 + \lambda)^2 - 2\lambda^2(\sigma_j^2 + \lambda)}{(\sigma_j^2 + \lambda)^4}(v_j^\top w_{\mathrm{true}})^2
               = 2\lambda \sum_{j=1}^D \frac{\sigma_j^2}{(\sigma_j^2 + \lambda)^3}(v_j^\top w_{\mathrm{true}})^2.
\end{align*}
在 $\lambda = 0$ 处求值：
\[
  V'(0) = -2\sigma^2 \sum_{j=1}^D \frac{1}{\sigma_j^4}, \qquad
  B'(0) = 0,
\]
从而总误差在原点处的一阶导数为
\[
  \left.\frac{d}{d\lambda}\mathrm{MSE}(\lambda)\right|_{\lambda=0}
  = -2\sigma^2 \sum_{j=1}^D \frac{1}{\sigma_j^4} < 0 \qquad (\text{只要 }\sigma^2 > 0).
\]
由导数保号性，存在 $\delta > 0$，使得对任意 $\lambda \in (0, \delta)$ 均有 $\mathrm{MSE}(\lambda) < \mathrm{MSE}(0) = \mathrm{MSE}_{\mathrm{OLS}}$。直觉上，在 $\lambda = 0$ 附近引入微小正则化只带来二阶小量 $O(\lambda^2)$ 的平方偏差，却立竿见影地削减了一阶量 $O(\lambda)$ 的方差，这是一场数学上稳赚不赔的交易。
\end{proof}

\subsection*{习题 4.3.4\quad 三类几何假设空间的精确 VC 维}
\begin{enumerate}[label=(\alph*), leftmargin=1.8em]
  \item \textbf{一维单侧阈值 $\mathcal{H}_1 = \{\mathbb{I}[x \ge \theta] : \theta \in \R\}$}：$\mathbf{d_{\mathrm{VC}}(\mathcal{H}_1) = 1}$。
  \begin{itemize}
    \item \emph{能打散 1 点}：取 $x_1 = 0$，取 $\theta = -1$ 得 $+1$，取 $\theta = 1$ 得 $0$，实现全部 $2^1 = 2$ 种标记。
    \item \emph{不能打散任意 2 点}：设 $x_1 < x_2$，任何阈值分类器均满足单调性 $h(x_1) \le h(x_2)$，故标记 $(y_1, y_2) = (1, 0)$ 永远无法实现。
  \end{itemize}

  \item \textbf{一维闭区间 $\mathcal{H}_2 = \{\mathbb{I}[a \le x \le b] : a \le b\}$}：$\mathbf{d_{\mathrm{VC}}(\mathcal{H}_2) = 2}$。
  \begin{itemize}
    \item \emph{能打散 2 点}：取 $x_1 = 0,\; x_2 = 2$。四种标记 $(0,0)$、$(1,0)$、$(0,1)$、$(1,1)$ 分别由闭区间 $[-2,-1]$、$[-0.5, 0.5]$、$[1.5, 2.5]$、$[0, 2]$ 精确实现。
    \item \emph{不能打散任意 3 点}：设 $x_1 < x_2 < x_3$，任何闭区间 $[a, b]$ 若同时包含两端点 $x_1, x_3$，由区间凸性必包含中间点 $x_2$，因此标记 $(1, 0, 1)$ 永远无法实现。
  \end{itemize}

  \item \textbf{二维轴对齐矩形 $\mathcal{H}_3$}：$\mathbf{d_{\mathrm{VC}}(\mathcal{H}_3) = 4}$。
  \begin{itemize}
    \item \emph{能打散 4 点}：取菱形四顶点 $A(-1,0),\, B(1,0),\, C(0,1),\, D(0,-1)$。任取其一子集 $S \subseteq \{A,B,C,D\}$：若 $S = \varnothing$ 取远离原点的空区矩形；否则取 $S$ 的最小轴对齐包围盒 $[\min_{p\in S} x_1,\, \max_{p\in S} x_1] \times [\min_{p\in S} x_2,\, \max_{p\in S} x_2]$。容易逐一检验（例如只选 $A, B$ 得 $[-1,1]\times[0,0]$ 不含 $C, D$；选 $A, B, C$ 得 $[-1,1]\times[0,1]$ 不含 $D$），全部 $2^4 = 16$ 种二分标记均可实现。
    \item \emph{不能打散任意 5 点}：任给平面上 5 个点，分别找出 $x_1$ 最小、$x_1$ 最大、$x_2$ 最小、$x_2$ 最大的至多 4 个极值点（若有并列任取其一），将其标为 $+1$，至少余下第 5 个点 $q$，将其标为 $0$。任何包含那 4 个极值点的轴对齐矩形必满足 $a_1 \le x_{1,\min} \le q_1 \le x_{1,\max} \le b_1$ 且 $a_2 \le x_{2,\min} \le q_2 \le x_{2,\max} \le b_2$，因而必然把 $q$ 也包入正类，该标记不可实现。
  \end{itemize}
\end{enumerate}

\subsection*{习题 4.3.5\quad 留一法（LOOCV）线性回归的 $O(1)$ 快捷公式}
\begin{proof}
设设计矩阵 $X \in \R^{N\times D}$ 列满秩，第 $i$ 行为 $x_i^\top$，记 $A = X^\top X$，则帽子矩阵对角元为 $h_{ii} = x_i^\top A^{-1} x_i \in (0, 1)$，全量 OLS 估计为 $\hat{w} = A^{-1} X^\top y$，全量拟合值为 $\hat{y}_i = x_i^\top \hat{w}$。

删去第 $i$ 个样本后，新正规方程的系数矩阵为 $A_{(-i)} = X_{(-i)}^\top X_{(-i)} = A - x_i x_i^\top$，右侧向量为 $X_{(-i)}^\top y_{(-i)} = X^\top y - x_i y_i$。应用 Sherman-Morrison 秩-1 逆矩阵公式：
\[
  (A - x_i x_i^\top)^{-1} = A^{-1} + \frac{A^{-1} x_i x_i^\top A^{-1}}{1 - x_i^\top A^{-1} x_i}
  = A^{-1} + \frac{A^{-1} x_i x_i^\top A^{-1}}{1 - h_{ii}}.
\]
将其代入留一参数 $\hat{w}^{(-i)} = (A - x_i x_i^\top)^{-1}(X^\top y - x_i y_i)$，并左乘 $x_i^\top$ 计算留一预测值 $\hat{y}_i^{(-i)} = x_i^\top \hat{w}^{(-i)}$：
\begin{align*}
  \hat{y}_i^{(-i)}
  &= \left(x_i^\top A^{-1} + \frac{h_{ii}\, x_i^\top A^{-1}}{1 - h_{ii}}\right)(X^\top y - x_i y_i)
   = \frac{1}{1 - h_{ii}}\, x_i^\top A^{-1}(X^\top y - x_i y_i) \\
  &= \frac{1}{1 - h_{ii}}\big(\hat{y}_i - h_{ii} y_i\big).
\end{align*}
因此第 $i$ 个样本的留一预测残差为
\[
  y_i - \hat{y}_i^{(-i)}
  = \frac{(1 - h_{ii}) y_i - (\hat{y}_i - h_{ii} y_i)}{1 - h_{ii}}
  = \frac{y_i - \hat{y}_i}{1 - h_{ii}}.
\]
于是整个 $N$ 折留一法均方误差仅需拟合一次全量模型即可精确算出：$\mathrm{LOOCV} = \frac{1}{N}\sum_{i=1}^N \left(\frac{y_i - \hat{y}_i}{1 - h_{ii}}\right)^2$。（Python 核验：在随机 $30\times 5$ 回归数据上与暴力重训 30 次的最大绝对误差为 $2.11\times 10^{-15}$。）
\end{proof}

\subsection*{习题 4.3.6\quad 单边可实现情形（Realizable PAC）的 $O(1/\varepsilon)$ 加速界}
\begin{proof}
定义“坏假设”集合 $\mathcal{H}_{\mathrm{bad}} := \{h \in \mathcal{H} : R(h) > \varepsilon\}$。因问题可实现（存在 $h^* \in \mathcal{H}$ 使 $R(h^*) = 0$，从而 $R_{\mathrm{emp}}(h^*) = 0$），一致学习器 $\hat{h}$ 总能找到满足 $R_{\mathrm{emp}}(\hat{h}) = 0$ 的假设。

事件 $\{R(\hat{h}) > \varepsilon\}$ 发生，意味着至少存在一个坏假设 $h \in \mathcal{H}_{\mathrm{bad}}$ 在全部 $N$ 个 i.i.d. 训练样本上恰好零失误（$R_{\mathrm{emp}}(h) = 0$）。对任意固定 $h \in \mathcal{H}_{\mathrm{bad}}$，其在单个样本上预测正确的概率为 $1 - R(h) < 1 - \varepsilon$，由样本独立性与初等不等式 $1 - \varepsilon \le e^{-\varepsilon}$ 得
\[
  \Prob\!\big(R_{\mathrm{emp}}(h) = 0\big) = \big(1 - R(h)\big)^N \le (1 - \varepsilon)^N \le e^{-N\varepsilon}.
\]
对所有坏假设取并集界（注意 $|\mathcal{H}_{\mathrm{bad}}| \le |\mathcal{H}|$，实际上因 $h^* \notin \mathcal{H}_{\mathrm{bad}}$ 甚至有 $|\mathcal{H}_{\mathrm{bad}}| \le |\mathcal{H}| - 1$）：
\[
  \Prob\!\big(R(\hat{h}) > \varepsilon\big)
  \le \Prob\!\left(\exists h \in \mathcal{H}_{\mathrm{bad}} : R_{\mathrm{emp}}(h) = 0\right)
  \le \sum_{h \in \mathcal{H}_{\mathrm{bad}}} \Prob\!\big(R_{\mathrm{emp}}(h) = 0\big)
  \le |\mathcal{H}|\, e^{-N\varepsilon}.
\]
令上界 $|\mathcal{H}|\, e^{-N\varepsilon} \le \delta$，两边取对数即得 $N \ge \frac{1}{\varepsilon}\ln\frac{|\mathcal{H}|}{\delta}$。在无噪声可实现情形下，由于单边零误差事件按指数 $(1-\varepsilon)^N$ 衰减（无需依赖均值集中在 $\sqrt{N}$ 尺度的 Hoeffding 界），样本复杂度从 $O(1/\varepsilon^2)$ 锐减至 $O(1/\varepsilon)$。
\end{proof}

\subsection*{习题 4.3.7\quad 岭回归压制高阶多项式方差实验}
\begin{enumerate}[label=(\arabic*), leftmargin=1.8em]
  \item \textbf{核心实现要点}：固定多项式次数 $d = 9$，构造惩罚矩阵 $P = \diag(0, 1, \dots, 1) \in \R^{10\times 10}$（不惩罚截距 $w_0$），用 \texttt{np.linalg.solve(A\_tr.T @ A\_tr + lam * P, A\_tr.T @ y\_tr)} 替代普通最小二乘，在相同 400 轮随机训练抽样（\texttt{seed=7}）上精确统计 $\Bias^2$、$\Var$ 与总 MSE。
  \item \textbf{400 轮实测数值自检表}（不可约噪声 $\sigma^2 = 0.0900$；无正则 $d=3$ 基准总 MSE $= 0.1073$）：
  \begin{center}\small
  \begin{tabular}{@{}lcccccc@{}}
    \toprule
    $\lambda$ & $0$（OLS） & $10^{-8}$ & $10^{-6}$ & $\mathbf{10^{-4}}$ & $10^{-2}$ & $10^{-1}$ \\
    \midrule
    $\Bias^2$ & $0.0000$ & $0.0000$ & $0.0001$ & $\mathbf{0.0007}$ & $0.0364$ & $0.0740$ \\
    $\Var$ & $0.0381$ & $0.0267$ & $0.0220$ & $\mathbf{0.0174}$ & $0.0128$ & $0.0101$ \\
    总 MSE & $0.1281$ & $0.1167$ & $0.1121$ & $\mathbf{0.1081}$ & $0.1392$ & $0.1740$ \\
    \bottomrule
  \end{tabular}
  \end{center}
  \item \textbf{结论解读}：随着 $\lambda$ 从 $0$ 增至 $10^{-4}$，高阶 $d=9$ 多项式的方差被腰斩一半以上（$0.0381 \to 0.0174$），而平方偏差仅微增 $0.0007$，总测试 MSE 从 $0.1281$ 降至 $\mathbf{0.1081}$，已与无正则最优低阶 $d=3$（$0.1073$）几乎持平！若进一步过大（$\lambda \ge 10^{-2}$），则过度压缩高次系数引发欠拟合，偏差飙升至 $0.0740$ 导致总误差反弹。
\end{enumerate}

\subsection*{习题 4.3.8\quad 亲手复现“特征选择数据泄露”假象}
\begin{enumerate}[label=(\arabic*), leftmargin=1.8em]
  \item \textbf{实验设定与实测基准}：使用 \texttt{rng = np.random.default\_rng(2026)} 生成 $N=100, D=5000$ 的纯标准正态噪声矩阵 $X$ 及独立均匀随机二分类标签 $y \in \{0,1\}$，配合 $L_2$ 逻辑回归（$\lambda = 0.1$）执行 5 折交叉验证：
  \begin{itemize}
    \item \textbf{错误流程（全量预筛 Top-20 特征后 CV）}：5 折平均验证准确率竟虚高至 $\mathbf{91.0\%}$！原因是在 5000 个纯噪声特征中，仅靠多重比较的随机巧合就总能挑出 20 列在全部 100 个样本（含验证折）上与 $y$ 高度伪相关的列；
    \item \textbf{正确流程（每折仅用 80 训练样本筛 Top-20 特征）}：5 折平均验证准确率诚实回归至 $\mathbf{58.0\%}$（受 $N=100$ 有限样本标准误 $\approx 5\%$ 波动，跨多随机种子均值严格收于 $50.0\%$）。
  \end{itemize}
  \item \textbf{工程防错铁律}：凡是\textbf{依赖样本统计量或标签}的任何预处理步骤（特征筛选、标准化均值/方差、PCA 主成分、SMOTE 过采样、目标编码），都必须封装进交叉验证的\textbf{单折内部 Pipeline}，严格遵循“只在训练折 \texttt{fit}，在验证折仅 \texttt{transform}”原则。
\end{enumerate}

% =========================================================================
\section*{第 4 章习题解答（共 8 题）}
\addcontentsline{toc}{section}{第 4 章习题解答（共 8 题）}
% =========================================================================

\subsection*{习题 4.1\quad 规范超平面与间隔计算}
已知超平面参数 $w = (3, -4)^\top,\; b = 5$，法向量模长为 $\norm{w}_2 = \sqrt{3^2 + (-4)^2} = 5$。
\begin{enumerate}[label=(\arabic*), leftmargin=1.8em]
  \item \textbf{函数间隔与几何间隔逐点计算}：由 $\hat{\gamma}_i = y_i(w^\top x_i + b)$ 与 $\gamma_i = \hat{\gamma}_i / \norm{w}_2$：
  \begin{itemize}
    \item $x_1 = (1, 1)^\top,\; y_1 = +1$：$w^\top x_1 + b = 3 - 4 + 5 = 4$，故 $\hat{\gamma}_1 = 4,\; \gamma_1 = \frac{4}{5} = 0.8$；
    \item $x_2 = (0, 2)^\top,\; y_2 = -1$：$w^\top x_2 + b = 0 - 8 + 5 = -3$，故 $\hat{\gamma}_2 = (-1)(-3) = 3,\; \gamma_2 = \frac{3}{5} = 0.6$；
    \item $x_3 = (3, 1)^\top,\; y_3 = +1$：$w^\top x_3 + b = 9 - 4 + 5 = 10$，故 $\hat{\gamma}_3 = 10,\; \gamma_3 = \frac{10}{5} = 2.0$；
    \item $x_4 = (-1, 3)^\top,\; y_4 = -1$：$w^\top x_4 + b = -3 - 12 + 5 = -10$，故 $\hat{\gamma}_4 = (-1)(-10) = 10,\; \gamma_4 = \frac{10}{5} = 2.0$。
  \end{itemize}
  \item \textbf{规范超平面缩放与隔离带总宽度}：最小函数间隔为 $\hat{\gamma} = \min_i \hat{\gamma}_i = 3$（在 $x_2$ 处取到）。将 $(w, b)$ 同除以 $3$ 得规范超平面参数：
  \[
    w_{\mathrm{can}} = \frac{w}{3} = \left(1,\; -\frac{4}{3}\right)^\top, \qquad
    b_{\mathrm{can}} = \frac{5}{3}.
  \]
  此时 $\norm{w_{\mathrm{can}}}_2 = \sqrt{1 + \frac{16}{9}} = \frac{5}{3}$，四个样本的几何间隔全部保持不变（最小几何间隔仍为 $\gamma = 1 / \norm{w_{\mathrm{can}}}_2 = \frac{3}{5} = 0.6$），对应的隔离带总宽度为
  \[
    \frac{2}{\norm{w_{\mathrm{can}}}_2} = 2 \times \frac{3}{5} = \frac{6}{5} = \mathbf{1.2}.
  \]
\end{enumerate}

\subsection*{习题 4.2\quad 合法核函数的快速甄别}
\begin{enumerate}[label=(\arabic*), leftmargin=1.8em]
  \item \textbf{$K(x, z) = 3 K_1(x, z) + 2 K_2(x, z)$}：\textbf{必为合法核}。
  对任意有限点集 $\{x_i\}_{i=1}^n$ 与任意向量 $c \in \R^n$，Gram 矩阵满足 $c^\top (3K_1 + 2K_2) c = 3(c^\top K_1 c) + 2(c^\top K_2 c) \ge 0$（正定核锥对非负线性组合封闭）。

  \item \textbf{$K(x, z) = K_1(x, z) - K_2(x, z)$}：\textbf{不一定为合法核}。
  \emph{显式反例}：取 $K_1(x, z) = x^\top z$ 与 $K_2(x, z) = 2 x^\top z$，两者均为合法线性核；但其差 $K(x, z) = -x^\top z$ 在单点 $x = (1, 0, \dots, 0)^\top$ 处就有 $K(x, x) = -1 < 0$，对应 $1\times 1$ Gram 矩阵 $[-1]$ 非半正定，故不是合法核。

  \item \textbf{$K(x, z) = \big(K_1(x, z) + 1\big)^3$}：\textbf{必为合法核}。
  展开得 $K(x, z) = K_1(x, z)^3 + 3 K_1(x, z)^2 + 3 K_1(x, z) + 1$。常数核 $1 = \inner{1}{1}$ 合法；由 Schur 乘积闭合性（习题 4.6），合法核的逐点乘积 $K_1^2, K_1^3$ 均合法；再由非负线性组合闭合性即证。

  \item \textbf{$K(x, z) = \cos(x - z)\;(x, z \in \R)$}：\textbf{必为合法核}。
  由三角和差公式直接分解为二维显式特征内积：
  \[
    \cos(x - z) = \cos x \cos z + \sin x \sin z = \inner{\phi(x)}{\phi(z)}, \qquad
    \phi(x) = \begin{pmatrix} \cos x \\ \sin x \end{pmatrix} \in \R^2.
  \]
  因可表为欧氏内积，其任意 Gram 矩阵 $K = \Phi\Phi^\top \succeq 0$，必为合法核。
\end{enumerate}

\subsection*{习题 4.3\quad 一维非线性可分数据的二次核手算}
样本：$x_1 = -1\ (y_1 = +1)$，$x_2 = 0\ (y_2 = -1)$，$x_3 = 1\ (y_3 = +1)$，特征映射 $\phi(x) = (x, x^2)^\top$：
\[
  \phi(x_1) = (-1, 1)^\top, \qquad \phi(x_2) = (0, 0)^\top, \qquad \phi(x_3) = (1, 1)^\top.
\]
\begin{enumerate}[label=(\arabic*), leftmargin=1.8em]
  \item \textbf{$3\times 3$ 核矩阵 $K$}：计算两两内积 $K_{ij} = \phi(x_i)^\top \phi(x_j) = x_i x_j + x_i^2 x_j^2$：
  $K_{11} = (-1)^2 + 1^2 = 2$，$K_{33} = 1^2 + 1^2 = 2$，$K_{13} = (-1)(1) + (1)(1) = 0$，且含 $x_2 = 0$ 的行列全为 $0$，故
  \[
    K = \begin{pmatrix}
      2 & 0 & 0 \\
      0 & 0 & 0 \\
      0 & 0 & 2
    \end{pmatrix}.
  \]
  \item \textbf{手解硬间隔对偶 QP}：对偶问题~\eqref{eq:svm-kernel-dual} 为
  \[
    \max_{\alpha_1, \alpha_2, \alpha_3 \ge 0}\;
    \alpha_1 + \alpha_2 + \alpha_3 - \frac{1}{2}\big(2\alpha_1^2 + 2\alpha_3^2\big)
    \quad\text{s.t.}\quad \alpha_1 - \alpha_2 + \alpha_3 = 0.
  \]
  将等式约束 $\alpha_2 = \alpha_1 + \alpha_3$ 代入目标函数，消去 $\alpha_2$ 得完全解耦的凹二次函数：
  \[
    W(\alpha_1, \alpha_3) = 2\alpha_1 + 2\alpha_3 - \alpha_1^2 - \alpha_3^2
    = -( \alpha_1 - 1 )^2 - ( \alpha_3 - 1 )^2 + 2.
  \]
  当且仅当 $\alpha_1^* = 1,\; \alpha_3^* = 1$ 时取得全局最大值 $W^* = 2$，回代等式约束得 $\alpha_2^* = \alpha_1^* + \alpha_3^* = 2$。因此
  \[
    \alpha^* = \mathbf{(1,\; 2,\; 1)^\top}
  \]
  （全部 3 个点均为支持向量！）。由 $w^* = \sum_i \alpha_i^* y_i \phi(x_i) = 1\cdot(-1,1)^\top - 2\cdot(0,0)^\top + 1\cdot(1,1)^\top = (0, 2)^\top$，代入任一支持向量（例如 $x_2 = 0, y_2 = -1$）求偏置：
  \[
    b^* = y_2 - \sum_{i=1}^3 \alpha_i^* y_i K(x_i, x_2) = -1 - 0 = \mathbf{-1}.
  \]
  检验其余两点：$y_1 f(x_1) = 1\cdot(2\cdot 1 - 1) = 1$，$y_3 f(x_3) = 1\cdot(2\cdot 1 - 1) = 1$，全部精确落在间隔边界上，且原始目标值 $\frac{1}{2}\norm{w^*}_2^2 = 2 = W^*$ 满足强对偶！

  \item \textbf{最终判别函数与一维决策边界点}：
  \[
    f(x) = \sum_{i=1}^3 \alpha_i^* y_i K(x_i, x) + b^*
    = 1\cdot(-x + x^2) + 1\cdot(x + x^2) - 1
    = \mathbf{2x^2 - 1}.
  \]
  令 $f(x) = 0 \implies 2x^2 - 1 = 0$，得原一维实轴上的两个决策边界点为 $\mathbf{x = -\frac{\sqrt{2}}{2}}$ 与 $\mathbf{x = +\frac{\sqrt{2}}{2}}$（约 $\pm 0.7071$）：区间 $(-\frac{\sqrt{2}}{2}, \frac{\sqrt{2}}{2})$ 内判为 $-1$，两侧判为 $+1$。
\end{enumerate}

\subsection*{习题 4.4\quad 软间隔支持向量状态诊断}
已知 $C = 2.0,\; w^* = (1, 0)^\top,\; b^* = -1$，故判别函数为 $f^*(x) = x_1 - 1$，且四个样本标签均为 $y = +1$。由最优松弛变量 $\xi_i^* = \max(0, 1 - y_i f^*(x_i))$ 及第~\ref{sec:svm-soft} 节 KKT 互补松弛条件（$\alpha_i^*[1 - \xi_i^* - y_i f^*(x_i)] = 0$ 且 $(C - \alpha_i^*)\xi_i^* = 0$），逐一诊断如下：
\begin{center}\small
\begin{tabular}{@{}lccccl@{}}
  \toprule
  \textbf{样本} & $y_i f^*(x_i)$ & $\xi_i^*$ & $\alpha_i^*$ \textbf{状态} & \textbf{几何位置} & \textbf{可否用于求 $b^*$？} \\
  \midrule
  $A(3, 1)$ & $2.0$ & $0$ & $\alpha_A^* = 0$ & 隔离带外侧（非支持向量） & 否 \\
  $B(2, -2)$ & $1.0$ & $0$ & $0 \le \alpha_B^* \le C$（典型为 $(0, C)$） & 恰好落在间隔边界上（自由 SV） & \textbf{是}（若 $0 < \alpha_B^* < C$） \\
  $C(1.5, 0)$ & $0.5$ & $0.5$ & $\alpha_C^* = C = 2.0$ & 侵入隔离带内但仍分类正确 & 否 \\
  $D(0.5, 2)$ & $-0.5$ & $1.5$ & $\alpha_D^* = C = 2.0$ & 越过超平面被误分类 & 否 \\
  \bottomrule
\end{tabular}
\end{center}
\textbf{说明}：对 $C, D$ 两点因 $\xi_i^* > 0$，互补松弛 $(C - \alpha_i^*)\xi_i^* = 0$ 强制 $\alpha_C^* = \alpha_D^* = C = 2.0$；对 $A$ 点因 $y_A f^*(x_A) = 2 > 1$ 且 $\xi_A^* = 0$，必有 $\alpha_A^* = 0$；只有处于间隔边界上的样本 $B$ 满足 $y_B f^*(x_B) = 1$（当其为自由支持向量 $0 < \alpha_B^* < C$ 时，可直接由 $b^* = y_B - (w^*)^\top x_B = 1 - 2 = -1$ 精确解出偏置 $b^*$）。

\subsection*{习题 4.5\quad 支持向量乘子之和与最大间隔的精确恒等式}
\begin{proof}
对于线性可分的硬间隔 SVM，原始问题为严格凸二次规划且存在严格可行点（Slater 条件满足），因此强对偶成立：$p^* = d^*$。
\begin{itemize}
  \item 原始最优值为 $p^* = \frac{1}{2}\norm{w^*}_2^2$；
  \item 由 KKT 平稳性条件 $w^* = \sum_{i=1}^N \alpha_i^* y_i x_i$，对偶目标中的二次项恰好等于
  \[
    \frac{1}{2}\sum_{i=1}^N\sum_{j=1}^N \alpha_i^* \alpha_j^* y_i y_j (x_i^\top x_j)
    = \frac{1}{2}\norm{\sum_{i=1}^N \alpha_i^* y_i x_i}_2^2
    = \frac{1}{2}\norm{w^*}_2^2.
  \]
  故对偶最优值为 $d^* = \sum_{i=1}^N \alpha_i^* - \frac{1}{2}\norm{w^*}_2^2$。
\end{itemize}
由强对偶 $p^* = d^*$ 立得
\[
  \frac{1}{2}\norm{w^*}_2^2 = \sum_{i=1}^N \alpha_i^* - \frac{1}{2}\norm{w^*}_2^2
  \implies \sum_{i=1}^N \alpha_i^* = \norm{w^*}_2^2 = \frac{1}{(\gamma^*)^2},
\]
其中 $\gamma^* = 1/\norm{w^*}_2$ 为最大几何间隔。

\textbf{几何意义}：正负两类数据越贴近、越难线性分开，所能容纳的最大几何间隔 $\gamma^*$ 就越小，因此 $\norm{w^*}_2 = 1/\gamma^*$ 必然越大，支撑起超平面的“对偶支撑力总和” $\sum_{i=1}^N \alpha_i^* = 1/(\gamma^*)^2$ 便按平方反比急剧增大。
\end{proof}

\subsection*{习题 4.6\quad 多项式核的特征维度与舒尔乘积合法性证明}
\begin{proof}
\begin{enumerate}[label=(\arabic*), leftmargin=1.8em]
  \item \textbf{Schur 乘积定理（$A, B \succeq 0 \implies A \had B \succeq 0$）}：因 $A \in \R^{n\times n}$ 对称半正定，做正交谱分解 $A = \sum_{r=1}^n \lambda_r u_r u_r^\top$，其中特征值 $\lambda_r \ge 0$。对任意向量 $c \in \R^n$，直接展开二次型：
  \begin{align*}
    c^\top (A \had B) c
    &= \sum_{i=1}^n \sum_{j=1}^n c_i c_j A_{ij} B_{ij}
     = \sum_{i=1}^n \sum_{j=1}^n c_i c_j \left(\sum_{r=1}^n \lambda_r u_{r,i} u_{r,j}\right) B_{ij} \\
    &= \sum_{r=1}^n \lambda_r \sum_{i=1}^n \sum_{j=1}^n (c_i u_{r,i}) B_{ij} (c_j u_{r,j})
     = \sum_{r=1}^n \lambda_r (c \had u_r)^\top B (c \had u_r).
  \end{align*}
  由于 $B \succeq 0$，每一项 $(c \had u_r)^\top B (c \had u_r) \ge 0$，又 $\lambda_r \ge 0$，故 $c^\top (A \had B) c \ge 0$ 对任意 $c \in \R^n$ 成立，且 $A \had B$ 显然对称，从而 $A \had B \succeq 0$。

  \item \textbf{核乘积与多项式核的合法性}：任取有限点集 $\{x_1, \dots, x_n\} \subset \R^d$，记 $K_1, K_2$ 对应的 Gram 矩阵为 $G^{(1)}, G^{(2)} \succeq 0$。逐点乘积函数 $K(x, z) = K_1(x, z) K_2(x, z)$ 的 Gram 矩阵恰好为 Hadamard 积 $G^{(1)} \had G^{(2)}$，由第 (1) 问立即得 $G^{(1)} \had G^{(2)} \succeq 0$，故 $K_1 K_2$ 必为合法正定核。

  对多项式核 $K(x, z) = (x^\top z + c)^p$（$c \ge 0, p \in \N^+$）：基础线性核 $K_{\mathrm{lin}}(x, z) = x^\top z$（Gram 阵 $XX^\top \succeq 0$）与非负常数核 $K_c(x, z) = c$（Gram 阵 $c\mathbf{1}\mathbf{1}^\top \succeq 0$）均合法，故其和 $K_{\mathrm{base}}(x, z) = x^\top z + c$ 合法；对其反复应用 $p - 1$ 次 Schur 乘积闭合性，立得 $(x^\top z + c)^p$ 对任意维度 $d$ 均为合法正定核（其对应的多项式单项式隐式特征空间维度为 $\binom{d + p}{p}$）。\qedhere
\end{enumerate}
\end{proof}

\subsection*{习题 4.7\quad 原始次梯度求解线性软间隔 SVM 与对偶解一致性校验}
\begin{enumerate}[label=(\arabic*), leftmargin=1.8em]
  \item \textbf{合页损失目标的次梯度公式}：记式~\eqref{eq:svm-hinge} 的目标为 $J(w, b) = \frac{1}{2}\norm{w}_2^2 + C\sum_{i=1}^N \max\!\big(0, 1 - y_i(w^\top x_i + b)\big)$，令活跃违约集合 $\mathcal{M}(w, b) = \{i : y_i(w^\top x_i + b) < 1\}$，则一个有效次梯度为：
  \[
    g_w = w - C\sum_{i \in \mathcal{M}(w, b)} y_i x_i, \qquad
    g_b = -C\sum_{i \in \mathcal{M}(w, b)} y_i.
  \]
  \item \textbf{核心实现与收敛技巧}：采用衰减步长 $\eta_t = \frac{0.05}{\sqrt{t}}$ 配合后半程 Polyak 参数平均（消除折点附近的锯齿抖动），与线性核 $K = X X^\top$ 版 SMO 求解器（$C = 1.0$，充分遍历直至对偶间隙 $< 2\times 10^{-3}$）进行交叉核验。
  \item \textbf{实测一致性自检基准}（\texttt{seed=2026}, $N=100$ 二维含轻度重叠高斯二分类数据，$C = 1.0$）：
  \begin{itemize}
    \item \textbf{对偶线性 SMO}：$w_{\mathrm{SMO}} = (1.31715,\; 0.82475)^\top,\; b_{\mathrm{SMO}} = 0.14974$，原始目标 $J = 7.783541$（对偶值 $7.781543$，间隙 $2.00\times 10^{-3}$）；
    \item \textbf{原始次梯度（Polyak 平均）}：$w_{\mathrm{SG}} = (1.31859,\; 0.82588)^\top,\; b_{\mathrm{SG}} = 0.14972$，原始目标 $J = 7.781608$；
    \item \textbf{相对误差校验}：权重相对误差 $\mathbf{0.118\%}$、偏置相对误差 $\mathbf{0.008\%}$、目标函数值相对误差 $\mathbf{0.025\%}$，全部远小于 $1\%$ 校验红线！
  \end{itemize}
\end{enumerate}

\subsection*{习题 4.8\quad RBF-SVM 的 $(C, \gamma)$ 二维网格搜索与过拟合诊断}
复用 \texttt{exp4\_rbf\_svm\_circles} 的含噪同心圆数据集（\texttt{seed=2026}, 160 训练 / 80 测试）与 SMO 算法（每次重置 \texttt{alpha=0, b=0}，固定 SMO 内核随机种子 \texttt{42}，35 轮扫描），测得完整 $3\times 5$ 对照矩阵（格内依次为 \textbf{测试准确率\% / 支持向量占比\% / 过拟合间隙(训练$-$测试)\%}）：
\begin{center}\footnotesize
\begin{tabular}{@{}lccccc@{}}
  \toprule
  $C \backslash \gamma$ & $\gamma = 0.05$ & $\gamma = 0.2$ & $\gamma = 0.8$ & $\gamma = 4.0$ & $\gamma = 12.0$ \\
  \midrule
  $C = 0.5$ & $73.8 / 88.1 / +6.9$ & $96.2 / 54.4 / -3.8$ & $93.8 / 48.1 / +1.9$ & $95.0 / 65.0 / +1.9$ & $93.8 / 92.5 / +3.8$ \\
  $C = 5.0$ & $96.2 / 49.4 / -3.8$ & $96.2 / 29.4 / -2.5$ & $95.0 / 26.2 / +1.9$ & $93.8 / 49.4 / +3.8$ & $92.5 / 78.1 / +6.2$ \\
  $C = 50.0$ & $96.2 / 31.2 / -0.6$ & $96.2 / 30.6 / +0.0$ & $95.0 / 27.5 / +1.2$ & $95.0 / 48.8 / +3.1$ & $\mathbf{92.5 / 79.4 / +6.9}$ \\
  \bottomrule
\end{tabular}
\end{center}
\textbf{偏差-方差机理剖析}：
\begin{itemize}
  \item \textbf{左下角（低 $C=0.5$ + 低 $\gamma=0.05$）——典型高偏差欠拟合}：核带宽极宽、曲率极平，叠加对违约惩罚很轻，模型画不出足够弯曲的闭环，测试准确率暴跌至 $73.8\%$，且高达 $88.1\%$ 的样本沦为带内支持向量；
  \item \textbf{黄金带（$C \in \{5, 50\},\; \gamma \in \{0.2, 0.8\}$）}：支持向量最稀疏（最低仅占 $\mathbf{26.2\%}$），测试准确率稳定达到 $95.0\%$--$96.2\%$，训练与测试间隙接近零；
  \item \textbf{右下角（高 $C=50$ + 高 $\gamma=12.0$）——最危险的高方差过拟合角落}：高 $\gamma$ 使每个样本的高斯核塌缩成极窄“针尖”（$\sigma = (2\gamma)^{-1/2} \approx 0.20$），失去跨样本共享信息的能力，导致支持向量比例反常飙升至 $79.4\%$；同时高 $C$ 强迫模型绝不放过任何噪声离群点，在噪声环交界处雕出细碎孤岛，使训练准确率飙至 $99.4\%$ 而测试准确率跌至 $92.5\%$（过拟合间隙扩至全场最高的 $+6.9\%$）。
\end{itemize}

% =========================================================================
\section*{第 5 章习题解答（共 8 题）}
\addcontentsline{toc}{section}{第 5 章习题解答（共 8 题）}
% =========================================================================

\subsection*{习题 5.1\quad 三分类节点不纯度计算}
父节点共 10 个样本，三类频数为 $(5, 3, 2)$，概率向量 $p = (0.5, 0.3, 0.2)$。
\begin{enumerate}[label=(\arabic*), leftmargin=1.8em]
  \item \textbf{父节点 Gini 不纯度与香农熵}：
  \begin{align*}
    G(p) &= 1 - (0.5^2 + 0.3^2 + 0.2^2) = 1 - (0.25 + 0.09 + 0.04) = \mathbf{0.6200}, \\
    H(p) &= -\big(0.5\log_2 0.5 + 0.3\log_2 0.3 + 0.2\log_2 0.2\big) \\
         &\approx 0.5000 + 0.5211 + 0.4644 = \mathbf{1.4855\;\text{bits}}.
  \end{align*}
  \item \textbf{分裂后的 Gini 增益与信息增益}：
  \begin{itemize}
    \item 左子节点 $(5, 0, 0)$（权重 $\frac{5}{10} = 0.5$）：完全纯净，故 $G_L = 0,\; H_L = 0$。
    \item 右子节点 $(0, 3, 2)$（权重 $\frac{5}{10} = 0.5$，比例 $0.6, 0.4$）：
    \begin{align*}
      G_R &= 1 - (0.6^2 + 0.4^2) = 1 - 0.52 = \mathbf{0.4800}, \\
      H_R &= -\big(0.6\log_2 0.6 + 0.4\log_2 0.4\big) \approx 0.4422 + 0.5288 = \mathbf{0.9710\;\text{bits}}.
    \end{align*}
  \end{itemize}
  按样本比例加权求差，得
  \begin{align*}
    \Delta G &= G(p) - \big(0.5 G_L + 0.5 G_R\big) = 0.6200 - 0.5 \times 0.4800 = \mathbf{0.3800}, \\
    \operatorname{IG} &= H(p) - \big(0.5 H_L + 0.5 H_R\big) = 1.4855 - 0.5 \times 0.9710 = \mathbf{1.0000\;\text{bits}}.
  \end{align*}
  （注：信息增益恰为精确的 $1.0000\;\text{bit}$！因为按二元分裂区分出占比 $0.5$ 的第 1 类，恰等价于揭示了 1 枚公平硬币的信息量：$H(0.5, 0.3, 0.2) = H(0.5, 0.5) + 0.5 H(0.6, 0.4)$。）
\end{enumerate}

\subsection*{习题 5.2\quad 随机森林 vs GBDT 诊断四连问}
\begin{enumerate}[label=(\arabic*), leftmargin=1.8em]
  \item \textbf{错误}。由集成方差公式~\eqref{eq:rf-var-formula} $\Var(\bar{T}_M) = \rho\sigma^2 + \frac{1-\rho}{M}\sigma^2$，随着随机森林树数 $M$ 增大，第二项 $\frac{1-\rho}{M}\sigma^2$ 严格单调趋近于 $0$，集成方差单调下降并收敛至下界 $\rho\sigma^2$，而偏差保持与单棵树一致不变。因此单纯增加 $M$ \textbf{不会}引起随机森林过拟合，只会使测试误差趋稳（代价是计算开销增加）。
  \item \textbf{正确}。GBDT 采用前向串行加法模型 $F_M(x) = F_0 + \eta\sum_{m=1}^M h_m(x)$，每棵新树都在刻意拟合前序模型的剩余残差。当树数 $M$ 过大且学习率 $\eta = 0.2$ 较大时，模型会不断在函数空间中记住训练噪声，导致方差持续攀升、测试误差掉头恶化（严重过拟合）。
  \item \textbf{正确}。CART 分类树在寻找切分时只依赖单个特征值的\textbf{相对排序次序}（比较 $x_{ij} \le s$）。因为 $x \mapsto \ln(1 + x)$ 在 $x \ge 0$ 上严格单调递增，它完好保留所有样本在该特征上的全序关系，每个候选切分的左右样本划分完全一致，故树结构、不纯度增益及预测结果均完全不变（仅分裂阈值 $s$ 对应变换为 $\ln(1 + s)$）。
  \item \textbf{错误，正好说反！}随机森林（Bagging）的核心任务是\textbf{降方差}，无法降低基模型的偏差，因此必须使用低偏差、高方差的\textbf{深树}（通常不剪枝）作为基学习器；而 GBDT（Boosting）的核心任务是串行\textbf{降偏差}，若单棵树过深（$d=20$）会一步就把训练噪声吸干导致剧烈过拟合，因此 GBDT 偏好高偏差、低方差的\textbf{浅树}（如 $d \in \{2, 3, \dots, 6\}$ 甚至决策树桩）。
\end{enumerate}

\subsection*{习题 5.3\quad 一维回归树手算与代价复杂度剪枝}
给定 4 个样本 $(1, 1),\, (2, 2),\, (3, 9),\, (4, 10)$。
\begin{enumerate}[label=(\arabic*), leftmargin=1.8em]
  \item \textbf{评估全部 3 个中点切分的残差平方和（RSS）}：
  \begin{itemize}
    \item $s = 1.5$：左 $\{1\}$（均值 $1$，RSS $= 0$），右 $\{2, 9, 10\}$（均值 $7$，RSS $= (2-7)^2 + (9-7)^2 + (10-7)^2 = 25 + 4 + 9 = 38$），总 RSS $= \mathbf{38}$；
    \item $s = 2.5$：左 $\{1, 2\}$（均值 $1.5$，RSS $= 0.25 + 0.25 = 0.5$），右 $\{9, 10\}$（均值 $9.5$，RSS $= 0.25 + 0.25 = 0.5$），总 RSS $= \mathbf{1.0}$；
    \item $s = 3.5$：左 $\{1, 2, 9\}$（均值 $4$，RSS $= 9 + 4 + 25 = 38$），右 $\{10\}$（均值 $10$，RSS $= 0$），总 RSS $= \mathbf{38}$。
  \end{itemize}
  因此最优第一次分裂为 $\mathbf{x \le 2.5}$（总误差从根节点骤降至 $1$）。

  \item \textbf{三棵候选树的代价复杂度 $R_\alpha(T) = R(T) + \alpha|T|$ 与最优 $\alpha$ 区间}：
  先算根节点均值 $\bar{y} = \frac{1+2+9+10}{4} = 5.5$，其总残差平方和为
  \[
    R(T_{\mathrm{root}}) = (-4.5)^2 + (-3.5)^2 + 3.5^2 + 4.5^2 = 20.25 + 12.25 + 12.25 + 20.25 = \mathbf{65}.
  \]
  三棵树的代价复杂度分别为：
  \begin{itemize}
    \item 单叶根节点 $T_{\mathrm{root}}$（$|T_{\mathrm{root}}| = 1$）：$R_\alpha(T_{\mathrm{root}}) = 65 + \alpha$；
    \item 1 次分裂树 $T_1$（$|T_1| = 2$）：$R_\alpha(T_1) = 1 + 2\alpha$；
    \item 4 叶完全树 $T_0$（$|T_0| = 4$）：$R_\alpha(T_0) = 0 + 4\alpha = 4\alpha$。
  \end{itemize}
  要求 $T_1$ 严格优于 $T_0$ 与 $T_{\mathrm{root}}$，同时解两个不等式：
  \[
    \begin{cases}
      R_\alpha(T_1) < R_\alpha(T_0) \iff 1 + 2\alpha < 4\alpha \iff \alpha > 0.5, \\
      R_\alpha(T_1) < R_\alpha(T_{\mathrm{root}}) \iff 1 + 2\alpha < 65 + \alpha \iff \alpha < 64.
    \end{cases}
  \]
  故当且仅当 $\mathbf{\alpha \in (0.5,\; 64)}$ 时，1 次分裂树 $T_1$ 严格胜出——几何上也最直观：只要叶节点惩罚适中，模型就会果断抓住从 $2$ 到 $9$ 的巨大跃迁，而拒绝为组内微小的 $1$ 单位噪声继续过度切分。
\end{enumerate}

\subsection*{习题 5.4\quad XGBoost 二阶叶权重与分裂增益手算}
已知 $y = (1, 1, 0, 0),\; p = (0.8, 0.4, 0.6, 0.2),\; \lambda = 1.0,\; \gamma = 0.05$。
\begin{enumerate}[label=(\arabic*), leftmargin=1.8em]
  \item \textbf{逐样本一阶导 $g_i = p_i - y_i$ 与二阶导 $h_i = p_i(1 - p_i)$}：
  \begin{center}\small
  \begin{tabular}{@{}ccccc@{}}
    \toprule
    样本 $i$ & $y_i$ & $p_i$ & $g_i = p_i - y_i$ & $h_i = p_i(1 - p_i)$ \\
    \midrule
    1 & $1$ & $0.8$ & $-0.2$ & $0.16$ \\
    2 & $1$ & $0.4$ & $-0.6$ & $0.24$ \\
    3 & $0$ & $0.6$ & $+0.6$ & $0.24$ \\
    4 & $0$ & $0.2$ & $+0.2$ & $0.16$ \\
    \bottomrule
  \end{tabular}
  \end{center}
  \item \textbf{分裂前父叶最优权重 $w^*$}：汇总 $G = \sum_{i=1}^4 g_i = -0.2 - 0.6 + 0.6 + 0.2 = 0$，$H = \sum_{i=1}^4 h_i = 0.16 + 0.24 + 0.24 + 0.16 = 0.80$。由式~\eqref{eq:xgb-leaf-gain} 得
  \[
    w^* = -\frac{G}{H + \lambda} = -\frac{0}{0.80 + 1.0} = \mathbf{0}.
  \]
  （因为父节点中正负样本的欠预测与过预测误差恰好对称抵消，不分裂时无法对全体给出有效统一修正。）

  \item \textbf{分裂 $I_L = \{1, 2\},\; I_R = \{3, 4\}$ 的增益 $\operatorname{Gain}$ 与决策}：
  汇总左右子叶统计量：
  \begin{align*}
    G_L &= -0.2 - 0.6 = -0.8, & H_L &= 0.16 + 0.24 = 0.40, \\
    G_R &= 0.6 + 0.2 = +0.8, & H_R &= 0.24 + 0.16 = 0.40.
  \end{align*}
  核验守恒：$G_L + G_R = 0 = G$，$H_L + H_R = 0.80 = H$。代入式~\eqref{eq:xgb-split-gain}：
  \begin{align*}
    \operatorname{Gain}
    &= \frac{1}{2}\left[\frac{(-0.8)^2}{0.40 + 1.0} + \frac{(0.8)^2}{0.40 + 1.0} - \frac{0^2}{0.80 + 1.0}\right] - 0.05 \\
    &= \frac{1}{2}\left[\frac{0.64}{1.4} + \frac{0.64}{1.4} - 0\right] - 0.05
    = \frac{16}{35} - \frac{1}{20} = \frac{57}{140} \approx \mathbf{0.4071}.
  \end{align*}
  因为 $\operatorname{Gain} \approx 0.4071 > 0$，\textbf{应当执行该分裂}！分裂后左叶输出正权重 $w_L^* = +\frac{0.8}{1.4} = +\frac{4}{7}$ 抬高正类概率，右叶输出负权重 $w_R^* = -\frac{4}{7}$ 压低负类概率。
\end{enumerate}

\subsection*{习题 5.5\quad 回归中位数叶预测与 $L_1$ 绝对值损失的 GBDT}
\begin{enumerate}[label=(\arabic*), leftmargin=1.8em]
  \item \textbf{证明最小化 $\sum_{i=1}^n |y_i - c|$ 的最优常数是样本中位数}：
  \begin{proof}
  目标函数 $L(c) = \sum_{i=1}^n |c - y_i|$ 是关于 $c \in \R$ 的连续凸函数，故 $c^*$ 为全局极小值点当且仅当 $0 \in \partial L(c^*)$。由次微分加法法则：
  \[
    \partial L(c) = \sum_{i : y_i < c} (+1) + \sum_{i : y_i > c} (-1) + \sum_{i : y_i = c} [-1, 1].
  \]
  记 $n_<(c) = \#\{i : y_i < c\}$、$n_>(c) = \#\{i : y_i > c\}$、$n_=(c) = \#\{i : y_i = c\}$，则
  \[
    0 \in \partial L(c^*)
    \iff \big|n_<(c^*) - n_>(c^*)\big| \le n_=(c^*).
  \]
  将样本排序为 $y_{(1)} \le y_{(2)} \le \dots \le y_{(n)}$：
  \begin{itemize}
    \item 若 $n = 2k + 1$ 为奇数，取 $c^* = y_{(k+1)}$ 时至多有 $k$ 个严格小于 $c^*$、至多 $k$ 个严格大于 $c^*$，满足 $|n_< - n_>| \le n_=$（反之若 $c < y_{(k+1)}$ 则 $n_> - n_< \ge 1 + n_=$）；
    \item 若 $n = 2k$ 为偶数，则闭区间 $[y_{(k)}, y_{(k+1)}]$ 上任一点均满足 $n_< = n_> = k$（端点处亦满足 $|n_< - n_>| \le n_=$），惯例取中点 $\frac{1}{2}(y_{(k)} + y_{(k+1)})$。
  \end{itemize}
  综上所述，$c^* = \operatorname{median}(y_1, \dots, y_n)$ 恰是全局最优常数预测。
  \end{proof}

  \item \textbf{$L_1$-GBDT 的伪残差与叶输出重估机制}：
  \begin{itemize}
    \item \textbf{第 $m$ 轮伪残差}：损失函数 $\ell_1(y_i, F) = |y_i - F|$ 关于预测值 $F$ 的负次梯度为符号函数：
    \[
      r_{im} = -\left.\frac{\partial \ell_1(y_i, F)}{\partial F}\right|_{F = F_{m-1}(x_i)} = \sign\!\big(y_i - F_{m-1}(x_i)\big) \in \{-1, 0, +1\}.
    \]
    \item \textbf{为何树分裂后必须重估叶输出为残差中位数}：拟合回归树 $h_m(x)$ 时仅利用了伪残差 $r_{im} \in \{-1, +1\}$（只包含误差\textbf{方向}信息、彻底丢弃了误差\textbf{幅度}信息！）来划分出若干叶区域 $\{R_{jm}\}$。若直接沿用回归树叶均值（在 $[-1, 1]$ 内），步长幅值完全失真。因此一旦划分好叶区域 $R_{jm}$，标准 $L_1$-GBDT 回到真实 $L_1$ 损失执行最优线搜索：
    \[
      \gamma_{jm} = \argmin_{\gamma \in \R} \sum_{i \in R_{jm}} \big|y_i - F_{m-1}(x_i) - \gamma\big|
      = \operatorname{median}_{i \in R_{jm}}\!\big(y_i - F_{m-1}(x_i)\big),
    \]
    恰好应用了第 (1) 问的结论！这使得 $L_1$-GBDT 既能用平方误差树高效切分特征空间，又能对重尾离群点保持极强鲁棒性。
  \end{itemize}
\end{enumerate}

\subsection*{习题 5.6\quad 随机森林特征子采样相关性极限分析}
由定理~\ref{thm:rf-variance}，记纯 Bagging（$m_{\mathrm{try}} = d$）与随机森林（$m_{\mathrm{try}} = \sqrt{d}$）的集成方差分别为：
\begin{align*}
  V_{\mathrm{bag}}(M) &= \left(\rho_0 + \frac{1 - \rho_0}{M}\right)\sigma^2 = \left(0.60 + \frac{0.40}{M}\right)\sigma^2, \\
  V_{\mathrm{rf}}(M)  &= \left(\rho_1 + \frac{1 - \rho_1}{M}\right)(1.25\sigma^2) = \left(0.3125 + \frac{0.9375}{M}\right)\sigma^2.
\end{align*}
\begin{enumerate}[label=(\arabic*), leftmargin=1.8em]
  \item \textbf{$M \to \infty$ 极限下的方差降幅}：取极限得 $V_{\mathrm{bag}}(\infty) = 0.60\sigma^2$，$V_{\mathrm{rf}}(\infty) = 0.3125\sigma^2$。相对降低百分比为
  \[
    \frac{0.60 - 0.3125}{0.60} \times 100\% = \frac{0.2875}{0.60} \times 100\% = \frac{23}{48} \times 100\% \approx \mathbf{47.92\%}.
  \]
  尽管单棵树方差略增 $25\%$，但由于打破强特征垄断使树间相关性腰斩以上，大森林极限方差反而骤降近一半！
  \item \textbf{随机森林开始严格优于纯 Bagging 的最小树数 $M_{\min}$}：解严格不等式 $V_{\mathrm{rf}}(M) < V_{\mathrm{bag}}(M)$：
  \[
    0.3125 + \frac{0.9375}{M} < 0.60 + \frac{0.40}{M}
    \iff \frac{0.5375}{M} < 0.2875
    \iff M > \frac{0.5375}{0.2875} = \frac{43}{23} \approx 1.8696.
  \]
  因 $M$ 为正整数，故仅需 $\mathbf{M_{\min} = 2}$ 棵树，随机森林即可严格击败纯 Bagging（检验：$M=1$ 时 $1.25 > 1.00$，单棵树 Bagging 略优；$M=2$ 时 $V_{\mathrm{rf}}(2) = 0.78125\sigma^2 < 0.80\sigma^2 = V_{\mathrm{bag}}(2)$）。
\end{enumerate}

\subsection*{习题 5.7\quad 手写随机森林袋外（OOB）验证器}
\begin{enumerate}[label=(\arabic*), leftmargin=1.8em]
  \item \textbf{核心实现代码}（嵌入 \texttt{exp5\_trees\_rf\_gbdt} 随机森林循环）：
  \begin{lstlisting}
oob_sum, oob_cnt = np.zeros(N_tr), np.zeros(N_tr)
for m in range(1, M_max + 1):
    b_idx = rng.integers(0, N_tr, size=N_tr)
    oob_mask = ~(np.bincount(b_idx, minlength=N_tr) > 0)
    tree = fit_reg_tree(X_tr[b_idx], y_tr[b_idx],
                        max_depth=8, min_leaf=2, max_features=3, rng_local=rng)
    oob_sum[oob_mask] += pred_tree(tree, X_tr[oob_mask])
    oob_cnt[oob_mask] += 1
    valid = oob_cnt > 0
    oob_mse = np.mean((oob_sum[valid] / oob_cnt[valid] - y_tr[valid]) ** 2)
  \end{lstlisting}
  \item \textbf{实测自检基准与有限样本现象解析}：
  \begin{itemize}
    \item \textbf{OOB 树数占比收敛于 $e^{-1} \approx 36.79\%$}：在 $N_{\mathrm{tr}} = 220$ 上实测 $M = 10, 20, 40, 60$ 时每样本平均 OOB 树数占比依次为 $\mathbf{37.05\%,\; 36.50\%,\; 36.83\%,\; 36.83\%}$，精准贴合理论极限 $36.8\%$。
    \item \textbf{OOB MSE 与含噪测试 MSE 一致性检验}：注意 OOB 评估对象是\textbf{含噪训练标签} $y_{\mathrm{tr}} = f(x) + \varepsilon$，因此它估计的是含噪测试 MSE $= \text{clean test MSE} + \sigma^2$（$\sigma^2 = 0.49$）。在 10 组不同随机种子（$N_{\mathrm{tr}}=220, N_{\mathrm{te}}=1000, M=60$）上跨种子平均，测得 $\overline{\mathrm{OOB\ MSE}} = \mathbf{1.3826}$，而独立含噪测试集均值为 $\mathbf{1.3762}$，相对偏差仅 $\mathbf{0.47\%}$（远小于 $10\%$），严格证实 OOB 是无偏的“免费交叉验证”！若只看单一小样本种子（如 \texttt{seed=2026}），因该批 220 个训练点恰好抽到偏低的经验噪声方差（$0.4279 < 0.49$）与略易拟合的区域，单次 OOB MSE（$1.0830$）与含噪测试 MSE（$1.5438$）会出现有限样本抽样起伏，增大 $N_{\mathrm{tr}}$ 或做多轮平均即迅速收敛到 $1\%$ 以内。
  \end{itemize}
\end{enumerate}

\subsection*{习题 5.8\quad GBDT 过拟合警报与早停（Early Stopping）实验}
基于 \texttt{exp5\_trees\_rf\_gbdt} 数据集（\texttt{seed=2026}, $N_{\mathrm{tr}}=220, N_{\mathrm{te}}=220$，测试 MSE 为无噪真值误差），将迭代拉长至 $M = 250$：
\begin{enumerate}[label=(\arabic*), leftmargin=1.8em]
  \item \textbf{两组超参数实测轨迹对比}：
  \begin{center}\footnotesize\setlength{\tabcolsep}{4pt}
  \begin{tabular}{@{}lcccc@{}}
    \toprule
    \textbf{超参数组} & 最优测试 MSE（轮数） & $M=50$（训/测） & $M=100$（训/测） & $M=250$（训/测） \\
    \midrule
    (A) $d=3,\; \eta=0.05$ & $\mathbf{0.4430}$（$M=250$） & $0.7198 / 1.0827$ & $0.3293 / 0.5970$ & $0.1146 / \mathbf{0.4430}$ \\
    (B) $d=6,\; \eta=0.30$ & $0.7112$（$M=22$） & $0.0002 / 0.7249$ & $0.0000 / 0.7277$ & $0.0000 / 0.7278$ \\
    \bottomrule
  \end{tabular}
  \end{center}
  \textbf{现象剖析}：(B) 组深树快步长早在 $\mathbf{M = 22}$ 轮就到达测试误差谷底（$0.7112$），随后在 $M = 50$ 时训练 MSE 已经几乎清零（$0.0002$），测试误差反转上扬并在高位停滞（若进一步取消叶样本限制至 \texttt{min\_leaf=1}，测试 MSE 会直接反弹至 $0.91$ 以上，叠加噪声方差 $0.49$ 后含噪测试 MSE 超过 $1.40$！）；反观 (A) 组“浅树 + 慢步长”虽然前期收敛从容，但 250 轮全程稳定下降至 $\mathbf{0.4430}$，不仅彻底避开过拟合，最终精度比 (B) 组最优时刻还低出近 $38\%$！

  \item \textbf{20\% 留出验证集早停策略（\texttt{patience = 15}）}：将 220 个训练样本切分为 176 训练 + 44 验证，每轮追踪验证集 MSE，若连续 15 轮未刷新历史最小值则终止并回滚至 \texttt{best\_M}：
  \begin{itemize}
    \item 对 (B) 组（$d=6, \eta=0.30$）：早停机制精准在 $M = 44$ 轮拉响刹车并回滚至 $\mathbf{\texttt{best\_M} = 29}$（验证 MSE $= 0.9414$，测试 MSE 锁定在 $0.6780$），避免后面 200 余轮无意义且有害的过拟合计算；
    \item 对 (A) 组（$d=3, \eta=0.05$）：验证 MSE 持续改善至 $M = 250$ 仍未触发停机（验证 MSE $= 0.7913$，测试 MSE $= 0.5059$），再次印证“浅树 + 小学习率 + 早停”是 GBDT 的黄金工程法则。
  \end{itemize}
\end{enumerate}

% =========================================================================
\section*{第 6 章习题解答（共 8 题）}
\addcontentsline{toc}{section}{第 6 章习题解答（共 8 题）}
% =========================================================================

\subsection*{习题 6.1\quad GMM 独立参数量计数与自由度对比}
设维度 $D = 10$、簇数 $K = 4$：
\begin{enumerate}[label=(\arabic*), leftmargin=1.8em]
  \item \textbf{K-Means}：只需估计 $K$ 个 $D$ 维簇中心 $\mu_1, \dots, \mu_K$，独立参数总数为
  \[
    p_{\mathrm{KM}} = K D = 4 \times 10 = \mathbf{40}.
  \]
  \item \textbf{对角协方差 GMM}：包括 $K - 1$ 个独立混合权重（因 $\sum_k \pi_k = 1$）、$KD$ 个均值参数、以及每个簇 $D$ 个对角方差参数：
  \[
    p_{\mathrm{diag}} = (K - 1) + K D + K D = 3 + 40 + 40 = \mathbf{83}.
  \]
  \item \textbf{满协方差 GMM}：每个对称 $D\times D$ 协方差矩阵 $\Sigma_k$ 有 $\frac{D(D+1)}{2}$ 个独立元素，总参数数为
  \[
    p_{\mathrm{full}} = (K - 1) + K D + K\frac{D(D+1)}{2} = 3 + 40 + 4 \times 55 = \mathbf{263}.
  \]
\end{enumerate}
\textbf{高维小样本选型启示}：满协方差 GMM 的参数量随特征维度 $D$ 按\textbf{平方阶 $O(KD^2)$} 爆炸。当样本量 $N$ 不大（例如某簇有效样本数 $N_k \le D$）时，经验协方差矩阵 $\hat{\Sigma}_k$ 必然秩亏奇异，导致行列式 $|\hat{\Sigma}_k| \to 0$、对数似然发散至 $+\infty$ 以及严重过拟合；改用对角协方差（$O(KD)$）或球形协方差（$O(KD + K)$）能将自由度压低一个数量级，换取数值可逆性与低方差泛化。

\subsection*{习题 6.2\quad K-Means++ 抽样概率手算}
一维样本为 $x = (0,\, 1,\, 4,\, 6)$。
\begin{enumerate}[label=(\arabic*), leftmargin=1.8em]
  \item \textbf{已选 $\mu_1 = 0$，抽取第二个质心 $\mu_2$}：计算各点到 $\mu_1 = 0$ 的平方距离 $D(x_i)^2 = (x_i - 0)^2$：
  \[
    D^2 = (0^2,\; 1^2,\; 4^2,\; 6^2) = (0,\; 1,\; 16,\; 36), \qquad \sum_{i=1}^4 D(x_i)^2 = 53.
  \]
  按平方距离归一化得精确抽样概率：
  \[
    P(\mu_2 = 1) = \frac{1}{53} \approx 0.0189, \qquad
    P(\mu_2 = 4) = \frac{16}{53} \approx 0.3019, \qquad
    P(\mu_2 = 6) = \frac{36}{53} \approx 0.6792
  \]
  （$x_1 = 0$ 的概率为 $0$）。可见最远点 $x_4 = 6$ 占据了超过三分之二的抽中机会。

  \item \textbf{已选 $\{\mu_1, \mu_2\} = \{0, 6\}$，抽取第三个质心 $\mu_3$}：计算各点到最近已选质心的距离 $D(x_i) = \min(|x_i - 0|,\, |x_i - 6|)$：
  \begin{itemize}
    \item $x_2 = 1$：$D(1) = \min(1, 5) = 1 \implies D(1)^2 = 1$；
    \item $x_3 = 4$：$D(4) = \min(4, 2) = 2 \implies D(4)^2 = 4$；
    \item $x_1 = 0$ 与 $x_4 = 6$ 均已有 $D = 0$。
  \end{itemize}
  平方距离总和为 $1 + 4 = 5$，故精确抽样概率为
  \[
    P(\mu_3 = 1) = \frac{1}{5} = \mathbf{0.2000}, \qquad
    P(\mu_3 = 4) = \frac{4}{5} = \mathbf{0.8000}.
  \]
\end{enumerate}

\subsection*{习题 6.3\quad 硬币投掷（伯努利混合）模型的 EM 推导}
设参数 $\theta = (\theta_A, \theta_B)$，先验 $\pi_A = \pi_B = 0.5$ 已知。给定硬币 $k \in \{A, B\}$，第 $i$ 轮（共抛 $L$ 次、出现 $h_i$ 次正面）的二项似然为
\[
  P(h_i \mid z_i = k; \theta) = \binom{L}{h_i}\theta_k^{h_i}(1 - \theta_k)^{L - h_i}.
\]
\begin{enumerate}[label=(\arabic*), leftmargin=1.8em]
  \item \textbf{E-步后验响应度}：由贝叶斯定理，分子分母中公共常数 $\frac{1}{2}\binom{L}{h_i}$ 全部精确约去：
  \[
    \gamma_{iA} = P\!\left(z_i = A \mid h_i; \theta^{(t)}\right)
    = \frac{\big(\theta_A^{(t)}\big)^{h_i}\big(1 - \theta_A^{(t)}\big)^{L - h_i}}
           {\big(\theta_A^{(t)}\big)^{h_i}\big(1 - \theta_A^{(t)}\big)^{L - h_i} + \big(\theta_B^{(t)}\big)^{h_i}\big(1 - \theta_B^{(t)}\big)^{L - h_i}},
  \]
  且 $\gamma_{iB} = 1 - \gamma_{iA}$。

  \item \textbf{M-步闭式更新}：写出完全数据期望对数似然 $Q(\theta \mid \theta^{(t)})$（省略与 $\theta$ 无关的常数项）：
  \[
    Q(\theta \mid \theta^{(t)}) = \sum_{i=1}^N \Big[
      \gamma_{iA}\big(h_i \ln\theta_A + (L - h_i)\ln(1 - \theta_A)\big)
      + \gamma_{iB}\big(h_i \ln\theta_B + (L - h_i)\ln(1 - \theta_B)\big)
    \Big] + \text{const}.
  \]
  关于 $\theta_A$ 求偏导并置零：
  \[
    \frac{\partial Q}{\partial \theta_A}
    = \frac{\sum_{i=1}^N \gamma_{iA} h_i}{\theta_A} - \frac{\sum_{i=1}^N \gamma_{iA}(L - h_i)}{1 - \theta_A} = 0
    \implies \theta_A^{(t+1)} = \frac{\sum_{i=1}^N \gamma_{iA} h_i}{\sum_{i=1}^N \gamma_{iA} L}.
  \]
  同理可得硬币 $B$ 的闭式更新：
  \[
    \theta_B^{(t+1)} = \frac{\sum_{i=1}^N \gamma_{iB} h_i}{\sum_{i=1}^N \gamma_{iB} L}.
  \]
  直觉极其自然：分子是归因于硬币 $k$ 的\textbf{期望正面总次数}，分母是归因于硬币 $k$ 的\textbf{期望总抛掷次数}，即“按软权重加权的经验正面频率”。
\end{enumerate}

\subsection*{习题 6.4\quad 例~\ref{ex:gmm-hand} 的对称检验与硬 K-Means 对比}
一维数据 $x = (-2, -1, 2, 3)$，初始中心 $\mu^{(0)} = (-1, 1)$。
\begin{enumerate}[label=(\arabic*), leftmargin=1.8em]
  \item \textbf{标准 K-Means（Lloyd 算法）完整迭代过程}：
  \begin{itemize}
    \item \textbf{第 1 轮指派}：计算各点到 $\mu_1^{(0)} = -1$ 与 $\mu_2^{(0)} = 1$ 的距离：
    $-2$ 距 $(-1, 1)$ 为 $(1, 3) \to$ 簇 1；
    $-1$ 距 $(-1, 1)$ 为 $(0, 2) \to$ 簇 1；
    $2$ 距 $(-1, 1)$ 为 $(3, 1) \to$ 簇 2；
    $3$ 距 $(-1, 1)$ 为 $(4, 2) \to$ 簇 2。
    故指派结果为 $c^{(1)} = (1, 1, 2, 2)$。
    \item \textbf{第 1 轮质心更新}：$\mu_1^{(1)} = \frac{-2 + (-1)}{2} = \mathbf{-1.5}$，$\mu_2^{(1)} = \frac{2 + 3}{2} = \mathbf{2.5}$。
    \item \textbf{第 2 轮验证收敛}：以 $\mu^{(1)} = (-1.5, 2.5)$ 重新指派（中垂线位于 $x = 0.5$），指派结果仍为 $(1, 1, 2, 2)$，质心不再变化，算法在 \textbf{1 轮}即达收敛！
    \item \textbf{最终目标值（惯量）}：
    \[
      J = (-2 + 1.5)^2 + (-1 + 1.5)^2 + (2 - 2.5)^2 + (3 - 2.5)^2 = 0.25 + 0.25 + 0.25 + 0.25 = \mathbf{1.0}.
    \]
  \end{itemize}
  \item \textbf{与 GMM-EM 第 1 轮结果 $(-1.4875, 2.2683)$ 的差异来源}：
  硬 K-Means 在第 1 轮做了“非黑即白”的 $0/1$ 切割，把 $x = -1$ 百分之百划归左簇；而在等方差 GMM-EM 的第 1 轮 E-步中，由于初始中心 $-1$ 与 $+1$ 距离 $x = -1$ 的方差尺度相当（$\sigma^2 = 1$），由软后验公式 $\gamma_{i2} = \frac{1}{1 + e^{-2x_i}}$，样本 $x = -1$ 仍有高达 $\gamma_{22} = \frac{1}{1 + e^{2}} \approx \mathbf{11.92\%}$ 的概率归属右簇（而 $x = 2$ 归属左簇的权重仅 $\frac{1}{1+e^4} \approx 1.80\%$，不对称！）。正是这 $11.92\%$ 的软权重在 M-步加权平均中把右簇均值 $\mu_2^{(1)}$ 从 $2.5$ 明显向左拽至 $2.2683$（左簇均值仅被 $x=2$ 极微弱地右拉至 $-1.4875$）。随着后续 EM 迭代推进方差缩小，软响应度迅速锐化，最终两簇中心也会收敛到非常接近 $(-1.5, 2.5)$。
\end{enumerate}

\subsection*{习题 6.5\quad ELBO 在切点处的一阶梯度一致性}
\begin{proof}
固定第 $t$ 轮的 E-步最优分布 $q^{(t+1)}(z) = p_{\theta^{(t)}}(z \mid x)$，定义边际对数似然与 ELBO 之差函数：
\[
  d(\theta) := \ln p_\theta(x) - \mathcal{L}\big(q^{(t+1)}, \theta\big).
\]
由精确 ELBO--KL 分解恒等式（定理~\ref{thm:elbo-decomp}），对任意参数 $\theta$ 恒有
\[
  d(\theta) = \DKL\!\Big(q^{(t+1)}(z) \,\big\|\, p_\theta(z \mid x)\Big) \ge 0.
\]
又在切点 $\theta = \theta^{(t)}$ 处，$p_{\theta^{(t)}}(z \mid x) = q^{(t+1)}(z)$，故
\[
  d(\theta^{(t)}) = \DKL\!\Big(q^{(t+1)} \,\big\|\, q^{(t+1)}\Big) = 0.
\]
这表明 $\theta = \theta^{(t)}$ 是非负可微函数 $d(\theta)$ 的无约束全局极小值点！由多元微积分 Fermat 一阶必要条件，必有
\[
  \nabla_\theta d(\theta)\Big|_{\theta = \theta^{(t)}} = 0
  \implies \nabla_\theta \mathcal{L}\big(q^{(t+1)}, \theta\big)\Big|_{\theta = \theta^{(t)}}
  = \nabla_\theta \ln p_\theta(x)\Big|_{\theta = \theta^{(t)}}.
\]
\textbf{推论（EM 不动点必为真实对数似然驻点）}：若 $\theta^*$ 是 EM 算法的不动点，且在参数空间开集内部达到 M-步极大值，则对固定 $q^* = p_{\theta^*}(z \mid x)$，函数 $\theta \mapsto \mathcal{L}(q^*, \theta)$ 在 $\theta = \theta^*$ 处取得极大值，因而 $\nabla_\theta \mathcal{L}(q^*, \theta)\big|_{\theta=\theta^*} = 0$。结合上述梯度一致性定理，立得 $\nabla_\theta \ln p_\theta(x)\big|_{\theta=\theta^*} = 0$，即 $\theta^*$ 必为真实边际对数似然的驻点。
\end{proof}

\subsection*{习题 6.6\quad 高斯混合马氏距离决策边界}
在等权重 $\pi_1 = \pi_2 = 0.5$ 下，后验相等边界 $\gamma_{i1} = \gamma_{i2} = 0.5$ 等价于两个高斯密度的对数似然相等：$\ln \mathcal{N}(x \mid \mu_1, \Sigma_1) = \ln \mathcal{N}(x \mid \mu_2, \Sigma_2)$。消去常数 $-\frac{D}{2}\ln(2\pi)$ 并乘以 $-2$，得精确边界方程：
\[
  (x - \mu_1)^\top \Sigma_1^{-1} (x - \mu_1) + \ln|\Sigma_1|
  = (x - \mu_2)^\top \Sigma_2^{-1} (x - \mu_2) + \ln|\Sigma_2|.
\]
按 $x$ 的次数整理为二次型方程：
\[
  x^\top \big(\Sigma_1^{-1} - \Sigma_2^{-1}\big) x
  - 2\big(\Sigma_1^{-1}\mu_1 - \Sigma_2^{-1}\mu_2\big)^\top x
  + \big(\mu_1^\top\Sigma_1^{-1}\mu_1 - \mu_2^\top\Sigma_2^{-1}\mu_2 + \ln|\Sigma_1| - \ln|\Sigma_2|\big) = 0.
\]
\begin{enumerate}[label=(\arabic*), leftmargin=1.8em]
  \item \textbf{共享协方差情形（$\Sigma_1 = \Sigma_2 = \Sigma$，对应 LDA）}：二次项矩阵 $\Sigma_1^{-1} - \Sigma_2^{-1} = 0$ 以及对数行列式差 $\ln|\Sigma_1| - \ln|\Sigma_2| = 0$ 同时精确消去！方程退化为一条直线（仿射超平面）：
  \[
    \underbrace{\big[\Sigma^{-1}(\mu_1 - \mu_2)\big]^\top}_{w^\top} x
    - \frac{1}{2}\big(\mu_1^\top \Sigma^{-1}\mu_1 - \mu_2^\top\Sigma^{-1}\mu_2\big) = 0,
  \]
  等价于 $w^\top\!\left(x - \frac{\mu_1 + \mu_2}{2}\right) = 0$。该直线穿过两质心中点，但其法向量为 $w = \Sigma^{-1}(\mu_1 - \mu_2)$——只要 $\Sigma$ 不是球形矩阵 $\sigma^2 I$（存在主轴拉伸与倾斜相关），$\Sigma^{-1}(\mu_1 - \mu_2)$ 便不再平行于 $\mu_1 - \mu_2$，决策直线会顺着椭圆长轴方向发生偏转（这正是第 6.4 节倾斜双椭圆上 GMM 完胜 K-Means 垂直平分线的几何奥秘！）。

  \item \textbf{异方差情形（$\Sigma_1 \neq \Sigma_2$，对应 QDA）}：对称矩阵 $Q := \Sigma_1^{-1} - \Sigma_2^{-1} \neq 0$，方程保留非零二次项 $x^\top Q x + b^\top x + c = 0$。在二维平面上，依 $Q$ 的两个特征值同号、一零一非零、或异号，该方程依次刻画出\textbf{椭圆、抛物线或双曲线}等一般二次曲线决策边界。
\end{enumerate}

\subsection*{习题 6.7\quad K-Means++ 实现与多重启稳定性对比}
\begin{enumerate}[label=(\arabic*), leftmargin=1.8em]
  \item \textbf{K-Means++ 初始化纯 NumPy 核心代码}：
  \begin{lstlisting}
def init_kmeans_pp(X, K, rng):
    N = len(X)
    idx = [int(rng.integers(0, N))]
    for _ in range(1, K):
        d2 = np.min(np.sum((X[:, None, :] - X[idx][None, :, :]) ** 2, axis=2), axis=1)
        idx.append(int(rng.choice(N, p=d2 / np.sum(d2))))
    return X[idx].copy()
  \end{lstlisting}
  \item \textbf{50 次独立重启实测自检表}（$N=300, K=5$ 个分离二维高斯团簇，\texttt{data\_seed=2026}, \texttt{trial\_seed=42}）：
  \begin{center}\footnotesize\setlength{\tabcolsep}{4pt}
  \begin{tabular}{@{}lccccc@{}}
    \toprule
    \textbf{初始化方法} & 惯量均值 & 惯量中位数 & 惯量最差值 & 平均轮数 & 劣质解（$J>550$） \\
    \midrule
    完全随机均匀初始化 & $1082.93$ & $532.96$ & $2975.93$ & $7.72$ 轮 & $16 / 50$（$32\%$） \\
    \textbf{K-Means++ 初始化} & $\mathbf{681.42}$ & $\mathbf{532.96}$ & $\mathbf{2054.95}$ & $\mathbf{3.82}$ \textbf{轮} & $\mathbf{5 / 50}$（$\mathbf{10\%}$） \\
    \bottomrule
  \end{tabular}
  \end{center}
  \textbf{实验结论}：两种方法的中位数均指向同一全局最优解 $J^* = 532.96$，但随机均匀初始化有近三分之一（$32\%$）的概率把两个质心挤进同一团簇而漏掉远处簇，平均惯量高达 $1082.93$；K-Means++ 将陷落劣质局部极小的频率压低至原来的三分之一（$10\%$），平均收敛步数同步砍半（$7.72 \to 3.82$ 轮）。
\end{enumerate}

\subsection*{习题 6.8\quad GMM-EM 协方差类型对比与 BIC 自动选 $K$ 实验}
复用 \texttt{exp6\_kmeans\_vs\_gmm\_em} 的倾斜平行双椭圆数据集（\texttt{seed=2026}, $N = 360, D = 2$）：
\begin{enumerate}[label=(\arabic*), leftmargin=1.8em]
  \item \textbf{对角协方差 GMM vs 满协方差 GMM 对比}：在 M-步中把协方差更新替换为仅保留对角方差向量 \texttt{v = (gamma[:, k] @ (diff ** 2)) / Nk[k] + 1e-6} 及 \texttt{Sigmas[k] = np.diag(v)}（5 次随机种子取最大似然）：
  \begin{itemize}
    \item \textbf{满协方差 GMM}：聚类纯度 $\mathbf{100.00\%}$，最终对数似然 $\ell = -1286.45$；
    \item \textbf{对角协方差 GMM}：聚类纯度断崖式暴跌至 $\mathbf{55.00\%}$（下降了整整 $45$ 个百分点，与 K-Means 的球形溃败几乎一样惨！），最终对数似然骤降至 $\ell = -1525.01$。
  \end{itemize}
  \textbf{几何成因}：对角协方差强行令非对角项 $\Sigma_{12} = 0$，把每个高斯等高线死锁在\textbf{水平/竖直轴对齐椭圆}内，完全丧失了旋转椭圆去贴合 $45^\circ$ 倾斜双带的能力；因此它被迫像 K-Means 一样从腰部横劈两条平行带（把左下两半合成一簇、右上两半合成另一簇）。

  \item \textbf{满协方差 GMM 遍历 $K \in \{1, 2, 3, 4, 5\}$ 的 BIC 自动选簇实测表}：
  二维满协方差 GMM 的独立参数量为 $p_K = (K-1) + 2K + 3K = 6K - 1$，$\operatorname{BIC}(K) = -2\ell(\hat{\theta}_K) + p_K \ln N$（$\ln 360 \approx 5.8861$）。每组运行 3 个随机种子（各 80 轮 EM）取最大似然，得实测结果：
  \begin{center}\small
  \begin{tabular}{@{}ccccc@{}}
    \toprule
    簇数 $K$ & 独立参数量 $p_K = 6K - 1$ & 最大对数似然 $\ell(\hat{\theta}_K)$ & 复杂度惩罚 $p_K\ln 360$ & $\operatorname{BIC}(K)$ \\
    \midrule
    $K = 1$ & $5$ & $-1520.36$ & $29.43$ & $3070.15$ \\
    $\mathbf{K = 2}$ & $\mathbf{11}$ & $\mathbf{-1286.45}$ & $\mathbf{64.75}$ & $\mathbf{2637.65}$ \textbf{（全局最小）} \\
    $K = 3$ & $17$ & $-1283.36$ & $100.06$ & $2666.77$ \\
    $K = 4$ & $23$ & $-1275.51$ & $135.38$ & $2686.40$ \\
    $K = 5$ & $29$ & $-1274.24$ & $170.70$ & $2719.19$ \\
    \bottomrule
  \end{tabular}
  \end{center}
  \textbf{核心洞察}：训练对数似然 $\ell(\hat{\theta}_K)$ 随 $K$ 从 $1$ 增至 $5$ \textbf{单调严格上升}（$-1520.36 \to -1274.24$），因而绝不能直接看似然选簇数；而 $\operatorname{BIC}(K)$ 在 $K=1 \to 2$ 时因似然暴涨 $+233.9$ 而陡降 $432.5$，但在 $K \ge 3$ 之后每多拆一个冗余簇带来的微小似然收益远不足以抵消 $+6\ln 360 \approx +35.3$ 的复杂度罚金，从而精准在真实簇数 $\mathbf{K = 2}$ 处锁定全局最小值！
\end{enumerate}


\end{document}
